% !TeX program = xelatex
% 数学 LaTeX 排版 Skill v4.0.0；版式保持 v3 金标准. XeLaTeX 编译至少两次. 不分发字体.
% WZ-LOCKED-CLASS-BEGIN
\begin{filecontents*}[overwrite]{elegantbook-original-adapter.cls}
\NeedsTeXFormat{LaTeX2e}
\ProvidesClass{elegantbook-original-adapter}[2026/09/23 v3.0 fixed mathematical book environments]
\LoadClass[10pt,letterpaper,oneside]{book}
\RequirePackage[UTF8,scheme=plain,fontset=fandol]{ctex}
\RequirePackage{fontspec}
\setmainfont{cmunrm.otf}[BoldFont=cmunbx.otf,ItalicFont=cmunti.otf,BoldItalicFont=cmunbi.otf]
\setCJKfamilyfont{sourceheader}[ItalicFont=FandolKai-Regular.otf]{FandolKai-Regular.otf}
\RequirePackage{amsmath,amssymb,mathrsfs}
\RequirePackage{amsthm,etoolbox}
\RequirePackage{xcolor,geometry,setspace,fancyhdr,enumitem,needspace,xurl}
\definecolor{structurecolor}{HTML}{880E4F}
\definecolor{main}{HTML}{9C27B0}
\geometry{paperwidth=612bp,paperheight=792bp,left=20mm,right=20mm,top=95.04bp,bottom=20mm,headheight=12pt,headsep=25pt,footskip=30pt}
\setstretch{1.6}
\setlength{\parindent}{2em}
\setlength{\parskip}{0pt}
\setlist{nosep,leftmargin=2em}
\AtBeginDocument{\setlength{\abovedisplayskip}{4pt plus 1pt minus 1pt}\setlength{\belowdisplayskip}{4pt plus 1pt minus 1pt}\setlength{\abovedisplayshortskip}{2pt}\setlength{\belowdisplayshortskip}{3pt}}
\fancyhf{}
\fancyhead[L]{{\itshape\CJKfamily{sourceheader}\rightmark}}
\fancyhead[R]{{\itshape\CJKfamily{sourceheader}\leftmark}}
\fancyfoot[C]{\thepage}
\renewcommand{\headrulewidth}{0.4pt}
\renewcommand{\headrule}{{\color{structurecolor}\hrule height\headrulewidth width\headwidth\vskip-\headrulewidth}}
\pagestyle{fancy}
\fancypagestyle{cover}{\fancyhf{}\fancyfoot[C]{\thepage}\renewcommand{\headrulewidth}{0pt}\renewcommand{\headrule}{}}
\RequirePackage{hyperref}
\hypersetup{unicode=true,colorlinks=true,linkcolor=structurecolor,urlcolor=structurecolor,citecolor=structurecolor,linktoc=section,bookmarksnumbered=true,bookmarksopen=true,pdfborder={0 0 0},pdfstartview=FitH}
\setcounter{tocdepth}{1}
\renewcommand*\l@section{\@dottedtocline{1}{1.5em}{2.4em}}
\renewcommand{\thesection}{\thechapter.\arabic{section}}
\newcommand{\originalchapter}[1]{\par\addvspace{14pt}\Needspace{5\baselineskip}\refstepcounter{chapter}\setcounter{section}{0}\addcontentsline{toc}{chapter}{\protect\numberline{\thechapter}#1}\markboth{\thechapter\quad #1}{}\begingroup\centering\color{structurecolor}\bfseries\fontsize{14.4}{20}\selectfont\thechapter\quad #1\par\endgroup\nobreak\vspace{8pt}}
\newcommand{\originalsection}[1]{\par\addvspace{9pt}\Needspace{4\baselineskip}\refstepcounter{section}\addcontentsline{toc}{section}{\protect\hspace*{1.5em}\protect\numberline{\protect\textcolor{structurecolor}{\thesection}}#1}\markright{\thesection\quad #1}\noindent{\fontsize{12}{17}\selectfont\color{structurecolor}\thesection\quad}{\bfseries\fontsize{12}{17}\selectfont #1}\par\nobreak\vspace{4pt}}

% Semantic environments are registered once here, not re-designed in manuscripts.
\newtheoremstyle{wzstatement}{6pt}{5pt}
  {\normalfont\kaishu}{0pt}{\normalfont\bfseries\color{structurecolor}}
  {}{.7em}{\thmname{#1}\thmnumber{ #2}\thmnote{\normalfont\ (#3)}}
\newtheoremstyle{wzdefinition}{6pt}{5pt}
  {\normalfont\songti}{0pt}{\normalfont\bfseries\color{structurecolor}}
  {}{.7em}{\thmname{#1}\thmnumber{ #2}\thmnote{\normalfont\ (#3)}}
\newtheoremstyle{wzproblem}{7pt}{5pt}
  {\normalfont\songti}{2em}{\normalfont\bfseries\color{main}}
  {}{.7em}{\thmname{#1}\thmnumber{ #2}}
\newtheoremstyle{wzremark}{4pt}{4pt}
  {\normalfont\songti}{0pt}{\normalfont\bfseries}
  {}{.7em}{\thmname{#1}}
\theoremstyle{wzstatement}
\newtheorem{theorem}{定理}[chapter]
\newtheorem{lemma}[theorem]{引理}
\newtheorem{proposition}[theorem]{命题}
\newtheorem{corollary}[theorem]{推论}
\theoremstyle{wzdefinition}
\newtheorem{definition}[theorem]{定义}
\theoremstyle{wzproblem}
\newtheorem{example}{例题}[chapter]
\newtheorem{exercise}{习题}[chapter]
\theoremstyle{wzremark}
\newtheorem*{remark}{注}
\renewcommand{\proofname}{证明}
% Retain the native amsthm QED stack; only adjust the fixed typography.
\expandafter\patchcmd\csname\string\proof\endcsname{\normalfont \topsep6\p@\@plus6\p@\relax}
  {\normalfont\songti \topsep5\p@\@plus1\p@\relax}
  {}{\ClassError{elegantbook-original-adapter}{Cannot patch proof spacing}{Check the amsthm version.}}
\expandafter\patchcmd\csname\string\proof\endcsname{\itshape}{\normalfont\bfseries}
  {}{\ClassError{elegantbook-original-adapter}{Cannot patch proof heading}{Check the amsthm version.}}
\expandafter\patchcmd\csname\string\proof\endcsname{\ignorespaces}
  {\setlength{\parindent}{2em}\ignorespaces}
  {}{\ClassError{elegantbook-original-adapter}{Cannot patch proof paragraphs}{Check the amsthm version.}}
\AtBeginEnvironment{proof}{\par\Needspace{3\baselineskip}}
\renewcommand{\qedsymbol}{\ensuremath{\square}}
\newenvironment{solution}{\begin{proof}[解答]}{\end{proof}}
\newcommand{\wzcheckchapter}{\ifnum\value{chapter}<1
  \ClassError{elegantbook-original-adapter}{Numbered mathematics before chapter 1}{Move it after the first originalchapter.}\fi}

\newcommand{\wzregisterguard}[1]{\AtBeginEnvironment{#1}{\wzcheckchapter\par\Needspace{3\baselineskip}}}
\forcsvlist{\wzregisterguard}{theorem,lemma,proposition,corollary,definition,example,exercise}
% Front matter and bibliography are not numbered mathematical chapters.
\newcommand{\originalpreface}{\par\addvspace{10pt}\Needspace{5\baselineskip}
  \noindent{\bfseries\fontsize{12}{17}\selectfont 前言}\par\nobreak\vspace{4pt}}
\newcommand{\originalcontents}{\par\addvspace{14pt}
  \begin{center}{\color{structurecolor}\bfseries\fontsize{14.4}{20}\selectfont 目录}\end{center}
  \vspace{-10pt}\@starttoc{toc}}
\renewcommand{\bibname}{参考文献}
\renewenvironment{thebibliography}[1]{
  \par\addvspace{12pt}\Needspace{3\baselineskip}\phantomsection
  \addcontentsline{toc}{chapter}{\bibname}
  \markboth{\bibname}{}
  \noindent{\color{structurecolor}\bfseries\fontsize{14.4}{20}\selectfont\bibname}\par\nobreak\vspace{6pt}
  \list{\@biblabel{\@arabic\c@enumiv}}
    {\settowidth\labelwidth{\@biblabel{#1}}\leftmargin\labelwidth
     \advance\leftmargin\labelsep\usecounter{enumiv}
     \let\p@enumiv\@empty\renewcommand\theenumiv{\@arabic\c@enumiv}
     \setlength{\itemsep}{3pt}\setlength{\parsep}{0pt}}
  \sloppy\clubpenalty4000\widowpenalty4000\sfcode`\.=1000\relax
}{\def\@noitemerr{\@latex@warning{Empty bibliography}}\endlist}

\numberwithin{equation}{chapter}
\renewcommand{\theequation}{\thechapter.\arabic{equation}}
\allowdisplaybreaks[2]
\emergencystretch=2em
\clubpenalty=6000
\widowpenalty=6000
\raggedbottom
\endinput
\end{filecontents*}
% WZ-LOCKED-CLASS-END
\documentclass{elegantbook-original-adapter}
% ===== 元数据内容槽 =====
\newcommand{\BookTitle}{量子力学：波函数、态空间与近似方法}
\newcommand{\BookAuthor}{}
\newcommand{\BookDate}{}
\newcommand{\PDFTitle}{量子力学中文学习册}
\newcommand{\PDFAuthor}{}
\hypersetup{pdftitle={\PDFTitle},pdfauthor={\PDFAuthor}}
% 数学记号可以增加；禁止重定义版式、环境或证毕符号.
\newcommand{\dd}{\,\mathrm d}
\newcommand{\R}{\mathbb R}
\newcommand{\norm}[1]{\lVert #1\rVert}
\newcommand{\inner}[2]{\langle #1,#2\rangle}
\usepackage{tikz,pgfplots,longtable,booktabs}
\usetikzlibrary{arrows.meta,calc,positioning}
\pgfplotsset{compat=1.18}
\newcommand{\ket}[1]{\lvert #1\rangle}
\newcommand{\bra}[1]{\langle #1\rvert}
\newcommand{\braket}[2]{\langle #1\vert #2\rangle}
\newcommand{\ii}{\mathrm i}
\newcommand{\e}{\mathrm e}
\newcommand{\Tr}{\operatorname{Tr}}
\begin{document}
\frontmatter
\thispagestyle{cover}
\begin{center}{\color{structurecolor}\bfseries\fontsize{17}{23}\selectfont 量子力学}\par
\vspace{5pt}{\bfseries 波函数、态空间与近似方法}\par
\vspace{5pt}根据 MIT OpenCourseWare 公开课程中文重编\par
2026年10月11日\end{center}
\originalpreface
本册以会初等微积分、尚未系统学习量子力学的自学者为读者. 复数振幅、概率密度、波函数和算符在首次使用时解释. 单变量微积分和常微分方程的基本运算作为先修; 希尔伯特空间的严格谱理论没有假定已经掌握. 无界算符的论证在文中说明所用边界和共同定义域, 不把形式计算当作完整的泛函分析证明.

材料主轴为 Barton Zwiebach 的 8.04 Quantum Physics I (Spring 2016)、8.05 Quantum Physics II (Fall 2013) 和 8.06 Quantum Physics III (Spring 2018). 本册是中文重新组织和补充推导的非官方改编, 没有 MIT 官方中文或认可的含义. 原始文件保留在仓库量子力学来源目录. 原件和中文改编遵循 \href{https://creativecommons.org/licenses/by-nc-sa/4.0/}{CC BY-NC-SA 4.0}, 署名 MIT OpenCourseWare、Barton Zwiebach 及原件另署的贡献者. 本册新增解释、例题和答案标为独立整理. 不使用 MIT 校徽或宣传性标志.

本册现有19章理论, 39份作业、测试与期末文档中227道可用大题的全部小问. 原17理论章与19考核篇保留, 新增谱与变分、含时辐射两章, 第8章接入线代补充, 20份新增作业共115题. 两张公式表另归档. 题解均为\textbf{AI独立解答}, 所选OCW页面没有取得官方答案. 23份先前被整份暂缓的原件现全部恢复; 本次按完整MIT任务逐题补入, 商业教材只有题号而没有题面的10道不编造. 改编信用保留, 表格与数值图自主推导和绘制. 独立模型复核不称为教师答案, 本册也不是三课每页原件的全文翻译.

\noindent\textbf{8.04 课程安排.}
每周两次90分钟讲授和两次60分钟习题课; 官方先修是8.03振动与波, 成绩至少C. 三部分为量子基本概念、一维势中的运动、一维散射与中心势. 两次课内考试和三小时期末. 总评为考试平均75\%、作业25\%; 考试平均中两次测试各25\%、期末50\%, 所以折合总评分别18.75\%、18.75\%、37.5\%. 作业最低或缺交的一套不计入平均. 2016页面实际提供2015秋测试, 年份不能合并.

\noindent\textbf{8.05 课程安排.}
每周两次90分钟讲授和两次60分钟习题课; 先修8.04, 至少C. 大纲包括量子力学一般结构、动力学、二态系统、角动量与自旋、径向方程、角动量合成、全同粒子. 一次90分钟闭卷期中和三小时期末. 总评: 期中25\%、期末45\%、作业30\%. 官网发布11套作业, 以及2011--2013年的期中与期末试卷和两张公式表.

\noindent\textbf{8.06 课程安排.}
每周两次90分钟讲授和两次60分钟习题课; 先修8.05, 至少C. 大纲以定态微扰、氢精细结构、WKB、含时微扰、带电粒子、绝热近似、散射和全同粒子展开. 总评: 作业25\%、期中15\%、学期论文25\%、期末35\%. 论文成绩须至少C才能通过, 另有提案、初稿与同伴编辑环节. 课纲述九套计分作业和末周一套不计分作业, OCW作业页只有九套. 90分钟期中和三小时期末\textbf{明确不向OCW公开}; 第5章带电粒子讲义也未提供. 本册没有伪造这些题面或章节原件.

下列课历保留官方讲次与主题顺序; 主题用中文表述, 完整官网链接见卷末来源对应. 作业序号和考试位置保留; 官网未给具体日历日期, 不另造日期.

\noindent\textbf{8.04 课历.}
\begin{longtable}{p{.10\textwidth}p{.67\textwidth}p{.14\textwidth}}
\toprule 讲次 & 主题 & 考核节点\\\midrule\endhead
1--2 & 线性、复数、随机性、叠加与纠缠; 双缝、干涉仪、炸弹探测 & \\
3--5 & 光电效应、康普顿散射、物质波; 伽利略变换、群速度; 动量与位置算符、薛定谔方程 & 作业1、2\\
6--9 & 归一化、概率流、厄米算符; 波包与傅里叶变换; 期望值、Ehrenfest方程; 测量和不确定性 & 作业3、4\\
10--12 & 定态、圆周粒子、无限和有限势阱; 转折点、半经典近似 & 测试1; 作业5\\
13--15 & 打靶法; $\delta$ 势、节点; 谐振子微分方程和升降算符 & 作业6\\
16--18 & 散射、反射和透射; 共振透射; 相移、Levinson定理 & 作业7; 测试2\\
19--21 & 共振、复波数; 角动量代数、勒让德函数; 径向方程与氢原子 & 作业8、9\\
22--24 & 氢原子尺度、级数与量子化; 维里关系、轨道与转折点; 自旋 & 作业10; 期末\\
\bottomrule
\end{longtable}
\noindent\textbf{8.05 课历.}
\begin{longtable}{p{.12\textwidth}p{.65\textwidth}p{.14\textwidth}}
\toprule 讲次 & 主题 & 作业\\\midrule\endhead
1--3 & 波动力学 & 1\\
3--4 & 自旋$1/2$、bra、ket与算符 & 2\\
5--8 & 向量空间与算符 & 3、4\\
9 & Dirac记号 & 5\\
10--11 & 不确定性与相容观测量 & 6\\
12--16 & 量子动力学 & 7\\
16--18 & 二态系统 & 8\\
18--20 & 多粒子态与张量积 & 9\\
20--23 & 角动量 & 10\\
24--26 & 角动量合成 & 11\\
\bottomrule
\end{longtable}
\noindent\textbf{8.06 课历.}
\begin{longtable}{p{.12\textwidth}p{.65\textwidth}p{.14\textwidth}}
\toprule 讲次 & 主题 & 节点\\\midrule\endhead
L1--L3 & 定态与简并微扰 & 作业1\\
L4--L6 & 氢精细结构、Zeeman效应、半经典近似引入 & 作业2\\
L7--L9 & WKB、Airy函数和连接式; 含时微扰 & 作业3、4\\
L10--L11 & Fermi黄金规则和谐波跃迁 & 作业5; E1\\
L12--L14 & 氢电离、光与原子、带电粒子与电磁场 & 论文提案\\
L15--L18 & 绝热近似、Berry相位、分子 & 作业6、7; 论文初稿\\
L19--L22 & 散射、全同粒子引入 & 作业8\\
L23--L24 & 全同粒子、论文返还与期末复习 & E2\\
\bottomrule
\end{longtable}

\clearpage
\originalcontents
\clearpage
\mainmatter
\originalchapter{复数振幅与物质波}
\label{ch:waves}
\originalsection{干涉现象与复数振幅}
先看一个物理问题. 在Mach--Zehnder干涉仪中, 入射光被分成两条路径, 在第二个分束器重新汇合. 改变一条路径的光程, 两个出射口的亮暗会改变. 把入射减弱到单光子情形, 一次只留下一个探测事件, 重复实验的计数仍随两路的相对相位改变. 因此要解释的首先是探测统计, 不是给光子预先指定一条未被测量的轨迹.

如果只给两条路径各分配一个非负概率, 再把概率相加, 就无法描述一处增强、另一处抵消的变化. 量子描述先给各路径一个可相加的\textbf{概率振幅}, 汇合时相加, 最后取模方得到概率. 振幅既有大小又有相位; 相位记录类似波峰与波谷之间的相对位置. 复数是同时记录这两个量的方便工具. 下面介绍所需运算, 而不是先假定已经学过抽象态空间.

复数 $z=a+\ii b$ 中 $\ii^2=-1$, 由两个实数指定, 共轭为 $z^*=a-\ii b$, 模为 $|z|=\sqrt{a^2+b^2}$. 当 $z\ne0$ 时可写成 $z=r\e^{\ii\theta}$, 其中 $r=|z|>0$, $\theta$ 只确定到 $2\pi$ 的整数倍. Euler公式 $\e^{\ii\theta}=\cos\theta+\ii\sin\theta$ 把旋转写成乘法. 它由指数、正弦和余弦的绝对收敛幂级数, 将偶次和奇次项分开得到. 绝对收敛保证这次分组合法.

振幅相加与概率相加有不同结果. 两条不可区分的传播路径若有振幅 $z_1,z_2$, 合振幅为 $z_1+z_2$, 对应概率与 $|z_1+z_2|^2$ 成比例. 展开为 $|z_1|^2+|z_2|^2+2\operatorname{Re}(z_1^*z_2)$. 最后一项取决于相对相位, 是干涉项. 如果实验记录区分了两条路径, 不能直接用两个纯振幅相加来描述未读取记录后的统计, 必须把记录系统纳入, 再对它作概率求和.

\begin{example}
两个等大振幅分别为 $A$ 和 $A\e^{\ii\delta}$. 求总强度与 $\delta$ 的关系.
\end{example}
\begin{solution}
利用共轭乘积, $|A+A\e^{\ii\delta}|^2=|A|^2(2+2\cos\delta)=4|A|^2\cos^2(\delta/2)$. $\delta=0$ 时增强, $\delta=\pi$ 时抵消. 强度改变不意味着单条路径的振幅模改变; 改变的是交叉项.
\end{solution}

整体相位与相对相位要分开. 把整个态乘以同一个 $\e^{\ii\chi}$, 所有概率不变. 把叠加中的一个分量单独乘以相位, 一般会改变后续测量. 全局相位不表示新的物理态, 但局部相位差可以被干涉实验读取.

\originalsection{光子和物质波的尺度}
光电效应表明, 单个光子的能量为 $E=h\nu=\hbar\omega$, 其中 $\omega=2\pi\nu$, $\hbar=h/(2\pi)$. 真空光子满足 $E=pc$, 因而 $p=\hbar k$, $k=2\pi/\lambda$. 电子等有质量粒子也满足 de Broglie关系 $p=h/\lambda$; 在非相对论自由运动中, $E=p^2/(2m)$, 所以 $\omega(k)=\hbar k^2/(2m)$.

这些关系并不说明电子是经典介质的微小振动. 波函数是一种给出实验概率的复振幅. 经典波的场强、光子的电磁场与单个电子的概率振幅属于不同对象, 不能因为都含正弦就认作同一个物理量.

量子波长是否明显, 取决于与装置尺度的比较. 质量 $m=10^{-2}\,\mathrm{kg}$、速度 $v=10^{-1}\,\mathrm{m/s}$ 的小球, 波长约为 $6.6\times10^{-31}\,\mathrm m$. 这种尺度远小于实际对位置的控制精度. 冷原子速度很低, 波长能达到装置可分辨的尺度, 干涉便能够出现. 波长小不是量子理论失效, 而是某些实验量可用经典近似描述.

\originalsection{相速度与群速度}
单一平面波 $\e^{\ii(kx-\omega t)}$ 的等相位点满足 $kx-\omega t=\text{常数}$, 因而相速度为 $v_{\rm ph}=\omega/k$. 平面波在整条实线上都存在, 没有局域中心. 要描述局域粒子, 必须叠加不同波数; 群速度近似先考虑频谱集中在 $k_0$ 附近的情形.

令两个分量的波数为 $k_0\pm\Delta k/2$, 频率为 $\omega_0\pm\Delta\omega/2$. 相加得到
\[
2\e^{\ii(k_0x-\omega_0t)}\cos\!\left(\frac{\Delta k\,x-\Delta\omega\,t}{2}\right).
\]
快速振荡由前因子控制, 包络由余弦控制. 包络速度为 $\Delta\omega/\Delta k$; 窄频谱极限中成为 $v_g=\omega'(k_0)$. 对非相对论粒子, $v_g=\hbar k_0/m=p_0/m$, 而 $v_{\rm ph}=p_0/(2m)$. 两波拍频只是包络示意, 还不是可归一化的局域态. 窄谱波包的主包络近似以群速度传播; 二阶色散会使包络展宽. 自由非相对论粒子的平均位置速度严格为 $\langle p\rangle/m$, 一般多峰波包却未必有单一的峰值速度. 对相对论自由粒子, $\mathrm dE/\mathrm dp=pc^2/E=v$, 因而群速度仍是粒子速度.

\begin{remark}
由自由粒子色散关系和叠加原理构造运动方程, 是建立薛定谔方程的动机. 它不能仅凭经典力学推出全部量子测量规则; Born概率解释和量子态结构仍是物理理论的输入.
\end{remark}

\originalchapter{波函数与概率}
\originalsection{概率密度和归一化}
对一维无自旋粒子, 先研究能由单一复振幅描述的制备, 称为\textbf{纯态}; 它与对多个制备作经典随机混合的情形不同, 后者见第12章. \textbf{位置表象}是把每个位置$x$作为振幅的坐标, 用复函数$\psi(x,t)$表示态. 在时刻$t$做位置测量, 落入区间$[a,b]$的概率为$\int_a^b|\psi(x,t)|^2\dd x$. 这里$\rho=|\psi|^2$是\textbf{概率密度}, 不是概率本身. 在一维中$\rho$的量纲是长度的倒数, $\psi$的量纲是长度的负二分之一次方. 单个点的概率通常是零, 这不妨碍附近小区间有非零概率.

\begin{definition}
函数 $\psi$ 若满足 $\int_{\R}|\psi|^2\dd x<\infty$, 称为平方可积. 归一化态满足 $\int_{\R}|\psi|^2\dd x=1$. 若 $0<\|\psi\|<\infty$ 而范数不为1, 可除以其范数归一化; 零函数平方可积, 但不能表示归一化态.
\end{definition}
不能对任意函数进行这一步. 平面波的 $|\psi|$ 是常数, 在实线上的积分发散; 它是连续谱的广义本征态, 不是可归一化单粒子态. 后文散射中的单色波也有同样限制. 可以把它放在周期盒子中归一化, 或用平面波叠加成波包.

\begin{example}
设 $a>0$, $\psi(x)=N\exp[-x^2/(4a^2)]$. 求归一化系数、$\langle x\rangle$ 和 $\Delta x$.
\end{example}
\begin{solution}
由 Gaussian积分 $\int_\R\e^{-x^2/(2a^2)}\dd x=\sqrt{2\pi}\,a$, 可取 $N=(2\pi)^{-1/4}a^{-1/2}$. 概率密度关于零对称, 所以 $\langle x\rangle=0$. 对Gaussian积分的参数求导, 或分部积分, 得 $\int_\R x^2\e^{-x^2/(2a^2)}\dd x=\sqrt{2\pi}\,a^3$. 除以归一化积分得 $\langle x^2\rangle=a^2$, 所以标准差 $\Delta x=a$.
\end{solution}

\originalsection{内积与叠加}
函数也可以像向量一样比较. 定义 $\langle\phi,\psi\rangle=\int\phi^*(x)\psi(x)\dd x$. 它对第二项线性, 对第一项共轭线性, 满足 $\langle\phi,\psi\rangle=\langle\psi,\phi\rangle^*$. 范数为 $\|\psi\|=\sqrt{\langle\psi,\psi\rangle}$. 内积为零的态叫正交态.

若 $u_1,u_2$ 是正交归一态, $\psi=c_1u_1+c_2u_2$ 的范数平方是 $|c_1|^2+|c_2|^2$, 因而归一化要求这个和为1. 若两个基函数不正交, 范数还包含 $c_1^*c_2\langle u_1,u_2\rangle$ 及其共轭. 不能在一般叠加中不检查正交性就把系数模方认作概率.

对于离散非简并自伴观测量的正交本征基, 系数 $c_n=\langle u_n,\psi\rangle$, 测得第$n$个本征值的概率为 $|c_n|^2$. 若本征值简并, 概率须对整个本征子空间的系数模方求和; 连续谱的概率密度在第4章讨论. 这条测量规则不是说粒子在测量前已经随机选定某个基态. 同一个态在不同测量基中有不同展开, 干涉可以检验其相干性.

\originalsection{概率、振幅和期望值}
当对应一阶、二阶矩存在时, 位置的期望值是 $\langle x\rangle=\int x|\psi(x)|^2\dd x$. 它是相同制备下大量独立测量的平均, 一次测量无需等于平均. 方差 $\langle(x-\langle x\rangle)^2\rangle$ 非负, 标准差用来量化分布的宽度.

\begin{center}
\begin{tikzpicture}
\begin{axis}[width=10cm,height=4.4cm,axis lines=middle,xlabel={$x/a$},ylabel={$a\rho$},xlabel style={at={(axis description cs:1,-.08)},anchor=north},ylabel style={at={(axis description cs:-.06,1)},anchor=south},xmin=-3.5,xmax=3.5,ymin=0,ymax=.48,xtick={-2,0,2},ytick={.2,.4},samples=120,domain=-3.4:3.4]
\addplot[thick,structurecolor] {exp(-x*x/2)/sqrt(2*pi)};
\draw[dashed] (axis cs:-1,0)--(axis cs:-1,.242);
\draw[dashed] (axis cs:1,0)--(axis cs:1,.242);
\end{axis}
\end{tikzpicture}
\par Gaussian概率密度; 虚线对应距离均值一个标准差, 不表示概率的硬边界.
\end{center}

$\psi$ 的正负和复相位不会直接显示在位置概率密度中. 但后续动力学或干涉测量可以显示相位差. 因而只给 $|\psi|^2$ 一般不能指定完整纯态. 例如同模的 $\psi$ 和 $\e^{\ii qx/\hbar}\psi$ 有相同位置分布, 若动量均值存在, 后者的动量均值比前者增加 $q$.

\originalchapter{薛定谔方程与概率流}
\originalsection{从平面波到线性方程}
已知初始位置分布和相位后, 如何预测稍后的探测统计? 我们需要振幅的演化规律. 记$\partial_x$为固定时间对$x$求导, $\partial_t$为固定位置对$t$求导, 帽号标记作用于函数的算符. $\hat p^2\psi$指连续作用两次$\hat p$, 而非一个测量数的任意平方.

对平面波 $\psi=\e^{\ii(kx-\omega t)}$, 求导会取出其波数或频率: $-\ii\hbar\partial_x\psi=\hbar k\psi$, $\ii\hbar\partial_t\psi=\hbar\omega\psi$. 因此动量算符取为 $\hat p=-\ii\hbar\partial_x$, 自由粒子的能量算符取为 $\hat p^2/(2m)=-\hbar^2\partial_x^2/(2m)$. 位置算符 $\hat x$ 的作用是乘以$x$: $(\hat x\psi)(x)=x\psi(x)$.

算符的输入和输出都是函数. 例如 $\hat p\psi$ 不是测量得到的一个数, 而是新的函数. 只有当 $\hat p\psi=p\psi$ 时, $\psi$ 才是动量本征态, $p$ 才是它确定的动量本征值. 类比矩阵作用于向量, 本征向量方向不变, 只乘上本征值.

每个自由平面波都满足$\ii\hbar\partial_t\psi_j=-\hbar^2\partial_x^2\psi_j/(2m)$; 因求导线性, 常系数$c_1\psi_1+c_2\psi_2$也满足同一方程. 连续Fourier叠加在允许微分进积分的条件下仍满足. 这一步从确定波数推广到局域自由波包, 加外势是进一步的物理输入.

把外部实势能 $V(x,t)$ 作为乘法算符, 得到
\begin{equation}\label{eq:sch}
\ii\hbar\frac{\partial\psi}{\partial t}
=\hat H\psi,\qquad
\hat H=-\frac{\hbar^2}{2m}\frac{\partial^2}{\partial x^2}+V(x,t).
\end{equation}
$\hat H$ 是 Hamilton算符. 这是时间一阶、空间二阶的偏微分方程. 给初始波函数和适当空间边界, 它决定之后的波函数. 时间一阶意味着不再独立指定 $\partial_t\psi(x,0)$. 它决定的是概率振幅的确定演化, 单次测量仍是随机的.

动量关系直接保证自由平面波满足此方程; 加势能是与实验相符的量子动力学规则, 不是把经典轨道强行改写为正弦. 对带磁场的粒子, Hamilton算符还需用机械动量与矢势, 不能无条件继续用仅含标量势的式\eqref{eq:sch}.

\originalsection{连续性方程}
小区间概率的增加应该等于左边流入减右边流出. 一维$J$的量纲为每时间的概率, 正号指向正$x$. 以下用撇号表示对$x$求导, 并从演化方程找出这种流.

假定 $V$ 实值且 $\psi$ 足够光滑. 先用乘积求导$\partial_t(\psi^*\psi)=\psi^*\partial_t\psi+\psi\partial_t\psi^*$, 再用\eqref{eq:sch}和其共轭求 $\rho=\psi^*\psi$ 的时间导数, 势能两项恰好抵消:
\begin{align*}
\partial_t\rho
&=\frac{\ii\hbar}{2m}(\psi^*\psi''-\psi\psi^{*\prime\prime})\\
&=-\partial_x\left[\frac{\hbar}{2m\ii}(\psi^*\psi'-\psi\psi^{*\prime})\right].
\end{align*}
所以定义概率流 $J=\hbar\operatorname{Im}(\psi^*\psi')/m$, 有 $\partial_t\rho+\partial_xJ=0$. 对区间$[a,b]$积分, $\frac{\mathrm d}{\mathrm dt}\int_a^b\rho\dd x=J(a,t)-J(b,t)$. 右边是流入减流出, 并不是概率被势能产生或消灭.

实线中若两端流为零, 总概率不变. 有限区间可以施加硬墙、周期或其他自伴边界条件保证适当的总流抵消. 若人为使用复吸收势, 连续性方程会多出源项; 它是开放系统的有效描述, 不能拿来证明上述闭合系统归一化守恒.

\begin{example}
求平面波 $A\e^{\ii kx}$ 的概率流. 再求实函数 $\sin kx$ 的概率流.
\end{example}
\begin{solution}
第一种情形 $\psi^*\psi'=\ii k|A|^2$, 所以 $J=\hbar k|A|^2/m$. 流向随$k$正负变化. 第二种情形 $\psi^*\psi'$ 实值, 因而$J=0$. 正弦是等权反向平面波叠加, 两股流互相抵消. 两者在无限实线上都不能按普通积分归一化, 这些流是单色波的形式结果.
\end{solution}

推广到三维时,$\nabla=(\partial_x,\partial_y,\partial_z)$表示三个空间导数, $\nabla\cdot\boldsymbol J=\partial_xJ_x+\partial_yJ_y+\partial_zJ_z$称散度. 三维中 $\rho=|\psi|^2$, $\boldsymbol J=(\hbar/m)\operatorname{Im}(\psi^*\nabla\psi)$, 满足 $\partial_t\rho+\nabla\cdot\boldsymbol J=0$. 这里密度量纲变为体积的倒数; 一个空间区域中的概率改变由其表面净流给出.

\originalchapter{动量表象与不确定性}
\originalsection{傅里叶展开}
相同的位置模方可配不同相位, 从而有不同的动量统计. 为找每个动量读数的振幅, 把局域波包分解成确定波数的平面波; 连续积分是离散展开求和的推广, $\phi(k)$记录各波数分量. 采用一致的变换约定,
\[
\psi(x)=\frac1{\sqrt{2\pi}}\int_\R\phi(k)\e^{\ii kx}\dd k,
\qquad
\phi(k)=\frac1{\sqrt{2\pi}}\int_\R\psi(x)\e^{-\ii kx}\dd x.
\]
对平方可积函数, Plancherel定理保证该变换在$L^2$（平方可积函数集合）意义下定义且保持范数: $\int|\psi|^2\dd x=\int|\phi|^2\dd k$. 这里调用标准傅里叶理论的$L^2$版本, 不把分布积分的形式运算当作它的证明. 本册Gaussian和光滑快速下降波包可直接核对所有积分.

$|\phi(k)|^2\dd k$ 是波数在小区间中的概率. 若改用动量$p=\hbar k$, 概率幅必须是 $\widetilde\psi(p)=\phi(p/\hbar)/\sqrt\hbar$, 才有 $|\widetilde\psi(p)|^2\dd p=|\phi(k)|^2\dd k$. 不可以只把$k$写成$p$却忽略变量代换因子$\dd p=\hbar\dd k$.

对光滑快速下降函数分部积分, 动量算符在$k$表象中变为乘以$\hbar k$, 位置算符变为$\ii\partial_k$. 从而 $\langle p\rangle=\int\hbar k|\phi(k)|^2\dd k$. 这解释了空间中快速的相位振荡为什么对应较大的动量.

两条表象规则可直接复算. 记上述变换为 $\mathcal F$, 对边界项消失的光滑快速下降函数, 分部积分给
$\mathcal F(\psi')=\ii k\phi$, 所以 $\mathcal F(-\ii\hbar\psi')=\hbar k\phi$. 对 $k$ 求导则有 $\phi'(k)=\mathcal F(-\ii x\psi)$, 因而 $\mathcal F(x\psi)=\ii\phi'$. 快速下降在这里表示函数及所需导数比任意倒数幂衰减得快; Gaussian满足这些条件. $L^2$收敛则指函数差的模方积分趋零, 未要求每一个点都收敛.

对第2章Gaussian, 不需要复变量积分也可完成变换. 令 $I(k)=\int\e^{-x^2/(4a^2)}\e^{-\ii kx}\dd x$. Gaussian衰减允许微分进积分. 由 $x\e^{-x^2/(4a^2)}=-2a^2\partial_x\e^{-x^2/(4a^2)}$, 分部积分得 $I'(k)=-2a^2kI(k)$. 又 $I(0)=2a\sqrt\pi$, 故 $I(k)=2a\sqrt\pi\e^{-a^2k^2}$. 乘上变换归一系数得
$\phi(k)=(2a^2/\pi)^{1/4}\e^{-a^2k^2}$. 它的密度仍是Gaussian, $\Delta k=1/(2a)$; 空间越窄, 波数分布越宽.


\originalsection{厄米性与边界}
测量均值应为实数, 这引出伴随运算. 伴随$A^\dagger$由$\inner\phi{A\psi}=\inner{A^\dagger\phi}\psi$定义; 有限维正交归一基中$(A^\dagger)_{ij}=(A_{ji})^*$, 即交换行列再取共轭. 定义域是允许输入使作用和内积有意义的函数集合. 观测量的期望值写为 $\langle A\rangle=\langle\psi,A\psi\rangle$. 在有限维, $A^\dagger=A$ 称为厄米, 等价于矩阵共轭转置不变. 对微分算符, 积分区间和定义域是算符的一部分. 例如
\[
\langle\phi,\hat p\psi\rangle-\langle\hat p\phi,\psi\rangle
=-\ii\hbar[\phi^*(x)\psi(x)]_{a}^{b}.
\]
只有边界项为零, 才在所选定义域上有对称性. 自伴还要求伴随算符的定义域相同, 它比这个分部积分等式更强. 后文有限维矩阵没有此区别; 实线上的标准动量实现, 定义域可表述为函数与其弱导数都平方可积; 弱导数表示用对光滑试探函数的分部积分等式来定义导数. 本册计算先在光滑快速下降的子类上进行. 在有限区间, $\psi(b)=\e^{\ii\theta}\psi(a)$的一族边界给一阶动量的自伴实现. 硬墙的两端Dirichlet（函数值固定为零）边界给二阶Hamilton算符的自伴实现, 却不能自动给同域的一阶动量自伴性. 每次形式对易计算限制在所需乘积都有意义的共同定义域上.

自伴本征态 $Au=a u$ 的本征值实, 因为 $\langle u,Au\rangle=\langle Au,u\rangle$ 给出 $a\|u\|^2=a^*\|u\|^2$. 不同本征值的本征态正交: 若$Au=au$, $Av=bv$, 则 $(b-a)\langle u,v\rangle=0$. 两个结论不能反过来独立证明一个未指定边界的微分式自伴.

\originalsection{不确定性关系的推导}
Gaussian已给出位置窄则动量宽的具体例子. 现在问任意态能否使这两个宽度同时任意小. 去掉均值后的算符作用$(A-\langle A\rangle)\psi$记录偏离确定平均的部分, 其范数就是标准差; 比较两种这样的变化能把两种宽度联系起来.

归一态的方差定义为 $(\Delta A)^2=\|(A-\langle A\rangle)\psi\|^2$. 因为是范数平方, 它非负, 并且仅在$\psi$为$A$的本征态时为零. 记 $[A,B]=AB-BA$, 算符乘积按右边先作用. 假定 $A,B$自伴, 当前态在所需定义域中且方差有限. 令 $a=(A-\langle A\rangle)\psi$, $b=(B-\langle B\rangle)\psi$. Cauchy--Schwarz不等式可由范数非负得到: 当$a\ne0$时, 取 $c=\langle a,b\rangle/\|a\|^2$, 展开$\norm{b-ca}^2=\norm b^2-c\inner ba-c^*\inner ab+|c|^2\norm a^2$, 代入该$c$后得 $|\langle a,b\rangle|^2\le\|a\|^2\|b\|^2$; $a=0$时结论直接成立. 所以
\[
(\Delta A)^2(\Delta B)^2\ge|\langle a,b\rangle|^2
\ge|\operatorname{Im}\langle a,b\rangle|^2.
\]
在$AB\psi$和$BA\psi$都有意义且伴随交换合法时,
$\langle a,b\rangle-\langle b,a\rangle=\langle\psi,(AB-BA)\psi\rangle$, 也就是 $2\ii\operatorname{Im}\langle a,b\rangle=\langle[A,B]\rangle$. 因而
\begin{equation}\label{eq:uncertainty}
\Delta A\,\Delta B\ge\frac12|\langle[A,B]\rangle|.
\end{equation}
对快速下降波包, $xp\psi-p(x\psi)=-\ii\hbar x\psi'+\ii\hbar(\psi+x\psi')=\ii\hbar\psi$, 得 $\Delta x\,\Delta p\ge\hbar/2$. 它描述同一个制备态的两种测量分布宽度, 不是一次测量仪器误差的简单相乘.

对第2章的实Gaussian, $p\psi=\ii\hbar x\psi/(2a^2)$, 所以$\langle p\rangle=0$, $\Delta p=\hbar/(2a)$. 因而它饱和此不等式. 乘以$\e^{\ii p_0x/\hbar}$会把动量均值平移到$p_0$, 但两种标准差不变. 波包可以有任意大的平均动量同时保持同样的不确定性.

\originalchapter{定态与一维势阱}
\originalsection{定态分离与叠加}
当$V$不含时间, 可寻找分离变量解$\psi(x,t)=u(x)T(t)$. 代入薛定谔方程, 除以$uT$后, 时间变量和空间变量相等, 因而都等于常数$E$. 于是 $\hat Hu=Eu$, $T(t)=\e^{-\ii Et/\hbar}$. 前者是\emph{定态薛定谔方程}, 后者规定这个本征态的时间相位.

定态的密度不随时间改变, 因为$|T|=1$. 多个能量不同的定态叠加一般不是定态. 若 $\psi=\sum_n c_nu_n\e^{-\ii E_nt/\hbar}$, 密度中的交叉项含 $\e^{\ii(E_m-E_n)t/\hbar}$. 动力学的可观测变化正来自相对相位变化. 只求$Hu=Eu$而不展开初始条件, 还没有解决任意初态的时间演化.

\originalsection{无限深势阱}
考虑$0<x<L$中$V=0$, 区间外有无限高墙. 硬墙边界取 $u(0)=u(L)=0$. 在区间内 $u''+k^2u=0$, $k^2=2mE/\hbar^2$. Hamilton算符非负, 因为分部积分给出 $\langle u,Hu\rangle=(\hbar^2/2m)\int|u'|^2\dd x$. $E=0$的线性解在两端为零后只能恒为零, 因此非零本征态有$E>0$.

通解$u=A\sin kx+B\cos kx$. 左边界给$B=0$, 右边界要求$\sin kL=0$, 所以$k=n\pi/L$, $n=1,2,\ldots$. 归一化后,
\[
u_n(x)=\sqrt{\frac2L}\sin\frac{n\pi x}{L},
\qquad E_n=\frac{n^2\pi^2\hbar^2}{2mL^2}.
\]
$n=0$不能作为态: 对应函数是零函数, 无法归一化. 最低能量大于零不是额外的边界约定; 它由局域化和微分方程共同决定. 无限墙处不要求$u'$和外侧零导数连续, 因为这里势能奇异; 若两者都强行设零, 解会消失.

\begin{example}
势阱中初态为$(u_1+u_2)/\sqrt2$. 求其能量测量概率及位置密度.
\end{example}
\begin{solution}
两个能量分别为$E_1,E_2$, 概率各$1/2$, $\langle H\rangle=(E_1+E_2)/2$. 随时间演化后
\[
|\psi(x,t)|^2=\frac{u_1(x)^2+u_2(x)^2}{2}
+u_1(x)u_2(x)\cos\frac{(E_2-E_1)t}{\hbar}.
\]
能量分布不变, 位置分布却振荡; 这两个结论对应不同观测量, 并不矛盾. 正交性使交叉项在整个区间积分为零, 所以总概率仍为1.
\end{solution}

\originalsection{有限深势阱}
取$V=-V_0$在$|x|<a$内, 外部$V=0$, $V_0>0$. 束缚态有$-V_0<E<0$. $E\ge0$时外区振荡、常数或线性解都无法在两端平方可积. 又 $H+V_0$的期望值是非负导数项加非负外区项, $E=-V_0$若要取等只能两项全零, 使解恒为零; $E<-V_0$更不可能. 阱内定义 $k=\sqrt{2m(E+V_0)}/\hbar$, 阱外定义 $\kappa=\sqrt{-2mE}/\hbar$. 外区要取衰减指数, 排除发散指数. 势能有限跳跃没有$\delta$项. 跨边界积分给 $u'(a+\epsilon)-u'(a-\epsilon)=(2m/\hbar^2)\int_{a-\epsilon}^{a+\epsilon}(V-E)u\dd x\to0$, 所以$u'$连续. 若$u$有跳跃,$u''$会出现$\delta'$分布而右端没有对应项, 因而$u$也须连续. 这些条件共同决定能级, 不能只在阱内要求正弦.

反射与对称势的Hamilton算符对易, 所以偶、奇子空间各自保持不变, 可分别寻找本征态. 偶态内部取$A\cos kx$, 右外侧取$B\e^{-\kappa(x-a)}$. 在$x=a$处连续给$B=A\cos ka$, 导数连续给$-\kappa B=-Ak\sin ka$, 所以 $k\tan ka=\kappa$. 奇态内部取$A\sin kx$, 右侧$B\e^{-\kappa(x-a)}$, 左侧须带对应负号. 右边界给 $B=A\sin ka$, $-\kappa B=Ak\cos ka$, 因而$-k\cot ka=\kappa$. 另有$k^2+\kappa^2=2mV_0/\hbar^2$. 这两类超越方程可以图解或数值求根.

令$z=ka$, $z_0=a\sqrt{2mV_0}/\hbar$. 偶态方程是$z\tan z=\sqrt{z_0^2-z^2}$. 在$0<z<\min(z_0,\pi/2)$中, 左边从零递增, 右边递减; 右端或趋零, 或左边趋无穷, 故存在交点. 任意非零深度的一维对称方阱至少有一个束缚态. 三维球阱有不同阈值, 不能把这个结论不加条件搬过去.

外部密度并非零, 而是指数衰减. 经典禁止区表示$E<V$时经典动能无法为正, 不表示量子位置概率必须为零. 束缚态指数尾部没有持续向外的概率流, 所以粒子并不因尾部非零而不断逃走.

\originalchapter{反射、透射与隧穿}
\originalsection{势阶与流量系数}
取$V=0$于$x<0$, $V=V_0>0$于$x>0$. 先考虑$E>V_0$和从左入射. 左右波数为$k=\sqrt{2mE}/\hbar$, $q=\sqrt{2m(E-V_0)}/\hbar$. 设
\[
u(x)=\begin{cases}
\e^{\ii kx}+r\e^{-\ii kx},&x<0,\\
t\e^{\ii qx},&x>0.
\end{cases}
\]
$1+r=t$及$k(1-r)=qt$来自两条边界连续条件. 解得$r=(k-q)/(k+q)$, $t=2k/(k+q)$. 左右速度不同, 因而透射概率必须取流量比:
\[
R=|r|^2,\qquad T=\frac qk|t|^2,
\qquad R+T=1.
\]
入射流为$\hbar k/m$, 反射流为$-\hbar k|r|^2/m$, 透射流为$\hbar q|t|^2/m$. 把$T$直接写成$|t|^2$会违反流守恒. 单色$R,T$首先是流量比; 对归一化、窄能谱入射波包, 远离势后两部分充分分离, 总反射和透射概率近似为$R(E_0),T(E_0)$. 一般波包要按入射能量概率密度加权积分. 即使经典粒子有足够能量越过势阶, 量子波也可以部分反射.

若$0<E<V_0$, 右侧只能取$t\e^{-\kappa x}$, $\kappa=\sqrt{2m(V_0-E)}/\hbar$. 该项流为零, 所以$T=0$, $R=1$, 但右侧密度非零. 单个半无限势阶没有远处透射流; 它与有限宽势垒的隧穿不同.

\originalsection{矩形势垒的匹配}
取$V=V_0$于$0<x<b$, 其他处为零, $0<E<V_0$. 中间波函数为$A\e^{\kappa x}+B\e^{-\kappa x}$, 两边分别为入射加反射和透射. 从右边界向左传播比较便于计算. 在$x=b$, $u(b)=t\e^{\ii kb}$, $u'(b)=\ii kt\e^{\ii kb}$. 先定义 $\cosh s=(\e^s+\e^{-s})/2$, $\sinh s=(\e^s-\e^{-s})/2$. 方程$u''=\kappa^2u$在$b$处给定函数与导数后的解是
$u(x)=u(b)\cosh[\kappa(x-b)]+u'(b)\sinh[\kappa(x-b)]/\kappa$.
代入$x=b$可以恢复两项初值, 再求二阶导数验证方程, 因而由常微分方程唯一性确定该解. 在$x=0$代入得到
\begin{align*}
u(0)&=t\e^{\ii kb}\left(\cosh\kappa b-\frac{\ii k}{\kappa}\sinh\kappa b\right),\\
u'(0)&=t\e^{\ii kb}\left(-\kappa\sinh\kappa b+\ii k\cosh\kappa b\right).
\end{align*}
左侧 $u(0)=1+r$, $u'(0)=\ii k(1-r)$. 用 $u(0)+u'(0)/(\ii k)=2$ 消去$r$, 得
\[
t=\frac{\e^{-\ii kb}}
{\cosh\kappa b+\frac{\ii}{2}(\kappa/k-k/\kappa)\sinh\kappa b}.
\]
两侧波数相同, 所以$T=|t|^2$. 利用$\cosh^2s=1+\sinh^2s$和$k^2+\kappa^2=2mV_0/\hbar^2$, 得
\begin{equation}\label{eq:tunnel}
T=\left[1+\frac{V_0^2\sinh^2(\kappa b)}{4E(V_0-E)}\right]^{-1}.
\end{equation}
它对有限$b$为正, $b\to0$时$T\to1$. 单独实指数的流为零, 但$A\e^{\kappa x}+B\e^{-\kappa x}$的交叉项给$J=-2\hbar\kappa\operatorname{Im}(A^*B)/m$, 可以承载与两侧相同的净流. 不能从禁止区的实波数变虚就断言势垒内部没有流. 厚势垒$\kappa b\gg1$下, $\sinh^2\kappa b\simeq\e^{2\kappa b}/4$, 因而$T\simeq16E(V_0-E)\e^{-2\kappa b}/V_0^2$. 指数依赖是隧穿对厚度特别敏感的原因.

\begin{center}
\begin{tikzpicture}[x=1.2cm,y=1cm]
\draw[-{Stealth}] (-2.8,0)--(4.7,0) node[right]{$x$};
\draw[-{Stealth}] (-2.3,-.15)--(-2.3,2.3) node[above]{$V$};
\draw[thick] (-2,0)--(0,0)--(0,1.6)--(2.4,1.6)--(2.4,0)--(4.2,0);
\draw[dashed] (-1.9,.8)--(4.1,.8);
\node[above] at (1.2,1.6){$V_0$};
\node[left] at (-2.3,.8){$E$};
\node[below] at (0,0){$0$};\node[below] at (2.4,0){$b$};
\draw[-{Stealth},structurecolor,thick] (-1.8,.35)--(-.6,.35);
\draw[-{Stealth},structurecolor,thick] (3,.35)--(3.9,.35);
\node[above] at (-1.2,.35){入射};\node[above] at (3.5,.35){透射};
\end{tikzpicture}
\par 矩形势垒示意; 图中长度和高度任取, 概率由匹配方程决定.
\end{center}

隧穿没有让粒子获得额外能量. 静态势中入射和透射有同一个$E$. 禁止区内不能沿用经典动能$mv^2/2$来给某条轨迹赋予负数; 波函数的指数解属于新的运动描述. 在$E>V_0$情形, 中间波数实, 相同匹配给含$\sin^2(qb)$的透射率, $qb=n\pi$时有完全透射, 这是共振而非能量从势垒中消失.

\Needspace{9\baselineskip}
\originalchapter{谐振子与Gaussian波包}
\originalsection{升降算符与能级}
在势能极小点附近, Taylor展开的一次项为零, 二次项通常主导, 因而出现 $V(x)=m\omega^2x^2/2$, $\omega>0$. 高次项小才允许此近似; 分子振动和局部束缚常因此近似为谐振子. 定义
\[
a=\sqrt{\frac{m\omega}{2\hbar}}x+\frac{\ii p}{\sqrt{2m\hbar\omega}},
\qquad a^\dagger=\sqrt{\frac{m\omega}{2\hbar}}x-\frac{\ii p}{\sqrt{2m\hbar\omega}}.
\]
在共同的光滑快速下降定义域, $[x,p]=\ii\hbar$给$[a,a^\dagger]=1$, 而 $H=\hbar\omega(a^\dagger a+1/2)$. 令$N=a^\dagger a$, 则对归一态有$\langle N\rangle=\|a\psi\|^2\ge0$. 对易式 $[N,a]=-a$, $[N,a^\dagger]=a^\dagger$ 表示$a$降低、$a^\dagger$升高本征值一个单位.

若$N\ket{\nu}=\nu\ket{\nu}$且$\nu$有限, 连续降阶后的非零态必须保持非负本征值. 降阶范数平方依次含$\nu,\nu-1,\nu-2,\ldots$; 如果$\nu$不是非负整数, 在某一步会出现负范数平方. 因而序列必须终止于$a\ket0=0$, 并有$\nu=n=0,1,\ldots$. 对标准实线谐振子自伴实现, 这些态构成完整基; 完备性是Hermite函数的$L^2$性质, 不由仅有对易式的形式代数自动保证.

基态方程$a\psi_0=0$是 $\psi_0'=-m\omega x\psi_0/\hbar$, 所以
\[
\psi_0(x)=\left(\frac{m\omega}{\pi\hbar}\right)^{1/4}
\e^{-m\omega x^2/(2\hbar)},\qquad E_n=\hbar\omega\left(n+\frac12\right).
\]
由$\ket n=(a^\dagger)^n\ket0/\sqrt{n!}$构造激发态, 可逐次验证$a\ket n=\sqrt n\ket{n-1}$, $a^\dagger\ket n=\sqrt{n+1}\ket{n+1}$. 能量间隔相等源自升降算符代数; 基态能量非零源自$H$中的$1/2$.

\originalsection{位置、动量与期望值}
反解得到 $x=\sqrt{\hbar/(2m\omega)}(a+a^\dagger)$ 和 $p=-\ii\sqrt{m\hbar\omega/2}(a-a^\dagger)$. 在$\ket n$中, 单个升降算符的期望值为零. 展开$x^2$后, $a^2$和$(a^\dagger)^2$也把态移到正交的能级; 只有$a a^\dagger$和$a^\dagger a$留下:
\[
\langle x^2\rangle_n=\frac{\hbar}{m\omega}\left(n+\frac12\right),
\qquad \langle p^2\rangle_n=m\hbar\omega\left(n+\frac12\right).
\]
所以$\Delta x\Delta p=\hbar(n+1/2)$, 仅基态在这些能量本征态中达到$\hbar/2$. 平均动能与势能各为$E_n/2$, 但不是说每次测量都得到这两个平均值.

\originalsection{自由Gaussian的扩散}
给宽度参数$a>0$的初态 $\psi(x,0)=(2\pi a^2)^{-1/4}\e^{-x^2/(4a^2)}$. 自由演化在波数表象中把每个分量乘以$\e^{-\ii\hbar k^2t/(2m)}$. 第4章已算初态波数振幅为$(2a^2/\pi)^{1/4}\e^{-a^2k^2}$. 因而需计算$\int\exp[-(a^2+\ii\hbar t/(2m))k^2+\ii kx]\dd k$. 对$\operatorname{Re}A>0$的Gaussian积分为$\sqrt{\pi/A}\exp[B^2/(4A)]$; 它从实正$A$的积分经解析延拓成立. 取$A=a^2+\ii\hbar t/(2m)$、$B=\ii x$代入. 记$\tau=\hbar t/(2ma^2)$, 得
\[
\psi(x,t)=\frac{(2\pi a^2)^{-1/4}}{\sqrt{1+\ii\tau}}
\exp\left[-\frac{x^2}{4a^2(1+\ii\tau)}\right],
\]
平方根取与$t=0$连续的分支. 取模平方得
\[
\rho(x,t)=\frac1{\sqrt{2\pi}a\sqrt{1+\tau^2}}
\exp\left[-\frac{x^2}{2a^2(1+\tau^2)}\right].
\]
因此$\Delta x(t)=a\sqrt{1+\tau^2}$, 而动量分布只有相位改变, $\Delta p(t)=\hbar/(2a)$保持不变. 这与不同动量分量具有不同速度相符. 初态不带线性相位时中心固定; 若乘以$\e^{\ii p_0x/\hbar}$, 中心变为$p_0t/m$, 宽度保持同样演化.

初态达到不确定性最小值, 后来$\Delta x\Delta p$变大. 量子演化没有违反不确定性关系: 它给下界, 并不要求所有时刻都取等. 演化还产生位置与动量的相关性, 仅用两个独立标准差不能完整描述波包.

\originalchapter{态空间、测量与动力学}
\originalsection{向量空间、基与维数}
这一补充对应8.05原讲义Chapter 3全章的七节主题, 原作者Barton Zwiebach（2013年10月21日）并信用Sheldon Axler的线性代数讨论. 采用独立中文解释, 没有复制或打包Axler商业教材. 先从两能级系数$(c_1,c_2)$理解向量: 可加、可乘复数, 运算封闭, 有零矢和加法逆元. 这些代数运算满足交换、结合和分配律. 用复数作标量与用实数作标量会改变维数: $\mathbb C$自身在复数域维数1, 在实数域基为$1,\ii$, 维数2. 因而说“二维态”必须说明采用复振幅的维数.

一组矢\textbf{张成}空间, 指每个矢可写为它们的有限线性组合; \textbf{线性独立}指只有全零系数组合得零. 同时满足两者的是基, 系数展开唯一. 若有$n$个基矢, 维数为$n$; 矩阵空间的标准基是每次只在一格放1的$E_{ij}$, 维数为行数乘列数. 多项式空间包含任意长独立列$1,x,\ldots,x^N$, 因而无限维. 量子平方可积函数还需要范数极限来容纳无限展开, 这就是Hilbert空间比纯代数向量空间多的完整性条件.

子空间必须含零且对加法和倍乘封闭. $U+W=\{u+w\}$; 若交集只有零, 每个和分解唯一, 记$U\oplus W$. 正交直和另要求两子空间内任意两矢内积0. 代数直和不自动正交, 例如平面中的斜轴给唯一坐标却不是正交坐标.
\originalsection{线性映射、矩阵与本征矢}
$T:V\to W$满足$T(au+bv)=aTu+bTv$时称线性映射, $\ker T$是被映为0的矢, $\operatorname{im}T$是全部输出. \textbf{单射}指$Tu=Tv$必有$u=v$, 即不同输入不合并为同一输出; \textbf{满射}指每个$w\in W$都有$v\in V$使$Tv=w$, 即输出覆盖整个目标空间. 对线性映射, 检验单射只需解$Tv=0$并确认只有零解, 检验满射则确认$\operatorname{im}T=W$. 有限维中取零核基$k_1,\ldots,k_r$, 扩成全空间基$k_1,\ldots,k_r,v_1,\ldots,v_s$. 后面$s$个像张成像空间; 若$\sum a_jTv_j=0$, 则$\sum a_jv_j$落在零核, 与扩充基独立性矛盾, 故它们独立. 所以$\dim V=\dim\ker T+\dim\operatorname{im}T$. 相同有限维的输入输出间, 单射、满射、可逆等价; 无限维左移序列是满射而非单射, 不能照搬.

对$T:V\to V$, 选基后$Tv_j=\sum_iT_{ij}v_i$, 第$j$列表示第$j$基矢的输出. 复合对应矩阵乘法, 恒等对应单位阵. 对$n\times n$方阵, \textbf{迹}是对角元素之和, \textbf{行列式}是带排列符号的乘积和:
\[
 \Tr T=\sum_{i=1}^nT_{ii},\qquad
 \det T=\sum_{\sigma\in S_n}(-1)^{N_\sigma}\prod_{i=1}^nT_{i,\sigma(i)}.
\]
这里$S_n$为$1,\ldots,n$的全部排列, $N_\sigma$是满足$i<j$而$\sigma(i)>\sigma(j)$的逆序对数; 例如二维行列式为$T_{11}T_{22}-T_{12}T_{21}$. 实际计算可用一行加另一行的倍数消元, 这不改行列式; 每次交换两行变号, 上三角阵的行列式为对角元素之积. \textbf{特征多项式}定义为$p_T(z)=\det(zI-T)$, 可按行列式计算并整理成$z$的多项式. 换基列矩阵$C$满足$[v]_{\rm old}=C[v]_{\rm new}$, 算符新矩阵$C^{-1}TC$, 其迹、行列式和特征多项式不变. 这说明一个矩阵表格依基, 其物理作用不依基.

$Tu=\lambda u,u\ne0$定义本征对. 求本征值先解$p_T(\lambda)=0$, 再解$(T-\lambda I)u=0$找非零本征矢; 有对角本征基时, 迹与行列式分别给本征值之和与积, 可用来核对计算. 不同本征值的本征矢独立: 若最短非平凡依赖$\sum_{j=1}^na_ju_j=0$, 作用$T-\lambda_nI$消去最后一项, 剩下更短非平凡依赖, 矛盾. 复空间的特征多项式有根保证存在一个本征矢; 实矩阵如二维旋转可无实本征值. 有$n$个独立本征矢才可对角化, 存在一个本征矢远远不够. 例如$\begin{pmatrix}1&1\\0&1\end{pmatrix}$的$(T-I)(x,y)=0$只给$y=0$, 本征矢仅为$(x,0)$且$x\ne0$, 故本征空间一维, 不能对角化.
\originalsection{内积、正交化与投影}
内积$(u,v)$对第二变量线性、对第一变量共轭线性, 共轭交换, $\norm v^2=(v,v)>0$对非零$v$成立. 复空间若不共轭,$(1,\ii)$的“平方和”会为0而非正长度. Schwarz不等式来自将$u$减去沿$v$的最佳倍数; 三角不等式随之成立, 作业3有完整推导.

给独立列$v_1,\ldots,v_n$, Gram--Schmidt先$e_1=v_1/\norm{v_1}$, 再$w_j=v_j-\sum_{i<j}e_i(e_i,v_j)$, $e_j=w_j/\norm{w_j}$. 正交性由逐项内积核验; 独立性保证$w_j\ne0$. 它不会改变每一步已张成的子空间. 对子空间$U$的正交基,$P_Uv=\sum_je_j(e_j,v)$, 自然有$P_U^2=P_U=P_U^\dagger$. $v-P_Uv$与每个$e_j$正交, 对有限维$U$因而$V=U\oplus U^\perp$; 一般Hilbert空间中需$U$闭, 投影展开按范数收敛理解, 分解唯一. 以两部分的Pythagoras式证明$P_Uv$是唯一最佳近似, 详见作业4的函数逼近应用. 单有$P^2=P$只是斜投影, 还要Hermitian才能作为正交测量投影.
\originalsection{对偶、伴随与幺正}
线性泛函$\varphi$把矢映成复数, 且保持线性. 正交基中设$\varphi(e_j)=a_j$, 则$\varphi(v)=\sum a_jv_j=(u,v)$, 唯一$u=\sum a_j^*e_j$. 这就是有限维Riesz对应, 在Hilbert空间需要泛函连续. bra是这种泛函, 其系数恰为ket系数的复共轭.

伴随由$(u,Tv)=(T^\dagger u,v)$定义. 正交基中矩阵为共轭转置, 因为逐基内积给$(T^\dagger)_{ji}=T_{ij}^*$. 可核$(ST)^\dagger=T^\dagger S^\dagger$, 和的伴随逐项相加, 常数取共轭. Hermitian满足$T=T^\dagger$, 本征值实与异值本征矢正交的证明在下文给出谱分解. 有限维幺正满足$U^\dagger U=UU^\dagger=I$并保持所有内积; 反过来对每对矢保持内积便有$U^\dagger U=I$, 同有限维时单射亦满射, 再得$UU^\dagger=I$; 无限维只保内积的非满射等距嵌入不称幺正. 正交归一基经幺正作用仍正交归一. 幺正本征值模1: 对$Uu=zu$比较范数得$|z|=1$.

任意有限维复矩阵还有Schur上三角化: 选一个单位本征矢作第一基矢, 第一列除顶格外为0, 下方余块再递归酉换基, 得上三角. Hermitian上三角必须同时下三角, 所以对角. 更一般normal矩阵满足$TT^\dagger=T^\dagger T$, 上三角normal的第一行除对角元素外须为0（比较第一格对角的行范数与列范数）, 归纳仍对角. 这解释何种算符拥有完备正交本征基, 不是把所有矩阵都假设为可测观测量.

\originalsection{把代数规则补到可复算的算符语言}
Hermitian矩阵对实倍乘封闭, 对复倍乘不封闭（$\ii I$不Hermitian）; 它们构成实向量空间, $n\times n$时维数$n^2$, 与全部复矩阵的复维数$n^2$要区分. 无限序列空间含任意长的标准独立列, 函数空间亦有多项式独立列; 这说明无限维而非声称有限和能表示每个函数. 任一独立列不长于有限生成列: 逐次把独立矢代入生成式, 用一个非零系数替换对应生成矢, 保持生成; 若独立列更长, 最后一矢便落在前面独立矢张成空间, 矛盾. 由此所有有限基长度相同.

线性映射的和、倍乘与复合仍线性; 复合通常不交换. 在多项式上$D=\partial_x,M=x$满足$(DM-MD)f=f$. 序列的右移$R(a_0,a_1,\ldots)=(0,a_0,a_1,\ldots)$单射非满射, 左移$L(a_0,a_1,\ldots)=(a_1,a_2,\ldots)$满射非单射, $LR=I$但$RL\ne I$. 这提供了无限维不能凭单侧逆推出双侧逆的具体检查.

子空间$W$不变指$TW\subset W$, 才能将$T$限制为$W$上的算符. 本征空间$\ker(T-\lambda I)$是不变子空间, 包含零矢; 本征矢则必须非零. Jordan例中本征值1的代数重数2、几何重数1, 因而这两个“重数”不能互换.

正交投影在$U$上为恒等、在$U^\perp$上为零, 本征值只有0、1; 取适应该分解的有限正交基, 得$\operatorname{rank}P=\Tr P=\dim U$. $U=V$时$P=I$可逆, 只有真子空间投影才不可逆. 理想测量的子空间概率是$\norm{P\psi}^2$, 条件态为$P\psi/\norm{P\psi}$, 若概率0则该条件态未定义.

对固定$u$, 泛函$v\mapsto(u,Tv)$由Riesz对应唯一矢$T^\dagger u$, 故伴随存在唯一; 泛函对$u$共轭线性而Riesz对应也共轭, 两者抵消使$T^\dagger$对$u$线性. 定义立即给$(T^\dagger)^\dagger=T$, $(cT)^\dagger=c^*T^\dagger$及逆序乘积. 非正交基的Gram矩阵$G_{ij}=(v_i,v_j)$使内积为$x^\dagger Gy$, 因而伴随矩阵是$G^{-1}T^*G$, 其中$T^*$为共轭转置, 不能无条件用$T^*$代替伴随.

复空间若$(v,Av)=0$对全体$v$, 令$B(u,v)=(u,Av)$. 先用$u+v$扣除两个对角项得$B(u,v)+B(v,u)=0$, 再用$u+\ii v$得$\ii B(u,v)-\ii B(v,u)=0$, 故所有$B(u,v)=0$, $A=0$. 若全态$(v,Tv)$实, 则$(v,(T-T^\dagger)v)=0$, 上式推出$T=T^\dagger$; 反向由共轭交换立即成立. 实空间探针不能得到第二个式子, 这正是实反对称矩阵零二次型例的原因.

满射线性算符保持范数, 由复极化恢复全部内积, 因而幺正; 有限维中范数保持已保证单射并进而满射. 若$v_j=Uu_j$为两正交基, 则$U$在两基中的矩阵相同: $(v_i,Uv_j)=(Uu_i,U^2u_j)=(u_i,Uu_j)$. 这是同一个算符在特定两基的性质, 不是任意换基都不改变矩阵.


范数的三条性质为正定、$\norm{cv}=|c|\norm v$及三角不等式. 完备指每个Cauchy序列（任意$\epsilon>0$后项两两距离小于$\epsilon$）有空间内的范数极限, 有限维自动完备而一般函数子空间未必. 核与像对线性组合封闭且含零, 故确为子空间. 两个双侧逆若存在,$S=S(TQ)=(ST)Q=Q$, 所以逆唯一. 有限生成列删除冗余项可成基, 独立列持续添加外部矢直至张成亦成基. 有限维$T-\lambda I$不单射等价行列式0, 将本征空间和特征多项式联系起来.

Riesz表示唯一性也直接可核: 若$(u,v)=(u',v)$对全体$v$, 取$v=u-u'$得$\norm{u-u'}^2=0$. 伴随逆序证明为$(u,STv)=(S^\dagger u,Tv)=(T^\dagger S^\dagger u,v)$, 唯一性给结论. 例如正交基中$T=\begin{pmatrix}0&2&\ii\\1&-\ii&0\\3\ii&1&7\end{pmatrix}$, 则$T^\dagger=\begin{pmatrix}0&1&-3\ii\\2&\ii&1\\-\ii&0&7\end{pmatrix}$; 每格为对称位置的复共轭. 对无界算符仍必须保留后文所说的定义域条件, 不能以有限矩阵论证替代.


\originalsection{Dirac记号和谱分解}
前面的势阱问题已给出物理入口: 同一个制备可以用位置波函数描述, 也可以展开成各能量态的振幅; 能量测量的概率由相应振幅给出. 线性代数在这里整理的是这些可叠加的振幅和可改变振幅的操作. 换一套坐标不应改变制备或探测结果, 因而需要不绑在位置坐标上的记号.

位置波函数是态的一种表示. 态写为$\ket\psi$, 读作ket. 与它对应的$\bra\psi$读作bra, 是把任意输入态映到复数内积的规则, 称为\textbf{线性泛函}. 内积为$\braket\phi\psi$. $\bra\psi$对输入ket线性, 即保持加法和复数倍乘; 但从$\ket\psi$到$\bra\psi$的对应是共轭线性的, 标量倍乘会变成其共轭. $\ket\psi$可在不同基中表示为不同的系数列, 物理态不会因换坐标而改变. 外积$\ket u\bra v$是算符: 它作用于$\ket\psi$后得到$\ket u\braket v\psi$.

先算熟悉的两盒能级$\psi=c_1u_1+c_2u_2$, 系数列$(c_1,c_2)^{\mathsf T}$. 恒等$I$不改任何态. 投影$P_1$只留下第一分量,$P_1\psi=c_1u_1$, 矩阵$\operatorname{diag}(1,0)$; $P_2$同理为$\operatorname{diag}(0,1)$, $P_1+P_2=I$. 能量读数$E_1$概率$|c_1|^2$, 若非零, 测后为$u_1$到整体相位. 投影是在算什么由此可直接核对.

下面的有限维谱证明作为进阶线代工具说明: 子空间是对加法和倍乘封闭的集合,$u^\perp$是全部与$u$正交的矢; 使用复多项式存在根的标准工具. 先考虑有限维复内积空间, 避免无界定义域问题. 厄米矩阵有一个本征向量, 因为其特征多项式在复数域有根. 本征值为实, 如第4章所证. 若$u$是归一化本征向量, $u^\perp$在$A$下不变: 对$v\perp u$, $\braket u{Av}=\braket{Au}v=0$. 对维数归纳, 便得到正交归一本征基$\ket{u_n}$. 所以
\[
I=\sum_n\ket{u_n}\bra{u_n},\qquad
A=\sum_n a_n\ket{u_n}\bra{u_n}.
\]
如果有简并, 同一本征值的所有本征向量组成子空间, 用其投影$P_a$统一写$A=\sum_a aP_a$.

对归一态, Born测量规则为$\Pr(A=a)=\bra\psi P_a\ket\psi=\|P_a\psi\|^2$. 理想投影测量后, 若结果$a$发生且概率非零, 条件态为$P_a\ket\psi/\|P_a\psi\|$. 简并时不指定简并子空间内某一个任意基向量; 测量$A$本身只确定子空间. 更一般的测量可以不同于投影测量, 本册这一节只使用理想投影模型.

\originalsection{相容性与同时本征基}
相容性问能否制备对两个读数都确定的态, 并用一套共同确定读数态展开. 两个对角测量例如盒的$H$和任意$f(H)$使用相同驻波, 所以读数可以同时明确. 不意味着任何具体装置的顺序实施都无额外扰动. 若$A,B$是对易的有限维厄米算符, 对$A$的本征子空间$V_a$, 有 $A(Bv)=B(Av)=aBv$, 所以$B$保持$V_a$. 在每个$V_a$内将$B$厄米对角化, 可得到同时本征基. 若$A$本征值不简并, $V_a$是一维, 所以每个$A$本征向量已是$B$本征向量.

对易是\emph{存在}同时基的条件, 不等于任意两个给定本征向量都相同. 连续谱和无界算符的情形须使用适当的谱投影相容性, 不能仅凭某个稠密子域上的形式$[A,B]=0$断言强对易.

\originalsection{幺正演化和期望值方程}
闭合系统演化保持总概率以及态之间的内积. 幺正就是$U^\dagger U=I$; 在能量基中每个$u_n$仅乘模1相位$\e^{-\ii E_nt/\hbar}$, 可逐项看到内积保持. 有限维且$H$定常, 令$U(t)=\e^{-\ii Ht/\hbar}$. 逐本征态定义指数, 便有$U^\dagger U=I$, 所以演化保持内积和所有态的范数. 它也说明两个态若受同一Hamilton算符演化, 重叠$\braket{\phi(t)}{\psi(t)}$不变. 位置概率密度可能分离, 复内积仍可以保持不变.

点号表示时间导数. 算符显含时间是测量规则本身随时间改变, 与态变化是两种不同来源. 对于显含时间的观测量$A(t)$,
\begin{align*}
\frac{\mathrm d}{\mathrm dt}\langle A\rangle
&=\langle\dot\psi,A\psi\rangle+\langle\psi,\dot A\psi\rangle+\langle\psi,A\dot\psi\rangle\\
&=\frac{\ii}{\hbar}\langle[H,A]\rangle+\left\langle\frac{\partial A}{\partial t}\right\rangle.
\end{align*}
上式代$\dot\psi=-\ii H\psi/\hbar$与其共轭, 两个态导数项分别给$\ii\langle HA\rangle/\hbar$及$-\ii\langle AH\rangle/\hbar$, 才得到对易子. 对$p^2/(2m)+V(x)$, 利用乘积对易式先算$[p^2,x]=p[p,x]+[p,x]p=-2\ii\hbar p$. 对试探函数$f$, $(Vp-pV)f=-\ii\hbar Vf\prime+\ii\hbar(V\prime f+Vf\prime)=\ii\hbar V\prime f$, 因而$[V,p]=\ii\hbar V\prime$, 对合适光滑态, $\frac{\mathrm d}{\mathrm dt}\langle x\rangle=\langle p\rangle/m$, $\frac{\mathrm d}{\mathrm dt}\langle p\rangle=-\langle V'(x)\rangle$. 这就是Ehrenfest关系. 一般$\langle V'(x)\rangle\ne V'(\langle x\rangle)$, 因而平均位置不总是精确满足经典轨道方程. 二次势因力为线性函数而使两者相同.

\begin{remark}
如果$H(t)$在不同时刻不对易, $U(t)$不能简单写成$\exp[-\ii\int_0^tH(s)\dd s/\hbar]$. 需要时间排序或直接解矩阵微分方程. 这条限制在驱动二态和含时微扰中会再次出现.
\end{remark}

\originalsection{Schrödinger、Heisenberg与相互作用表象}
三种表象分配时间依赖不同, 实验概率相同. Schrödinger态$\ket{\psi_S(t)}=U(t)\ket{\psi_0}$, 固定观测量写$A_S$. Heisenberg态固定为$\ket{\psi_0}$, 算符$A_H=U^\dagger A_SU$. 将它放进均值$\bra{\psi_0}A_H\ket{\psi_0}$便恢复$\bra{\psi_S}A_S\ket{\psi_S}$. 乘积、对易子也在同一共轭下保持, 即便$A_S$显含时间也如此.

例如自由粒子$\dot p_H=0,\dot x_H=p_H/m$, 因而$p_H=p,x_H=x+pt/m$. 恒力$V=gx$时$p_H=p-gt$, $x_H=x+pt/m-gt^2/(2m)$. 谐振子给$x_H=x\cos\omega t+p\sin\omega t/(m\omega)$,$p_H=p\cos\omega t-m\omega x\sin\omega t$. 这些是算符等式, 对任意合适初态取均值即可预测均值轨迹, 也可平方来计算宽度.

若$H=H_0+V(t)$且$H_0$不含时间, 用已会解的$H_0$先除去快速演化: $\ket{\psi_I}=\e^{\ii H_0t/\hbar}\ket{\psi_S}$, $V_I=\e^{\ii H_0t/\hbar}V\e^{-\ii H_0t/\hbar}$, 直接求导给$\ii\hbar\dot\psi_I=-H_0\psi_I+(H_0+V_I)\psi_I=V_I\psi_I$, 显示两个$H_0$项如何抵消. 积分一次并把右侧态再代回, 产生
\[
U_I(t)=I-\frac\ii\hbar\int_0^tV_I(s)\dd s+
\left(-\frac\ii\hbar\right)^2\int_0^t\dd s\int_0^s\dd s' V_I(s)V_I(s')+\cdots.
\]
早时刻算符在右, 积分只取$s'>0,s'<s$, 这就是时间排序的具体意义. 在弱驱动下截在第一阶, 得第15章的转移振幅; 并非先要求读者知道一个抽象时间排序符号.

\originalchapter{自旋与二态系统}
\originalsection{自旋二分之一}
Stern--Gerlach装置让原子束穿过不均匀磁场, 再在屏上记录偏转. 磁矩的势能$H_{\rm mag}=-\boldsymbol\mu\cdot\boldsymbol B$, 沿$z$的场梯度产生$F_z\simeq\mu_z\partial_zB_z$, 把磁矩分量转换为位置带. 均匀场主要使方向进动, 本身不提供这个分束力. 用自旋$1/2$的有效自由度描述双输出, 磁矩与自旋成正比. Stern--Gerlach实验中的分束表明某些内禀角动量的分量取离散值. 自旋不是一个有确定表面速度的经典小球自转模型; 它是独立于轨道坐标的量子自由度. 对自旋$1/2$, 选$S_z$本征基$\ket{+},\ket{-}$, 本征值为$\pm\hbar/2$. 任意归一纯态是$\alpha\ket{+}+\beta\ket{-}$, $|\alpha|^2+|\beta|^2=1$.

先选$z+$, 再沿$x$测量得到两种结果各半; 筛选$x+$后再测$z$仍各半. 因此$x+$不是原$z+$标签, 可取$\ket{x;\pm}=(\ket+\pm\ket-)/\sqrt2$. 从$S_x=(\hbar/2)(\ket{x;+}\bra{x;+}-\ket{x;-}\bra{x;-})$恢复非对角矩阵, 它是测量基变化而非多加一个经典随机标签. $y$向上态为$(\ket++\ii\ket-)/\sqrt2$, 相对$\ii$相位使它与$x$测量不同. 三种自旋分量可表示为$S_j=(\hbar/2)\sigma_j$, 其中
\[
\sigma_x=\begin{pmatrix}0&1\\1&0\end{pmatrix},\quad
\sigma_y=\begin{pmatrix}0&-\ii\\\ii&0\end{pmatrix},\quad
\sigma_z=\begin{pmatrix}1&0\\0&-1\end{pmatrix}.
\]
令$\epsilon_{ijk}$为Levi--Civita符号: $\epsilon_{123}=1$, 交换任意两指标变号, 重复指标时为零. 直接矩阵乘法给$\sigma_i\sigma_j=\delta_{ij}I+\ii\sum_k\epsilon_{ijk}\sigma_k$, 所以$[S_i,S_j]=\ii\hbar\sum_k\epsilon_{ijk}S_k$. 对$S_z$确定的态, $S_x,S_y$一般不确定. 例如$\ket+$中$\langle S_x\rangle=0$, 但$\langle S_x^2\rangle=\hbar^2/4$, 所以$\Delta S_x=\hbar/2$, 不是零.

对单位向量$\boldsymbol n=(\sin\theta\cos\phi,\sin\theta\sin\phi,\cos\theta)$, $\theta$是相对$z$轴的极角,$\phi$是$xy$平面方位角; 纯态的方向指均值向量方向, 三分量并非都确定. 沿该方向的自旋是$S_{\boldsymbol n}=\boldsymbol n\cdot\boldsymbol S$. $+\hbar/2$本征态可取
$\ket{\boldsymbol n,+}=\cos(\theta/2)\ket+ +\e^{\ii\phi}\sin(\theta/2)\ket-$. 代入$\boldsymbol n\cdot\boldsymbol\sigma$验证其本征值为1. 例如$\theta=\pi/2,\phi=0$恢复$x+$, $\phi=\pi/2$恢复$y+$, 直接矩阵乘法给各对应本征值$+\hbar/2$. 从$z$向上态测量沿$\boldsymbol n$自旋, 向上概率为$\cos^2(\theta/2)$.

\originalsection{一般二态Hamilton算符}
二态不是“只有两个可能的纯态”, 而是态空间维数为二. 例如双势阱左右局域制备$L,R$中矩阵为$\left(\begin{smallmatrix}E_L&K\\K^*&E_R\end{smallmatrix}\right)$. 对角能量表左右偏置, 非对角$K$使局域态可改变. 它给$E_c=(E_L+E_R)/2$, $h_z=(E_L-E_R)/2$, $h_x=\operatorname{Re}K,h_y=-\operatorname{Im}K$, 各量都为能量. 厄米$2\times2$矩阵有四个实参数, 可唯一写为$H=E_cI+\boldsymbol h\cdot\boldsymbol\sigma$, $\boldsymbol h\in\R^3$. 由于$(\boldsymbol h\cdot\boldsymbol\sigma)^2=|\boldsymbol h|^2I$, 本征值为$E_c\pm|\boldsymbol h|$. 假定$E_c,\boldsymbol h$不随时间改变. 指数的偶次、奇次项分开求和, 得
\[
U(t)=\e^{-\ii E_ct/\hbar}\left[
\cos\frac{|\boldsymbol h|t}{\hbar}I
-\ii\sin\frac{|\boldsymbol h|t}{\hbar}\frac{\boldsymbol h\cdot\boldsymbol\sigma}{|\boldsymbol h|}\right],
\]
$\boldsymbol h=0$时按连续极限取$U=\e^{-\ii E_ct/\hbar}I$. 公共能量$E_c$只贡献整体相位, 能级差$2|\boldsymbol h|$决定可观测振荡.

\begin{example}
在正交基$\ket L,\ket R$中, $H=E_cI-\Delta\sigma_x$, $\Delta>0$. $\Delta$为能量耦合而非频率, 初态$\ket L$. 求在时间$t$测得$\ket R$的概率.
\end{example}
\begin{solution}
演化给$\ket{\psi(t)}=\e^{-\ii E_ct/\hbar}[\cos(\Delta t/\hbar)\ket L+\ii\sin(\Delta t/\hbar)\ket R]$, 所以$P_R(t)=\sin^2(\Delta t/\hbar)$. 能量本征態为$(\ket L\pm\ket R)/\sqrt2$, 初态同时包含两者. 转移不是沿某条指定经典路径发生; 它由两个本征态的相对相位产生.
\end{solution}

加入不对称项$\epsilon\sigma_z$后, 令$\Omega=\sqrt{\Delta^2+\epsilon^2}/\hbar$, 转移概率为$P_R=\Delta^2\sin^2(\Omega t)/(\Delta^2+\epsilon^2)$. 概率振荡的角频率为$2\Omega$, 周期$\pi/\Omega$, 最大值为$\Delta^2/(\Delta^2+\epsilon^2)$. 能级失谐会阻止完全转移, 这从演化矩阵的非对角元直接得到. 磁共振和氨分子二态模型使用同一代数, 但具体的$\boldsymbol h$须从各自物理相互作用识别.

\originalchapter{角动量与中心势}
\originalsection{角动量代数}
中心力沿半径, 经典力矩$\boldsymbol r\times\boldsymbol F=0$; 量子中心势也不选空间方向, 旋转不改变Hamiltonian. 因而应寻找旋转生成的守恒量. $L^2$测总轨道角动量大小的平方, $L_z$测相对选定轴的投影; 不同分量不对易, 不能同时赋确定$x,y,z$读数. $\boldsymbol L$为轨道,$\boldsymbol S$为内禀自旋,$\boldsymbol J$可为任一种或它们之和. $j,m$无量纲, 实际读数带$\hbar$单位. 轨道角动量定义为$\boldsymbol L=\boldsymbol r\times\boldsymbol p$, $\boldsymbol p=-\ii\hbar\nabla$. 在光滑函数的共同定义域, 利用$[x_i,p_j]=\ii\hbar\delta_{ij}$逐项展开, 可得$[L_x,L_y]=\ii\hbar L_z$及循环式. $L^2=L_x^2+L_y^2+L_z^2$与三个分量都对易. 因而可选$L^2,L_z$共同本征态. 以下对一般角动量$\boldsymbol J$使用同一代数.

设$J^2\ket{j,m}=\hbar^2j(j+1)\ket{j,m}$, $J_z\ket{j,m}=\hbar m\ket{j,m}$. 初始可把$J^2$本征值写为非负参数$\hbar^2\lambda$; 下面的终止条件说明$\lambda=j(j+1)$的来源. 定义$J_\pm=J_x\pm\ii J_y$, 则$[J_z,J_\pm]=\pm\hbar J_\pm$且$[J^2,J_\pm]=0$. 升降只改变$m$, 不改变$J^2$.

由 $J_-J_+=J^2-J_z^2-\hbar J_z$ 和 $J_+J_-=J^2-J_z^2+\hbar J_z$, 范数平方为$\hbar^2[\lambda-m(m\pm1)]$. 它必须非负, 所以$m$不能无限升降. 若最高和最低值为$m_+,m_-$, 终止给$\lambda=m_+(m_++1)=m_-(m_--1)$. 相减因式分解得$(m_+-m_-+1)(m_++m_-)=0$, 前一因子正, 所以$m_-=-m_+$; “表示”在此指由升降相连的一组态, 不是粒子绕过多少圈. 记$m_+=j\ge0$. 两端之间的整数步数为$2j$, 因而$j$为非负整数或半整数, $m=-j,-j+1,\ldots,j$. 归一化系数为
\[
J_\pm\ket{j,m}=\hbar\sqrt{(j\mp m)(j\pm m+1)}\ket{j,m\pm1}.
\]
这是角动量量子数和系数的代数推导, 不能仅从“整数轨道圈数”推出自旋半整数.

\originalsection{球谐函数与径向方程}
球坐标定义为$x=r\sin\theta\cos\phi,y=r\sin\theta\sin\phi,z=r\cos\theta$. 中心势角方向相同, 可先找乘积$R(r)Y(\theta,\phi)$, 用角函数分离旋转性质. $L^2=-\hbar^2[\partial_\theta^2+\cot\theta\partial_\theta+\csc^2\theta\partial_\phi^2]$, 与$L_z$一起给角向本征方程. 在南北极正则、方位单值的解限制$\ell,m$, 本册不要求预先会全部关联Legendre函数. 以下$\mu$表示相对运动的约化质量, 固定中心单粒子时取粒子质量; 两体约化在下一章展开. 对于轨道角动量, 球坐标中$L_z=-\ii\hbar\partial_\phi$, 所以$L_z$本征函数的方位因子为$\e^{\ii m\phi}$. 位置波函数对$\phi\mapsto\phi+2\pi$单值, 强制$m$为整数. 因而轨道角动量的$j=\ell$也为整数. 角向本征函数称为球谐函数$Y_{\ell m}(\theta,\phi)$, 归一化为$\int|Y_{\ell m}|^2\dd\Omega=1$, $\dd\Omega=\sin\theta\dd\theta\dd\phi$. 最简单$Y_{00}=1/\sqrt{4\pi}$给每立体角等概率, 因$\int\dd\Omega=4\pi$而归一. $Y_{10}=\sqrt{3/(4\pi)}\cos\theta$, 代$L^2$得$2\hbar^2Y_{10}$, 代$L_z$得0; $\int\cos^2\theta\dd\Omega=4\pi/3$检验归一. 这分别是$s$与一个$p$角分布, $p$不是沿确定轴的经典轨道. 这与自旋量子数的半整数可能性并不冲突, 因为自旋不以这种单值标量角函数表示.

若$V=V(r)$, 三维方程分离为$\psi=R(r)Y_{\ell m}$. 使用
$\nabla^2=r^{-2}\partial_r(r^2\partial_r)-L^2/(\hbar^2r^2)$, 得
\[
-\frac{\hbar^2}{2\mu}\left(R''+\frac2rR'\right)
+\left[V(r)+\frac{\hbar^2\ell(\ell+1)}{2\mu r^2}\right]R=ER.
\]
取$u=rR$, 实际$R\prime=u\prime/r-u/r^2$, $R\prime\prime+2R\prime/r=u\prime\prime/r$, 因而一阶导数项消失:
\begin{equation}\label{eq:radial}
-\frac{\hbar^2}{2\mu}u''+
\left[V(r)+\frac{\hbar^2\ell(\ell+1)}{2\mu r^2}\right]u=Eu,
\qquad r>0.
\end{equation}
体积元$\dd^3r=r^2\dd r\dd\Omega$, 所以$\int|\psi|^2\dd^3r=\int|R|^2r^2\dd r\int|Y|^2\dd\Omega=\int|u|^2\dd r$. $|u|^2\dd r$是薄球壳概率, 不是某个空间点的密度. 归一化是$\int_0^\infty|u|^2\dd r=1$, 不是$\int|R|^2\dd r=1$. 对通常正则中心势, 物理边界为$u(0)=0$. 额外的离心项把三维问题化为半轴上的有效一维问题. 它不能与全实线的任意一维势阱混同.

\originalsection{角动量合成}
先把联合测量标签写清: $\ket{+-}$表示系统1读$+\hbar/2$、系统2读$-\hbar/2$的乘积制备. $\otimes$把两个态组合, 满足$(a u+bv)\otimes w=a(u\otimes w)+b(v\otimes w)$, 一侧算符只作用该侧. 四个自旋组合形成四维基, 第12章系统讨论一般联合态. 对两个独立系统, $\boldsymbol J=\boldsymbol J^{(1)}+\boldsymbol J^{(2)}$, 不同系统的分量对易, 所以总角动量仍满足同一代数. 乘积基$\ket{j_1,m_1}\otimes\ket{j_2,m_2}$使$J_z$本征值为$\hbar(m_1+m_2)$, 但一般不是$J^2$本征态.

两个自旋$1/2$的最高总磁量子数态是$\ket{++}$, $j=1,m=1$. $J_-=S_{1-}+S_{2-}$, 每个分量分别给$\hbar\ket{-+}$和$\hbar\ket{+-}$, 合态范数为$\sqrt2\hbar$. 用$J_-$降阶, 得三重态$\ket{1,0}=(\ket{+-}+\ket{-+})/\sqrt2$, $\ket{1,-1}=\ket{--}$. 剩下的正交态是$\ket{0,0}=(\ket{+-}-\ket{-+})/\sqrt2$, 直接验证被$J_z$和$J_\pm$湮灭. 四维空间分为三维$j=1$和一维$j=0$子空间.

又因为$J^2=(J^{(1)})^2+(J^{(2)})^2+2\boldsymbol J^{(1)}\cdot\boldsymbol J^{(2)}$, 若相互作用为$\lambda\boldsymbol J^{(1)}\cdot\boldsymbol J^{(2)}$, 两类态的能量移分别为$\lambda\hbar^2/4$与$-3\lambda\hbar^2/4$. 这给自旋耦合的分裂, 不需要对原来的$4\times4$矩阵盲目求特征多项式.

\originalsection{轨道与自旋二分之一的合成}
为下一章分裂准备一个一般结果. 固定轨道$\ell$, 乘自旋$1/2$共$2(2\ell+1)$维. 最大$m=\ell+1/2$只有$\ket{\ell,\ell}\ket+$, 被$J_+$湮灭, 所以它生成$j_+=\ell+1/2$的$2\ell+2$维阶梯. 其正交补在$J_i$作用下保持, 还剩$2\ell$维. 当$\ell\geq1$, $m=\ell-1/2$有两候选$\ket{\ell,\ell-1}\ket+$和$\ket{\ell,\ell}\ket-$, 已有阶梯占一个, 与它正交的一个最高态必被$J_+$湮灭. 因而补空间生成$j_-=\ell-1/2$的$2\ell$维阶梯, 维数已全用完. $\ell=0$只剩$j=1/2$, 没有负$j$.

具体下支最高态写$a\ket{\ell,\ell-1}\ket++b\ket{\ell,\ell}\ket-$, 条件$J_+=0$给$\sqrt{2\ell}a+b=0$, 归一取$a=1/\sqrt{2\ell+1},b=-\sqrt{2\ell/(2\ell+1)}$. 这同时给末态自旋测量概率, 后面2012期末衰变题直接使用. “最高权态”只指一组中最大$m$的态, 并非另外一种物理粒子.

\originalchapter{氢原子的量子化}
\originalsection{两体约化与尺度}
氢原子发光谱有离散频率, 提示内部能量差离散. 现在将中心势方法用于解释束缚能量与位置分布; 质心的自由运动能量不计为内部能级. 电子和原子核都可运动. 设质量为$m_e,M$, 质心坐标为$\boldsymbol X=(m_e\boldsymbol r_e+M\boldsymbol r_N)/(m_e+M)$, 相对坐标为$\boldsymbol r=\boldsymbol r_e-\boldsymbol r_N$. 对只依赖相对坐标的相互作用, 动能分为质心项和相对项:
\[
H=\frac{P^2}{2(m_e+M)}+\frac{p^2}{2\mu}-\frac{g}{r},
\quad \mu=\frac{m_eM}{m_e+M},\quad g=\frac{Ze^2}{4\pi\varepsilon_0}.
\]
$Z$为核电荷整数,$e>0$为基本电荷大小, 核$+Ze$、电子$-e$,$\varepsilon_0$为SI真空介电常数. $P=p_e+p_N$共轭于质心$X$, 相对共轭动量$p=(Mp_e-m_ep_N)/(m_e+M)$. 反解$p_e=m_eP/(m_e+M)+p,p_N=MP/(m_e+M)-p$, 平方相加交叉项抵消, 就得到上式动能. $r$是两粒子间距, 不假定核在实验室固定. 这里使用SI库仑常数; 官方作业的Gaussian单位中写$g=Ze^2$. 两种写法不能混用. 对氢$Z=1$, 且$M\gg m_e$, 所以$\mu$接近$m_e$, 但谱线精密计算不能永远忽略差别.

定义长度$a_Z=\hbar^2/(\mu g)$, 它使动能尺度$\hbar^2/(\mu a_Z^2)$和势能尺度$g/a_Z$相当. 这个量纲估计说明原子大小的来源, 还没有得到精确能级. 精确量子化须解径向方程和两端边界.

\originalsection{径向级数与能量}
$E=0$是无限远分离阈值, 束缚态$E<0$. 远处$u\prime\prime\simeq\kappa^2u$, 所以$\kappa^{-1}$为指数尾衰减长度; 用$\rho=2\kappa r$把这个尺度移除, $\eta$则衡量库仑项强度. 束缚态$E<0$, 设$\kappa=\sqrt{-2\mu E}/\hbar$, $\rho=2\kappa r$, $\eta=\mu g/(\hbar^2\kappa)$. 式\eqref{eq:radial}变为
\[
\frac{\mathrm d^2u}{\mathrm d\rho^2}
+\left[-\frac14+\frac\eta\rho-\frac{\ell(\ell+1)}{\rho^2}\right]u=0.
\]
在零附近, 标准三维Coulomb实现取$u(0)=0$, 代$u\sim\rho^s$到最奇异两项给$s(s-1)=\ell(\ell+1)$, 两根$\ell+1,-\ell$, 正则解为$\rho^{\ell+1}$; $\ell=0$时另一支常数虽然半轴平方可积, 仍被此边界排除. 在无穷处, 可积解含$\e^{-\rho/2}$. 因而写$u=\rho^{\ell+1}\e^{-\rho/2}F(\rho)$. 令$g_0=\rho^{\ell+1}\e^{-\rho/2}$, 则$g_0\prime/g_0=(\ell+1)/\rho-1/2$, $u\prime\prime=g_0[F\prime\prime+2(g_0\prime/g_0)F\prime+(g_0\prime\prime/g_0)F]$, 代入并除掉公共因子后,
\[
\rho F''+(2\ell+2-\rho)F'+(\eta-\ell-1)F=0.
\]
设$F=\sum_{s\ge0}c_s\rho^s$, 比较$\rho^s$系数得到
$c_{s+1}=(s+\ell+1-\eta)c_s/[(s+1)(s+2\ell+2)]$.

先看递推的实际用法: $c_1=(\ell+1-\eta)c_0/(2\ell+2)$, $c_2=(\ell+2-\eta)c_1/[2(2\ell+3)]$. 例如$\ell=0,\eta=1$给$F=1$, $\ell=0,\eta=2$给$F\propto1-\rho/2$, $\ell=1,\eta=2$又给$F=1$. 特殊函数增长定理在这里作为明确的引用工具: 该方程原点正则解若不在有限阶终止, 无穷端含$\e^\rho$增长支, 因而乘$\e^{-\rho/2}$后仍不正规. 它通常名为合流超几何函数的大$\rho$渐近性质, 本册不让读者先学复Gamma函数再理解量子化. 避免这一项需要分子在某个$s=n_r$处为零, 即$\eta=n_r+\ell+1=:n$. 因而$n=1,2,\ldots$, $\ell=0,1,\ldots,n-1$, 且
\[
E_n=-\frac{\mu g^2}{2\hbar^2n^2},\qquad
u_{n\ell}\propto r^{\ell+1}\e^{-r/(na_Z)}
L_{n-\ell-1}^{2\ell+1}\!\left(\frac{2r}{na_Z}\right).
\]
$n_r=n-\ell-1$是径向多项式次数, 也为内部径向节点数, $n$标内部能量,$\ell$标角分布,$m$标轴投影. 这里$L$表示关联Laguerre多项式; 多项式形式正是有限递推的结果. 非终止解的增长判据可由上述方程的合流超几何大参数渐近验证, 不是仅凭“大$s$像指数”对任意级数作一般定理. 正则起点和可积终点共同排除其他能量.

\originalsection{基态分布和维里关系}
氢样原子基态为$\psi_{100}=(\pi a_Z^3)^{-1/2}\e^{-r/a_Z}$. 位置密度$|\psi|^2$在$r=0$最大, 但半径落在$[r,r+\dd r]$的概率为$4r^2\e^{-2r/a_Z}\dd r/a_Z^3$. 对其取导, 得最可能的\emph{半径}是$a_Z$. 密度最大位置与半径分布最大值不同, 原因是球壳体积随$r^2$增长.

用$\int_0^\infty r^s\e^{-\lambda r}\dd r=s!/\lambda^{s+1}$, 得$\langle r\rangle=3a_Z/2$, $\langle r^2\rangle=3a_Z^2$, $\langle1/r\rangle=1/a_Z$. 若核半径$R\ll a_Z$, 核内概率约为$4R^3/(3a_Z^3)$, 忽略指数的相对误差为$O(R/a_Z)$.

这也解释光谱: 忽略精细结构,$n=2\to1$能差$13.6(1-1/4)=10.2$ eV, 发出光子$\hbar\omega=10.2$ eV; 还需选择定则判断具体态是否允许. 伸缩生成元使位置与动量反向缩放, 表示压窄原子会提高动能而增强吸引势的平衡. 取伸缩生成元$D=(\boldsymbol r\cdot\boldsymbol p+\boldsymbol p\cdot\boldsymbol r)/2$. 在适当束缚态定义域,$[D,H]=\ii\hbar(2T-\boldsymbol r\cdot\nabla V)$. 定态中任意无显含时算符的期望值不变, 所以$2\langle T\rangle=\langle\boldsymbol r\cdot\nabla V\rangle$. 对$V=-g/r$, 有$\boldsymbol r\cdot\nabla V=g/r=-V$, 因而$\langle T\rangle=-\langle V\rangle/2=-E$. 此关系检验符号和能量尺度, 不证明电子沿一条Bohr圆轨道运动.

\originalchapter{复合系统、纠缠与全同粒子}
\originalsection{张量积和纠缠}
对两枚自旋, 联合测量有$++,+-,-+,--$四种结果; 每种配一个复振幅$c_{ij}$. 独立乘积制备时$c_{ij}=a_i b_j$, 一般联合制备却不必能这样分解. 整理这四个组合所需的联合态空间不是把维数相加, 而是张量积$\mathcal H_A\otimes\mathcal H_B$. 若两侧基为$\ket{i}_A$和$\ket{j}_B$, 联合基为$\ket{i}_A\otimes\ket{j}_B$, 一般归一态$\ket\Psi=\sum_{ij}c_{ij}\ket{i,j}$, $\sum_{ij}|c_{ij}|^2=1$. 两个二态系统有四个基向量, 四个复系数在归一化后描述联合纯态.

只作用在$A$上的算符为$A\otimes I_B$, 它在不同$B$分量中使用同一个$A$, 即$(A\otimes I)\ket{i,j}=(A\ket i)\otimes\ket j$. 例如$S_z\otimes I$在$\ket{+,-}$读$+\hbar/2$, 不改第二标签. 此处下标$A$也作系统名, 算符$A$是一个测量规则, 应由上下文分清. 若不存在$\ket\alpha\otimes\ket\beta$的乘积表达, 纯态称为纠缠态. 例如$\ket\Psi=(\ket{00}+\ket{11})/\sqrt2$. 假定它能写为$(a\ket0+b\ket1)\otimes(c\ket0+d\ket1)$, 则要求$ac,bd$非零而$ad,bc$为零, 彼此矛盾. 因而它不能分解.

在该态中测量$A$的$z$基, 结果0和1概率各半; 给定$A$结果0, $B$结果确定为0. 这说明关联, 并不允许$A$主动选择其随机测量结果来发送信息. 忽略$A$的结果时, $B$仍有相同的一半一半分布.

\originalsection{约化密度算符}
先区别两种同样在$z$上各半的制备: 随机选$z+$或$z-$给矩阵$\rho_{\rm mix}=\left(\begin{smallmatrix}1/2&0\\0&1/2\end{smallmatrix}\right)$; 相干$x+$给$\rho_{x+}=\left(\begin{smallmatrix}1/2&1/2\\1/2&1/2\end{smallmatrix}\right)$. 在$x$测量前者仍各半、后者必为$+$, 非对角项保留了相位干涉信息. 经典概率$p_s$不是可相加的复概率幅. 记$\Tr$为在任意正交归一基上的对角矩阵元之和, 称为迹. 纯态密度算符是$\rho=\ket\psi\bra\psi$, 统计混合则为$\rho=\sum_s p_s\ket{\psi_s}\bra{\psi_s}$, $p_s\ge0$, $\sum_sp_s=1$. $\rho\geq0$指每个矢的$\bra v\rho\ket v\geq0$, 不是要求每个矩阵元非负. 它满足$\rho\ge0$, $\Tr\rho=1$, 期望值写为$\Tr(\rho A)$. 混合态表示对制备结果或其他子系统未保留的信息, 一般不能用单个波函数代替.

不查看系统$B$的结果时, 把所有$B$标签的可能性求和, 仍留下$A$上的矩阵; 这称偏迹, 全迹则只给一个数. 将$\ket\Psi\bra\Psi=\sum_{iji\prime j\prime}c_{ij}c_{i\prime j\prime}^*\ket i\bra{i\prime}\otimes\ket j\bra{j\prime}$展开, 对$B$取迹$\Tr(\ket j\bra{j\prime})=\delta_{jj\prime}$, 就保留相同$B$标签的项. 联合纯态中只观测$A$时, 定义$\rho_A=\Tr_B\rho_{AB}$, 矩阵元为$(\rho_A)_{ii'}=\sum_j c_{ij}c_{i'j}^*$. 这样$\bra\Psi A\otimes I\ket\Psi=\Tr(\rho_AA)$. 对上面的纠缠态, $\rho_A=I/2$, 局部$S_z$均值0, $S_z^2$均值$\hbar^2/4$, 局部$x$也各半. 虽然联合态纯, 单个子系统的约化态混合. 相反,$(\ket0+\ket1)/\sqrt2$的纯态密度矩阵有非零非对角元, 与$I/2$不同; 在$x$基测量即可区别.

\originalsection{交换对称性}
电子是费米子, 光子是玻色子. 本节$x_i$的空间坐标与自旋标签在完整态中一同交换. 全同粒子的内禀属性相同, 用标签1、2不能增加可测的物理差别. 交换算符$P_{12}$把$\psi(x_1,x_2)$变为$\psi(x_2,x_1)$, 满足$P_{12}^2=I$. 标准三维非相对论量子理论中, 玻色子用对称态, 费米子用反对称态. 自旋与统计的联系是更深的相对论量子场论结论; 本册把相应交换规则作为物理输入, 不从$P^2=I$独自证明它.

两个正交单粒子态$a,b$的规范化联合态为
\[
\ket{a,b}_\pm=\frac{\ket a\otimes\ket b\pm\ket b\otimes\ket a}{\sqrt2}.
\]
若$a=b$, 反对称式为零, 无法归一化, 因而两个全同费米子不能占据同一个\emph{完整单粒子态}. 空间波函数相同但自旋不同, 不属于同一完整态, 可以占据同一空间轨道. 对电子的总态, 空间对称态须与反对称自旋单态搭配, 空间反对称态须与对称自旋三重态搭配, 使总态反对称.

\begin{example}
两个无相互作用、无自旋的全同粒子在同一无限深势阱中. 单粒子能量为$E_n=n^2E_1$. 比较两玻色子与两费米子的最低总能量.
\end{example}
\begin{solution}
玻色子可共同占据$n=1$, 所以总能量$2E_1$. 无自旋费米子必须占据两个不同单粒子态, 最低选择$n=1,2$, 反对称归一态为$(u_1(x_1)u_2(x_2)-u_2(x_1)u_1(x_2))/\sqrt2$, 总能量$5E_1$. “无自旋”是本数学模型的限制; 对真实自旋$1/2$电子, 可以用相反自旋占据同一空间基态, 最低能量就不同.
\end{solution}

\originalchapter{变分原理与非简并微扰}
\originalsection{变分上界}
许多势无法解析求解. 近似方法的目标是用已知可解模型得到有条件、可检验的估计. 若自伴$H$下界为离散基态$E_0$, 用完备本征基展开归一试探态$\psi=\sum_nc_nu_n$, 得 $\langle H\rangle=\sum_n|c_n|^2E_n\ge E_0\sum_n|c_n|^2=E_0$. 连续谱贡献也高于谱下界, 用谱测度可得到相同不等式. 因而试探态给\emph{基态能量的上界}; 试探族越小, 最优化结果未必越准确.

压窄波包会增大动量宽度和动能, 放宽会使束缚势能升高; 优化是在这两种代价之间找平衡, 不是任意把函数调得很尖. 取谐振子Hamilton算符和宽度为$a$的Gaussian试探态. 前文已算 $\langle x^2\rangle=a^2$, $\langle p^2\rangle=\hbar^2/(4a^2)$, 所以
$E(a)=\hbar^2/(8ma^2)+m\omega^2a^2/2$. 求导令零给$a^2=\hbar/(2m\omega)$, 最小值$\hbar\omega/2$. 此次恰等于精确基态, 因为试探族包含真正基态; 对一般势并不保证取等.

变分法的参数控制形状, 不是新的物理可观测量. 在实际问题中应先满足边界与归一化, 再优化; 不满足算符定义域的尖锐函数可能导致动能发散, 不能只算势能后报告一个很低的数.

\originalsection{非简并微扰展开}
定常微扰跟踪原有能级在小改变后的能量与态. $\lambda$为无量纲强度或记账参数, $W$为能量算符, 上标$(0),(1),(2)$是各阶系数, 真实能量改变量要乘$\lambda,\lambda^2$, 不是把系数当成已乘强度的位移. 设$H(\lambda)=H_0+\lambda W$, $W$厄米. 假定$E_n^{(0)}$是孤立非简并本征值, 且在所讨论$\lambda$邻域有足够光滑的本征对. 这是应用前提, 并非任意小系数都自动保证级数收敛. 写
\[
E_n=E_n^{(0)}+\lambda E_n^{(1)}+\lambda^2E_n^{(2)}+\cdots,
\quad
\ket n=\ket{n^{(0)}}+\lambda\ket{n^{(1)}}+\cdots.
\]
固定相位和尺度为中间归一化$\braket{n^{(0)}}n=1$, 因而$\braket{n^{(0)}}{n^{(1)}}=0$. 最后若计算概率须再归一化全态.

比较本征方程的一阶系数:
$(H_0-E_n^{(0)})\ket{n^{(1)}}=-(W-E_n^{(1)})\ket{n^{(0)}}$. 左乘$\bra{n^{(0)}}$给$E_n^{(1)}=W_{nn}$. 左乘不同$m$的本征bra, 得
\[
\ket{n^{(1)}}=\sum_{m\ne n}
\frac{W_{mn}}{E_n^{(0)}-E_m^{(0)}}\ket{m^{(0)}}.
\]
这里$W_{mn}=\bra{m^{(0)}}W\ket{n^{(0)}}$. 二阶本征方程完整为$(H_0-E_n^{(0)})\ket{n^{(2)}}=-W\ket{n^{(1)}}+E_n^{(1)}\ket{n^{(1)}}+E_n^{(2)}\ket{n^{(0)}}$. 左乘$\bra{n^{(0)}}$得$E_n^{(2)}=\bra{n^{(0)}}W\ket{n^{(1)}}$, 因中间归一化把另一项消去. 二阶方程左乘$\bra{n^{(0)}}$后, 中间归一化消去一阶能量乘态修正项, 得
\[
E_n^{(2)}=\sum_{m\ne n}\frac{|W_{mn}|^2}{E_n^{(0)}-E_m^{(0)}}.
\]
有限维中求和自动有限; 无限维中需要求和收敛及相应算符条件, 有连续谱时还需积分. 对基态, 所有分母负, 二阶能量修正非正. 这为计算提供符号检查.

\originalsection{扰动大小与一个算例}
重要的单通道混合诊断是$|\lambda W_{mn}/(E_n^{(0)}-E_m^{(0)})|$, 不是仅看$|\lambda|$. 多通道时还应检查$\lambda^2\sum_{m\ne n}|W_{mn}|^2/(E_n^{(0)}-E_m^{(0)})^2\ll1$或相应连续谱积分; 每个小比值不自动保证总态修正小. 当两个能级很近, 很小的矩阵元也会造成大混合, 应使用下一章的简并或近简并处理.

\begin{example}
谐振子加上弱线性势$\lambda f x$, 求基态能量到二阶并与精确解比较.
\end{example}
\begin{solution}
$f$量纲为力, $+fx$的力为$-f$, 基态的$\langle x\rangle=0$, 一阶修正为零. $x=\sqrt{\hbar/(2m\omega)}(a+a^\dagger)$只把$\ket0$连到$\ket1$. 因而二阶系数$E_0^{(2)}$为
$|f\langle1|x|0\rangle|^2/(E_0^{(0)}-E_1^{(0)})=-f^2/(2m\omega^2)$.
真实二阶能量改变量为$\lambda^2E_0^{(2)}$. 配方可得
\[
\frac12m\omega^2x^2+\lambda f x
=\frac12m\omega^2\left(x+\frac{\lambda f}{m\omega^2}\right)^2
-\frac{\lambda^2f^2}{2m\omega^2}.
\]
所以精确基态能量为$\hbar\omega/2-\lambda^2f^2/(2m\omega^2)$, 与微扰一致. 这里二阶以后能量修正为零, 因为势的变化只平移中心并加常数, 不是一般扰动都如此.
\end{solution}

\originalsection{从近似系数理解能级排斥}
扰动对角元$W_{nn}$是原态对新能量项的直接平均, 非对角元则把原来不同态连接起来. 为观察反馈, 取$H=\left(\begin{smallmatrix}E_a&\lambda w\\\lambda w^*&E_b\end{smallmatrix}\right)$,$E_b>E_a$. 精确能量$(E_a+E_b)/2\pm\sqrt{(E_b-E_a)^2/4+\lambda^2|w|^2}$, 小$|\lambda w|/(E_b-E_a)$下低能下降$-\lambda^2|w|^2/(E_b-E_a)$、高能升高同量. 态的混合幅约$\lambda w/(E_b-E_a)$, 能隙很小则这个展开失效. 基态二阶下降与变分上界并不矛盾: 微扰比较的是原$H_0$能量, 变分比较的是同一个新$H$的精确基态与试探态.

对真正非谐$V=m\omega^2x^2/2+g x^4$,$g>0$, 用第2章Gaussian宽度$a$, $\langle x^4\rangle=3a^4$, 上界$E(a)=\hbar^2/(8ma^2)+m\omega^2a^2/2+3ga^4$. 设$y=a^2$, 最佳参数由$-\hbar^2/(8my^2)+m\omega^2/2+6gy=0$确定, 左边严格递增所以唯一正根. $g$增加会缩窄最佳态, 但Gaussian不含精确四次势基态, 上界不等于精确答案. 上界保证能量方向, 不自动保证波函数误差或激发能误差同样小.

\originalchapter{简并微扰与能级分裂}
\originalsection{简并子空间中的对角化}
若$H_0$有多维本征子空间$\mathcal D$, 非简并公式会出现零分母. 在一阶方程上用完整$E_0$本征空间投影$P$投到$\mathcal D$, 左边$P(H_0-E_0)=0$, 因而
$PWP\ket{\chi}=E^{(1)}\ket\chi$. 零阶态不应事先任取, 而要选$W$在简并子空间内的本征向量. 首先对角化一个较小的厄米矩阵, 才能决定一阶分裂与正确起始态.

\begin{example}
$H_0=\operatorname{diag}(E_0,E_0,E_3)$, $E_3\ne E_0$. 微扰在前两维为 $\begin{pmatrix}a&b\\b^*&a\end{pmatrix}$, $b=|b|\e^{\ii\varphi}\ne0$. 求一阶能量和零阶态.
\end{example}
\begin{solution}
特征方程是$(a-w)^2-|b|^2=0$, 所以$w_\pm=a\pm|b|$. 对应归一向量可取 $(\e^{\ii\varphi}\ket1\pm\ket2)/\sqrt2$, 能量为$E_0+\lambda(a\pm|b|)+O(\lambda^2)$. 与第3态的耦合先在子空间外修正中出现, 由非零能隙$E_0-E_3$控制. 当$E_3$也接近$E_0$, 应把它一并纳入近简并空间.
\end{solution}

若$PWP$仍有简并, 一阶没有完全区分态. 可以进一步研究二阶有效算符$PWQ(E_0-QH_0Q)^{-1}QWP$, $Q=I-P$. 此式要求补空间上对应逆算符存在, 不可在仍有零能隙时直接取逆. 一阶未分裂不等于扰动完全没有作用.

\originalsection{自旋轨道耦合}
在固定$n,\ell$的氢样子空间, 考虑$W=\xi(r)\boldsymbol L\cdot\boldsymbol S$. 若径向矩阵元已知, 用总角动量$\boldsymbol J=\boldsymbol L+\boldsymbol S$对角化角向部分最自然. $L^2,S^2,J^2,J_z$相容, 而$L_z,S_z$通常分别不守恒.

利用 $2\boldsymbol L\cdot\boldsymbol S=J^2-L^2-S^2$, 自旋$1/2$时, $j=\ell+1/2$的角向本征值为$\hbar^2\ell/2$, $j=\ell-1/2$为$-\hbar^2(\ell+1)/2$. 后者只在$\ell\ge1$存在. 所以一阶能量移为对应角向值乘$\langle\xi(r)\rangle_{n\ell}$. 这清楚分开了径向积分与角动量代数; 尚未求径向积分时, 不能把系数写成一个虚构数值.

氢精细结构还包含相对论动能修正与Darwin项, 并非仅有自旋轨道耦合. 完整结果的推导需这些有效Hamilton算符及其适用阶数; 本册此处完成角向分裂方法, 完整精细结构计算留作来源中的8.06第2章对应专题, 覆盖表标为尚未完整重编. 这样不会把“一个分裂项已算”误记成“整套精细结构已推导”.

\originalsection{弱场Zeeman效应}
弱外磁场相对于精细结构能隙很小, 仍可用$\ket{n,\ell,j,m_j}$为零阶态. 残余简并中,$H_Z$与$L^2,J_z$对易, 所以不同$\ell,m_j$不混合; 相同$\ell,m_j$的其他$j$分量由非零精细结构能隙分开. 忽略核磁矩, 电子磁扰动为$H_Z=\mu_B B(L_z+2S_z)/\hbar$, $\mu_B=e\hbar/(2m_e)>0$, 磁场沿$z$轴. 正号来自电子磁矩为负且能量为$-\boldsymbol\mu\cdot\boldsymbol B$.

在固定$j>0$的单个不可约角动量空间, 向量算符满足$[J_i,V_j]=\ii\hbar\epsilon_{ijk}V_k$. 标准投影性质给$P_j\boldsymbol S P_j=c_j\boldsymbol J$, 其中$c_j$与$m_j$无关. 两侧点乘$\boldsymbol J$并利用固定$J^2$, 就由$\boldsymbol S\cdot\boldsymbol J$确定$c_j$. 因此向量算符$\boldsymbol S$的对角期望沿$\boldsymbol J$的方向, 比例由$\boldsymbol S\cdot\boldsymbol J$决定. 取$2\boldsymbol S\cdot\boldsymbol J=J^2+S^2-L^2$, 得
\[
\langle S_z\rangle=\hbar m_j
\frac{j(j+1)+s(s+1)-\ell(\ell+1)}{2j(j+1)}.
\]
其比例关系是固定$j$空间中的向量投影性质, 可由角动量升降矩阵元验证; 这里引用此明确性质, 不把它视为对任意向量场成立. 因而$\Delta E_Z=\mu_BBg_jm_j$, 其中
$g_j=1+[j(j+1)+s(s+1)-\ell(\ell+1)]/[2j(j+1)]$, $s=1/2$. 若磁场超过精细结构尺度, 固定$j$近似失效, 需换用更适合的基并重新对角化. 使用公式必须先比较能隙.

\Needspace{9\baselineskip}
\originalchapter{含时微扰与绝热演化}
\originalsection{相互作用表象与跃迁幅}
设$H(t)=H_0+\lambda W(t)$, $H_0\ket n=E_n\ket n$. 写$\ket{\psi(t)}=\sum_n c_n(t)\e^{-\ii E_nt/\hbar}\ket n$. 代入薛定谔方程, $H_0$项与指数导数抵消, 得
\[
\ii\hbar\dot c_m(t)=\lambda\sum_nW_{mn}(t)
\e^{\ii(E_m-E_n)t/\hbar}c_n(t).
\]
若初态为$\ket i$, 零阶取$c_n=\delta_{ni}$. 代入右侧并从0到$t$积分, 得一阶跃迁幅
\[
c_f^{(1)}(t)=-\frac{\ii\lambda}{\hbar}\int_0^t
W_{fi}(s)\e^{\ii\omega_{fi}s}\dd s,
\qquad \omega_{fi}=(E_f-E_i)/\hbar.
\]
此时$f\ne i$的概率最低阶为$|c_f^{(1)}|^2$, 是$\lambda^2$阶, 不是$\lambda$阶. 只保留跃迁概率却不给初态概率作相应二阶修正, 总概率会轻微偏离1; 这是截断误差, 不是精确动力学破坏幺正性.

\originalsection{共振与黄金规则的限制}
若一段时间内 $W_{fi}(s)=w\e^{-\ii\omega s}$, 定义失谐$\delta=\omega_{fi}-\omega$, 积分给
\[
P_{i\to f}(t)\simeq\frac{4\lambda^2|w|^2}{\hbar^2}
\frac{\sin^2(\delta t/2)}{\delta^2}.
\]
$\delta=0$时取连续极限$P\simeq\lambda^2|w|^2t^2/\hbar^2$. 它在长时间会大于1, 说明一阶幅近似已超出适用范围; 不能把它当精确的永久增长规律. 对离散二态共振, 完整解产生有界的Rabi振荡.

若末态形成稠密连续谱, 记$\operatorname{sinc}z=\sin z/z$（零处连续取1）. 状态密度和矩阵元在共振窄窗内变化缓慢, 积分中的$sinc^2$核在长于相关时间且短于显著耗尽时间的窗口集中在能量守恒点. 分布意义下
$4\sin^2(\delta t/2)/\delta^2\to2\pi t\delta_{\rm D}(\delta)$.
把末态求和换为能量积分, 得Fermi黄金规则, 单位时间的跃迁概率即跃迁率为 $\Gamma=(2\pi/\hbar)\lambda^2|W_{fi}|^2\rho_E(E_i+\hbar\omega)$. 此处$W_{fi}$指驱动的对应Fourier分量$w$, 多个末态通道要分别加权求和或积分. $\rho_E$按单位能量计数, 不能与按角频率的密度混用而漏掉$\hbar$. 一个孤立的有限二态系统不满足连续谱条件, 不能直接用常数跃迁率描述全部时间.

\originalsection{绝热条件和Berry相位}
令$H(t)\ket{n(t)}=E_n(t)\ket{n(t)}$, 假定目标本征空间是一维, 与其余能级在整个路径上有非零间隙, 本征态及Hamilton算符足够平滑. 对$m\ne n$, 对本征方程求导并左乘$\bra m$, 得
$\braket m{\dot n}=\bra m\dot H\ket n/(E_n-E_m)$. 展开态后, 非对角项因快速的能隙相位振荡而抵消. 一个常用的局部小量是
\[
\frac{\hbar|\bra m\dot H\ket n|}{|E_n-E_m|^2}\ll1.
\]
它是有用的诊断, 不是不加全程光滑性、间隙和时间积累条件就成立的充分定理. 特别在避免交叉的最小能隙处, 要检查最差情况.

忽略跃迁后, 从$n$态出发的近似解为
\[
\ket{\psi(t)}\simeq
\exp\left[-\frac{\ii}{\hbar}\int_0^tE_n(s)\dd s\right]
\e^{\ii\gamma_n(t)}\ket{n(t)},
\quad \gamma_n(t)=\ii\int_0^t\braket{n(s)}{\dot n(s)}\dd s.
\]
归一化导数保证$\braket n{\dot n}$纯虚, 所以$\gamma_n$实. 第一项是动力学相位, 第二项是几何相位. 变换$\ket n\mapsto\e^{\ii\chi(t)}\ket n$会把$\gamma_n$改变$-\chi(t)+\chi(0)$, 恰好抵消本征态相位变化. 对闭合参数回路, 若采用周期规范$\ket{n(T)}=\ket{n(0)}$, Berry相位就是$\gamma_n(T)$模$2\pi$. 对任意连续规范, 规范不变的闭环量应写为 $\gamma_{\rm loop}=\gamma_n(T)+\arg\braket{n(0)}{n(T)}$模$2\pi$, 两项的规范变化恰好抵消. 这个闭环量可以通过干涉读取, 不能把任意开路径下单独的$\gamma_n$当作规范不变观测量.

\originalchapter{半经典近似与三维散射}
\originalsection{WKB的局部构造}
定态一维方程为$-\hbar^2u''/(2m)+Vu=Eu$. 写$u=\exp[\ii S(x)/\hbar]$, 代入得到$(S')^2-\ii\hbar S''=2m(E-V)$. 在允许区令$p(x)=\sqrt{2m(E-V)}>0$, 展开$S=S_0+\hbar S_1+\cdots$. 零阶给$S_0'=\pm p$, 一阶给$2S_0'S_1'=\ii S_0''$, 所以$S_1=(\ii/2)\log p+\text{常数}$. 于是
\[
u_\pm(x)\simeq\frac C{\sqrt{p(x)}}
\exp\left[\pm\frac{\ii}{\hbar}\int^x p(s)\dd s\right].
\]
$1/\sqrt p$不是装饰性的归一化, 它使单向流近似为常数. 近似要求$\hbar|p'|/p^2\ll1$, 即动量在一段局部波长内变化小. 在$E=V$的转折点,$p=0$, 该式发散, 必须另作匹配.

在禁止区令$\kappa(x)=\sqrt{2m(V-E)}/\hbar$, 两个局部解为$\kappa^{-1/2}\exp(\pm\int\kappa\dd x)$. 对一个平滑简单转折点, 将$V-E$线性化后, 方程化为Airy方程. 其衰减解的允许区渐近包含$\pi/4$相位. 两个简单转折点的束缚态匹配给
$\int_{x_1}^{x_2}p(x)\dd x=\pi\hbar(n+1/2)$, $n=0,1,\ldots$.
这里明确使用Airy函数的标准渐近连接公式; 本册完成局部WKB展开, 没有在此证明Airy渐近理论. 硬墙或高阶转折点不能继续使用同一$1/2$修正. 对右转折点$x_2$且右侧禁止, 标准Airy渐近给连接
\[
\frac{C}{\sqrt{|p(x)|}}\exp\!\left[-\frac1\hbar\int_{x_2}^{x}|p(s)|\dd s\right]
\longleftrightarrow
\frac{2C}{\sqrt{p(x)}}\sin\!\left[\frac1\hbar\int_x^{x_2}p(s)\dd s+\frac\pi4\right].
\]
左端同样连接. 两个允许区正弦要表示同一解, 两端$\pi/4$合成$\pi/2$, 给上式$n+1/2$. 此处$|p|=\sqrt{2m(V-E)}$, 与禁止区的$\hbar\kappa$相同.

对谐振子, 转折点$x_\pm=\pm\sqrt{2E/(m\omega^2)}$, 积分是椭圆面积的一半,$\int_{x_-}^{x_+}\sqrt{2mE-m^2\omega^2x^2}\dd x=\pi E/\omega$. 因而WKB给$E=\hbar\omega(n+1/2)$, 恰为精确能级. 这次取等来自二次势的特殊性, 不意味着WKB对任意最低能级都精确.

\originalsection{散射振幅和截面}
短程势的三维弹性散射远处可写为$u(\boldsymbol r)\simeq\e^{\ii kz}+f(\theta,\phi)\e^{\ii kr}/r$, 能量$E=\hbar^2k^2/(2\mu)$. 入射波流密度是$\hbar k/\mu$, 散射球面波的径向流最低阶是$(\hbar k/\mu)|f|^2/r^2$. 在固体角$\dd\Omega$上乘球面面积$r^2\dd\Omega$, 再除入射流密度, 得
$\mathrm d\sigma/\mathrm d\Omega=|f|^2$. 截面量纲为面积, 振幅$f$量纲为长度. 长程库仑势的渐近相位不同, 不能不加修正沿用此短程边界形式.

对中心势, 不同$\ell$通道不混合. $P_\ell$表示勒让德多项式, 径向相移约定为$u_\ell(r)\sim\sin(kr-\ell\pi/2+\delta_\ell)$, 与自由正则波比较. 径向远区相对于自由解多一个相移$\delta_\ell$, 弹性通道幺正性给$S_\ell=\e^{2\ii\delta_\ell}$. 用平面波的球谐展开, 可得
\[
f(\theta)=\frac1k\sum_{\ell=0}^\infty
(2\ell+1)\e^{\ii\delta_\ell}\sin\delta_\ell\,P_\ell(\cos\theta),
\quad
\sigma=\frac{4\pi}{k^2}\sum_\ell(2\ell+1)\sin^2\delta_\ell.
\]
总截面由勒让德多项式正交性积分得到: $\int_{-1}^1P_\ell P_{\ell'}\dd z=2\delta_{\ell\ell'}/(2\ell+1)$. 此处调用平面波的标准球谐展开, 详细相移求解与共振分析仍需对应径向边界问题.

\originalsection{Born近似}
自由外向Green函数为$G(\boldsymbol r-\boldsymbol r')=-\e^{\ii k|\boldsymbol r-\boldsymbol r'|}/(4\pi|\boldsymbol r-\boldsymbol r'|)$, 满足$(\nabla^2+k^2)G=\delta^{(3)}$. 这个负号可由在原点附近积分拉普拉斯检验. 将方程写成$(\nabla^2+k^2)u=(2\mu/\hbar^2)Vu$, 得积分方程
\[
u(\boldsymbol r)=\e^{\ii\boldsymbol k\cdot\boldsymbol r}
-\frac{\mu}{2\pi\hbar^2}\int
\frac{\e^{\ii k|\boldsymbol r-\boldsymbol r'|}}{|\boldsymbol r-\boldsymbol r'|}
V(\boldsymbol r')u(\boldsymbol r')\dd^3r'.
\]
第一Born近似在积分内部用入射波代替完整$u$. 定义末态动量方向$\boldsymbol k'=k\widehat{\boldsymbol r}$. 远区$|\boldsymbol r-\boldsymbol r'|\simeq r-\widehat{\boldsymbol r}\cdot\boldsymbol r'$, 因而
\[
f_{\rm B}(\boldsymbol k',\boldsymbol k)
=-\frac{\mu}{2\pi\hbar^2}\int
\e^{\ii(\boldsymbol k-\boldsymbol k')\cdot\boldsymbol r'}V(\boldsymbol r')\dd^3r',
\quad |\boldsymbol k'|=|\boldsymbol k|.
\]
所以散射测量与势能的Fourier变换相关. 这也说明短程、弱势或合适高能条件为什么有助于近似. 是否足够弱需比较下一次迭代或实际误差; 不能只看到“势很小”就保证在阈值和共振附近有效.


\begin{example}
短程Gaussian势$V(r)=V_*\e^{-r^2/a^2}$, $a>0$. 用Born近似求微分截面.
\end{example}
\begin{solution}
三维Fourier积分可分成三个一维Gaussian, 得$\int\e^{\ii\boldsymbol Q\cdot\boldsymbol r}\e^{-r^2/a^2}\dd^3r=\pi^{3/2}a^3\e^{-a^2Q^2/4}$. 散射角$\theta$对应$Q=|\boldsymbol k-\boldsymbol k'|=2k\sin(\theta/2)$. 所以
$f_B=-\mu\sqrt\pi V_*a^3\e^{-a^2k^2\sin^2(\theta/2)}/(2\hbar^2)$,
$\mathrm d\sigma/\mathrm d\Omega=\mu^2\pi V_*^2a^6\e^{-2a^2k^2\sin^2(\theta/2)}/(4\hbar^4)$.
$\mu V_*a^2/\hbar^2$是本模型的无量纲势强诊断, 但具体误差还依赖$ka$及高阶迭代.
\end{solution}

\originalchapter{带电粒子与氢原子的量子对称性}
本章为补足已选进阶内容的独立推导. 8.06 Spring 2018没有公开第5章带电粒子PDF, 所以这里没有冒称翻译该缺失原件. 公式按前文算符、概率流和角动量工具组织, 带电粒子例子与因子分解方法也为考核篇提供所需支持.
\originalsection{电磁势、机械动量与规范}
在电场中粒子可受力、在磁场中方向可偏转. 用标量势$\phi$和矢势$\boldsymbol A$表示$\boldsymbol E=-\nabla\phi-\partial_t\boldsymbol A$,$\boldsymbol B=\nabla\times\boldsymbol A$. 对质量$m$、电荷$q$的无自旋粒子, 最小耦合Hamiltonian为$H=(\boldsymbol p-q\boldsymbol A)^2/(2m)+q\phi$, $\boldsymbol p=-\ii\hbar\nabla$为正则动量, $\boldsymbol\pi=\boldsymbol p-q\boldsymbol A$为机械动量. 平方有算符顺序, 要按$(p_i-qA_i)(p_i-qA_i)$作用, 不能先把$p_i,A_i$交换. $q\boldsymbol A$量纲为动量.

同一电磁场可以用不同势: $\boldsymbol A' =\boldsymbol A+\nabla\chi$, $\phi'=\phi-\partial_t\chi$. 波函数同时变$\psi'=\e^{\ii q\chi/\hbar}\psi$, 直接求导得$\boldsymbol\pi'\psi'=\e^{\ii q\chi/\hbar}\boldsymbol\pi\psi$, 时间导数中额外$-q\dot\chi$与标势变化抵消, 所以方程协变. 此处$\chi$量纲作用量/电荷. 密度不变, 流为$\boldsymbol J=m^{-1}\operatorname{Re}(\psi^*\boldsymbol\pi\psi)$, 也保持不变. 规范选择改变描述中的局部相位, 不改变电磁场或计数统计.

机械动量分量一般不对易: $[\pi_i,\pi_j]=\ii q\hbar\epsilon_{ijk}B_k$, 由$[p_i,A_j]=-\ii\hbar\partial_iA_j$求得. 所以不能在磁场中继续将三个机械动量当成自由粒子相容读数. 用Heisenberg方程对位置求导给速度$\boldsymbol\pi/m$, 可与经典Lorentz力联系; 非均匀场时算符乘积要对称排列, 均值也不总能以经典点轨迹代替.
\originalsection{均匀磁场的Landau能级}
取$B\hat z$, 无电势, 选$\boldsymbol A=(0,Bx,0)$. $p_y,p_z$与$H$对易, 对其广义本征值$\hbar k_y,p_z$将波函数写成$\e^{\ii k_yy+\ii p_zz/\hbar}u(x)$. 剩下
\[
H_x=\frac{p_x^2}{2m}+\frac{(qB)^2}{2m}(x-x_c)^2,\quad
x_c=\frac{\hbar k_y}{qB},\quad\omega_c=\frac{|qB|}{m},
\]
故$E_{n,p_z}=\hbar\omega_c(n+1/2)+p_z^2/(2m)$,$n\geq0$. 新能量不依赖中心$x_c$, 同一能级有很多轨道中心制备; 这不是漏掉变量而是磁平移对称导致简并. 波函数宽度尺度$\ell_B=\sqrt{\hbar/|qB|}$, 能级间隔$\hbar\omega_c$对应经典回旋频率, 但最低能仍有零点能量.

若矩形面积$L_xL_y$, $y$周期使$k_y=2\pi j/L_y$, 允许中心在样品内, 每个Landau能级约有$|qB|L_xL_y/(2\pi\hbar)$个态, 每面积态数$|qB|/(2\pi\hbar)$. 边界附近能量可变化, 该数忽略边缘有限尺寸修正. 若为电子还需自旋项$+\mu_BB\sigma_z$（$g\simeq2$, 不计辐射异常磁矩）, 与轨道能量相加; 无自旋结果不是完整电子磁谱.
\originalsection{无局部磁场区域的Aharonov--Bohm相位}
细长磁通管外部$\boldsymbol B=0$, 但绕管闭合路径$\oint\boldsymbol A\cdot\dd\boldsymbol l=\Phi$. 两条绕行路径的振幅相对相位为$q\Phi/\hbar$, 可改变干涉计数. 在单连通的小区域可用规范移除$\boldsymbol A$, 围住被排除磁通管的全区域却未必可用单值规范同时移除; 边界拓扑参与条件, 不代表粒子在局部无场区域受经典磁力.

半径$b$环上的粒子是可算模型. 取$A_\varphi=\Phi/(2\pi b)$, 单值边界$\psi(\varphi+2\pi)=\psi(\varphi)$给$\psi_n=\e^{\ii n\varphi}/\sqrt{2\pi}$,$n\in\mathbb Z$（按角度归一）. 能量
\[
E_n(\Phi)=\frac{\hbar^2}{2mb^2}\left(n-\frac{q\Phi}{h}\right)^2.
\]
磁通改变机械角动量, 能量谱作为集合的周期为$h/|q|$, 因磁通加一个量子只重标$n$. 固定$n$曲线不是周期函数, 需区分单支与全谱. 若把矢势变零, 波函数边界变扭转$\psi'(2\pi)=\e^{-\ii q\Phi/\hbar}\psi'(0)$; 同时偷偷保留原周期边界会改变物理问题.
\originalsection{径向因子分解与氢谱}
第11章由边界和级数求氢谱; 此处给另一条不依赖非终止级数渐近的代数路线. 本节取固定核$m=m_e$,$a_Z=\hbar^2/(mg)$; 若计核运动则同时将$m$改为约化质量$\mu$, 不混用两套长度和质量. 对固定$\ell$的半轴Coulomb径向算符, 取实超势（只是因子分解的函数名, 不是另一个物理势）$W_\ell=-(\ell+1)/r+1/[(\ell+1)a_Z]$, $A_\ell=\mathrm d/\mathrm dr+W_\ell$, $A_\ell^\dagger=-\mathrm d/\mathrm dr+W_\ell$, 使用正规径向边界使伴随积分项消失. 展开得
\[
H_\ell=\frac{\hbar^2}{2m}A_\ell^\dagger A_\ell-\frac{\hbar^2}{2ma_Z^2(\ell+1)^2}.
\]
正性给基态下界, $A_\ell u=0$解$u\propto r^{\ell+1}\e^{-r/[(\ell+1)a_Z]}$确实正规且取到下界. 伙伴$A_\ell A_\ell^\dagger$将离心项变为$(\ell+1)(\ell+2)/r^2$, 对应$H_{\ell+1}$加已知常数. 用$H^{(1)}A^\dagger=A^\dagger H^{(2)}$逐次从更大$\ell$的零态提升, 得$E_{n_r,\ell}=-\hbar^2/[2ma_Z^2(n_r+\ell+1)^2]$, 与第11章一致. 因子分解不是每个势都具有的结构; 一旦零态不正规或伙伴不属于同一边界域, 形式阶梯就不能直接报告物理谱.
\originalsection{氢的额外简并与Runge--Lenz量}
一般中心势只有$m$的$2\ell+1$简并, 氢同$n$还包含不同$\ell$, 不由普通空间旋转单独解释. 取固定核$H=p^2/(2m)-g/r$, 定义无量纲$\boldsymbol R=(\boldsymbol p\times\boldsymbol L-\ii\hbar\boldsymbol p)/(mg)-\boldsymbol r/r$. 等价用$(\boldsymbol p\times\boldsymbol L-\boldsymbol L\times\boldsymbol p)/2$使Hermitian性明确. 在合适共同域的Coulomb算符计算给$[H,R_i]=0$, $[L_i,R_j]=\ii\hbar\epsilon_{ijk}R_k$, $[R_i,R_j]=-2\ii\hbar H\epsilon_{ijk}L_k/(mg^2)$以及
\[
\boldsymbol R\cdot\boldsymbol L=0,\qquad R^2=1+\frac{2H}{mg^2}(L^2+\hbar^2).
\]
这是量子修正后的算符恒等式, 可用$[p_i,r_j]$及径向导数逐项检验; $\hbar^2$项不在经典公式中. 垂直性的一段具体计算见2012期末第3题. 这些多项对易恒等式作为本节引用的Coulomb代数工具, 后面的谱推理从它们明确展开, 不冒称已经逐项证明所有长代数式.

固定$E<0$能量空间, 令$\boldsymbol A=\sqrt{-mg^2/(2E)}\boldsymbol R$, $\boldsymbol J_\pm=(\boldsymbol L\pm\boldsymbol A)/2$. 代上式给两套互相对易的角动量代数, $J_+^2=J_-^2=[-mg^2/(2E)-\hbar^2]/4$. 两套量子数相同为$j$, 所以$-mg^2/(2E)=\hbar^2(2j+1)^2$. 记$n=2j+1$, 得氢能量与忽略自旋时简并$n^2$. 不同$\ell=0,\ldots,n-1$维数之和$\sum(2\ell+1)=n^2$, 正是这套隐藏对称性的同一空间. 精细结构或外场改变$H$, 不再自动保持$\boldsymbol R$守恒, 因而能够打破这些简并.

\originalchapter{8.04 作业1及独立解答}
\label{ch:ps1}
本章题面按 \href{https://ocw.mit.edu/courses/8-04-quantum-physics-i-spring-2016/resources/mit8_04s16_ps1_2016/}{8.04 Spring 2016 Problem Set 1} 中文转写, 发布于2016年2月4日, 截止2016年2月11日17时. 保留7题和全部子问. 原件由MIT课程组提供; 未取得官方解答. 以下均为\textbf{AI独立解答}. 第1、2题使用原题Gaussian单位, $e^2$已经包含SI中相应库仑常数; 不直接将SI电荷数值代入这些式子.

\begin{exercise}
经典氢原子中, 电子绕固定质子作近圆周运动. 原题给$V=-e^2/r$, 库仑力大小$e^2/r^2$, 及Larmor辐射功率$\mathrm dE/\mathrm dt=-2e^2a^2/(3c^3)$.
(a) 求每周损失能量$\Delta E$与动能$K$之比, 说明何时可把轨道视为逐渐缩小的圆.
(b) 在半径50 pm、1 pm、1 fm时比较$v$与$c$.
(c) 从50 pm缩到1 pm需多久? 非相对论近似是否合理? 形式上把终点改为1 fm有多大改变?
(d) $r\to0$时能量如何变化? 有无下界? （10分）
\end{exercise}
\begin{solution}
圆周平衡给$mv^2/r=e^2/r^2$, 所以$v^2=e^2/(mr)$、$K=e^2/(2r)$、$E=-e^2/(2r)$. 一周时间为$2\pi r/v$, 加速度为$v^2/r$. 辐射功率乘周期后,
$\Delta E=4\pi e^2v^3/(3c^3r)$, 故$\Delta E/K=(8\pi/3)(v/c)^3$. 当$v\ll c$, 每周能量损失很小, 近圆轨道慢变是自洽的局部近似.

令$r_e=e^2/(mc^2)=2.818\,\mathrm{fm}$, 则$v/c=\sqrt{r_e/r}$. 三个半径给约$0.00751$、$0.0531$和$1.68$. 最后一个超过1, 是非相对论模型失效的信号, 不能解释为电子实际超光速.

由$\dot E=(e^2/2r^2)\dot r$和辐射式, $\dot r=-4e^4/(3m^2c^3r^2)$. 积分得
\[
t=\frac{m^2c^3}{4e^4}(r_i^3-r_f^3)
=\frac{r_i^3-r_f^3}{4cr_e^2}.
\]
取50 pm和1 pm得$t\simeq1.31\times10^{-11}\,\mathrm s$. 到1 pm速度仍约$c$的5\%, 相对论修正小. 把终点形式上改为1 fm, 时间仅增加约原值的$8\times10^{-6}$, 因为积分受起始大半径主导. 但很小半径处的辐射与轨道模型已不自洽, 这个形式估值不能验证最后阶段的真实运动. 经典$E=-e^2/(2r)$无下界, $r\to0$趋$-\infty$, 正是稳定性失败.
\end{solution}

\begin{exercise}
电子绕固定重质子作圆周运动, $V=-e^2/r$.
(a) 给$K,V,E$间关系.
(b) 若角动量大小量子化为$L=n\hbar$, $n$为正整数, 求$E_n,r_n$, 用电子静能$E_e=mc^2$、约化Compton长度$\bar\lambda=\hbar/(mc)$、$\alpha=e^2/(\hbar c)$表达. （5分）
\end{exercise}
\begin{solution}
上题的圆周条件给$K=-V/2$, $E=V/2=-K$. 又$L^2=m^2v^2r^2=me^2r$, 故$r_n=n^2\hbar^2/(me^2)=n^2\bar\lambda/\alpha$, $E_n=-mc^2\alpha^2/(2n^2)=-E_e\alpha^2/(2n^2)$. 这是Bohr旧量子模型, 能级与非相对论氢谱同形, 但把真实电子态等同确定圆轨道不是现代量子力学的结论.
\end{solution}

\begin{exercise}
取$h=6.6\times10^{-34}\,\mathrm{J\,s}$, $k_B=1.38\times10^{-23}\,\mathrm{J/K}$.
(a) 求下列de Broglie波长并解释日常为何不易直接看到它: 2000 kg汽车以22 m/s运动; 10 g弹珠以10 cm/s运动; 直径100 nm、质量1 fg烟粒在300 K具有平均平动动能$3k_BT/2$; $^{87}$Rb原子冷却到100 $\mu$K且同样取$3k_BT/2$.
(b) 光子$E=h\nu$: (i) 求400--700 nm可见光的频率和能量范围; (ii) 求2.5 GHz、300 W微波炉, 633 nm、10 mW激光, 850 MHz、0.25 W手机每秒发出光子数; (iii) 200 ml水升温10 K需多少该微波光子, 水密度1 g/ml、比热4 J/(g K)? (iv) 同功率下无线电与X射线何者更适合经典波近似? （10分）
\end{exercise}
\begin{solution}
非相对论$p=mv$, 热动能条件给$p=\sqrt{3mk_BT}$; 1 fg为$10^{-18}$ kg. 四种波长依次约为$1.50\times10^{-38}$ m、$6.6\times10^{-31}$ m、$5.92\times10^{-15}$ m和$2.70\times10^{-8}$ m. 最后一项使用$^{87}$Rb质量约$87\times1.66\times10^{-27}$ kg. 可见烟粒的几何直径与其中心质点波长是不同概念. 宏观物体波长很小, 又难保持不同路径相干, 日常难以分辨其量子干涉.

取$c=3.0\times10^8$ m/s. 可见光频率约$4.29\times10^{14}$到$7.50\times10^{14}$ Hz, 能量$2.83\times10^{-19}$到$4.95\times10^{-19}$ J, 约1.77--3.09 eV. 低波长端对应高频高能. 光子率为$\dot N=P/(h\nu)$, 三装置分别约$1.82\times10^{26}$、$3.20\times10^{16}$、$4.46\times10^{23}\,\mathrm{s^{-1}}$. 激光频率须先用$c/\lambda$计算.

水的吸收能量为$200\times4\times10=8000$ J, 每个微波光子$1.65\times10^{-24}$ J, 所以需约$4.85\times10^{27}$个, 假设这些能量全部被水吸收. 同功率和相同观察时间下, 低频光每光子能量更低、光子数更多, 相对的离散光子涨落较小, 通常更适合经典波统计近似. 实际经典性还依赖制备和相干性, 不能仅由频率对所有量子光态下断言.
\end{solution}

\begin{exercise}
写$z=a+\ii b=r\e^{\ii\theta}$.
(a) 用Taylor级数证明Euler公式.
(b) 写笛卡尔与极坐标参数的转换.
(c) 证明乘$\ii$等于转$90^\circ$, 给$\ii z$的实部, 并证明乘$\e^{\ii\phi}$等于转$\phi$.
(d) 对$z^*=a-\ii b$: (i) 是否有数同时为实数和纯虚数? (ii) 求$(z^*)^*$及$z^*$的极形式; (iii) 用$z,z^*$表达实虚部; (iv) 求$zz^*$的多种表达.
(e) 推出二倍角、三倍角及和角公式. （15分）
\end{exercise}
\begin{solution}
绝对收敛指数级数中偶次项为$\sum_j(-1)^j\theta^{2j}/(2j)!=\cos\theta$, 奇次项为$\ii\sum_j(-1)^j\theta^{2j+1}/(2j+1)!=\ii\sin\theta$, 故Euler公式成立. $a=r\cos\theta$, $b=r\sin\theta$, $r=\sqrt{a^2+b^2}$, 当$r>0$时$\theta=\operatorname{atan2}(b,a)\pmod{2\pi}$, 即按象限取角; $r=0$时角度任意.

$\ii=\e^{\ii\pi/2}$, 所以$\ii z=r\e^{\ii(\theta+\pi/2)}=-b+\ii a$, 实部$-b$. 一般乘相位只将$\theta$加$\phi$. 同时实且纯虚的唯一数为0; $(z^*)^*=z$, $z^*=r\e^{-\ii\theta}$, $\operatorname{Re}z=(z+z^*)/2$, $\operatorname{Im}z=(z-z^*)/(2\ii)$, $zz^*=a^2+b^2=r^2=|z|^2$.

由$\e^{2\ii\theta}=(\cos\theta+\ii\sin\theta)^2$比较实虚部, 得$\cos2\theta=\cos^2\theta-\sin^2\theta$, $\sin2\theta=2\sin\theta\cos\theta$. 立方比较给$\cos3\theta=\cos^3\theta-3\cos\theta\sin^2\theta$, $\sin3\theta=3\sin\theta\cos^2\theta-\sin^3\theta$. 最后$\e^{\ii(A+B)}=\e^{\ii A}\e^{\ii B}$给$\cos(A+B)=\cos A\cos B-\sin A\sin B$, $\sin(A+B)=\sin A\cos B+\cos A\sin B$.
\end{solution}

\begin{exercise}
光子与自由电子碰撞, 解释光子为何不能被完全吸收. （5分）
\end{exercise}
\begin{solution}
在电子初始静止系, 设光子能量$\epsilon>0$. 若最终只有一个电子, 动量守恒给$p_f=\epsilon/c$, 能量守恒给$E_f=mc^2+\epsilon$. 但电子质量壳条件要求$E_f^2=p_f^2c^2+m^2c^4$, 代入得到$2mc^2\epsilon=0$, 矛盾. 因而至少还须有散射光子或其他物体带走动量; 束缚电子可以与原子核共同分担反冲, 不受这个自由二体约束.
\end{solution}

\begin{exercise}
Mach--Zehnder干涉仪入射态为$(\alpha,\beta)^{\mathsf T}$, $|\alpha|^2+|\beta|^2=1$, 探测器为$D_0,D_1$.
(a) 在下臂放相位$\delta_l$, 求$P_0,P_1$.
(b) 只在上臂放$\delta_u$.
(c) 两个相位器同时放入. （10分）
\end{exercise}
\begin{solution}
使用原讲义Lecture 2的两个分束器约定
$B_1=\frac1{\sqrt2}\begin{pmatrix}-1&1\\1&1\end{pmatrix}$,
$B_2=\frac1{\sqrt2}\begin{pmatrix}1&1\\1&-1\end{pmatrix}$.
上、下臂相位矩阵为$D=\operatorname{diag}(\e^{\ii\delta_u},\e^{\ii\delta_l})$. 输出为$B_2DB_1(\alpha,\beta)^{\mathsf T}$. 令$\delta=\delta_l-\delta_u$, 去掉不影响概率的整体$\e^{\ii\delta_u}$, 得
\begin{align*}
P_0(\delta)&=\frac14|\alpha(\e^{\ii\delta}-1)+\beta(1+\e^{\ii\delta})|^2,\\
P_1(\delta)&=\frac14|-\alpha(1+\e^{\ii\delta})+\beta(1-\e^{\ii\delta})|^2.
\end{align*}
(a)代$\delta=\delta_l$, (b)代$\delta=-\delta_u$, (c)代$\delta=\delta_l-\delta_u$. 因而两相位都相同不改变结果. $\delta=0$时$P_0=|\beta|^2$, $P_1=|\alpha|^2$, 与原讲义输出$(\beta,-\alpha)$一致. 不能把后一测试中两个相同分束器的约定直接套入本题.
\end{solution}

\begin{exercise}
Elitzur--Vaidman思想实验采用原讲义的平衡干涉仪, 从下方入射, 下臂放待检光子吸收器.
(a) 重复测试直至有足够把握判断状态, 有效吸收器有多少比例可在未触发吸收的情况下被认证?
(b) 80\%为失效吸收器. 随机取一个, 发入一个光子, 在$D_0$探测到. 该吸收器失效的条件概率是多少? 原题以“炸弹”命名这些理想化吸收器. （10分）
\end{exercise}
\begin{solution}
失效时两臂畅通,$D_0$概率1,$D_1$概率0. 有效时一次测试中, 吸收概率$1/2$, $D_0$与$D_1$各$1/4$. $D_1$事件认证有效且没有吸收; $D_0$不确定, 可重试. 对有效者, 认证比例是$\sum_{j\ge0}(1/4)^j(1/4)=1/3$. 这是重复次数趋大的理想极限; 有限次连续$D_0$不能逻辑上完全证明失效, 只能按先验降低有效概率.

Bayes规则给$\Pr(\text{失效}\mid D_0)=0.8\times1/[0.8\times1+0.2\times(1/4)]=16/17$. 将条件概率直接写成80\%会忽略此次测量提供的信息.
\end{solution}

\originalchapter{8.04 2015秋测试1及独立解答}
本章中文转写 \href{https://ocw.mit.edu/courses/8-04-quantum-physics-i-spring-2016/resources/mit8_04s16_test1_fall2015/}{Test 1, Fall 2015}, 原题日期2015年10月20日, 90分钟, 五题共100分, 禁止书籍、笔记和计算器. 它由2016春OCW课程页面发布, 不是2016春新试卷. 本章保留全部五题子问; 原件公式表随来源PDF保存, 本册此前已解释解题所需关系. 下列均为\textbf{AI独立解答}, 无官方答案身份.

\begin{exercise}
质子的电荷半径约0.9 fm. Compton波长与此尺度相比是“大得多、稍大、稍小、小得多”中的哪一种? （10分）
\end{exercise}
\begin{solution}
原题公式表使用$\lambda_C=h/(mc)$, 不是约化波长$\bar\lambda_C=\hbar/(mc)$. 用$\hbar c=197.3\,\mathrm{MeV\,fm}$及$m_pc^2=938\,\mathrm{MeV}$, 得$\lambda_C=2\pi\times197.3/938\simeq1.32$ fm, 所以是“稍大”. 如果漏掉$2\pi$会算成0.210 fm并选错方向.
\end{solution}

\begin{exercise}
相对论粒子满足$E^2=p^2c^2+m^2c^4$. 用$E=\hbar\omega$、$p=\hbar k$求群速度$\mathrm d\omega/\mathrm dk$, 用粒子速度$v$表示且不作近似. （10分）
\end{exercise}
\begin{solution}
链式法则给$v_g=\mathrm dE/\mathrm dp=pc^2/E$. 因$p=\gamma mv$, $E=\gamma mc^2$, 比值恰为$v$. 这依赖群速度导数, 若拿$\omega/k=E/p=c^2/v$代替, 会得到相速度并误以为信息超光速.
\end{solution}

\begin{exercise}
自由粒子质量$m$, 初始波函数$\psi(x,0)=\sin(k_0x)$, $k_0\in\R$.
(a) 求初始概率流.
(b) 测量动量可能得什么?
(c) 求$\psi(x,t)$. （15分）
\end{exercise}
\begin{solution}
(a)实函数使$\operatorname{Im}(\psi^*\psi')=0$, 所以$J=0$. (b)$\sin k_0x=(\e^{\ii k_0x}-\e^{-\ii k_0x})/(2\ii)$, 动量可能为$\pm\hbar k_0$. (c)两者都有自由能量$E=\hbar^2k_0^2/(2m)$, 所以$\psi(x,t)=\e^{-\ii\hbar k_0^2t/(2m)}\sin k_0x$. 此单色波在实线上不可普通归一化, 这些是原题的广义本征态计算; 周期盒子中可正规归一化. 原题允许$k_0=0$的形式值, 但此时给零函数, 不能代表物理态.
\end{solution}

\begin{exercise}
改造干涉仪, 两个分束器相同. 反射和透射概率为$R,T$, $R+T=1$, 正振幅$r=\sqrt R$, $t=\sqrt T$. 分束器矩阵除右下角外全为正实, 光子从上方入射.
(a) 构造$U$并验证幺正.
(b) 下臂放失效理想吸收器, 求$D_0,D_1$概率.
(c) 有效吸收器时, 求吸收、$D_0,D_1$概率.
(d) 重复不确定结果, 求有效者未被吸收而得到认证的比例$f$及最大可能值. （35分）
\end{exercise}
\begin{solution}
(a)从上入射反射留在上臂, 透射进入下臂, 两列因此为$(r,t)^{\mathsf T}$和$(t,-r)^{\mathsf T}$. 所以$U=\begin{pmatrix}r&t\\t&-r\end{pmatrix}$, 两列模为1且内积$rt-tr=0$, $U^\dagger U=I$. 此约定两个器件相同, 与上一作业不同.

(b)失效时两臂畅通,$U^2=I$, 上方输入回到$D_0$, 所以$P_0=1$, $P_1=0$. (c)第一分束器后是$(r,t)$, 吸收概率$T$. 未吸收的\textbf{未归一化}振幅为$(r,0)$, 第二分束器给$(r^2,rt)$, 所以$P_0=R^2$, $P_1=RT$, 与吸收项相加为1. 若先把幸存态归一化再忘记幸存概率, 会误算总概率.

(d)一次$D_0$不确定概率$R^2$, 取有效器件无限重试的理想极限, 认证比例
$f=\sum_{j\ge0}(R^2)^jRT=RT/(1-R^2)=R/(1+R)$, $0<R<1$. 随$R$增大趋$1/2$. 这里$1/2$是上确界, 在$R=1$端点并不能认证: 光子永远不进下臂, $D_1$永不发生. 越接近此极限, 有效器件直至被吸收或由$D_1$认证的平均总试次$1/(1-R^2)$越大.
\end{solution}

\begin{center}
\begin{tikzpicture}[x=1cm,y=1cm,>=Stealth]
\coordinate (B1) at (0,0);\coordinate (MU) at (2,1.3);\coordinate (ML) at (2,-1.3);\coordinate (B2) at (4,0);
\draw[->] (-1,.6)--(B1);\draw[->] (B1)--(MU);\draw[->] (MU)--(B2);
\draw[->] (B1)--(ML);\draw[->] (ML)--(B2);
\draw[->] (B2)--(5.3,.7);\draw[->] (B2)--(5.3,-.7);
\node[below left] at (B1){$U$};\node[above] at (MU){上臂};
\node[below] at (ML){下臂, 理想吸收器};\node[above] at (B2){$U$};
\node[right] at (5.3,.7){$D_0$};\node[right] at (5.3,-.7){$D_1$};
\end{tikzpicture}
\par 测试题拓扑重绘; 振幅基与原题一致, 臂长没有作为参数.
\end{center}

\begin{exercise}
同一自由Hamilton算符下的两个态重叠不随时间变. 一个似乎矛盾的例子: 两个Gaussian在$t=0$位置重合, 以远大于动量不确定性的相反动量传播. 声称“初始重叠大, 中心分开几个位置标准差后重叠小”. 找出错误.
令$\widehat\psi(x,0)=(2\pi)^{-1/4}a^{-1/2}\e^{-x^2/(4a^2)}$, $a>0$, $q\in\R$且量纲为动量, $\psi_1=\e^{\ii qx/\hbar}\widehat\psi$, $\psi_2=\e^{-\ii qx/\hbar}\widehat\psi$.
(a)求$\widehat\psi$的动量均值及其是否变化.
(b)求两个态动量均值及其是否变化.
(c)求初始重叠$\gamma(0)=\int\psi_1^*\psi_2\dd x$.
(d)用不等式表达平均动量大于不确定性, 判断例子的错误. （30分）
\end{exercise}
\begin{solution}
(a)实Gaussian的平均动量为0, $[H,p]=0$所以不随时间变化. (b)求导相位给$\langle p\rangle_1=q$, $\langle p\rangle_2=-q$, 同样守恒. 两者标准差均为$\Delta p=\hbar/(2a)$.

(c)初始内积的被积函数包含$\e^{-2\ii qx/\hbar}$, 不是仅两个密度相乘. 用前文Gaussian变换,
\[
\gamma(0)=\frac1{\sqrt{2\pi}a}\int
\e^{-x^2/(2a^2)-2\ii qx/\hbar}\dd x
=\e^{-2a^2q^2/\hbar^2}.
\]
(d)$|q|\gg\hbar/(2a)$表示平均动量大于标准差, 此时$\gamma(0)\ll1$, 初始位置密度重合也不意味着态内积大. 两态本来就因快速相对相位振荡而近正交, 后来内积保持这一小值. 错误在第一个声称, 不是幺正演化违反了重叠守恒.
\end{solution}

\originalchapter{8.05 作业6：演化速度与谐振子}
来源为8.05 Fall 2013的Assignment 6, 2013年10月10日发布, 10月18日15时截止. 以下完整转述六道题的各小问; 保留题号, 答案均为\textbf{AI独立解答}, 未取得官方答案. 该作业共有65分, 前五题各10分, 第六题15分.

\originalsection{第1题：初态重叠的短时变化}
\begin{exercise}
归一化态在不含时间的$H$下演化. 令$S(t)=|\braket{\psi(0)}{\psi(t)}|^2$. (a)说明它不能超过1, 并求初态为能量本征态时的值. (b)展开到$t^2$项, 用能量不确定性表达结果.
\end{exercise}
\begin{solution}
Cauchy--Schwarz不等式给$S\leq\norm{\psi(0)}^2\norm{\psi(t)}^2=1$. 能量本征态仅获得相位$\e^{-\ii Et/\hbar}$, 故$S=1$. 对一般初态, 在有关能量矩存在而足以作下述展开时,
\[
\braket{\psi(0)}{\psi(t)}=1-\frac{\ii t}{\hbar}\langle H\rangle-\frac{t^2}{2\hbar^2}\langle H^2\rangle+O(t^3).
\]
乘以其复共轭, 一阶项相消, 得$S(t)=1-(\Delta H)^2t^2/\hbar^2+O(t^3)$. 这说明重叠的最初下降由能量分散控制, 与平均能量的整体相位无关. 有限维系统直接满足展开条件; 对无界$H$不能只凭归一化就假定任意阶矩存在.
\end{solution}

\originalsection{第2题：到达正交态需要多久}
\begin{exercise}
已知$\Delta H\Delta Q\geq(\hbar/2)|\mathrm d\langle Q\rangle/\mathrm dt|$, 且$\Delta H$不随时间变化. 取$Q=\ket{\psi(0)}\bra{\psi(0)}$, 写$S(t)=\cos^2\phi(t)$, $\phi(0)=0$, $0\leq\phi\leq\pi/2$. (a)证明$\dot\phi\leq\Delta H/\hbar$, 推出$\Delta Ht_\perp\geq h/4$. (b)证明在$t\leq\pi\hbar/(2\Delta H)$时, $S(t)\geq\cos^2(\Delta Ht/\hbar)$.
\end{exercise}
\begin{solution}
$Q^2=Q$, 故$\langle Q\rangle=S$, $\Delta Q=\sqrt{S(1-S)}=\sin\phi\cos\phi$. 在$0<\phi<\pi/2$时代入不等式并约去该因子, 得更强的$|\dot\phi|\leq\Delta H/\hbar$. 端点处不能直接除以零, 但可由相邻内点积分和连续性得到同一角距离界. 因而$\phi(t)\leq\Delta Ht/\hbar$. 正交态要求$\phi=\pi/2$, 所以$t_\perp\geq\pi\hbar/(2\Delta H)=h/(4\Delta H)$. 在右侧角度不超过$\pi/2$的范围, $\cos^2$递减, 就得到(b). 若$\Delta H=0$, 态在单一能量子空间中仅变相位, 永不到达正交态; 此时$S=1$.
\end{solution}

\originalsection{第3题：自旋使时间下界取等号}
\begin{exercise}
自旋$1/2$粒子有$H=-\gamma B S_z$, 磁场沿$z$, 初态自旋沿正$x$. (a)求$S_z$基中的演化态并描述方向变化. (b)证明第2题的正交时间下界取等号.
\end{exercise}
\begin{solution}
以$S_z$本征值$\pm\hbar/2$的态为$\ket{\pm}$, 初态为$(\ket{+}+\ket{-})/\sqrt2$. 演化为
\[
\ket{\psi(t)}=\frac{\e^{\ii\gamma Bt/2}\ket{+}+\e^{-\ii\gamma Bt/2}\ket{-}}{\sqrt2},\qquad
\langle\boldsymbol S\rangle=\frac\hbar2(\cos\gamma Bt,-\sin\gamma Bt,0).
\]
因此自旋在赤道平面绕$z$以有向角速度$-\gamma B$进动. 重叠为$S(t)=\cos^2(\gamma Bt/2)$, $\Delta H=\hbar|\gamma B|/2$. 若$\gamma B\neq0$, 首次正交时间为$\pi/|\gamma B|$, 两者乘积正好是$\pi\hbar/2=h/4$. 无场或零耦合时不存在有限正交时间.
\end{solution}

\originalsection{第4题：求和规则与维里定理}
\begin{exercise}
$H=p^2/(2m)+V(x)$有离散完备能量基$\ket a$, $H\ket a=E_a\ket a$. (a)由双重对易子证明$\sum_{a'}|x_{aa'}|^2(E_{a'}-E_a)=\hbar^2/(2m)$. (b)证明$p_{aa'}=(\ii m/\hbar)(E_a-E_{a'})x_{aa'}$, 从而$\sum_{a'}|x_{aa'}|^2(E_a-E_{a'})^2=(\hbar^2/m^2)\langle a|p^2|a\rangle$. (c)推导定态的$2\langle T\rangle=\langle xV'(x)\rangle$, 并专门讨论$V=\alpha x^n$.
\end{exercise}
\begin{solution}
这里还需有关态在对易运算的共同定义域中, 求和收敛, 且积分边界项消失. $[x,H]=\ii\hbar p/m$, 所以$[[x,H],x]=\hbar^2/m$. 另一方面, 插入完备基给
\[
\langle a|[[x,H],x]|a\rangle=2\sum_{a'}(E_{a'}-E_a)|x_{aa'}|^2,
\]
即(a)的Thomas--Reiche--Kuhn规则. $[H,x]=-\ii\hbar p/m$的矩阵元给(b)第一式; 再用$\sum_{a'}|p_{aa'}|^2=\langle p^2\rangle_a$即得第二式. 对(c), 定态中$\langle[xp,H]\rangle=0$, 而$[xp,H]=2\ii\hbar T-\ii\hbar xV'$, 得所求关系. 用对称的$(xp+px)/2$也得同样结果. 若$V=\alpha x^n$, 则$2\langle T\rangle=n\langle V\rangle$, 前提仍是所讨论束缚态和边界允许上述运算.
\end{solution}

\originalsection{第5题：谐振子的矩阵与截断}
\begin{exercise}
(a)由基态唯一性证明激发态不简并. (b)证明$[a,(a^\dagger)^n]=n(a^\dagger)^{n-1}$. (c)给$a,a^\dagger,x,p,x^2,p^2,N$的一般矩阵元及$m,n=0,1,2,3$的截断矩阵. (d)计算截断$x,p$的对易子, 解释它为何不是$\ii\hbar I$. (e)求$\ket n$的$\Delta x,\Delta p$, 判断是否取不确定性等号, 并比较同能量经典最大位移.
\end{exercise}
\begin{solution}
若两个独立的第$n$激发态存在, $a^n$将它们映到基态子空间. 在该能量空间中$(a^\dagger)^na^n=n!I$, 故映射无核且保留线性独立, 与唯一基态矛盾. 对(b), 使用$[a,BC]=[a,B]C+B[a,C]$归纳, 每个因子贡献一次$(a^\dagger)^{n-1}$.

设$c_x=\sqrt{\hbar/(2m\omega)}$, $c_p=\sqrt{m\hbar\omega/2}$. 完整无限矩阵的非零元素由
\[
\begin{aligned}
a_{mn}&=\sqrt n\,\delta_{m,n-1},\\
(a^\dagger)_{mn}&=\sqrt{n+1}\,\delta_{m,n+1},\qquad N_{mn}=n\delta_{mn},\\
x_{mn}&=c_x\bigl(\sqrt n\,\delta_{m,n-1}+\sqrt{n+1}\,\delta_{m,n+1}\bigr),\\
p_{mn}&=\ii c_p\bigl(\sqrt{n+1}\,\delta_{m,n+1}-\sqrt n\,\delta_{m,n-1}\bigr),\\
(x^2)_{mn}&=c_x^2\bigl((2n+1)\delta_{mn}+\sqrt{n(n-1)}\delta_{m,n-2}\\
&\hspace{3em}+\sqrt{(n+1)(n+2)}\delta_{m,n+2}\bigr),\\
(p^2)_{mn}&=c_p^2\bigl((2n+1)\delta_{mn}-\sqrt{n(n-1)}\delta_{m,n-2}\\
&\hspace{3em}-\sqrt{(n+1)(n+2)}\delta_{m,n+2}\bigr)
\end{aligned}
\]
给出. 下标仍从0起. 四维截断令
\[
A=\begin{pmatrix}0&1&0&0\\0&0&\sqrt2&0\\0&0&0&\sqrt3\\0&0&0&0\end{pmatrix},
\quad a_4=A,\quad a^\dagger_4=A^{\mathsf T},\quad
X_4=c_x\begin{pmatrix}0&1&0&0\\1&0&\sqrt2&0\\0&\sqrt2&0&\sqrt3\\0&0&\sqrt3&0\end{pmatrix},
\]
\[
P_4=\ii c_p\begin{pmatrix}0&-1&0&0\\1&0&-\sqrt2&0\\0&\sqrt2&0&-\sqrt3\\0&0&\sqrt3&0\end{pmatrix},\qquad N_4=\operatorname{diag}(0,1,2,3).
\]
完整$x^2,p^2$的四维截块则分别为
\[
(x^2)_4=c_x^2\begin{pmatrix}1&0&\sqrt2&0\\0&3&0&\sqrt6\\\sqrt2&0&5&0\\0&\sqrt6&0&7\end{pmatrix},\qquad
(p^2)_4=c_p^2\begin{pmatrix}1&0&-\sqrt2&0\\0&3&0&-\sqrt6\\-\sqrt2&0&5&0\\0&-\sqrt6&0&7\end{pmatrix}.
\]
它们不等于$X_4^2,P_4^2$: 后两者丢失经$\ket4$返回$\ket3$的路径. 同理$[A,A^{\mathsf T}]=\operatorname{diag}(1,1,1,-3)$, 因此$[X_4,P_4]=\ii\hbar\operatorname{diag}(1,1,1,-3)$. 有限矩阵对易子的迹恒为0, 不可能等于有非零迹的$\ii\hbar I_4$.

$\langle x\rangle=\langle p\rangle=0$, 故$\Delta x=\sqrt{\hbar(n+1/2)/(m\omega)}$, $\Delta p=\sqrt{m\hbar\omega(n+1/2)}$. 乘积$\hbar(n+1/2)$仅在$n=0$等于$\hbar/2$. 同能量经典运动的转折点为$x_{\max}=\sqrt{2E_n/(m\omega^2)}=\sqrt2\Delta x$; 均方宽度与最远位移本来就是不同量.
\end{solution}

\originalsection{第6题：二维各向异性谐振子}
\begin{exercise}
$H=(p_x^2+p_y^2)/(2m)+m(\omega_x^2x^2+\omega_y^2y^2)/2$. (a)定义两方向升降和数算符, 求本征态与能量. (b)对$\omega_x\approx\omega_y$, $\omega_x>\omega_y$, 画前三组能级并标$(n_x,n_y)$. (c)改用$N=N_x+N_y$, $n=N_x-N_y$, 写$E_{N,n}$并重新标图. 判断$\{N\}$、$\{N,n\}$、$\{N_x,N_y\}$、$\{H\}$是否是完备相容观测量组; 讨论等频和有理频率比. (d)等频时把$L_z=xp_y-yp_x$用升降算符表达并证明与$H$对易. (e)利用$a_L=(a_x+\ii a_y)/\sqrt2$, $a_R=(a_x-\ii a_y)/\sqrt2$求每个能级的角动量本征基, 证明$\{H,L_z\}$完备.
\end{exercise}
\begin{solution}
对$j=x,y$, 定义$a_j=\sqrt{m\omega_j/(2\hbar)}x_j+\ii p_j/\sqrt{2m\hbar\omega_j}$, $N_j=a_j^\dagger a_j$. 两方向对易, 所以
\[
H=\hbar\omega_x(N_x+\tfrac12)+\hbar\omega_y(N_y+\tfrac12),\qquad
\ket{n_x,n_y}=\frac{(a_x^\dagger)^{n_x}(a_y^\dagger)^{n_y}}{\sqrt{n_x!n_y!}}\ket{0,0}.
\]
$E=\hbar\omega_x(n_x+1/2)+\hbar\omega_y(n_y+1/2)$. 图取$\omega_x=1.1\omega_y$作原题近等频情形的示例. 纵轴减去基态能量$E_0=1.05\hbar\omega_y$, 横向错开只是为了标签可读, 没有额外量子数. 第一图标$(n_x,n_y)$, 第二图按(c)改标$(N,n)$.
\begin{center}
\begin{tikzpicture}[x=1cm,y=2.2cm,>=Stealth]
\draw[->](0,-.12)--(0,2.6)node[above]{$(E-E_0)/(\hbar\omega_y)$};
\foreach \y/\n/\xx in {0/{(0,0)}/1,1/{(0,1)}/1,1.1/{(1,0)}/3,2/{(0,2)}/1,2.1/{(1,1)}/3,2.2/{(2,0)}/5}{
\draw (\xx,\y)--++(.8,0);\node[right]at(\xx+.8,\y){$\n$};}
\foreach \y in {0,1,2}{\draw(-.08,\y)--(.08,\y);\node[left]at(-.08,\y){$\y$};}
\end{tikzpicture}
\end{center}
\begin{center}
\begin{tikzpicture}[x=1cm,y=2.2cm,>=Stealth]
\draw[->](0,-.12)--(0,2.6)node[above]{$(E-E_0)/(\hbar\omega_y)$};
\foreach \y/\n/\xx in {0/{(0,0)}/1,1/{(1,-1)}/1,1.1/{(1,1)}/3,2/{(2,-2)}/1,2.1/{(2,0)}/3,2.2/{(2,2)}/5}{
\draw (\xx,\y)--++(.8,0);\node[right]at(\xx+.8,\y){$\n$};}
\foreach \y in {0,1,2}{\draw(-.08,\y)--(.08,\y);\node[left]at(-.08,\y){$\y$};}
\end{tikzpicture}
\end{center}
\begin{center}\begin{tabular}{ccc}
\toprule $(n_x,n_y)$ & $(N,n)$ & $(E-E_0)/(\hbar\omega_y)$\\\midrule
$(0,0)$ & $(0,0)$ & $0$\\
$(0,1)$, $(1,0)$ & $(1,-1)$, $(1,1)$ & $1$, $1.1$\\
$(0,2)$, $(1,1)$, $(2,0)$ & $(2,-2)$, $(2,0)$, $(2,2)$ & $2$, $2.1$, $2.2$\\\bottomrule
\end{tabular}\end{center}
由于$n_x=(N+n)/2$, $n_y=(N-n)/2$, 允许的$n$为$-N,-N+2,\ldots,N$. 能量是
\[
E_{N,n}=\frac\hbar2\bigl[(\omega_x+\omega_y)(N+1)+(\omega_x-\omega_y)n\bigr].
\]
完备组要求联合本征值唯一标记一个态. $\{N\}$不能区分同一$N\geq1$中的多个态; $\{N,n\}$和$\{N_x,N_y\}$总是完备. $\{H\}$仅在$\omega_x/\omega_y$无理时完备: 若两态能量相同, 整数差满足$\omega_x\Delta n_x+\omega_y\Delta n_y=0$, 无理比迫使两个差均为0. 有理比存在不同整数态的简并, $H$单独不完备. 等频就是特殊的有理比, $H$只决定$N$, 简并度$N+1$; 前三组分别1、2、3重. 另外两套双算符组仍完备, $\{N\}$仍不完备.

从此取等频. 代入$x_j=c_x(a_j+a_j^\dagger)$和$p_j=\ii c_p(a_j^\dagger-a_j)$, 交叉对易后得$L_z=\ii\hbar(a_y^\dagger a_x-a_x^\dagger a_y)$. 每项增减的总数抵消, 故$[N,L_z]=0$, 进而$[H,L_z]=0$. 新算符满足$[a_L,a_L^\dagger]=[a_R,a_R^\dagger]=1$, 交叉对易为0, 且
\[
H=\hbar\omega(N_L+N_R+1),\qquad L_z=\hbar(N_R-N_L).
\]
因此$\ket{n_L,n_R}=(a_L^\dagger)^{n_L}(a_R^\dagger)^{n_R}\ket{0,0}/\sqrt{n_L!n_R!}$为共同本征基. 固定$N=n_L+n_R$时角动量为$\hbar(-N,-N+2,\ldots,N)$, 每个值出现一次. 能量给总数, 角动量给差数, 唯一反解$n_L,n_R$, 所以$\{H,L_z\}$对整个空间完备.
\end{solution}

\originalchapter{8.06 作业1：微扰的适用范围}
来源为8.06 Spring 2018的Assignment 1, 发布日2018年2月7日. 以下完整转述五题各小问, 分值依次10、15、25、10、20分. 所有题解为\textbf{AI独立解答}, 没有取得官方答案. 只转述该公开作业自身给出的题面, 不复制其阅读书单中的商业教材.

\originalsection{第1题：二阶厄米矩阵}
\begin{exercise}
(a)求$(\boldsymbol a\cdot\boldsymbol\sigma)^2$及其本征值, 用迹求$\boldsymbol a\cdot\boldsymbol\sigma$与$a_0I+\boldsymbol a\cdot\boldsymbol\sigma$的本征值. (b)令$\boldsymbol a=(\alpha,0,\beta)$, 求精确本征值; 分别在$|\alpha|\ll|\beta|$和$|\alpha|\gg|\beta|$时给主项与前两个修正项, 对照二阶微扰. (c)矩阵内积为$\langle A,B\rangle=\Tr(A^\dagger B)$, 对$A=(a_0I+\boldsymbol a\cdot\boldsymbol\sigma)/2$, 求厄米$Q_i$使$a_i=\langle A,Q_i\rangle$.
\end{exercise}
\begin{solution}
Pauli乘法中的反对称项与$a_ia_j$收缩为0, 故平方为$|\boldsymbol a|^2I$, 两个平方本征值均为$|\boldsymbol a|^2$. 原矩阵厄米且迹为0, 本征值为$\pm|\boldsymbol a|$; 加$a_0I$平移为$a_0\pm|\boldsymbol a|$. $\boldsymbol a=0$时两者都为0.

对(b), 精确值为$\pm\sqrt{\alpha^2+\beta^2}$. 小$\alpha$而$\beta\neq0$时,
\[
E_\pm=\pm\left(|\beta|+\frac{\alpha^2}{2|\beta|}-\frac{\alpha^4}{8|\beta|^3}+\cdots\right).
\]
以$\beta\sigma_z$为未扰动矩阵, $\alpha\sigma_x$只有非对角元, 一阶为0, 二阶为$\pm\alpha^2/(2|\beta|)$, 正好给前一修正项; 四阶项不能由二阶公式产生. 反过来以$\alpha\sigma_x$为主项, 小$\beta$的展开为$\pm(|\alpha|+\beta^2/(2|\alpha|)-\beta^4/(8|\alpha|^3)+\cdots)$, 其二阶结果同理. 两者都为0时不能套有非零能隙的非简并公式.

$\Tr I^2=2$, $\Tr\sigma_i\sigma_j=2\delta_{ij}$, $\Tr\sigma_i=0$, 所以(c)取$Q_0=I$, $Q_i=\sigma_i$即可. 内积中的$A^\dagger$在此等于$A$.
\end{solution}

\originalsection{第2题：三次势的形式基态}
\begin{exercise}
令$d^2=\hbar/(m\omega)$,
\[
H=\frac{p^2}{2m}+\frac12m\omega^2x^2+\lambda\sqrt2\hbar\omega\frac{x^3}{d^3}.
\]
(a)求未扰动基态的一阶和二阶能量变化. (b)把一阶波函数写成谐振子态的和. (c)画小$\lambda$的势, 判断所得态是否为真正基态, 说明微扰遗漏的效应.
\end{exercise}
\begin{solution}
$x=d(a+a^\dagger)/\sqrt2$, 所以扰动为$\lambda\hbar\omega(a+a^\dagger)^3/2$. 逐次作用给$(a+a^\dagger)^3\ket0=3\ket1+\sqrt6\ket3$, 不可把不对易的因子按普通二项式展开. 奇偶性使一阶能量为0, 而
\[
\Delta E^{(2)}=\frac{\lambda^2(\hbar\omega)^2}{4}
\left(\frac9{-\hbar\omega}+\frac6{-3\hbar\omega}\right)
=-\frac{11}{4}\lambda^2\hbar\omega.
\]
态为$\ket\psi=\ket0-(3\lambda/2)\ket1-(\sqrt6\lambda/6)\ket3+O(\lambda^2)$, 已在一阶归一化. 这是围绕局部阱的形式近似, 不能命名为实际全局基态.

取$\lambda=0.08$仅作图, 横轴$s=x/d$, 纵轴$v=V/(\hbar\omega)=s^2/2+\sqrt2\lambda s^3$. 势在左端趋向$-\infty$, 原点只是局部最低点. 对任意非零$\lambda$都不存在最低能量: 把归一化波包平移到负无穷(正$\lambda$)可使势能任意低而保持固定动能. 图中的局部势垒使原点附近的态可能成为寿命较长的亚稳态; 穿过势垒逃逸的隧穿衰变不会被有限阶实能量微扰描述. 负$\lambda$时图左右镜像.
\begin{center}\begin{tikzpicture}
\begin{axis}[width=.65\textwidth,height=5cm,axis lines=middle,xmin=-5.5,xmax=3,ymin=-4,ymax=7,xlabel={$s=x/d$},ylabel={$v$},samples=180,domain=-5.5:3]
\addplot[thick,structurecolor]{.5*x^2+sqrt(2)*.08*x^3};
\end{axis}\end{tikzpicture}\end{center}
\end{solution}

\originalsection{第3题：三维谐振子的简并空间}
\begin{exercise}
$H_0=\boldsymbol p^2/(2m)+m\omega^2\boldsymbol x^2/2$, $\delta H=\lambda\omega L_2$, 其中$L_2=x_3p_1-x_1p_3$是$y$方向分量, 不是角动量平方. (a)定义三方向升降和数算符, 求$\ket{n_1,n_2,n_3}$能量及$5\hbar\omega/2$简并度. (b)用升降算符写$\delta H$. (c)求该简并空间的扰动矩阵. (d)求扰动和全$H$的本征值、态. (e)在扰动本征基中写全$H$矩阵, 解释任意小非零扰动对基的选择. (f)证明不同未扰动能级之间的扰动矩阵元为0, 给出相应算符结论, 说明何以本题可以不要求$\lambda$小.
\end{exercise}
\begin{solution}
$a_j=\sqrt{m\omega/(2\hbar)}x_j+\ii p_j/\sqrt{2m\hbar\omega}$, $N_j=a_j^\dagger a_j$, 所以$H_0=\hbar\omega(N_1+N_2+N_3+3/2)$. $5\hbar\omega/2$的空间为$\ket{100},\ket{010},\ket{001}$三维. 直接代入得$L_2=\ii\hbar(a_1^\dagger a_3-a_3^\dagger a_1)$, 因而在此顺序中
\[
\delta H=\lambda\hbar\omega\begin{pmatrix}0&0&\ii\\0&0&0\\-\ii&0&0\end{pmatrix}.
\]
零本征值的态为$\ket{010}$, $+\lambda\hbar\omega$和$-\lambda\hbar\omega$分别对应$(\ket{100}-\ii\ket{001})/\sqrt2$和$(\ket{100}+\ii\ket{001})/\sqrt2$. 全能量在各扰动本征值上加$5\hbar\omega/2$. 按正、零、负的次序, (e)的矩阵是$\hbar\omega\operatorname{diag}(5/2+\lambda,5/2,5/2-\lambda)$. 无论非零$\lambda$多小, 选出的第一、第三态都含等幅的原基分量; 这种有限角度的基变换不随$\lambda$趋零而缩小. 在$\lambda=0$恰好重新有基选择自由.

扰动中的每项保持总激发数, 故$[H_0,\delta H]=0$. 对$E_\varphi\neq E_\psi$, $(E_\varphi-E_\psi)\bra\varphi\delta H\ket\psi=\bra\varphi[H_0,\delta H]\ket\psi=0$. 因此各能级块互不混合, 本题对这个块的对角化是精确的, 不需忽略块外高阶修正. 一般扰动并不具有此性质. 这些块内谱式对所有实$\lambda$成立; 若还要求整个三维振子有受限的最低能量, 应注意圆偏振模的频率为$\omega(1\pm\lambda)$, 稳定束缚情形要求$|\lambda|<1$.
\end{solution}

\originalsection{第4题：圆周粒子的极化率与三重势}
\begin{exercise}
质量$m$的粒子在半径$a$圆周上运动, $\psi(\varphi+2\pi)=\psi(\varphi)$且$\int_0^{2\pi}|\psi|^2\dd\varphi=1$. (a)求$H_0=L_z^2/(2ma^2)$的本征态、能量和简并. (b)粒子电荷$q$, 沿$x$加电场$\varepsilon$, 扰动$-q\varepsilon a\cos\varphi$. 求基态一阶波函数、诱导偶极矩$\langle qx\rangle$及极化率. (c)撤去场, 用$b\cos3\varphi$模拟乙烷两端$\mathrm{CH}_3$基团的相对转动, 求一阶基态和二阶能量, 解释三重对称效应. 原题以$2ma^2$代替实际分子转动惯量系数, 本题保留该简化.
\end{exercise}
\begin{solution}
周期边界使$u_n(\varphi)=\e^{\ii n\varphi}/\sqrt{2\pi}$, $n\in\mathbb Z$, $E_n=Bn^2$, 其中$B=\hbar^2/(2ma^2)$是能量尺度, 此处不是磁场. 基态$n=0$不简并, 非零$\pm n$成对. 因为$\cos\varphi=(\e^{\ii\varphi}+\e^{-\ii\varphi})/2$, 只有$n=\pm1$与0耦合. 两个一阶系数为$q\varepsilon a/(2B)$, 所以
\[
\psi(\varphi)=\frac1{\sqrt{2\pi}}\left(1+\frac{q\varepsilon a}{B}\cos\varphi\right)+O(\varepsilon^2),\qquad
\langle qx\rangle=\frac{q^2a^2}{B}\varepsilon+O(\varepsilon^2).
\]
用了$\int_0^{2\pi}\cos^2\varphi\dd\varphi=\pi$. 极化率为$2mq^2a^4/\hbar^2$, 为正, 与电荷符号无关. 近似要求$|q\varepsilon a|\ll B$.

三重势只耦合$0$与$\pm3$, 系数为$-b/(18B)$, 得
\[
\psi=\frac1{\sqrt{2\pi}}\left(1-\frac b{9B}\cos3\varphi\right)+O(b^2),\qquad
E_0=0-\frac{b^2}{18B}+O(b^3).
\]
一阶能量为0. 二阶项来自两条相同的$|b/2|^2/(-9B)$路径. $b>0$时低势方向$\cos3\varphi=-1$的概率增加, 三个等价相对朝向出现; $b<0$时转过$\pi/3$. 能量下降反映基态通过调整分布集中到有利朝向, 而非粒子完全静止在经典最低点.
\end{solution}

\originalsection{第5题：有限核半径的能量修正}
\begin{exercise}
质子电荷均匀分布在半径$R$球内, 忽略氢精细结构、Lamb移和超精细分裂. 原题使用Gaussian单位, 球外势能为$-e^2/r$. (a)求所有$r$的电子势能并验证连续. (b)利用核内未扰动基态波函数几乎不变, 求最低阶能量位移; $R=10^{-13}\,\mathrm{cm}$时求其相对13.6eV结合能的大小. (c)解释为何$\ell=0$最显著, 估计$\ell=1$的额外抑制因子. (d)原题引述2010年缪子氢测量$R=0.84184(67)\times10^{-13}\,\mathrm{cm}$及当时CODATA值$0.8768(69)\times10^{-13}\,\mathrm{cm}$, 要求解释重206.8倍的缪子何以更敏感. 这些是原题的\textbf{历史数据}, 不是对当前常数或争议状态的断言.
\end{exercise}
\begin{solution}
球内包围电荷为$e(r/R)^3$, Gauss定律给电场$er/R^3$, 球外$e/r^2$. 电子势能导数为$eE$, 用$V(\infty)=0$并在$R$连续匹配得
\[
V(r)=\begin{cases}-\dfrac{e^2}{2R}(3-r^2/R^2),&0\leq r<R,\\-e^2/r,&r\geq R.\end{cases}
\]
相对点核的改变量仅在球内, 为$\delta V=e^2(1/r-3/(2R)+r^2/(2R^3))$. 以$|\psi_{100}(0)|^2=1/(\pi a_0^3)$代替核内密度,
\[
\Delta E\simeq\frac{4e^2}{a_0^3}\int_0^R\left(r-\frac{3r^2}{2R}+\frac{r^4}{2R^3}\right)\dd r
=\frac{2e^2R^2}{5a_0^3}.
\]
位移为正, 有限核减弱最深处吸引. 因为结合能大小为$e^2/(2a_0)$, 比例为$(4/5)(R/a_0)^2$. 用$a_0\simeq5.29\times10^4\,\mathrm{fm}$, $R=1\,\mathrm{fm}$, 得$2.86\times10^{-10}$, 即约$3.9\times10^{-9}\,\mathrm{eV}$. 略去密度随$r$变化的相对误差为$O(R/a_0)$.

正则径向解在小$r$有$R_\ell(r)\propto r^\ell$, 所以$\ell=1$的核内密度相比同尺度的$s$态额外带$(R/a_0)^2$量级, 约$3.6\times10^{-10}$; 不同主量子数另有系数, 这里只是尺度估计. $s$态在原点有非零密度, 对有限核最敏感.

Bohr长度实际用约化质量$\mu=m_{\rm lepton}m_p/(m_{\rm lepton}+m_p)$, 因而$\Delta E\propto a^{-3}\propto\mu^3$, 相对结合能的位移则正比$\mu^2$. 若把质子当无限重, 缪子使长度缩小206.8倍. 保留$m_p\simeq1836.15m_e$时, 缪子氢与普通氢的约化质量比约186, 仍使绝对有限核位移增强约$6.4\times10^6$倍. 这解释了半径灵敏度; 实际谱线提取还需原题暂时忽略的其他效应, 不能把本简化式当作完整实验分析.
\end{solution}

\originalchapter{8.04 作业2：波长与干涉仪}
来源为8.04 Spring 2016 Assignment 2, 2016年2月11日发布, 2月18日17时截止, 共65分. 以下转述全部五题及小问, 解答均为\textbf{AI独立解答}. 题中的加速器预期属于2016年题面背景, 不代表当前建设进度.
\originalsection{第1题：非相对论与相对论波长（20分）}
\begin{exercise}
(a)写电子非相对论波长为$\lambda_{\rm nr}=\delta/\sqrt{E_{\rm kin}}$埃, 能量的数值用eV, 求无量纲$\delta$. (b)对$E=\gamma m_ec^2$, 写$\lambda_{\rm r}=\ell/\sqrt{\gamma^2-1}$, 求$\ell$的fm值和名称. (c)用$\ell,\gamma$重写非相对论式, 比较两波长. (d)计算波长等于电子Compton波长时的能量及速度; 1 TeV电子和7 TeV质子的波长; 非相对论波长相对真实波长误差至多10\%时的最大动能与速度.
\end{exercise}
\begin{solution}
由$p=\sqrt{2m_eE_{\rm kin}}$, $\lambda=h/p$得到$\delta=h/[10^{-10}\sqrt{2m_e(1\,\mathrm{eV})}]=12.264$. 相对论关系$E^2=p^2c^2+m_e^2c^4$给$\ell=h/(m_ec)=2426.31$ fm, 为\textbf{非约化}Compton波长. 这里不是$\hbar/(m_ec)$.

用动能$(\gamma-1)m_ec^2$代入(a), 得$\lambda_{\rm nr}=\ell/\sqrt{2(\gamma-1)}$, 因而$\lambda_{\rm nr}/\lambda_{\rm r}=\sqrt{(\gamma+1)/2}>1$（非零动能）. 若$\lambda=\ell$, 则$p=m_ec$, $\gamma=\sqrt2$, 总能量$722.66$ keV, 动能$211.66$ keV, $\beta=1/\sqrt2=0.7071>0.2$, 显然需要相对论.

高能时$pc\simeq E$, $hc=1.23984\,\mathrm{GeV\,fm}$, 所以1 TeV电子$\lambda\simeq1.23984\times10^{-3}$ fm, 7 TeV质子$\lambda\simeq1.77120\times10^{-4}$ fm, 后者约短7倍. 保留质量项只需用$pc=\sqrt{E^2-m^2c^4}$. 按误差定义$(\lambda_{\rm nr}-\lambda_{\rm r})/\lambda_{\rm r}\leq0.1$, 有$\gamma\leq1.42$, 最大动能$0.42m_ec^2=214.62$ keV, $\beta=\sqrt{1-1/1.42^2}=0.7100$. 若改以近似值作误差分母, 阈值会不同; 本解明确采用真实值为分母.
\end{solution}
\originalsection{第2题：三个电子长度（10分）}
\begin{exercise}
题面采用Gaussian单位$\alpha=e^2/(\hbar c)$、$r_0=e^2/(m_ec^2)$、$\bar\lambda_C=\hbar/(m_ec)$. (a)由$e^2,\hbar,m_e$构造Bohr半径并给fm值. (b)证明三个长度之比为$1:\alpha:\alpha^2$, 并给数值.
\end{exercise}
\begin{solution}
$e^2$的量纲为能量乘长度. 设长度为$\hbar^a m_e^b(e^2)^d$, 比较质量、长度、时间指数得$a=2,b=-1,d=-1$, 故$a_0=\hbar^2/(m_ee^2)=52917.7$ fm. $\bar\lambda_C=\alpha a_0=386.159$ fm, $r_0=\alpha\bar\lambda_C=2.81794$ fm. SI中各$e^2$应替换成$e^2/(4\pi\epsilon_0)$; 这些不同单位制的电荷公式不能混用.
\end{solution}
\originalsection{第3题：概率守恒为何要求幺正（10分）}
\begin{exercise}
双端口归一化入射振幅$u$经过线性矩阵$R$成为$Ru$. 证明对任意入射态保持概率意味着$R^\dagger R=I$.
\end{exercise}
\begin{solution}
条件为$u^\dagger(R^\dagger R-I)u=0$. 令Hermitian矩阵$M=R^\dagger R-I$. 分别取两个单位基矢得$M_{11}=M_{22}=0$; 取$(1,1)/\sqrt2$得$\operatorname{Re}M_{12}=0$, 取$(1,\ii)/\sqrt2$得$\operatorname{Im}M_{12}=0$. 因而$M=0$. 反过来$R^\dagger R=I$直接保证$\norm{Ru}^2=\norm u^2$. 只检验单个输入的概率守恒不能推出幺正.
\end{solution}
\originalsection{第4题：不对称分束器测弹（15分）}
\begin{exercise}
分束器反射、透射概率为$R,T$, $R+T=1$, 振幅$r=\sqrt R,t=\sqrt T$; 矩阵只有右下元素为负. 干涉仪用两个相同分束器, 从上端口输入, 弹在下支路. (a)求幺正矩阵. (b)求坏弹与好弹的两探测概率及爆炸概率. (c)重复测试直到认证或爆炸, 求好弹被无损认证的比例及上界.
\end{exercise}
\begin{solution}
$U=\left(\begin{smallmatrix}r&t\\t&-r\end{smallmatrix}\right)$, $U^\dagger U=U^2=I$. 坏弹不阻断, 输出$U^2(1,0)^{\mathsf T}=(1,0)^{\mathsf T}$, 所以$P_0=1,P_1=0$. 好弹时第一次分束后振幅为$(r,t)$: 下支路概率$T$爆炸, 未爆炸的\textbf{未归一化}振幅为$(r,0)$. 再分束得$(r^2,rt)$, 故$P_{\rm boom}=T$, $P_0=R^2$, $P_1=RT$, 三者和为1. $D_1$认证, $D_0$结果不确定可重试. 几何级数给$f=RT/(1-R^2)=R/(1+R)$（$0<R<1$）, 上确界$1/2$在$R\to1^-$逼近. $R=1$时既不爆炸也不认证, 不能代入说有一半认证成功; 需要的测试次数也随极限增大.
\end{solution}
\originalsection{第5题：平面物质波为何为复指数（10分）}
\begin{exercise}
令$\Psi=\cos\theta+\gamma\sin\theta$, $\theta=kx-\omega t$, 要求任意相移$\epsilon$只乘常数$a(\epsilon)$. 求$\gamma,a$, 指出通常使用的分支.
\end{exercise}
\begin{solution}
逐项比较$\cos\theta,\sin\theta$系数:
$a=\cos\epsilon+\gamma\sin\epsilon$, $a\gamma=-\sin\epsilon+\gamma\cos\epsilon$. 消去$a$得$(1+\gamma^2)\sin\epsilon=0$. 对任意$\epsilon$成立只允许$\gamma=\pm\ii$. 对应$a=\e^{\pm\ii\epsilon}$, $\Psi=\e^{\pm\ii\theta}$. 常规$p=\hbar k,E=\hbar\omega$的正动量正能量平面波选择$+\ii$, 即$\e^{\ii(kx-\omega t)}$; 共轭分支改变导数本征值.
\end{solution}

\originalchapter{8.04 作业4：自由波包与概率流}
来源为8.04 Spring 2016 Assignment 4, 2016年2月25日发布, 3月4日中午截止. 全部七题共65分, 解答为\textbf{AI独立解答}. 教材阅读条目仅为参考, 以下题面取自实际OCW作业PDF.
\originalsection{第1题：波包展宽尺度（5分）}
\begin{exercise}
(a)估算自由质子宽度$10^{-10}$ m及1 cm的显著展宽时间. (b)若$\Delta x\Delta p\sim\hbar$, 证明$\Delta p\ll p$时波包通过固定位置期间不显著展宽.
\end{exercise}
\begin{solution}
速度分散$\Delta v\sim\Delta p/m$, 展宽达到原宽度需要$t_s\sim\Delta x/\Delta v\sim m(\Delta x)^2/\hbar$. 质子质量$1.67262\times10^{-27}$ kg给$t_s\sim1.586\times10^{-13}$ s和$1.586\times10^3$ s（约26.4分钟）. 这里是量级估算; 按Gaussian宽度增大$\sqrt2$定义会有2的系数差. 通过时间$t_p\sim\Delta x/(p/m)$, 因而$t_p/t_s\sim\Delta p/p\ll1$.
\end{solution}
\originalsection{第2题：平面入射加球面散射（10分）}
\begin{exercise}
远场定态为$\psi=\e^{\ii kz}+f(\theta)\e^{\ii kr}/r$, 时间整体相位省略. (a)求入射流及其穿过大球的净通量. (b)求散射流径向分量及总通量. (c)保留$1/r$主阶求交叉流径向分量.
\end{exercise}
\begin{solution}
用$\boldsymbol J=(\hbar/m)\operatorname{Im}(\psi^*\boldsymbol\nabla\psi)$. 入射$\boldsymbol J_1=(\hbar k/m)\hat z$; 径向分量乘$\cos\theta$, 在球面积分为0（有流入也有流出）. 球面波给$J_{2r}=\hbar k|f(\theta)|^2/(mr^2)$, 总向外通量为$(\hbar k/m)\int|f|^2\dd\Omega$. $f$只依赖角度, 对$r$求导时视为常数. 交叉项为
\[
J_{12,r}=\frac{\hbar k}{mr}\operatorname{Im}\!\left\{\ii\left[f\e^{\ii kr(1-\cos\theta)}+\cos\theta f^*\e^{-\ii kr(1-\cos\theta)}\right]\right\}+O(r^{-2}).
\]
它来自$\psi_1^*\partial_r\psi_2+\psi_2^*\partial_r\psi_1$, 用$\partial_r\psi_1=\ii k\cos\theta\psi_1$即可逐项核对. 不能把交叉项先逐点丢掉再算总球面通量; 前向小角区域贡献导致光学定理, 原题不要求完成该积分.
\end{solution}
\originalsection{第3题：Gaussian的精确演化（15分）}
\begin{exercise}
初态$\psi_a(x,0)=(2\pi)^{-1/4}a^{-1/2}\e^{-x^2/(4a^2)}$. (a)验证归一化. (b)以$1/\sqrt{2\pi}$的波数Fourier约定求$\Phi_a(k)$. (c)求自由演化, 令$\tau=2ma^2/\hbar$. (d)写时刻$t$的概率密度及宽度$a(t)$.
\end{exercise}
\begin{solution}
$|\psi|^2=\e^{-x^2/(2a^2)}/(\sqrt{2\pi}a)$, 积分为1. 实Gaussian积分给$\Phi_a(k)=(2/\pi)^{1/4}\sqrt a\,\e^{-a^2k^2}$. 每个波数分量乘$\e^{-\ii\hbar k^2t/(2m)}$, 积分中$A=a^2+\ii\hbar t/(2m)$,$B=\ii x$, 代$\int\e^{-Ak^2+Bk}\dd k=\sqrt{\pi/A}\e^{B^2/(4A)}$（$\operatorname{Re}A>0$）得
\[
\psi_a(x,t)=\frac{\exp[-x^2/(4a^2(1+\ii t/\tau))]}{(2\pi)^{1/4}\sqrt a\sqrt{1+\ii t/\tau}},\quad
a(t)=a\sqrt{1+(t/\tau)^2}.
\]
平方根从$t=0$的正值连续选取. 取模方利用$\operatorname{Re}(1/(1+\ii u))=1/(1+u^2)$, 定义$G(x;b)=\exp[-x^2/(2b^2)]/(\sqrt{2\pi}b)$, 得$|\psi_a(x,t)|^2=G(x;a(t))$. 总概率始终1, 相位中的$x^2$项记录位置与动量相关, 并非热扩散.
\end{solution}
\originalsection{第4题：概率在Fourier变换下保持（10分）}
\begin{exercise}
(a)证明一维及三维$\int|\psi|^2\dd x=\int|\phi|^2\dd p$. (b)氢基态$\psi=N\e^{-r/a_0}$, 动量态$\phi=N'/(1+a_0^2p^2/\hbar^2)^2$, 求两常数及$p>\hbar/a_0$的概率.
\end{exercise}
\begin{solution}
在$d$维采用$\psi(x)=(2\pi\hbar)^{-d/2}\int\phi(p)\e^{\ii p\cdot x/\hbar}\dd^dp$. 展开模方后先保留两动量积分:
\[
\int|\psi(x)|^2\dd^dx=\frac1{(2\pi\hbar)^d}\int\dd^dp\,\dd^dq\,\phi(p)\phi(q)^*\int\e^{\ii(p-q)\cdot x/\hbar}\dd^dx.
\]
内层为$(2\pi\hbar)^d\delta^{(d)}(p-q)$, 对$q$积分留下$\int|\phi(p)|^2\dd^dp$. 可先对快速衰减函数证明再按平方范数连续推广, 不要求任意函数积分都绝对收敛.

位置归一化$4\pi N^2\int_0^\infty r^2\e^{-2r/a_0}\dd r=\pi a_0^3N^2=1$, 所以$N=(\pi a_0^3)^{-1/2}$. 令$u=a_0p/\hbar$, 则$\int_0^\infty u^2/(1+u^2)^4\dd u=\pi/32$, 从而$N'=2\sqrt2\,a_0^{3/2}/(\pi\hbar^{3/2})$（整体相位可选）. 大动量概率为
\[
P=\frac{32}{\pi}\int_1^\infty\frac{u^2\dd u}{(1+u^2)^4}=\frac12-\frac2{3\pi}=0.28779.
\]
最后一步令$u=\tan\theta$, 积分变成$\sin^2\theta\cos^4\theta$在$[\pi/4,\pi/2]$上的积分, 即可不用数值软件求出.
\end{solution}
\originalsection{第5、6题：均值运动与动量宽度（各10、5分）}
\begin{exercise}
第5题: 对$H=p^2/(2m)+V(x)$证明$\mathrm d\langle x\rangle/\mathrm dt=\langle p\rangle/m$, $\mathrm d\langle p\rangle/\mathrm dt=-\langle V'\rangle$. 第6题: 证明自由波包的$\Delta p$不变.
\end{exercise}
\begin{solution}
不显含时间的$A$满足$\dot{\langle A\rangle}=(\ii/\hbar)\langle[H,A]\rangle$. $[p^2,x]=-2\ii\hbar p$给第一式, $[V,p]=\ii\hbar V'$给第二式. 这不是通常的$-V'(\langle x\rangle)$: 只有力近似线性或分布足够窄时才能以均值代入. 自由时$[H,p]=[H,p^2]=0$, 所以$\langle p\rangle$和$\langle p^2\rangle$恒定, $\Delta p$亦恒定. 这些操作假定所需矩存在且边界允许分部积分.
\end{solution}
\originalsection{第7题：空间相位的可测作用（10分）}
\begin{exercise}
归一化$\psi_0$有均值$x_0,p_0$. 定义$(B_qf)(x)=\e^{\ii qx/\hbar}f(x)$, $q$为实动量, 新态$B_q\psi_0$. (a,b)求位置、动量均值. (c)解释该相位. (d)算$[p,B_q],[x,B_q]$.
\end{exercise}
\begin{solution}
模方不变, 所以$\langle x\rangle=x_0$. 求导给$pB_qf=B_qpf+qB_qf$, 故$\langle p\rangle=p_0+q$, $[p,B_q]=qB_q$, $[x,B_q]=0$. 这使整个动量分布平移$q$, 是动量助推; 空间依赖相位有可测后果, 与所有位置都相同的整体相位不同.
\end{solution}

\originalchapter{8.04 作业5：不确定性与势阱}
来源为8.04 Spring 2016 Assignment 5, 页面与标题为2016春, PDF发布行印为2015年3月11日, 截止行是2016年3月18日中午; 保留该日期差异, 不默改原件. 全部八题共65分, 以下为\textbf{AI独立解答}.
\originalsection{第1题：实Gaussian（5分）}
\begin{exercise}
令$\psi=N\e^{-x^2/(2a^2)}$, $a>0$. (a)在位置表象计算$\Delta x,\Delta p$, 检验不确定性是否饱和. (b)求动量波函数, 验证归一化并重新求$\Delta p$.
\end{exercise}
\begin{solution}
$N=(\pi a^2)^{-1/4}$, 偶性给两均值0. 由$I(b)=\int\e^{-bx^2}\dd x=\sqrt{\pi/b}$,$-I\prime(b)=\int x^2\e^{-bx^2}\dd x=\sqrt\pi/(2b^{3/2})$, 取$b=1/a^2$并除归一积分, 得$\langle x^2\rangle=a^2/2$. 用$\langle p^2\rangle=\hbar^2\int|\psi'|^2\dd x=\hbar^2/(2a^2)$, 得$\Delta x=a/\sqrt2,\Delta p=\hbar/(\sqrt2a)$, 乘积$\hbar/2$. Fourier积分给$\phi(p)=\sqrt{a/(\sqrt\pi\hbar)}\e^{-a^2p^2/(2\hbar^2)}$. 积分$|\phi|^2$为1, 其二阶矩为$\hbar^2/(2a^2)$, 与位置计算一致.
\end{solution}
\originalsection{第2题：复Gaussian与相关性（10分）}
\begin{exercise}
把$a$替换为复$\Delta$, $\operatorname{Re}(\Delta^2)>0$, $N$取实正. (a)求宽度, 以$|\Delta|$和$\operatorname{Re}(\Delta^2)$表示. (b)Fourier变换, 检验Parseval归一化, 并在动量表象重算$\Delta p$. (c)令$\Delta=|\Delta|\e^{\ii\phi_\Delta}$, 求不确定性乘积并讨论$\phi_\Delta=0,\pi/4$. (d)用于作业4自由Gaussian, 验证动量宽度不随时间变.
\end{exercise}
\begin{solution}
设$D=\Delta^2$, $A=\operatorname{Re}(1/D)=\operatorname{Re}D/|D|^2>0$, $N=(A/\pi)^{1/4}$. 共轭分式$1/z=z^*/|z|^2$证明题给恒等式. $\langle x^2\rangle=1/(2A)$, 而$\psi'=-x\psi/D$, 故
\[
\Delta x=\frac{|\Delta|^2}{\sqrt{2\operatorname{Re}(\Delta^2)}},\qquad
\Delta p=\frac{\hbar}{\sqrt{2\operatorname{Re}(\Delta^2)}},\qquad
\phi(p)=\frac{N\sqrt D}{\sqrt\hbar}\e^{-\Delta^2p^2/(2\hbar^2)}.
\]
固定Fourier变换取$\operatorname{Re}\sqrt D>0$的平方根; $D=\Delta^2$允许原参数$\Delta$正负两代表, 不能未经选择直接把根写成$\Delta$. 这个选择使变换等号精确, 而不仅到整体相位. 模方动量密度为$N^2|\Delta|^2\e^{-\operatorname{Re}D\,p^2/\hbar^2}/\hbar$, 用前题$-I\prime(b)/I(b)=1/(2b)$并取$b=\operatorname{Re}D/\hbar^2$, 积分为1（亦为Parseval检验）, 两均值由偶性为0, 动量方差$\hbar^2/[2\operatorname{Re}D]$与位置表象计算相同. 乘积为$\hbar/[2\cos(2\phi_\Delta)]$（要求分母正）. $\phi_\Delta=0$取等号; $\phi_\Delta\to\pi/4$时位置宽度发散, 端点本身不可归一, 不能当成有限宽度态. 作业4有$D(t)=2a^2(1+\ii t/\tau)$, 实部恒为$2a^2$, 所以$\Delta p=\hbar/(2a)$始终不变. 位置展宽来自相关相位, 不来自动量分布变宽.
\end{solution}
\originalsection{第3题：盒中定态的宽度（15分）}
\begin{exercise}
区间$[0,a]$有无限高墙, $\Psi_n=N\sin(n\pi x/a)\e^{-\ii\phi_n(t)}$. (a)求相位和归一化. (b)求位置一、二阶矩及宽度. (c)求$\langle p\rangle,\langle p^2\rangle$及宽度,$n=1,2,\ldots$. (d)判断不确定性是否满足及饱和. (e)说明上述量是否随时间变.
\end{exercise}
\begin{solution}
$N=\sqrt{2/a}$, $E_n=\hbar^2n^2\pi^2/(2ma^2)$, $\phi_n(t)=E_nt/\hbar+\phi_n(0)$. 用$\sin^2u=(1-\cos2u)/2$. 令$b=2n\pi/a$, 两次分部积分给$\int_0^a x^2\cos bx\dd x=2a/b^2=a^3/(2n^2\pi^2)$; 一阶余弦矩为0. 将这两项代回归一密度积分得
\[
\langle x\rangle=\frac a2,\quad\langle x^2\rangle=a^2\left(\frac13-\frac1{2n^2\pi^2}\right),\quad
\Delta x=a\sqrt{\frac1{12}-\frac1{2n^2\pi^2}}.
\]
$\int\psi\psi'\dd x=0$给$\langle p\rangle=0$, $-\hbar^2\psi\prime\prime=(n\pi\hbar/a)^2\psi$给$\langle p^2\rangle=(n\pi\hbar/a)^2$, 因而$\Delta p=n\pi\hbar/a$. 乘积$\hbar\sqrt{n^2\pi^2/12-1/2}$, 最小$n=1$约$0.568\hbar>\hbar/2$, 不饱和. 定态仅乘整体相位, 全部矩不随时间变. 这里的动量矩可由全线延零波函数的动量分布解释; 不把盒内零边界的一阶导数算符误称为自伴动量观测量.
\end{solution}
\originalsection{第4、5题：半线硬墙与上台阶（各5、10分）}
\begin{exercise}
第4题: $x\leq0$为硬墙,$x>0$自由, 求不可归一化定态与能量. 第5题: 全线$V=0$（$x\leq0$）、$V_0$（$x>0$）, 求$0<E<V_0$的定态.
\end{exercise}
\begin{solution}
第4题零边界迫使$u_k(x)=A\sin kx$（$x>0$）, $E=\hbar^2k^2/(2m),k>0$. 不存在负能正规解; $E=0$线性解不提供可归一态. 第5题以单位左入射写$u=\e^{\ii kx}+r\e^{-\ii kx}$（左）,$u=t\e^{-\kappa x}$（右）, $k=\sqrt{2mE}/\hbar$, $\kappa=\sqrt{2m(V_0-E)}/\hbar$. 连续条件$1+r=t$, $\ii k(1-r)=-\kappa t$给$r=(k-\ii\kappa)/(k+\ii\kappa)$,$t=2k/(k+\ii\kappa)$. $|r|=1$, 右侧有概率尾巴却无净流. 乘任意复常数得到同一能量散射解, 没有从右衰减支路来的独立入射态.
\end{solution}
\originalsection{第6题：硬墙旁的半势阱（10分）}
\begin{exercise}
$x<0$硬墙,$0<x<a$势为$-V_0$,$x>a$为0. 求$E<0$束缚态, 判断是否总有束缚态; 令$z_0^2=2ma^2V_0/\hbar^2$, 求三束缚态阈值并联系宽$2a$全线势阱.
\end{exercise}
\begin{solution}
令$k=\sqrt{2m(E+V_0)}/\hbar$, $\kappa=\sqrt{-2mE}/\hbar$, 则内部$A\sin kx$, 外部$A\sin(ka)\e^{-\kappa(x-a)}$. 匹配导数得$k\cot ka=-\kappa$. 写$z=ka,w=\kappa a$, 有$z\cot z=-\sqrt{z_0^2-z^2}$, $E=\hbar^2z^2/(2ma^2)-V_0$. 根分别出现在$((n-1/2)\pi,n\pi)$与$z<z_0$的交集. 第一根需$z_0>\pi/2$, 所以并非总有束缚态. 第三个阈值为$5\pi/2$, 严格有三个需$z_0>5\pi/2$且$z_0\leq7\pi/2$（下一态在上端恰为零能、不正规）. 在等号处新态$E=0$的外尾为常数, 不可归一; 因而没有取到的最小值, 只有下确界. 将半线态奇延拓到负轴, 正是宽$2a$有限阱的\textbf{奇}态; 偶态没有对应, 不能把全阱所有束缚态数搬过来.
\end{solution}
\originalsection{第7题：模拟氢能量的方阱（5分）}
\begin{exercise}
全线$|x|<a_0$势$-V_0$, 外为0, 要使基态为$E_0=-e^2/(2a_0)=-13.6$ eV, 求$V_0$.
\end{exercise}
\begin{solution}
氢关系$|E_0|=\hbar^2/(2ma_0^2)$给$\kappa a_0=1$. 偶基态条件$z\tan z=1$, $0<z<\pi/2$, 数值根$z=0.860334$, 故$V_0=|E_0|(1+z^2)\simeq23.67$ eV. 此处只匹配一个能量与长度, 一维方阱没有氢的整个谱或角动量结构.
\end{solution}
\originalsection{第8题：为何不能处处禁阻（5分）}
\begin{exercise}
实定态满足$-\hbar^2\psi''/(2m)+(V-E)\psi=0$. (a)从$\langle H\rangle$证明能量高于势最小值. (b)用波函数曲率说明处处$E<V$不能形成正规束缚态.
\end{exercise}
\begin{solution}
在正规归一化态和消失边界项下,$E=\hbar^2\int|\psi'|^2/(2m)+\int V|\psi|^2\geq V_{\min}$. 全线正规态动能不可能为0, 因为那会要求非零常函数, 不可归一, 所以严格大于$V_{\min}$. 若处处$V-E>0$, 则$\psi''$与$\psi$同号: 正的段向上凸, 负的段向下凸. 一条两端衰减的非零连续曲线必在正最大处有非正曲率, 或在负最小处有非负曲率, 与方程冲突. 这个几何论证与积分正性是同一限制, 不能由禁阻区指数解在单个区间存在推出整个束缚态存在.
\end{solution}

\originalchapter{8.04 作业9：相移、共振与角动量}
来源为8.04 Spring 2016 Assignment 9, 2016年4月21日发布, 4月29日中午截止. 全部六题共65分. 以下题解、数值积分及重绘曲线为\textbf{AI独立解答}; 数值代码与记录随源码保存. 曲线展示所给模型, 不是证明的替代品.
\originalsection{第1题：检验定相位法（10分）}
\begin{exercise}
令$f(\lambda)=\int_{-\infty}^\infty\e^{-100(x-2)^2}\e^{\ii\phi}\dd x$, $\phi=50(x-\lambda x^4/32)$,$0\leq\lambda\leq1$. (a)求Gaussian全半高宽并判断积分$[1,3]$截断. (b)在$x=2$令相位驻定, 求$\lambda_*$及二阶展开. (c)求半高区域相位变化, 忽略变化求积分. (d)保留二阶求积分. (e)数值绘制$|f|$, 比较临界点及实际最大值.
\end{exercise}
\begin{solution}
半高条件$100(x-2)^2=\ln2$给全宽$\Delta=2\sqrt{\ln2}/10=0.16651$. $[1,3]$外包络至多$\e^{-100}$级, 截断误差绝对值由Gaussian尾积分控制, 足够安全. $\partial_x\phi|_2=50(1-\lambda)$给$\lambda_*=1$. 令$s=x-2$, 则$\phi(1,x)=75-37.5s^2+O(s^3)$. 半高区域相位主变化$37.5\ln2/100=0.25993$弧度, 为$0.08274\pi$. 若保留三四次项, 半高区精确最大变化约0.26722弧度, 二阶估计与它接近. 忽略变化给$f\simeq(\sqrt\pi/10)\e^{75\ii}$, 幅值0.177245. 保留二阶给
\[
f\simeq\sqrt{\frac\pi{100+37.5\ii}}\e^{75\ii}
=\frac{\sqrt\pi}{(100^2+37.5^2)^{1/4}}\e^{\ii[75-\tfrac12\arctan(0.375)]}.
\]
幅值约0.17151. 直接积分$f(1)=0.14367073-0.09363253\ii$, 幅值0.17148857; 实际最大在$\lambda=0.99161278$, 幅值0.17155529. 驻定相位给峰位置的近似, 包络有有限宽度、相位有三四阶项, 不要求精确最大恰在1. 下图取$[1,3]$复积分分别对实虚部自适应积分, 绝对容差$10^{-12}$; 峰值另作一维极值搜索.
\begin{center}\begin{tikzpicture}
\begin{axis}[width=0.78\textwidth,height=4.5cm,xlabel={$\lambda$},ylabel={$|f(\lambda)|$},grid=major,legend style={font=\small},]
\addplot+[mark=none] coordinates{(0,0.00034216409) (0.0066666667,0.00037145084) (0.013333333,0.00040306783) (0.02,0.00043718292) (0.026666667,0.00047397437) (0.033333333,0.00051363138) (0.04,0.00055635457) (0.046666667,0.0006023565) (0.053333333,0.00065186219) (0.06,0.00070510966) (0.066666667,0.00076235046) (0.073333333,0.00082385019) (0.08,0.00088988909) (0.086666667,0.0009607625) (0.093333333,0.0010367815) (0.1,0.0011182733) (0.10666667,0.0012055819) (0.11333333,0.0012990686) (0.12,0.0013991123) (0.12666667,0.0015061104) (0.13333333,0.0016204786) (0.14,0.0017426521) (0.14666667,0.0018730855) (0.15333333,0.0020122532) (0.16,0.0021606501) (0.16666667,0.0023187916) (0.17333333,0.0024872139) (0.18,0.0026664742) (0.18666667,0.0028571511) (0.19333333,0.0030598444) (0.2,0.0032751753) (0.20666667,0.0035037863) (0.21333333,0.0037463412) (0.22,0.0040035248) (0.22666667,0.0042760428) (0.23333333,0.0045646214) (0.24,0.0048700066) (0.24666667,0.0051929641) (0.25333333,0.0055342786) (0.26,0.0058947527) (0.26666667,0.0062752064) (0.27333333,0.0066764759) (0.28,0.0070994129) (0.28666667,0.0075448831) (0.29333333,0.0080137646) (0.3,0.0085069474) (0.30666667,0.0090253306) (0.31333333,0.0095698218) (0.32,0.010141335) (0.32666667,0.010740787) (0.33333333,0.011369098) (0.34,0.012027189) (0.34666667,0.012715976) (0.35333333,0.013436372) (0.36,0.014189282) (0.36666667,0.0149756) (0.37333333,0.015796209) (0.38,0.016651973) (0.38666667,0.017543739) (0.39333333,0.01847233) (0.4,0.019438546) (0.40666667,0.020443156) (0.41333333,0.021486897) (0.42,0.022570472) (0.42666667,0.023694541) (0.43333333,0.024859726) (0.44,0.026066598) (0.44666667,0.027315681) (0.45333333,0.028607446) (0.46,0.029942303) (0.46666667,0.031320604) (0.47333333,0.032742637) (0.48,0.034208621) (0.48666667,0.035718705) (0.49333333,0.037272961) (0.5,0.038871386) (0.50666667,0.040513896) (0.51333333,0.042200321) (0.52,0.043930408) (0.52666667,0.045703812) (0.53333333,0.047520098) (0.54,0.049378737) (0.54666667,0.051279104) (0.55333333,0.053220477) (0.56,0.055202034) (0.56666667,0.057222855) (0.57333333,0.059281918) (0.58,0.061378097) (0.58666667,0.063510166) (0.59333333,0.065676798) (0.6,0.06787656) (0.60666667,0.070107922) (0.61333333,0.07236925) (0.62,0.074658814) (0.62666667,0.076974782) (0.63333333,0.07931523) (0.64,0.08167814) (0.64666667,0.084061401) (0.65333333,0.086462817) (0.66,0.088880104) (0.66666667,0.091310899) (0.67333333,0.093752762) (0.68,0.09620318) (0.68666667,0.09865957) (0.69333333,0.10111929) (0.7,0.10357963) (0.70666667,0.10603783) (0.71333333,0.1084911) (0.72,0.11093659) (0.72666667,0.1133714) (0.73333333,0.11579264) (0.74,0.11819736) (0.74666667,0.12058262) (0.75333333,0.12294545) (0.76,0.1252829) (0.76666667,0.12759199) (0.77333333,0.12986977) (0.78,0.13211333) (0.78666667,0.13431974) (0.79333333,0.13648612) (0.8,0.13860965) (0.80666667,0.14068752) (0.81333333,0.14271699) (0.82,0.14469537) (0.82666667,0.14662004) (0.83333333,0.14848845) (0.84,0.15029814) (0.84666667,0.15204669) (0.85333333,0.15373182) (0.86,0.1553513) (0.86666667,0.15690304) (0.87333333,0.15838501) (0.88,0.15979532) (0.88666667,0.16113218) (0.89333333,0.16239391) (0.9,0.16357896) (0.90666667,0.16468591) (0.91333333,0.16571343) (0.92,0.16666036) (0.92666667,0.16752564) (0.93333333,0.16830837) (0.94,0.16900774) (0.94666667,0.16962312) (0.95333333,0.17015399) (0.96,0.17059996) (0.96666667,0.1709608) (0.97333333,0.17123638) (0.98,0.17142673) (0.98666667,0.171532) (0.99333333,0.17155248) (1,0.17148857)};
\end{axis}\end{tikzpicture}\end{center}

\end{solution}
\originalsection{第2题：半势阱的Levinson关系（10分）}
\begin{exercise}
半线硬墙加$0<x<a$的$-V_0$势阱, 令$z_0^2=2mV_0a^2/\hbar^2$, $u=ka$, $v=k\prime a=\sqrt{u^2+z_0^2}$. 给定$\tan\delta=(1-(k'/k)\cot(k'a)\tan ka)/(\tan ka+(k'/k)\cot(k'a))$. (a,b)讨论零能及高能极限. (c)写无量纲函数并连续选相移. (d)画$z_0=2,5,9$曲线, 终点分别15、20、30, 检验束缚态数.
\end{exercise}
\begin{solution}
匹配本质为$k\cot(ka+\delta)=k'\cot(k'a)$, 更稳定地写
\[
f(u;z_0)=\frac{u\sin v\cos u-v\cos v\sin u}{u\sin v\sin u+v\cos v\cos u}.
\]
零能$v\to z_0$, 普通非阈值时$\tan\delta\sim u(\tan z_0/z_0-1)\to0$. 题中三例均可取$\delta(0)=0$. 若$z_0=(n+1/2)\pi$恰在零能共振阈值, 此结论需改为半束缚态版本, 不能忽略该例外. 高能$v=u+z_0^2/(2u)+O(u^{-3})$, 匹配与自由波差趋0, 所以$\tan\delta\to0$; 但相移自身可以趋整数倍$\pi$, 不是必须趋0.

计算连续角度$\alpha(u)=\operatorname{unwrap}\arg[v\cos v+\ii u\sin v]$, 令$\delta=\alpha-u-\alpha(0)$, 即避免主值$\arctan$跳跃. 束缚态在$z_0>(n-1/2)\pi$时增加一个, 因此三例$N=1,2,3$. 所选分支应有$\delta(\infty)=-N\pi$, 满足$[\delta(0)-\delta(\infty)]/\pi=N$. 下图终点数值分别$-0.95647\pi,-1.80221\pi,-2.57827\pi$, 距极限仍有$0.04353\pi,0.19779\pi,0.42173\pi$, 大$z_0$需要更高能量才接近极限. 相移加任意固定整数$\pi$不改散射, 两端差不变.
\begin{center}\begin{tikzpicture}
\begin{axis}[width=0.78\textwidth,height=4.5cm,xlabel={$u=ka$},ylabel={$\delta/\pi$},grid=major,legend style={font=\small},legend pos=south west]
\addplot+[mark=none] coordinates{(1e-07,-3.1830989e-08) (0.0500001,-0.033252259) (0.1000001,-0.066201336) (0.1500001,-0.098565366) (0.2000001,-0.13010134) (0.2500001,-0.16061662) (0.3000001,-0.18997346) (0.3500001,-0.21808734) (0.4000001,-0.24492075) (0.4500001,-0.27047465) (0.5000001,-0.29477921) (0.5500001,-0.31788507) (0.6000001,-0.33985597) (0.6500001,-0.36076275) (0.7000001,-0.3806789) (0.75000009,-0.39967743) (0.80000009,-0.41782868) (0.85000009,-0.43519901) (0.90000009,-0.45185009) (0.95000009,-0.46783856) (1.0000001,-0.48321597) (1.0500001,-0.49802896) (1.1000001,-0.51231951) (1.1500001,-0.52612522) (1.2000001,-0.53947969) (1.2500001,-0.55241286) (1.3000001,-0.56495131) (1.3500001,-0.57711862) (1.4000001,-0.58893564) (1.4500001,-0.60042076) (1.5000001,-0.61159015) (1.5500001,-0.622458) (1.6000001,-0.63303667) (1.6500001,-0.64333693) (1.7000001,-0.65336804) (1.7500001,-0.66313798) (1.8000001,-0.67265348) (1.8500001,-0.68192021) (1.9000001,-0.69094283) (1.9500001,-0.69972512) (2.0000001,-0.70827001) (2.0500001,-0.71657971) (2.1000001,-0.72465575) (2.1500001,-0.73249905) (2.2000001,-0.74010998) (2.2500001,-0.74748842) (2.3000001,-0.75463382) (2.3500001,-0.76154524) (2.4000001,-0.76822143) (2.4500001,-0.77466087) (2.5000001,-0.78086181) (2.5500001,-0.78682238) (2.6000001,-0.79254059) (2.6500001,-0.79801446) (2.7000001,-0.80324203) (2.7500001,-0.80822145) (2.8000001,-0.81295108) (2.8500001,-0.81742955) (2.9000001,-0.82165582) (2.9500001,-0.82562931) (3.0000001,-0.82934997) (3.0500001,-0.83281836) (3.1000001,-0.83603574) (3.1500001,-0.83900417) (3.2000001,-0.84172661) (3.2500001,-0.84420695) (3.3000001,-0.84645012) (3.3500001,-0.84846213) (3.4000001,-0.85025015) (3.4500001,-0.85182247) (3.5000001,-0.85318859) (3.5500001,-0.85435916) (3.6000001,-0.85534594) (3.6500001,-0.85616177) (3.7000001,-0.85682047) (3.7500001,-0.85733672) (3.8000001,-0.85772594) (3.8500001,-0.85800415) (3.9000001,-0.85818778) (3.9500001,-0.85829348) (4.0000001,-0.85833794) (4.0500001,-0.85833771) (4.1000001,-0.85830899) (4.1500001,-0.85826743) (4.2000001,-0.85822799) (4.2500001,-0.85820476) (4.3000001,-0.85821085) (4.3500001,-0.85825824) (4.4000001,-0.85835771) (4.4500001,-0.85851879) (4.5000001,-0.85874971) (4.5500001,-0.85905736) (4.6000001,-0.85944733) (4.6500001,-0.85992392) (4.7000001,-0.86049017) (4.7500001,-0.86114795) (4.8000001,-0.86189797) (4.8500001,-0.86273992) (4.9000001,-0.86367249) (4.9500001,-0.8646935) (5.0000001,-0.86579996) (5.0500001,-0.86698816) (5.1000001,-0.86825376) (5.1500001,-0.86959186) (5.2000001,-0.8709971) (5.2500001,-0.8724637) (5.3000001,-0.87398556) (5.3500001,-0.87555631) (5.4000001,-0.87716937) (5.4500001,-0.87881801) (5.5000001,-0.88049541) (5.5500001,-0.88219469) (5.6000001,-0.88390896) (5.6500001,-0.88563138) (5.7000001,-0.88735518) (5.7500001,-0.8890737) (5.8000001,-0.89078041) (5.8500001,-0.89246899) (5.9000001,-0.8941333) (5.9500001,-0.89576746) (6.0000001,-0.89736585) (6.0500001,-0.89892314) (6.1000001,-0.90043433) (6.1500001,-0.90189476) (6.2000001,-0.90330016) (6.2500001,-0.90464666) (6.3000001,-0.90593079) (6.3500001,-0.90714954) (6.4000001,-0.90830037) (6.4500001,-0.90938121) (6.5000001,-0.9103905) (6.5500001,-0.91132718) (6.6000001,-0.91219074) (6.6500001,-0.91298116) (6.7000001,-0.91369899) (6.7500001,-0.91434529) (6.8000001,-0.91492166) (6.8500001,-0.91543019) (6.9000001,-0.91587349) (6.9500001,-0.91625463) (7.0000001,-0.91657714) (7.0500001,-0.91684494) (7.1000001,-0.91706236) (7.1500001,-0.91723403) (7.2000001,-0.91736488) (7.2500001,-0.91746007) (7.3000001,-0.91752494) (7.3500001,-0.91756496) (7.4000001,-0.91758564) (7.4500001,-0.91759253) (7.5000001,-0.91759109) (7.55,-0.91758671) (7.6,-0.91758459) (7.65,-0.91758973) (7.7,-0.91760685) (7.75,-0.91764039) (7.8,-0.91769443) (7.85,-0.91777269) (7.9,-0.91787849) (7.95,-0.91801471) (8,-0.91818381) (8.05,-0.9183878) (8.1,-0.91862825) (8.15,-0.91890628) (8.2,-0.91922255) (8.25,-0.9195773) (8.3,-0.91997038) (8.35,-0.9204012) (8.4,-0.92086881) (8.45,-0.92137191) (8.5,-0.92190887) (8.55,-0.92247776) (8.6,-0.92307635) (8.65,-0.92370222) (8.7,-0.92435267) (8.75,-0.92502487) (8.8,-0.9257158) (8.85,-0.92642233) (8.9,-0.92714124) (8.95,-0.92786924) (9,-0.92860298) (9.05,-0.92933913) (9.1,-0.93007437) (9.15,-0.93080542) (9.2,-0.93152906) (9.25,-0.93224219) (9.3,-0.9329418) (9.35,-0.93362502) (9.4,-0.93428917) (9.45,-0.93493172) (9.5,-0.93555036) (9.55,-0.93614296) (9.6,-0.93670767) (9.65,-0.93724285) (9.7,-0.93774712) (9.75,-0.93821939) (9.8,-0.93865882) (9.85,-0.93906486) (9.9,-0.93943725) (9.95,-0.939776) (10,-0.94008141) (10.05,-0.94035406) (10.1,-0.94059482) (10.15,-0.94080478) (10.2,-0.94098532) (10.25,-0.94113804) (10.3,-0.94126476) (10.35,-0.9413675) (10.4,-0.94144847) (10.45,-0.94151001) (10.5,-0.94155462) (10.55,-0.94158487) (10.6,-0.94160344) (10.65,-0.94161303) (10.7,-0.94161639) (10.75,-0.94161623) (10.8,-0.94161524) (10.85,-0.94161606) (10.9,-0.94162121) (10.95,-0.94163313) (11,-0.94165411) (11.05,-0.94168627) (11.1,-0.94173158) (11.15,-0.94179181) (11.2,-0.94186852) (11.25,-0.94196307) (11.3,-0.94207658) (11.35,-0.94220995) (11.4,-0.94236384) (11.45,-0.9425387) (11.5,-0.94273473) (11.55,-0.9429519) (11.6,-0.94318997) (11.65,-0.94344848) (11.7,-0.94372677) (11.75,-0.94402399) (11.8,-0.94433912) (11.85,-0.94467094) (11.9,-0.94501812) (11.95,-0.94537918) (12,-0.94575252) (12.05,-0.94613644) (12.1,-0.94652916) (12.15,-0.94692882) (12.2,-0.94733355) (12.25,-0.94774141) (12.3,-0.94815046) (12.35,-0.94855878) (12.4,-0.94896445) (12.45,-0.9493656) (12.5,-0.94976042) (12.55,-0.95014716) (12.6,-0.95052415) (12.65,-0.95088984) (12.7,-0.95124278) (12.75,-0.95158163) (12.8,-0.9519052) (12.85,-0.95221244) (12.9,-0.95250245) (12.95,-0.95277449) (13,-0.95302799) (13.05,-0.95326254) (13.1,-0.9534779) (13.15,-0.95367401) (13.2,-0.95385098) (13.25,-0.9540091) (13.3,-0.95414879) (13.35,-0.95427066) (13.4,-0.95437549) (13.45,-0.95446415) (13.5,-0.9545377) (13.55,-0.9545973) (13.6,-0.95464421) (13.65,-0.9546798) (13.7,-0.95470554) (13.75,-0.95472294) (13.8,-0.95473357) (13.85,-0.95473903) (13.9,-0.95474097) (13.95,-0.954741) (14,-0.95474073) (14.05,-0.95474176) (14.1,-0.95474562) (14.15,-0.95475378) (14.2,-0.95476765) (14.25,-0.95478853) (14.3,-0.95481765) (14.35,-0.9548561) (14.4,-0.95490486) (14.45,-0.9549648) (14.5,-0.95503664) (14.55,-0.95512095) (14.6,-0.95521819) (14.65,-0.95532865) (14.7,-0.95545249) (14.75,-0.95558972) (14.8,-0.95574021) (14.85,-0.95590372) (14.9,-0.95607983) (14.95,-0.95626805) (15,-0.95646773)};
\addlegendentry{$z_0=2$}
\addplot+[mark=none] coordinates{(1e-07,-3.1830989e-08) (0.066666766,-0.035533708) (0.13333343,-0.070863755) (0.2000001,-0.10579805) (0.26666677,-0.14016595) (0.33333343,-0.17382546) (0.4000001,-0.20666719) (0.46666676,-0.23861497) (0.53333343,-0.26962395) (0.6000001,-0.29967677) (0.66666676,-0.32877877) (0.73333343,-0.35695299) (0.8000001,-0.38423544) (0.86666676,-0.41067083) (0.93333343,-0.43630919) (1.0000001,-0.46120304) (1.0666668,-0.48540535) (1.1333334,-0.50896799) (1.2000001,-0.53194073) (1.2666668,-0.55437055) (1.3333334,-0.57630125) (1.4000001,-0.59777325) (1.4666668,-0.61882357) (1.5333334,-0.63948589) (1.6000001,-0.65979067) (1.6666668,-0.67976533) (1.7333334,-0.69943443) (1.8000001,-0.71881987) (1.8666668,-0.73794112) (1.9333334,-0.75681539) (2.0000001,-0.7754578) (2.0666668,-0.7938816) (2.1333334,-0.81209829) (2.2000001,-0.83011778) (2.2666668,-0.84794854) (2.3333334,-0.8655977) (2.4000001,-0.88307113) (2.4666668,-0.90037361) (2.5333334,-0.91750886) (2.6000001,-0.93447959) (2.6666668,-0.95128764) (2.7333334,-0.96793395) (2.8000001,-0.98441869) (2.8666668,-1.0007412) (2.9333334,-1.0169001) (3.0000001,-1.0328934) (3.0666668,-1.0487181) (3.1333334,-1.0643708) (3.2000001,-1.0798474) (3.2666668,-1.0951428) (3.3333334,-1.1102516) (3.4000001,-1.1251675) (3.4666667,-1.1398834) (3.5333334,-1.1543917) (3.6000001,-1.1686838) (3.6666667,-1.1827505) (3.7333334,-1.1965817) (3.8000001,-1.2101665) (3.8666667,-1.2234931) (3.9333334,-1.2365489) (4.0000001,-1.2493202) (4.0666667,-1.2617925) (4.1333334,-1.2739505) (4.2000001,-1.2857775) (4.2666667,-1.2972564) (4.3333334,-1.3083687) (4.4000001,-1.3190953) (4.4666667,-1.3294164) (4.5333334,-1.3393113) (4.6000001,-1.3487589) (4.6666667,-1.3577381) (4.7333334,-1.3662275) (4.8000001,-1.3742064) (4.8666667,-1.381655) (4.9333334,-1.3885548) (5.0000001,-1.39489) (5.0666667,-1.4006472) (5.1333334,-1.4058175) (5.2000001,-1.4103966) (5.2666667,-1.4143863) (5.3333334,-1.4177951) (5.4000001,-1.4206398) (5.4666667,-1.4229458) (5.5333334,-1.4247477) (5.6000001,-1.4260895) (5.6666667,-1.4270245) (5.7333334,-1.4276145) (5.8000001,-1.4279285) (5.8666667,-1.4280407) (5.9333334,-1.4280288) (6.0000001,-1.4279715) (6.0666667,-1.4279457) (6.1333334,-1.4280245) (6.2000001,-1.4282751) (6.2666667,-1.4287568) (6.3333334,-1.4295203) (6.4000001,-1.4306063) (6.4666667,-1.432046) (6.5333334,-1.4338611) (6.6000001,-1.4360641) (6.6666667,-1.4386595) (6.7333334,-1.4416445) (6.8000001,-1.4450103) (6.8666667,-1.4487425) (6.9333334,-1.452823) (7.0000001,-1.4572301) (7.0666667,-1.4619395) (7.1333334,-1.4669253) (7.2000001,-1.4721599) (7.2666667,-1.4776152) (7.3333334,-1.4832622) (7.4000001,-1.4890722) (7.4666667,-1.4950159) (7.5333334,-1.5010647) (7.6000001,-1.5071898) (7.6666667,-1.5133631) (7.7333334,-1.5195568) (7.8000001,-1.5257435) (7.8666667,-1.5318962) (7.9333334,-1.5379886) (8.0000001,-1.5439949) (8.0666667,-1.5498897) (8.1333334,-1.5556483) (8.2000001,-1.5612469) (8.2666667,-1.5666623) (8.3333334,-1.5718723) (8.4000001,-1.5768557) (8.4666667,-1.5815926) (8.5333334,-1.5860646) (8.6000001,-1.5902552) (8.6666667,-1.5941499) (8.7333334,-1.5977364) (8.8000001,-1.6010053) (8.8666667,-1.6039503) (8.9333334,-1.6065688) (9.0000001,-1.6088616) (9.0666667,-1.610834) (9.1333334,-1.6124957) (9.2000001,-1.6138607) (9.2666667,-1.614948) (9.3333334,-1.6157807) (9.4000001,-1.6163866) (9.4666667,-1.616797) (9.5333334,-1.6170468) (9.6000001,-1.6171732) (9.6666667,-1.6172153) (9.7333334,-1.6172129) (9.8000001,-1.6172057) (9.8666667,-1.6172322) (9.9333334,-1.6173288) (10,-1.6175291) (10.066667,-1.6178632) (10.133333,-1.6183571) (10.2,-1.6190325) (10.266667,-1.6199063) (10.333333,-1.620991) (10.4,-1.6222943) (10.466667,-1.6238199) (10.533333,-1.6255672) (10.6,-1.6275321) (10.666667,-1.6297071) (10.733333,-1.6320817) (10.8,-1.6346431) (10.866667,-1.6373763) (10.933333,-1.6402644) (11,-1.6432892) (11.066667,-1.6464314) (11.133333,-1.6496707) (11.2,-1.6529864) (11.266667,-1.6563572) (11.333333,-1.6597618) (11.4,-1.6631786) (11.466667,-1.6665865) (11.533333,-1.6699645) (11.6,-1.6732921) (11.666667,-1.6765495) (11.733333,-1.6797175) (11.8,-1.6827782) (11.866667,-1.6857143) (11.933333,-1.6885102) (12,-1.6911516) (12.066667,-1.6936258) (12.133333,-1.6959221) (12.2,-1.6980317) (12.266667,-1.6999481) (12.333333,-1.7016674) (12.4,-1.7031881) (12.466667,-1.7045115) (12.533333,-1.7056417) (12.6,-1.7065858) (12.666667,-1.7073536) (12.733333,-1.7079578) (12.8,-1.7084138) (12.866667,-1.7087392) (12.933333,-1.708954) (13,-1.7090796) (13.066667,-1.709139) (13.133333,-1.7091557) (13.2,-1.7091538) (13.266667,-1.709157) (13.333333,-1.7091883) (13.4,-1.7092696) (13.466667,-1.7094209) (13.533333,-1.7096604) (13.6,-1.710004) (13.666667,-1.7104648) (13.733333,-1.7110536) (13.8,-1.7117781) (13.866667,-1.7126433) (13.933333,-1.7136517) (14,-1.714803) (14.066667,-1.7160946) (14.133333,-1.7175215) (14.2,-1.7190768) (14.266667,-1.7207518) (14.333333,-1.7225361) (14.4,-1.7244179) (14.466667,-1.7263844) (14.533333,-1.7284218) (14.6,-1.7305155) (14.666667,-1.7326506) (14.733333,-1.7348115) (14.8,-1.7369829) (14.866667,-1.7391491) (14.933333,-1.7412949) (15,-1.7434054) (15.066667,-1.745466) (15.133333,-1.7474631) (15.2,-1.7493837) (15.266667,-1.7512159) (15.333333,-1.752949) (15.4,-1.7545735) (15.466667,-1.7560812) (15.533333,-1.7574659) (15.6,-1.7587227) (15.666667,-1.7598487) (15.733333,-1.7608429) (15.8,-1.7617061) (15.866667,-1.7624413) (15.933333,-1.7630534) (16,-1.7635492) (16.066667,-1.7639375) (16.133333,-1.7642288) (16.2,-1.764435) (16.266667,-1.7645696) (16.333333,-1.764647) (16.4,-1.7646824) (16.466667,-1.7646917) (16.533333,-1.7646909) (16.6,-1.7646958) (16.666667,-1.7647217) (16.733333,-1.7647832) (16.8,-1.7648939) (16.866667,-1.765066) (16.933333,-1.7653104) (17,-1.7656362) (17.066667,-1.7660507) (17.133333,-1.7665594) (17.2,-1.767166) (17.266667,-1.7678723) (17.333333,-1.7686783) (17.4,-1.7695822) (17.466667,-1.7705808) (17.533333,-1.771669) (17.6,-1.7728407) (17.666667,-1.7740885) (17.733333,-1.7754038) (17.8,-1.7767772) (17.866667,-1.7781986) (17.933333,-1.7796573) (18,-1.781142) (18.066667,-1.7826415) (18.133333,-1.7841441) (18.2,-1.7856384) (18.266667,-1.7871132) (18.333333,-1.7885574) (18.4,-1.7899607) (18.466667,-1.7913132) (18.533333,-1.7926056) (18.6,-1.7938299) (18.666667,-1.7949787) (18.733333,-1.7960459) (18.8,-1.7970265) (18.866667,-1.7979168) (18.933333,-1.7987146) (19,-1.7994189) (19.066667,-1.8000303) (19.133333,-1.8005509) (19.2,-1.800984) (19.266667,-1.8013347) (19.333333,-1.8016091) (19.4,-1.8018147) (19.466667,-1.8019601) (19.533333,-1.8020547) (19.6,-1.802109) (19.666667,-1.8021338) (19.733333,-1.8021404) (19.8,-1.8021401) (19.866667,-1.8021443) (19.933333,-1.802164) (20,-1.8022096)};
\addlegendentry{$z_0=5$}
\addplot+[mark=none] coordinates{(1e-07,-3.1830989e-08) (0.1000001,-0.033428282) (0.2000001,-0.066841683) (0.3000001,-0.10022549) (0.4000001,-0.13356518) (0.5000001,-0.16684651) (0.6000001,-0.20005564) (0.7000001,-0.2331792) (0.8000001,-0.26620435) (0.9000001,-0.29911884) (1.0000001,-0.33191104) (1.1000001,-0.36456993) (1.2000001,-0.39708517) (1.3000001,-0.42944703) (1.4000001,-0.46164639) (1.5000001,-0.49367471) (1.6000001,-0.52552395) (1.7000001,-0.55718654) (1.8000001,-0.5886553) (1.9000001,-0.61992335) (2.0000001,-0.65098403) (2.1000001,-0.68183081) (2.2000001,-0.71245721) (2.3000001,-0.74285664) (2.4000001,-0.77302235) (2.5000001,-0.80294729) (2.6000001,-0.83262396) (2.7000001,-0.86204433) (2.8000001,-0.89119968) (2.9000001,-0.92008046) (3.0000001,-0.9486761) (3.1000001,-0.97697489) (3.2000001,-1.0049638) (3.3000001,-1.0326281) (3.4000001,-1.0599514) (3.5000001,-1.0869153) (3.6000001,-1.113499) (3.7000001,-1.139679) (3.8000001,-1.1654289) (3.9000001,-1.1907188) (4.0000001,-1.2155151) (4.1000001,-1.2397795) (4.2000001,-1.2634692) (4.3000001,-1.2865357) (4.4000001,-1.3089244) (4.5000001,-1.3305739) (4.6000001,-1.3514156) (4.7000001,-1.3713728) (4.8000001,-1.3903611) (4.9000001,-1.4082875) (5.0000001,-1.4250522) (5.1000001,-1.4405491) (5.2000001,-1.4546694) (5.3000001,-1.4673064) (5.4000001,-1.4783621) (5.5000001,-1.4877576) (5.6000001,-1.4954461) (5.7000001,-1.5014282) (5.8000001,-1.5057685) (5.9000001,-1.5086099) (6.0000001,-1.5101818) (6.1000001,-1.5107972) (6.2000001,-1.5108373) (6.3000001,-1.5107218) (6.4000001,-1.5108718) (6.5000001,-1.5116725) (6.6000001,-1.5134425) (6.7000001,-1.5164168) (6.8000001,-1.5207427) (6.9000001,-1.5264873) (7.0000001,-1.5336513) (7.1000001,-1.5421854) (7.2000001,-1.5520058) (7.3000001,-1.5630074) (7.4000001,-1.5750742) (7.5000001,-1.5880866) (7.6000001,-1.601926) (7.7000001,-1.6164779) (7.8000001,-1.6316338) (7.9000001,-1.6472912) (8.0000001,-1.6633538) (8.1000001,-1.6797315) (8.2000001,-1.6963388) (8.3000001,-1.713095) (8.4000001,-1.7299226) (8.5000001,-1.7467468) (8.6000001,-1.7634943) (8.7000001,-1.7800928) (8.8000001,-1.7964702) (8.9000001,-1.8125535) (9.0000001,-1.8282686) (9.1000001,-1.8435394) (9.2000001,-1.8582879) (9.3000001,-1.8724335) (9.4000001,-1.8858936) (9.5000001,-1.898584) (9.6000001,-1.9104201) (9.7000001,-1.9213187) (9.8000001,-1.9312008) (9.9000001,-1.9399957) (10,-1.9476459) (10.1,-1.9541135) (10.2,-1.9593874) (10.3,-1.9634909) (10.4,-1.9664888) (10.5,-1.9684915) (10.6,-1.9696572) (10.7,-1.9701871) (10.8,-1.9703165) (10.9,-1.9703003) (11,-1.9703954) (11.1,-1.9708425) (11.2,-1.9718512) (11.3,-1.9735887) (11.4,-1.9761744) (11.5,-1.9796802) (11.6,-1.9841343) (11.7,-1.9895275) (11.8,-1.995821) (11.9,-2.0029533) (12,-2.0108472) (12.1,-2.019415) (12.2,-2.0285625) (12.3,-2.0381926) (12.4,-2.0482066) (12.5,-2.0585063) (12.6,-2.0689942) (12.7,-2.0795744) (12.8,-2.0901522) (12.9,-2.1006345) (13,-2.1109297) (13.1,-2.1209477) (13.2,-2.1306) (13.3,-2.1398008) (13.4,-2.1484671) (13.5,-2.1565205) (13.6,-2.1638892) (13.7,-2.1705101) (13.8,-2.1763323) (13.9,-2.1813208) (14,-2.1854602) (14.1,-2.1887594) (14.2,-2.1912545) (14.3,-2.1930119) (14.4,-2.1941284) (14.5,-2.1947295) (14.6,-2.1949646) (14.7,-2.1949994) (14.8,-2.1950066) (14.9,-2.1951557) (15,-2.1956027) (15.1,-2.1964819) (15.2,-2.1979002) (15.3,-2.1999342) (15.4,-2.2026295) (15.5,-2.2060029) (15.6,-2.2100454) (15.7,-2.2147268) (15.8,-2.2199989) (15.9,-2.2258006) (16,-2.2320608) (16.1,-2.2387011) (16.2,-2.245639) (16.3,-2.2527888) (16.4,-2.2600638) (16.5,-2.2673769) (16.6,-2.2746417) (16.7,-2.2817733) (16.8,-2.2886892) (16.9,-2.29531) (17,-2.3015611) (17.1,-2.3073733) (17.2,-2.3126853) (17.3,-2.3174452) (17.4,-2.3216132) (17.5,-2.3251641) (17.6,-2.3280896) (17.7,-2.3304009) (17.8,-2.3321304) (17.9,-2.3333323) (18,-2.3340821) (18.1,-2.3344746) (18.2,-2.3346197) (18.3,-2.3346376) (18.4,-2.3346518) (18.5,-2.3347831) (18.6,-2.3351429) (18.7,-2.3358274) (18.8,-2.3369145) (18.9,-2.3384607) (19,-2.3405006) (19.1,-2.343048) (19.2,-2.3460971) (19.3,-2.3496251) (19.4,-2.3535951) (19.5,-2.3579582) (19.6,-2.3626568) (19.7,-2.3676263) (19.8,-2.3727974) (19.9,-2.3780979) (20,-2.383454) (20.1,-2.3887915) (20.2,-2.3940376) (20.3,-2.3991216) (20.4,-2.4039761) (20.5,-2.4085389) (20.6,-2.4127537) (20.7,-2.4165725) (20.8,-2.4199569) (20.9,-2.4228799) (21,-2.4253281) (21.1,-2.4273027) (21.2,-2.4288214) (21.3,-2.4299182) (21.4,-2.4306435) (21.5,-2.4310627) (21.6,-2.4312538) (21.7,-2.4313041) (21.8,-2.4313055) (21.9,-2.4313506) (22,-2.4315274) (22.1,-2.4319153) (22.2,-2.4325818) (22.3,-2.4335797) (22.4,-2.4349457) (22.5,-2.4367005) (22.6,-2.4388486) (22.7,-2.4413802) (22.8,-2.4442725) (22.9,-2.4474914) (23,-2.4509941) (23.1,-2.4547303) (23.2,-2.4586445) (23.3,-2.4626774) (23.4,-2.4667678) (23.5,-2.4708536) (23.6,-2.4748733) (23.7,-2.4787674) (23.8,-2.4824796) (23.9,-2.4859583) (24,-2.4891578) (24.1,-2.49204) (24.2,-2.4945754) (24.3,-2.4967452) (24.4,-2.4985419) (24.5,-2.4999705) (24.6,-2.5010492) (24.7,-2.5018092) (24.8,-2.5022942) (24.9,-2.5025588) (25,-2.5026664) (25.1,-2.5026864) (25.2,-2.502691) (25.3,-2.5027514) (25.4,-2.5029348) (25.5,-2.5033013) (25.6,-2.503901) (25.7,-2.5047726) (25.8,-2.5059425) (25.9,-2.507424) (26,-2.5092182) (26.1,-2.5113145) (26.2,-2.5136922) (26.3,-2.5163213) (26.4,-2.5191646) (26.5,-2.522179) (26.6,-2.525317) (26.7,-2.5285284) (26.8,-2.5317613) (26.9,-2.5349638) (27,-2.5380851) (27.1,-2.5410766) (27.2,-2.5438938) (27.3,-2.5464969) (27.4,-2.5488523) (27.5,-2.5509339) (27.6,-2.5527243) (27.7,-2.5542154) (27.8,-2.5554095) (27.9,-2.5563196) (28,-2.5569692) (28.1,-2.5573919) (28.2,-2.5576303) (28.3,-2.557734) (28.4,-2.5577581) (28.5,-2.55776) (28.6,-2.5577972) (28.7,-2.5579244) (28.8,-2.558191) (28.9,-2.5586394) (29,-2.5593032) (29.1,-2.5602058) (29.2,-2.5613609) (29.3,-2.5627714) (29.4,-2.5644309) (29.5,-2.5663239) (29.6,-2.5684271) (29.7,-2.5707105) (29.8,-2.5731389) (29.9,-2.5756728) (30,-2.57827)};
\addlegendentry{$z_0=9$}
\end{axis}\end{tikzpicture}\end{center}

\end{solution}
\originalsection{第3题：台阶与硬墙（10分）}
\begin{exercise}
半线$0<x<a$势$V_0>0$, $x>a$势0, 左端硬墙. (a)对$E>V_0$求$\cot\delta$. (b)对$E<V_0$求同式. (c)画$z_0=5$时从零能到高能的相移.
\end{exercise}
\begin{solution}
外部取$\sin(kx+\delta)$, 内部高能取$A\sin(qx)$, $q=\sqrt{2m(E-V_0)}/\hbar$. 令$C=(q/k)\cot qa$, 匹配$\cot(u+\delta)=C$, 解得$\cot\delta=(1+C\cot u)/(\cot u-C)$. 低能内部取$A\sinh\kappa x$, 因而同式只需$C=(\kappa/k)\coth\kappa a$, $\kappa=\sqrt{2m(V_0-E)}/\hbar$. 在$E=V_0$内解为线性, 连续极限$C=1/u$. 绘图先取$u+\delta$的连续辐角再减$u$, 不把不同主值支路拼接. $\delta(0)=0$, 中间为负（排斥使外波等效向外平移）, $u\to\infty$时从负侧趋0, 无束缚态. 下图画至35; 无限区间的末段由$\delta\sim-z_0^2/(2u)$描述.
\begin{center}\begin{tikzpicture}
\begin{axis}[width=0.78\textwidth,height=4.5cm,xlabel={$u=ka$},ylabel={$\delta/\pi$},grid=major,legend style={font=\small},]
\addplot+[mark=none] coordinates{(1e-07,-2.5465369e-08) (0.1000001,-0.025464971) (0.2000001,-0.050927374) (0.3000001,-0.076384683) (0.4000001,-0.10183433) (0.5000001,-0.12727371) (0.6000001,-0.1527002) (0.7000001,-0.17811111) (0.8000001,-0.2035037) (0.9000001,-0.22887517) (1.0000001,-0.25422263) (1.1000001,-0.27954308) (1.2000001,-0.30483346) (1.3000001,-0.33009054) (1.4000001,-0.35531098) (1.5000001,-0.38049128) (1.6000001,-0.40562775) (1.7000001,-0.43071652) (1.8000001,-0.45575349) (1.9000001,-0.4807343) (2.0000001,-0.50565433) (2.1000001,-0.53050863) (2.2000001,-0.55529188) (2.3000001,-0.57999837) (2.4000001,-0.60462193) (2.5000001,-0.62915585) (2.6000001,-0.65359284) (2.7000001,-0.67792493) (2.8000001,-0.70214334) (2.9000001,-0.72623841) (3.0000001,-0.75019942) (3.1000001,-0.77401443) (3.2000001,-0.79767007) (3.3000001,-0.82115129) (3.4000001,-0.84444103) (3.5000001,-0.86751986) (3.6000001,-0.89036548) (3.7000001,-0.91295211) (3.8000001,-0.93524972) (3.9000001,-0.95722306) (4.0000001,-0.9788303) (4.1000001,-1.0000214) (4.2000001,-1.0207358) (4.3000001,-1.0408992) (4.4000001,-1.0604196) (4.5000001,-1.0791811) (4.6000001,-1.0970353) (4.7000001,-1.1137891) (4.8000001,-1.129186) (4.9000001,-1.1428783) (5.0000001,-1.1543824) (5.1000001,-1.1630092) (5.2000001,-1.1677497) (5.3000001,-1.1670883) (5.4000001,-1.1587085) (5.5000001,-1.1390941) (5.6000001,-1.1033069) (5.7000001,-1.0462851) (5.8000001,-0.96855434) (5.9000001,-0.88372452) (6.0000001,-0.81130121) (6.1000001,-0.75989541) (6.2000001,-0.72700145) (6.3000001,-0.70718603) (6.4000001,-0.69592257) (6.5000001,-0.69007325) (6.6000001,-0.68752774) (6.7000001,-0.68682988) (6.8000001,-0.68692258) (6.9000001,-0.6869888) (7.0000001,-0.68635531) (7.1000001,-0.68443691) (7.2000001,-0.68070953) (7.3000001,-0.67470843) (7.4000001,-0.66605168) (7.5000001,-0.65449022) (7.6000001,-0.63998048) (7.7000001,-0.62276289) (7.8000001,-0.60341167) (7.9000001,-0.58281278) (8.0000001,-0.5620456) (8.1000001,-0.54219354) (8.2000001,-0.52415439) (8.3000001,-0.5085219) (8.4000001,-0.49556428) (8.5000001,-0.48527636) (8.6000001,-0.47746246) (8.7000001,-0.47181594) (8.8000001,-0.46797973) (8.9000001,-0.46558552) (9.0000001,-0.46427582) (9.1000001,-0.46371443) (9.2000001,-0.46359013) (9.3000001,-0.46361694) (9.4000001,-0.4635332) (9.5000001,-0.46310103) (9.6000001,-0.46210709) (9.7000001,-0.46036523) (9.8000001,-0.45772139) (9.9000001,-0.45406089) (10,-0.44931732) (10.1,-0.44348207) (10.2,-0.43661245) (10.3,-0.42883575) (10.4,-0.42034653) (10.5,-0.41139544) (10.6,-0.40226953) (10.7,-0.39326645) (10.8,-0.3846672) (10.9,-0.37671212) (11,-0.36958422) (11.1,-0.36340154) (11.2,-0.35821766) (11.3,-0.35402852) (11.4,-0.3507826) (11.5,-0.34839209) (11.6,-0.34674356) (11.7,-0.34570703) (11.8,-0.34514332) (11.9,-0.34490972) (12,-0.34486426) (12.1,-0.34486916) (12.2,-0.34479357) (12.3,-0.34451618) (12.4,-0.34392786) (12.5,-0.34293441) (12.6,-0.34145969) (12.7,-0.33944883) (12.8,-0.33687145) (12.9,-0.33372444) (13,-0.33003377) (13.1,-0.32585491) (13.2,-0.32127111) (13.3,-0.31638952) (13.4,-0.3113349) (13.5,-0.30624166) (13.6,-0.30124482) (13.7,-0.29647102) (13.8,-0.29203076) (13.9,-0.28801241) (14,-0.28447865) (14.1,-0.28146525) (14.2,-0.27898189) (14.3,-0.27701451) (14.4,-0.27552862) (14.5,-0.27447308) (14.6,-0.27378391) (14.7,-0.27338794) (14.8,-0.27320614) (14.9,-0.2731566) (15,-0.27315716) (15.1,-0.27312787) (15.2,-0.27299318) (15.3,-0.27268408) (15.4,-0.2721402) (15.5,-0.27131176) (15.6,-0.27016155) (15.7,-0.26866654) (15.8,-0.26681933) (15.9,-0.26462894) (16,-0.26212095) (16.1,-0.25933674) (16.2,-0.25633176) (16.3,-0.25317281) (16.4,-0.24993445) (16.5,-0.2466948) (16.6,-0.24353112) (16.7,-0.24051551) (16.8,-0.23771113) (16.9,-0.23516926) (17,-0.23292731) (17.1,-0.23100784) (17.2,-0.22941859) (17.3,-0.22815324) (17.4,-0.22719277) (17.5,-0.22650726) (17.6,-0.22605785) (17.7,-0.22579881) (17.8,-0.22567961) (17.9,-0.22564694) (18,-0.22564652) (18.1,-0.22562501) (18.2,-0.22553162) (18.3,-0.2253198) (18.4,-0.22494874) (18.5,-0.2243848) (18.6,-0.22360278) (18.7,-0.22258697) (18.8,-0.22133197) (18.9,-0.21984312) (19,-0.21813641) (19.1,-0.216238) (19.2,-0.21418311) (19.3,-0.21201435) (19.4,-0.20977967) (19.5,-0.20752988) (19.6,-0.20531604) (19.7,-0.2031869) (19.8,-0.20118643) (19.9,-0.19935189) (20,-0.19771228) (20.1,-0.19628734) (20.2,-0.19508718) (20.3,-0.19411235) (20.4,-0.19335439) (20.5,-0.19279667) (20.6,-0.19241556) (20.7,-0.19218165) (20.8,-0.1920612) (20.9,-0.19201747) (21,-0.1920122) (21.1,-0.19200693) (21.2,-0.19196433) (21.3,-0.19184953) (21.4,-0.19163126) (21.5,-0.19128296) (21.6,-0.19078378) (21.7,-0.19011939) (21.8,-0.18928258) (21.9,-0.18827365) (22,-0.18710049) (22.1,-0.18577829) (22.2,-0.18432895) (22.3,-0.18278013) (22.4,-0.18116395) (22.5,-0.17951544) (22.6,-0.17787091) (22.7,-0.17626614) (22.8,-0.17473472) (22.9,-0.17330649) (23,-0.1720063) (23.1,-0.17085301) (23.2,-0.16985893) (23.3,-0.16902953) (23.4,-0.16836351) (23.5,-0.16785321) (23.6,-0.16748514) (23.7,-0.16724082) (23.8,-0.16709767) (23.9,-0.16702999) (24,-0.16701004) (24.1,-0.1670091) (24.2,-0.1669985) (24.3,-0.16695064) (24.4,-0.16684001) (24.5,-0.16664406) (24.6,-0.16634404) (24.7,-0.16592569) (24.8,-0.16537983) (24.9,-0.16470272) (25,-0.16389627) (25.1,-0.16296797) (25.2,-0.16193066) (25.3,-0.16080191) (25.4,-0.15960332) (25.5,-0.1583595) (25.6,-0.15709697) (25.7,-0.15584287) (25.8,-0.15462377) (25.9,-0.15346443) (26,-0.15238674) (26.1,-0.15140878) (26.2,-0.15054419) (26.3,-0.14980167) (26.4,-0.14918481) (26.5,-0.14869215) (26.6,-0.14831737) (26.7,-0.14804978) (26.8,-0.14787485) (26.9,-0.14777493) (27,-0.14772997) (27.1,-0.14771838) (27.2,-0.14771783) (27.3,-0.14770609) (27.4,-0.14766184) (27.5,-0.14756546) (27.6,-0.14739974) (27.7,-0.1471505) (27.8,-0.14680715) (27.9,-0.1463631) (28,-0.14581601) (28.1,-0.14516796) (28.2,-0.1444253) (28.3,-0.14359845) (28.4,-0.14270146) (28.5,-0.14175136) (28.6,-0.14076742) (28.7,-0.13977032) (28.8,-0.13878117) (28.9,-0.13782055) (29,-0.13690769) (29.1,-0.13605956) (29.2,-0.1352902) (29.3,-0.13461024) (29.4,-0.13402648) (29.5,-0.13354175) (29.6,-0.1331549) (29.7,-0.13286099) (29.8,-0.13265158) (29.9,-0.13251519) (30,-0.13243779) (30.1,-0.13240344) (30.2,-0.13239493) (30.3,-0.1323944) (30.4,-0.13238409) (30.5,-0.13234695) (30.6,-0.13226728) (30.7,-0.13213134) (30.8,-0.1319278) (30.9,-0.13164826) (31,-0.1312875) (31.1,-0.13084373) (31.2,-0.13031865) (31.3,-0.1297174) (31.4,-0.12904835) (31.5,-0.12832272) (31.6,-0.12755413) (31.7,-0.12675799) (31.8,-0.12595081) (31.9,-0.12514949) (32,-0.12437055) (32.1,-0.12362946) (32.2,-0.12293995) (32.3,-0.12231347) (32.4,-0.12175877) (32.5,-0.12128154) (32.6,-0.12088433) (32.7,-0.12056646) (32.8,-0.12032419) (32.9,-0.12015088) (33,-0.12003739) (33.1,-0.11997247) (33.2,-0.11994323) (33.3,-0.11993566) (33.4,-0.11993522) (33.5,-0.11992733) (33.6,-0.11989799) (33.7,-0.11983425) (33.8,-0.11972474) (33.9,-0.11956005) (34,-0.11933313) (34.1,-0.11903953) (34.2,-0.11867762) (34.3,-0.11824861) (34.4,-0.11775652) (34.5,-0.11720801) (34.6,-0.11661211) (34.7,-0.11597984) (34.8,-0.11532369) (34.9,-0.11465712) (35,-0.11399398)};
\end{axis}\end{tikzpicture}\end{center}

\end{solution}
\originalsection{第4题：\texorpdfstring{$\delta$}{δ}障碍与共振（15分）}
\begin{exercise}
硬墙$x\leq0$, $x>0$有$g\delta(x-a)$, $g>0$, 令$\lambda=mag/\hbar^2$. (a)求$\tan\delta=-\sin^2u/h(u;\lambda)$及内振幅$A$. (b)讨论$u\ll1$和$\lambda\to\infty$. (c)在$\lambda=5$, $0<u<10$画相移、$\mathrm d\delta/\mathrm du$及$|A|$, 识别共振并检查Levinson关系.
\end{exercise}
\begin{solution}
令内波$A\sin kx$, 外波$\e^{\ii\delta}\sin(kx+\delta)$. 连续及导数跳跃$\psi\prime(a^+)-\psi\prime(a^-)=2mg\,\psi(a)/\hbar^2$给
\[
\cot(u+\delta)=\cot u+\frac{2\lambda}{u},\quad
\tan\delta=-\frac{\sin^2u}{u/(2\lambda)+\sin u\cos u},\quad
A=\e^{\ii\delta}\frac{\sin(u+\delta)}{\sin u}.
\]
零点处由连续极限解释最后一式. 小$u$时$\tan\delta\sim-[2\lambda/(1+2\lambda)]u$; 强障碍趋$-\tan u$, 对应外波像在$x=a$遇到硬墙, 共振附近极限不一致. 幅值可稳定计算为$|A|=[1+(2\lambda/u)\sin2u+(4\lambda^2/u^2)\sin^2u]^{-1/2}$. 下列曲线用连续$\delta=\operatorname{atan2}(-2\lambda\sin^2u,u+\lambda\sin2u)$, 其值在$[-\pi,0]$内. 在$|A|$峰处得到下表, 时间延迟峰位置略有不同, 不声称三种峰严格同点.
\begin{center}\begin{tabular}{ccc}\toprule $u$（幅值峰）&$|A|$&$\mathrm d\delta/\mathrm du$\\\midrule
2.87894&3.74007&14.12582\\5.84652&2.17045&3.95735\\8.88845&1.70956&2.01919\\\bottomrule\end{tabular}\end{center}
\begin{center}\begin{tikzpicture}
\begin{axis}[width=0.78\textwidth,height=4.5cm,xlabel={$u=ka$},ylabel={$\delta$},grid=major,legend style={font=\small},]
\addplot+[mark=none] coordinates{(1e-05,-9.0909091e-06) (0.0222322,-0.020211063) (0.0444544,-0.040412871) (0.0666766,-0.060614348) (0.0888988,-0.08081533) (0.111121,-0.10101565) (0.1333432,-0.12121514) (0.1555654,-0.14141364) (0.1777876,-0.16161097) (0.2000098,-0.18180697) (0.222232,-0.20200147) (0.2444542,-0.22219429) (0.2666764,-0.24238526) (0.2888986,-0.26257421) (0.3111208,-0.28276095) (0.333343,-0.3029453) (0.3555652,-0.3231271) (0.3777874,-0.34330614) (0.4000096,-0.36348224) (0.4222318,-0.38365521) (0.444454,-0.40382485) (0.4666762,-0.42399098) (0.4888984,-0.44415338) (0.5111206,-0.46431185) (0.5333428,-0.48446619) (0.555565,-0.50461617) (0.5777872,-0.52476158) (0.6000094,-0.54490219) (0.6222316,-0.56503777) (0.6444538,-0.5851681) (0.666676,-0.60529292) (0.6888982,-0.62541199) (0.7111204,-0.64552505) (0.7333426,-0.66563184) (0.7555648,-0.68573209) (0.777787,-0.70582553) (0.8000092,-0.72591186) (0.8222314,-0.74599079) (0.8444536,-0.76606201) (0.8666758,-0.78612521) (0.888898,-0.80618006) (0.9111202,-0.82622622) (0.9333424,-0.84626334) (0.9555646,-0.86629105) (0.9777868,-0.88630898) (1.000009,-0.90631673) (1.0222312,-0.9263139) (1.0444534,-0.94630005) (1.0666756,-0.96627476) (1.0888978,-0.98623755) (1.11112,-1.0061879) (1.1333422,-1.0261254) (1.1555644,-1.0460495) (1.1777866,-1.0659596) (1.2000088,-1.0858552) (1.222231,-1.1057356) (1.2444532,-1.1256002) (1.2666754,-1.1454484) (1.2888976,-1.1652795) (1.3111198,-1.1850926) (1.333342,-1.2048872) (1.3555642,-1.2246622) (1.3777864,-1.244417) (1.4000086,-1.2641505) (1.4222308,-1.2838619) (1.444453,-1.3035501) (1.4666752,-1.323214) (1.4888974,-1.3428526) (1.5111196,-1.3624645) (1.5333418,-1.3820486) (1.555564,-1.4016035) (1.5777862,-1.4211277) (1.6000084,-1.4406196) (1.6222306,-1.4600777) (1.6444528,-1.4795001) (1.666675,-1.498885) (1.6888972,-1.5182304) (1.7111194,-1.537534) (1.7333416,-1.5567936) (1.7555638,-1.5760067) (1.777786,-1.5951704) (1.8000082,-1.614282) (1.8222304,-1.6333382) (1.8444526,-1.6523356) (1.8666748,-1.6712705) (1.888897,-1.6901388) (1.9111192,-1.7089362) (1.9333414,-1.7276578) (1.9555636,-1.7462983) (1.9777858,-1.764852) (2.000008,-1.7833126) (2.0222302,-1.8016729) (2.0444524,-1.8199254) (2.0666746,-1.8380614) (2.0888968,-1.8560714) (2.111119,-1.8739449) (2.1333412,-1.8916701) (2.1555634,-1.9092337) (2.1777856,-1.9266208) (2.2000078,-1.9438146) (2.22223,-1.9607962) (2.2444522,-1.9775438) (2.2666744,-1.9940327) (2.2888966,-2.0102346) (2.3111188,-2.0261168) (2.333341,-2.0416413) (2.3555632,-2.056764) (2.3777854,-2.071433) (2.4000076,-2.0855872) (2.4222298,-2.0991538) (2.444452,-2.1120459) (2.4666742,-2.1241582) (2.4888964,-2.1353626) (2.5111186,-2.1455013) (2.5333408,-2.1543781) (2.555563,-2.1617461) (2.5777852,-2.1672898) (2.6000074,-2.1706013) (2.6222296,-2.1711435) (2.6444518,-2.1681967) (2.666674,-2.1607769) (2.6888962,-2.1475082) (2.7111184,-2.1264197) (2.7333406,-2.0946137) (2.7555628,-2.047722) (2.777785,-1.9790307) (2.8000072,-1.8782261) (2.8222294,-1.730364) (2.8444516,-1.5184891) (2.8666738,-1.2389701) (2.888896,-0.92798322) (2.9111182,-0.64952821) (2.9333404,-0.43892661) (2.9555626,-0.29204707) (2.9777848,-0.19188816) (3.000007,-0.12358984) (3.0222292,-0.076917882) (3.0444514,-0.045212803) (3.0666736,-0.024140512) (3.0888958,-0.010824761) (3.111118,-0.0033080128) (3.1333402,-0.00022322345) (3.1555624,-0.00059219055) (3.1777846,-0.0036991682) (3.2000068,-0.0090097881) (3.222229,-0.016117799) (3.2444512,-0.024709277) (3.2666734,-0.034538057) (3.2888956,-0.045408536) (3.3111178,-0.057163406) (3.33334,-0.069674774) (3.3555622,-0.082837618) (3.3777844,-0.096564903) (3.4000066,-0.1107839) (3.4222288,-0.12543334) (3.444451,-0.14046127) (3.4666732,-0.1558233) (3.4888954,-0.17148129) (3.5111176,-0.18740227) (3.5333398,-0.20355757) (3.555562,-0.21992216) (3.5777842,-0.23647401) (3.6000064,-0.25319372) (3.6222286,-0.27006407) (3.6444508,-0.28706972) (3.666673,-0.30419698) (3.6888952,-0.32143351) (3.7111174,-0.33876821) (3.7333396,-0.35619102) (3.7555618,-0.3736928) (3.777784,-0.3912652) (3.8000062,-0.40890056) (3.8222284,-0.42659185) (3.8444506,-0.44433257) (3.8666728,-0.46211668) (3.888895,-0.47993857) (3.9111172,-0.49779298) (3.9333394,-0.51567499) (3.9555616,-0.53357993) (3.9777838,-0.55150341) (4.000006,-0.56944123) (4.0222282,-0.58738939) (4.0444504,-0.60534403) (4.0666726,-0.62330146) (4.0888948,-0.64125808) (4.111117,-0.6592104) (4.1333392,-0.677155) (4.1555614,-0.69508852) (4.1777836,-0.71300765) (4.2000058,-0.73090911) (4.222228,-0.74878964) (4.2444502,-0.76664597) (4.2666724,-0.78447481) (4.2888946,-0.80227286) (4.3111168,-0.82003678) (4.333339,-0.83776315) (4.3555612,-0.8554485) (4.3777834,-0.87308929) (4.4000056,-0.89068184) (4.4222278,-0.9082224) (4.44445,-0.92570707) (4.4666722,-0.9431318) (4.4888944,-0.96049239) (4.5111166,-0.97778443) (4.5333388,-0.99500334) (4.555561,-1.0121443) (4.5777832,-1.0292022) (4.6000054,-1.0461718) (4.6222276,-1.0630473) (4.6444498,-1.0798228) (4.666672,-1.096492) (4.6888942,-1.113048) (4.7111164,-1.1294838) (4.7333386,-1.1457915) (4.7555608,-1.161963) (4.777783,-1.1779894) (4.8000052,-1.1938613) (4.8222274,-1.2095684) (4.8444496,-1.2250998) (4.8666718,-1.2404436) (4.888894,-1.2555868) (4.9111162,-1.2705156) (4.9333384,-1.2852147) (4.9555606,-1.2996676) (4.9777828,-1.3138562) (5.000005,-1.3277608) (5.0222272,-1.3413597) (5.0444494,-1.3546293) (5.0666716,-1.3675432) (5.0888938,-1.3800727) (5.111116,-1.3921858) (5.1333382,-1.4038471) (5.1555604,-1.4150172) (5.1777826,-1.4256523) (5.2000048,-1.4357034) (5.222227,-1.4451157) (5.2444492,-1.4538275) (5.2666714,-1.4617694) (5.2888936,-1.4688632) (5.3111158,-1.4750201) (5.333338,-1.4801394) (5.3555602,-1.4841064) (5.3777824,-1.4867901) (5.4000046,-1.4880408) (5.4222268,-1.4876866) (5.444449,-1.48553) (5.4666712,-1.4813438) (5.4888934,-1.4748663) (5.5111156,-1.4657965) (5.5333378,-1.4537885) (5.55556,-1.4384464) (5.5777822,-1.4193203) (5.6000044,-1.3959048) (5.6222266,-1.3676413) (5.6444488,-1.3339297) (5.666671,-1.2941517) (5.6888932,-1.2477142) (5.7111154,-1.194119) (5.7333376,-1.1330623) (5.7555598,-1.0645627) (5.777782,-0.98909971) (5.8000042,-0.90772469) (5.8222264,-0.82209097) (5.8444486,-0.73435634) (5.8666708,-0.64694949) (5.888893,-0.5622539) (5.9111152,-0.48230606) (5.9333374,-0.4085967) (5.9555596,-0.34201163) (5.9777818,-0.28288985) (6.000004,-0.23114706) (6.0222262,-0.18641386) (6.0444484,-0.148157) (6.0666706,-0.11577071) (6.0888928,-0.088638239) (6.111115,-0.066168994) (6.1333372,-0.047818348) (6.1555594,-0.033095921) (6.1777816,-0.021566928) (6.2000038,-0.012849615) (6.222226,-0.006610774) (6.2444482,-0.0025605102) (6.2666704,-0.00044696394) (6.2888926,-5.1328429e-05) (6.3111148,-0.0011833416) (6.333337,-0.0036773017) (6.3555592,-0.0073886018) (6.3777814,-0.012190744) (6.4000036,-0.017972782) (6.4222258,-0.024637147) (6.444448,-0.032097784) (6.4666702,-0.040278579) (6.4888924,-0.049112016) (6.5111146,-0.058538025) (6.5333368,-0.068503015) (6.555559,-0.078959033) (6.5777812,-0.089863049) (6.6000034,-0.10117635) (6.6222256,-0.11286399) (6.6444478,-0.12489437) (6.66667,-0.13723882) (6.6888922,-0.14987123) (6.7111144,-0.16276784) (6.7333366,-0.17590687) (6.7555588,-0.18926837) (6.777781,-0.202834) (6.8000032,-0.21658686) (6.8222254,-0.23051131) (6.8444476,-0.24459287) (6.8666698,-0.25881808) (6.888892,-0.27317439) (6.9111142,-0.28765008) (6.9333364,-0.30223417) (6.9555586,-0.31691632) (6.9777808,-0.3316868) (7.000003,-0.3465364) (7.0222252,-0.36145638) (7.0444474,-0.37643842) (7.0666696,-0.3914746) (7.0888918,-0.40655729) (7.111114,-0.42167919) (7.1333362,-0.43683325) (7.1555584,-0.45201261) (7.1777806,-0.46721066) (7.2000028,-0.48242091) (7.222225,-0.49763701) (7.2444472,-0.51285274) (7.2666694,-0.52806195) (7.2888916,-0.54325855) (7.3111138,-0.55843648) (7.333336,-0.57358971) (7.3555582,-0.58871219) (7.3777804,-0.60379784) (7.4000026,-0.61884053) (7.4222248,-0.63383405) (7.444447,-0.64877208) (7.4666692,-0.6636482) (7.4888914,-0.67845583) (7.5111136,-0.69318823) (7.5333358,-0.70783844) (7.555558,-0.72239931) (7.5777802,-0.73686342) (7.6000024,-0.75122306) (7.6222246,-0.76547022) (7.6444468,-0.77959654) (7.666669,-0.79359328) (7.6888912,-0.80745128) (7.7111134,-0.8211609) (7.7333356,-0.83471202) (7.7555578,-0.84809394) (7.77778,-0.86129536) (7.8000022,-0.8743043) (7.8222244,-0.88710808) (7.8444466,-0.89969318) (7.8666688,-0.91204524) (7.888891,-0.92414893) (7.9111132,-0.93598789) (7.9333354,-0.9475446) (7.9555576,-0.9588003) (7.9777798,-0.96973486) (8.000002,-0.98032666) (8.0222242,-0.99055242) (8.0444464,-1.0003871) (8.0666686,-1.0098036) (8.0888908,-1.0187729) (8.111113,-1.0272632) (8.1333352,-1.0352405) (8.1555574,-1.0426678) (8.1777796,-1.049505) (8.2000018,-1.0557086) (8.222224,-1.0612313) (8.2444462,-1.0660218) (8.2666684,-1.0700245) (8.2888906,-1.0731789) (8.3111128,-1.0754191) (8.333335,-1.076674) (8.3555572,-1.0768662) (8.3777794,-1.0759121) (8.4000016,-1.0737215) (8.4222238,-1.0701971) (8.444446,-1.0652348) (8.4666682,-1.0587237) (8.4888904,-1.0505462) (8.5111126,-1.0405792) (8.5333348,-1.0286949) (8.555557,-1.0147632) (8.5777792,-0.99865406) (8.6000014,-0.9802417) (8.6222236,-0.95940922) (8.6444458,-0.93605516) (8.666668,-0.91010122) (8.6888902,-0.8815014) (8.7111124,-0.8502521) (8.7333346,-0.81640271) (8.7555568,-0.78006569) (8.777779,-0.74142473) (8.8000012,-0.70073966) (8.8222234,-0.65834643) (8.8444456,-0.61465099) (8.8666678,-0.57011665) (8.88889,-0.52524537) (8.9111122,-0.48055475) (8.9333344,-0.43655314) (8.9555566,-0.39371581) (8.9777788,-0.35246501) (9.000001,-0.31315562) (9.0222232,-0.27606734) (9.0444454,-0.24140321) (9.0666676,-0.20929337) (9.0888898,-0.17980254) (9.111112,-0.15293974) (9.1333342,-0.12866884) (9.1555564,-0.10691886) (9.1777786,-0.087593429) (9.2000008,-0.070578879) (9.222223,-0.055750989) (9.2444452,-0.042980314) (9.2666674,-0.032136275) (9.2888896,-0.023090172) (9.3111118,-0.015717317) (9.333334,-0.009898457) (9.3555562,-0.005520661) (9.3777784,-0.0024777847) (9.4000006,-0.00067064016) (9.4222228,-6.9480267e-06) (9.444445,-0.00040114091) (9.4666672,-0.0017740699) (9.4888894,-0.0040526512) (9.5111116,-0.0071694813) (9.5333338,-0.01106244) (9.555556,-0.015674295) (9.5777782,-0.020952319) (9.6000004,-0.026847921) (9.6222226,-0.033316303) (9.6444448,-0.040316133) (9.666667,-0.047809251) (9.6888892,-0.055760385) (9.7111114,-0.064136903) (9.7333336,-0.072908577) (9.7555558,-0.082047371) (9.777778,-0.091527245) (9.8000002,-0.10132398) (9.8222224,-0.11141501) (9.8444446,-0.1217793) (9.8666668,-0.13239717) (9.888889,-0.1432502) (9.9111112,-0.15432112) (9.9333334,-0.16559372) (9.9555556,-0.17705271) (9.9777778,-0.18868369) (10,-0.20047304)};
\end{axis}\end{tikzpicture}\end{center}

\begin{center}\begin{tikzpicture}
\begin{axis}[width=0.78\textwidth,height=4.5cm,xlabel={$u=ka$},ylabel={$\mathrm d\delta/\mathrm du$},grid=major,legend style={font=\small},]
\addplot+[mark=none] coordinates{(1e-05,-0.90909091) (0.013898875,-0.90908946) (0.02778775,-0.90908511) (0.041676625,-0.90907785) (0.0555655,-0.9090677) (0.069454375,-0.90905464) (0.08334325,-0.90903866) (0.097232125,-0.90901977) (0.111121,-0.90899795) (0.12500987,-0.9089732) (0.13889875,-0.9089455) (0.15278763,-0.90891485) (0.1666765,-0.90888123) (0.18056538,-0.90884464) (0.19445425,-0.90880505) (0.20834313,-0.90876246) (0.222232,-0.90871684) (0.23612088,-0.90866818) (0.25000975,-0.90861645) (0.26389862,-0.90856165) (0.2777875,-0.90850374) (0.29167637,-0.90844271) (0.30556525,-0.90837853) (0.31945413,-0.90831118) (0.333343,-0.90824062) (0.34723188,-0.90816684) (0.36112075,-0.9080898) (0.37500963,-0.90800946) (0.3888985,-0.90792581) (0.40278738,-0.9078388) (0.41667625,-0.90774839) (0.43056512,-0.90765456) (0.444454,-0.90755726) (0.45834288,-0.90745646) (0.47223175,-0.9073521) (0.48612063,-0.90724415) (0.5000095,-0.90713256) (0.51389837,-0.90701728) (0.52778725,-0.90689827) (0.54167612,-0.90677546) (0.555565,-0.90664882) (0.56945387,-0.90651828) (0.58334275,-0.90638379) (0.59723162,-0.90624528) (0.6111205,-0.9061027) (0.62500937,-0.90595598) (0.63889825,-0.90580505) (0.65278712,-0.90564985) (0.666676,-0.9054903) (0.68056487,-0.90532632) (0.69445375,-0.90515785) (0.70834263,-0.9049848) (0.7222315,-0.90480708) (0.73612037,-0.90462462) (0.75000925,-0.90443732) (0.76389812,-0.90424509) (0.777787,-0.90404783) (0.79167587,-0.90384545) (0.80556475,-0.90363784) (0.81945362,-0.9034249) (0.8333425,-0.90320651) (0.84723138,-0.90298256) (0.86112025,-0.90275293) (0.87500912,-0.9025175) (0.888898,-0.90227614) (0.90278687,-0.90202871) (0.91667575,-0.90177509) (0.93056462,-0.90151512) (0.9444535,-0.90124866) (0.95834237,-0.90097555) (0.97223125,-0.90069564) (0.98612013,-0.90040875) (1.000009,-0.90011473) (1.0138979,-0.89981338) (1.0277867,-0.89950453) (1.0416756,-0.89918798) (1.0555645,-0.89886352) (1.0694534,-0.89853096) (1.0833423,-0.89819008) (1.0972311,-0.89784064) (1.11112,-0.89748243) (1.1250089,-0.89711518) (1.1388978,-0.89673866) (1.1527866,-0.89635259) (1.1666755,-0.89595671) (1.1805644,-0.89555072) (1.1944533,-0.89513433) (1.2083421,-0.89470723) (1.222231,-0.8942691) (1.2361199,-0.8938196) (1.2500088,-0.89335837) (1.2638976,-0.89288506) (1.2777865,-0.89239927) (1.2916754,-0.89190061) (1.3055642,-0.89138867) (1.3194531,-0.890863) (1.333342,-0.89032316) (1.3472309,-0.88976866) (1.3611198,-0.889199) (1.3750086,-0.88861368) (1.3888975,-0.88801213) (1.4027864,-0.8873938) (1.4166753,-0.88675808) (1.4305641,-0.88610435) (1.444453,-0.88543195) (1.4583419,-0.88474018) (1.4722308,-0.88402834) (1.4861196,-0.88329566) (1.5000085,-0.88254133) (1.5138974,-0.88176454) (1.5277863,-0.88096438) (1.5416751,-0.88013995) (1.555564,-0.87929026) (1.5694529,-0.87841428) (1.5833417,-0.87751094) (1.5972306,-0.8765791) (1.6111195,-0.87561755) (1.6250084,-0.87462503) (1.6388973,-0.87360019) (1.6527861,-0.87254163) (1.666675,-0.87144784) (1.6805639,-0.87031724) (1.6944528,-0.86914817) (1.7083416,-0.86793883) (1.7222305,-0.86668737) (1.7361194,-0.86539177) (1.7500083,-0.86404993) (1.7638971,-0.86265959) (1.777786,-0.86121837) (1.7916749,-0.85972373) (1.8055638,-0.85817295) (1.8194526,-0.85656316) (1.8333415,-0.85489129) (1.8472304,-0.85315405) (1.8611193,-0.85134793) (1.8750081,-0.84946921) (1.888897,-0.84751386) (1.9027859,-0.8454776) (1.9166748,-0.84335582) (1.9305636,-0.8411436) (1.9444525,-0.83883564) (1.9583414,-0.83642622) (1.9722303,-0.83390922) (1.9861191,-0.83127803) (2.000008,-0.82852551) (2.0138969,-0.82564396) (2.0277858,-0.82262506) (2.0416746,-0.8194598) (2.0555635,-0.81613839) (2.0694524,-0.81265023) (2.0833413,-0.80898379) (2.0972301,-0.80512653) (2.111119,-0.80106476) (2.1250079,-0.79678356) (2.1388968,-0.7922666) (2.1527856,-0.787496) (2.1666745,-0.78245216) (2.1805634,-0.77711354) (2.1944522,-0.77145641) (2.2083411,-0.76545464) (2.22223,-0.75907931) (2.2361189,-0.75229845) (2.2500077,-0.74507657) (2.2638966,-0.7373742) (2.2777855,-0.7291474) (2.2916744,-0.72034705) (2.3055633,-0.71091817) (2.3194521,-0.70079905) (2.333341,-0.68992023) (2.3472299,-0.67820333) (2.3611188,-0.66555962) (2.3750076,-0.65188835) (2.3888965,-0.63707477) (2.4027854,-0.62098776) (2.4166743,-0.60347692) (2.4305631,-0.58436913) (2.444452,-0.56346434) (2.4583409,-0.54053041) (2.4722297,-0.51529679) (2.4861186,-0.48744676) (2.5000075,-0.4566077) (2.5138964,-0.42233903) (2.5277852,-0.38411694) (2.5416741,-0.34131514) (2.555563,-0.29318018) (2.5694519,-0.23879983) (2.5833408,-0.17706196) (2.5972296,-0.10660095) (2.6111185,-0.025726842) (2.6250074,0.067668698) (2.6388963,0.17623873) (2.6527851,0.30335206) (2.666674,0.45332671) (2.6805629,0.63175131) (2.6944518,0.84592978) (2.7083406,1.1054987) (2.7222295,1.4232819) (2.7361184,1.8164571) (2.7500072,2.3080961) (2.7638961,2.9290362) (2.777785,3.7196959) (2.7916739,4.7304623) (2.8055627,6.0168401) (2.8194516,7.6204633) (2.8333405,9.5200833) (2.8472294,11.540588) (2.8611183,13.263958) (2.8750071,14.109314) (2.888896,13.708893) (2.9027849,12.24409) (2.9166738,10.271873) (2.9305626,8.2959338) (2.9444515,6.5749965) (2.9583404,5.1747203) (2.9722293,4.0688443) (2.9861181,3.2035947) (3.000007,2.5256365) (3.0138959,1.9907483) (3.0277847,1.5647282) (3.0416736,1.2218599) (3.0555625,0.94298406) (3.0694514,0.71382064) (3.0833403,0.52366804) (3.0972291,0.36444208) (3.111118,0.22997921) (3.1250069,0.11553437) (3.1388958,0.017418793) (3.1527846,-0.067262541) (3.1666735,-0.14080268) (3.1805624,-0.20503382) (3.1944513,-0.26143219) (3.2083401,-0.31119652) (3.222229,-0.35530753) (3.2361179,-0.39457316) (3.2500068,-0.42966342) (3.2638956,-0.4611374) (3.2777845,-0.48946439) (3.2916734,-0.51504043) (3.3055622,-0.53820154) (3.3194511,-0.5592342) (3.33334,-0.57838381) (3.3472289,-0.59586147) (3.3611178,-0.61184958) (3.3750066,-0.62650633) (3.3888955,-0.63996945) (3.4027844,-0.65235928) (3.4166733,-0.66378131) (3.4305621,-0.67432833) (3.444451,-0.68408219) (3.4583399,-0.69311531) (3.4722288,-0.70149196) (3.4861176,-0.7092693) (3.5000065,-0.71649832) (3.5138954,-0.7232246) (3.5277843,-0.72948898) (3.5416731,-0.73532814) (3.555562,-0.74077511) (3.5694509,-0.74585963) (3.5833397,-0.75060862) (3.5972286,-0.75504642) (3.6111175,-0.75919509) (3.6250064,-0.76307468) (3.6388953,-0.76670344) (3.6527841,-0.77009797) (3.666673,-0.77327343) (3.6805619,-0.77624366) (3.6944508,-0.77902134) (3.7083396,-0.78161805) (3.7222285,-0.78404443) (3.7361174,-0.78631024) (3.7500063,-0.78842444) (3.7638951,-0.79039527) (3.777784,-0.79223032) (3.7916729,-0.79393657) (3.8055618,-0.79552044) (3.8194506,-0.79698788) (3.8333395,-0.79834435) (3.8472284,-0.79959489) (3.8611172,-0.80074415) (3.8750061,-0.80179644) (3.888895,-0.80275571) (3.9027839,-0.80362562) (3.9166728,-0.80440954) (3.9305616,-0.80511056) (3.9444505,-0.80573156) (3.9583394,-0.80627515) (3.9722283,-0.80674374) (3.9861171,-0.80713956) (4.000006,-0.80746463) (4.0138949,-0.80772078) (4.0277837,-0.80790971) (4.0416726,-0.80803294) (4.0555615,-0.80809183) (4.0694504,-0.80808762) (4.0833392,-0.80802141) (4.0972281,-0.80789417) (4.111117,-0.80770673) (4.1250059,-0.80745983) (4.1388947,-0.80715408) (4.1527836,-0.80678998) (4.1666725,-0.80636793) (4.1805614,-0.80588822) (4.1944503,-0.80535104) (4.2083391,-0.80475648) (4.222228,-0.80410453) (4.2361169,-0.8033951) (4.2500057,-0.80262799) (4.2638946,-0.8018029) (4.2777835,-0.80091945) (4.2916724,-0.79997715) (4.3055612,-0.79897543) (4.3194501,-0.79791362) (4.333339,-0.79679095) (4.3472279,-0.79560656) (4.3611167,-0.79435947) (4.3750056,-0.79304862) (4.3888945,-0.79167284) (4.4027834,-0.79023085) (4.4166722,-0.78872126) (4.4305611,-0.78714256) (4.44445,-0.78549315) (4.4583389,-0.78377126) (4.4722278,-0.78197504) (4.4861166,-0.78010249) (4.5000055,-0.77815146) (4.5138944,-0.77611968) (4.5277832,-0.77400471) (4.5416721,-0.77180397) (4.555561,-0.76951468) (4.5694499,-0.76713394) (4.5833387,-0.76465861) (4.5972276,-0.76208539) (4.6111165,-0.75941078) (4.6250054,-0.75663102) (4.6388942,-0.75374218) (4.6527831,-0.75074004) (4.666672,-0.74762013) (4.6805609,-0.74437773) (4.6944497,-0.7410078) (4.7083386,-0.73750501) (4.7222275,-0.73386367) (4.7361164,-0.73007777) (4.7500053,-0.72614089) (4.7638941,-0.72204624) (4.777783,-0.71778656) (4.7916719,-0.71335416) (4.8055607,-0.70874082) (4.8194496,-0.70393783) (4.8333385,-0.69893587) (4.8472274,-0.69372501) (4.8611162,-0.68829468) (4.8750051,-0.68263358) (4.888894,-0.67672963) (4.9027829,-0.67056996) (4.9166717,-0.66414076) (4.9305606,-0.65742729) (4.9444495,-0.65041372) (4.9583384,-0.64308313) (4.9722272,-0.63541735) (4.9861161,-0.62739686) (5.000005,-0.61900073) (5.0138939,-0.61020642) (5.0277828,-0.6009897) (5.0416716,-0.59132445) (5.0555605,-0.58118253) (5.0694494,-0.57053359) (5.0833382,-0.55934483) (5.0972271,-0.5475808) (5.111116,-0.53520317) (5.1250049,-0.52217039) (5.1388937,-0.50843745) (5.1527826,-0.49395549) (5.1666715,-0.47867146) (5.1805604,-0.46252764) (5.1944492,-0.44546126) (5.2083381,-0.4274039) (5.222227,-0.40828094) (5.2361159,-0.38801096) (5.2500047,-0.36650494) (5.2638936,-0.34366556) (5.2777825,-0.31938624) (5.2916714,-0.29355023) (5.3055603,-0.26602946) (5.3194491,-0.23668339) (5.333338,-0.20535766) (5.3472269,-0.17188263) (5.3611157,-0.13607176) (5.3750046,-0.097719919) (5.3888935,-0.056601408) (5.4027824,-0.012467984) (5.4166712,0.034953323) (5.4305601,0.085962468) (5.444449,0.14088868) (5.4583379,0.20009274) (5.4722267,0.26396907) (5.4861156,0.33294753) (5.5000045,0.40749464) (5.5138934,0.48811383) (5.5277822,0.57534437) (5.5416711,0.66975824) (5.55556,0.77195404) (5.5694489,0.88254696) (5.5833377,1.0021532) (5.5972266,1.1313671) (5.6111155,1.2707287) (5.6250044,1.4206796) (5.6388932,1.5815031) (5.6527821,1.7532484) (5.666671,1.935634) (5.6805599,2.1279332) (5.6944487,2.3288423) (5.7083376,2.5363413) (5.7222265,2.7475621) (5.7361154,2.9586868) (5.7500042,3.1649109) (5.7638931,3.3605078) (5.777782,3.539033) (5.7916709,3.6936881) (5.8055597,3.8178394) (5.8194486,3.9056412) (5.8333375,3.9526742) (5.8472264,3.9564784) (5.8611152,3.9168691) (5.8750041,3.8359607) (5.888893,3.7178942) (5.9027819,3.5683312) (5.9166707,3.393824) (5.9305596,3.2011831) (5.9444485,2.9969377) (5.9583374,2.7869454) (5.9722262,2.5761642) (5.9861151,2.3685675) (6.000004,2.1671684) (6.0138929,1.9741127) (6.0277817,1.7908081) (6.0416706,1.6180637) (6.0555595,1.4562231) (6.0694484,1.3052829) (6.0833372,1.1649912) (6.0972261,1.0349273) (6.111115,0.91456355) (6.1250039,0.80331186) (6.1388927,0.70055792) (6.1527816,0.60568532) (6.1666705,0.5180922) (6.1805594,0.43720204) (6.1944482,0.36247027) (6.2083371,0.29338785) (6.222226,0.22948256) (6.2361149,0.17031897) (6.2500037,0.11549733) (6.2638926,0.064651845) (6.2777815,0.017448604) (6.2916704,-0.026416648) (6.3055592,-0.067220822) (6.3194481,-0.10521575) (6.333337,-0.14063037) (6.3472259,-0.17367276) (6.3611147,-0.20453209) (6.3750036,-0.23338038) (6.3888925,-0.2603741) (6.4027814,-0.2856557) (6.4166702,-0.30935494) (6.4305591,-0.33159008) (6.444448,-0.35246902) (6.4583369,-0.3720903) (6.4722257,-0.39054395) (6.4861146,-0.40791235) (6.5000035,-0.42427091) (6.5138924,-0.43968877) (6.5277812,-0.45422932) (6.5416701,-0.46795081) (6.555559,-0.48090675) (6.5694479,-0.49314638) (6.5833367,-0.50471507) (6.5972256,-0.51565459) (6.6111145,-0.52600353) (6.6250034,-0.53579749) (6.6388922,-0.54506938) (6.6527811,-0.55384965) (6.66667,-0.56216648) (6.6805589,-0.57004599) (6.6944477,-0.57751239) (6.7083366,-0.58458815) (6.7222255,-0.59129414) (6.7361144,-0.59764975) (6.7500032,-0.603673) (6.7638921,-0.60938067) (6.777781,-0.61478837) (6.7916699,-0.61991065) (6.8055587,-0.62476107) (6.8194476,-0.62935225) (6.8333365,-0.63369599) (6.8472254,-0.63780329) (6.8611142,-0.6416844) (6.8750031,-0.64534892) (6.888892,-0.64880578) (6.9027809,-0.65206334) (6.9166697,-0.6551294) (6.9305586,-0.65801124) (6.9444475,-0.66071567) (6.9583364,-0.66324903) (6.9722252,-0.66561725) (6.9861141,-0.66782585) (7.000003,-0.66987997) (7.0138919,-0.67178442) (7.0277807,-0.67354364) (7.0416696,-0.6751618) (7.0555585,-0.67664274) (7.0694474,-0.67799001) (7.0833362,-0.67920693) (7.0972251,-0.68029654) (7.111114,-0.68126162) (7.1250029,-0.68210476) (7.1388917,-0.68282829) (7.1527806,-0.68343435) (7.1666695,-0.68392487) (7.1805584,-0.68430158) (7.1944472,-0.68456602) (7.2083361,-0.68471956) (7.222225,-0.68476336) (7.2361139,-0.68469843) (7.2500027,-0.68452562) (7.2638916,-0.6842456) (7.2777805,-0.68385888) (7.2916694,-0.6833658) (7.3055582,-0.68276656) (7.3194471,-0.6820612) (7.333336,-0.68124961) (7.3472249,-0.68033151) (7.3611137,-0.67930649) (7.3750026,-0.67817396) (7.3888915,-0.6769332) (7.4027804,-0.67558333) (7.4166692,-0.67412331) (7.4305581,-0.67255193) (7.444447,-0.67086784) (7.4583359,-0.66906953) (7.4722247,-0.66715529) (7.4861136,-0.66512328) (7.5000025,-0.66297147) (7.5138914,-0.66069764) (7.5277802,-0.6582994) (7.5416691,-0.65577417) (7.555558,-0.65311916) (7.5694469,-0.65033139) (7.5833357,-0.64740766) (7.5972246,-0.64434457) (7.6111135,-0.64113847) (7.6250024,-0.63778548) (7.6388912,-0.63428148) (7.6527801,-0.63062208) (7.666669,-0.62680263) (7.6805579,-0.62281818) (7.6944468,-0.61866351) (7.7083356,-0.61433305) (7.7222245,-0.60982095) (7.7361134,-0.60512098) (7.7500022,-0.60022655) (7.7638911,-0.5951307) (7.77778,-0.58982607) (7.7916689,-0.58430487) (7.8055577,-0.57855885) (7.8194466,-0.5725793) (7.8333355,-0.56635701) (7.8472244,-0.55988224) (7.8611132,-0.55314468) (7.8750021,-0.54613343) (7.888891,-0.53883699) (7.9027799,-0.53124316) (7.9166687,-0.52333904) (7.9305576,-0.51511101) (7.9444465,-0.50654463) (7.9583354,-0.49762463) (7.9722243,-0.48833486) (7.9861131,-0.47865821) (8.000002,-0.46857658) (8.0138909,-0.45807079) (8.0277798,-0.44712056) (8.0416686,-0.43570439) (8.0555575,-0.42379952) (8.0694464,-0.41138186) (8.0833352,-0.39842587) (8.0972241,-0.38490454) (8.111113,-0.37078924) (8.1250019,-0.35604966) (8.1388907,-0.34065371) (8.1527796,-0.32456745) (8.1666685,-0.30775493) (8.1805574,-0.29017814) (8.1944463,-0.27179691) (8.2083351,-0.25256878) (8.222224,-0.23244892) (8.2361129,-0.21139003) (8.2500017,-0.18934227) (8.2638906,-0.16625316) (8.2777795,-0.14206756) (8.2916684,-0.11672757) (8.3055572,-0.090172595) (8.3194461,-0.062339275) (8.333335,-0.033161612) (8.3472239,-0.0025710531) (8.3611128,0.029503328) (8.3750016,0.063134575) (8.3888905,0.098397498) (8.4027794,0.1353682) (8.4166682,0.17412349) (8.4305571,0.21474008) (8.444446,0.25729362) (8.4583349,0.30185754) (8.4722237,0.3485015) (8.4861126,0.39728971) (8.5000015,0.44827874) (8.5138904,0.501515) (8.5277793,0.55703183) (8.5416681,0.61484604) (8.555557,0.67495399) (8.5694459,0.73732711) (8.5833347,0.80190701) (8.5972236,0.86859998) (8.6111125,0.93727116) (8.6250014,1.0077385) (8.6388902,1.0797664) (8.6527791,1.1530599) (8.666668,1.2272591) (8.6805569,1.3019345) (8.6944457,1.3765841) (8.7083346,1.4506325) (8.7222235,1.5234323) (8.7361124,1.5942693) (8.7500013,1.6623714) (8.7638901,1.7269213) (8.777779,1.7870746) (8.7916679,1.8419806) (8.8055567,1.8908081) (8.8194456,1.932773) (8.8333345,1.9671672) (8.8472234,1.9933866) (8.8611122,2.0109574) (8.8750011,2.0195565) (8.88889,2.0190269) (8.9027789,2.0093848) (8.9166678,1.9908191) (8.9305566,1.9636823) (8.9444455,1.9284744) (8.9583344,1.8858208) (8.9722232,1.8364463) (8.9861121,1.7811459) (9.000001,1.7207565) (9.0138899,1.6561297) (9.0277787,1.5881064) (9.0416676,1.5174968) (9.0555565,1.4450632) (9.0694454,1.3715079) (9.0833343,1.2974655) (9.0972231,1.2234977) (9.111112,1.1500931) (9.1250009,1.077668) (9.1388897,1.00657) (9.1527786,0.93708271) (9.1666675,0.86943133) (9.1805564,0.80378881) (9.1944452,0.740282) (9.2083341,0.67899774) (9.222223,0.61998858) (9.2361119,0.5632782) (9.2500007,0.50886624) (9.2638896,0.45673265) (9.2777785,0.40684154) (9.2916674,0.35914451) (9.3055563,0.3135835) (9.3194451,0.2700933) (9.333334,0.22860351) (9.3472229,0.18904034) (9.3611117,0.15132796) (9.3750006,0.11538973) (9.3888895,0.081149035) (9.4027784,0.048530115) (9.4166672,0.017458591) (9.4305561,-0.012138073) (9.444445,-0.040330237) (9.4583339,-0.067185846) (9.4722228,-0.092770271) (9.4861116,-0.1171462) (9.5000005,-0.1403736) (9.5138894,-0.16250968) (9.5277782,-0.18360894) (9.5416671,-0.20372321) (9.555556,-0.22290168) (9.5694449,-0.24119101) (9.5833337,-0.25863543) (9.5972226,-0.2752768) (9.6111115,-0.29115469) (9.6250004,-0.30630657) (9.6388893,-0.32076779) (9.6527781,-0.33457179) (9.666667,-0.34775012) (9.6805559,-0.36033257) (9.6944447,-0.37234728) (9.7083336,-0.38382079) (9.7222225,-0.39477814) (9.7361114,-0.40524297) (9.7500002,-0.41523759) (9.7638891,-0.42478305) (9.777778,-0.4338992) (9.7916669,-0.44260479) (9.8055557,-0.4509175) (9.8194446,-0.45885402) (9.8333335,-0.46643009) (9.8472224,-0.47366058) (9.8611113,-0.48055953) (9.8750001,-0.48714016) (9.888889,-0.49341497) (9.9027779,-0.49939575) (9.9166667,-0.50509363) (9.9305556,-0.51051911) (9.9444445,-0.51568211) (9.9583334,-0.52059196) (9.9722222,-0.52525748) (9.9861111,-0.52968699) (10,-0.53380654)};
\end{axis}\end{tikzpicture}\end{center}

\begin{center}\begin{tikzpicture}
\begin{axis}[width=0.78\textwidth,height=4.5cm,xlabel={$u=ka$},ylabel={$|A|$},grid=major,legend style={font=\small},]
\addplot+[mark=none] coordinates{(1e-05,0.090909091) (0.013898875,0.090912478) (0.02778775,0.090922629) (0.041676625,0.090939548) (0.0555655,0.09096324) (0.069454375,0.090993713) (0.08334325,0.091030977) (0.097232125,0.091075042) (0.111121,0.091125924) (0.12500987,0.091183638) (0.13889875,0.091248203) (0.15278763,0.09131964) (0.1666765,0.09139797) (0.18056538,0.09148322) (0.19445425,0.091575416) (0.20834313,0.091674588) (0.222232,0.091780768) (0.23612088,0.091893991) (0.25000975,0.092014292) (0.26389862,0.092141711) (0.2777875,0.092276288) (0.29167637,0.092418068) (0.30556525,0.092567098) (0.31945413,0.092723424) (0.333343,0.0928871) (0.34723188,0.093058178) (0.36112075,0.093236716) (0.37500963,0.093422771) (0.3888985,0.093616407) (0.40278738,0.093817688) (0.41667625,0.09402668) (0.43056512,0.094243455) (0.444454,0.094468085) (0.45834288,0.094700646) (0.47223175,0.094941219) (0.48612063,0.095189884) (0.5000095,0.095446727) (0.51389837,0.095711837) (0.52778725,0.095985306) (0.54167612,0.096267229) (0.555565,0.096557705) (0.56945387,0.096856837) (0.58334275,0.097164729) (0.59723162,0.097481492) (0.6111205,0.097807239) (0.62500937,0.098142087) (0.63889825,0.098486159) (0.65278712,0.098839578) (0.666676,0.099202475) (0.68056487,0.099574985) (0.69445375,0.099957244) (0.70834263,0.1003494) (0.7222315,0.10075159) (0.73612037,0.10116398) (0.75000925,0.10158672) (0.76389812,0.10201997) (0.777787,0.10246391) (0.79167587,0.1029187) (0.80556475,0.10338452) (0.81945362,0.10386157) (0.8333425,0.10435002) (0.84723138,0.10485008) (0.86112025,0.10536195) (0.87500912,0.10588584) (0.888898,0.10642197) (0.90278687,0.10697056) (0.91667575,0.10753184) (0.93056462,0.10810606) (0.9444535,0.10869344) (0.95834237,0.10929426) (0.97223125,0.10990877) (0.98612013,0.11053725) (1.000009,0.11117997) (1.0138979,0.11183722) (1.0277867,0.1125093) (1.0416756,0.11319652) (1.0555645,0.11389921) (1.0694534,0.11461768) (1.0833423,0.11535228) (1.0972311,0.11610336) (1.11112,0.11687128) (1.1250089,0.11765643) (1.1388978,0.11845919) (1.1527866,0.11927996) (1.1666755,0.12011915) (1.1805644,0.12097721) (1.1944533,0.12185458) (1.2083421,0.12275171) (1.222231,0.12366908) (1.2361199,0.1246072) (1.2500088,0.12556657) (1.2638976,0.12654772) (1.2777865,0.1275512) (1.2916754,0.12857759) (1.3055642,0.12962747) (1.3194531,0.13070146) (1.333342,0.13180019) (1.3472309,0.13292432) (1.3611198,0.13407455) (1.3750086,0.13525157) (1.3888975,0.13645612) (1.4027864,0.13768898) (1.4166753,0.13895094) (1.4305641,0.14024283) (1.444453,0.14156551) (1.4583419,0.14291987) (1.4722308,0.14430684) (1.4861196,0.1457274) (1.5000085,0.14718255) (1.5138974,0.14867335) (1.5277863,0.15020089) (1.5416751,0.1517663) (1.555564,0.15337079) (1.5694529,0.1550156) (1.5833417,0.156702) (1.5972306,0.15843137) (1.6111195,0.16020511) (1.6250084,0.1620247) (1.6388973,0.16389168) (1.6527861,0.16580766) (1.666675,0.16777433) (1.6805639,0.16979345) (1.6944528,0.17186689) (1.7083416,0.17399656) (1.7222305,0.17618451) (1.7361194,0.17843287) (1.7500083,0.18074387) (1.7638971,0.18311986) (1.777786,0.1855633) (1.7916749,0.18807681) (1.8055638,0.19066309) (1.8194526,0.19332502) (1.8333415,0.19606563) (1.8472304,0.1988881) (1.8611193,0.2017958) (1.8750081,0.20479227) (1.888897,0.20788125) (1.9027859,0.21106671) (1.9166748,0.21435283) (1.9305636,0.21774404) (1.9444525,0.22124503) (1.9583414,0.22486079) (1.9722303,0.22859658) (1.9861191,0.23245801) (2.000008,0.23645103) (2.0138969,0.24058199) (2.0277858,0.24485763) (2.0416746,0.24928514) (2.0555635,0.25387219) (2.0694524,0.25862697) (2.0833413,0.26355823) (2.0972301,0.26867537) (2.111119,0.2739884) (2.1250079,0.27950812) (2.1388968,0.28524609) (2.1527856,0.29121474) (2.1666745,0.2974275) (2.1805634,0.30389879) (2.1944522,0.31064421) (2.2083411,0.31768065) (2.22223,0.32502635) (2.2361189,0.33270114) (2.2500077,0.34072654) (2.2638966,0.34912596) (2.2777855,0.35792496) (2.2916744,0.36715142) (2.3055633,0.37683585) (2.3194521,0.38701172) (2.333341,0.39771578) (2.3472299,0.40898853) (2.3611188,0.42087463) (2.3750076,0.43342351) (2.3888965,0.44668998) (2.4027854,0.46073496) (2.4166743,0.47562641) (2.4305631,0.4914403) (2.444452,0.50826181) (2.4583409,0.5261868) (2.4722297,0.54532346) (2.4861186,0.56579435) (2.5000075,0.58773882) (2.5138964,0.61131591) (2.5277852,0.63670788) (2.5416741,0.66412452) (2.555563,0.69380839) (2.5694519,0.72604131) (2.5833408,0.76115231) (2.5972296,0.79952763) (2.6111185,0.84162311) (2.6250074,0.88797986) (2.6388963,0.93924396) (2.6527851,0.99619149) (2.666674,1.0597602) (2.6805629,1.1310895) (2.6944518,1.2115711) (2.7083406,1.3029115) (2.7222295,1.407206) (2.7361184,1.5270219) (2.7500072,1.6654749) (2.7638961,1.8262593) (2.777785,2.0135375) (2.7916739,2.2314736) (2.8055627,2.4829673) (2.8194516,2.7667906) (2.8333405,3.0721146) (2.8472294,3.3706079) (2.8611183,3.611549) (2.8750071,3.7335359) (2.888896,3.6992242) (2.9027849,3.5241332) (2.9166738,3.263311) (2.9305626,2.9738908) (2.9444515,2.6933109) (2.9583404,2.4390407) (2.9722293,2.216043) (2.9861181,2.0232322) (3.000007,1.8572243) (3.0138959,1.7141351) (3.0277847,1.5903231) (3.0416736,1.4826329) (3.0555625,1.3884261) (3.0694514,1.3055312) (3.0833403,1.2321716) (3.0972291,1.1668949) (3.111118,1.1085104) (3.1250069,1.0560378) (3.1388958,1.0086663) (3.1527846,0.96572096) (3.1666735,0.9266371) (3.1805624,0.89093943) (3.1944513,0.85822563) (3.2083401,0.82815332) (3.222229,0.80042961) (3.2361179,0.77480262) (3.2500068,0.75105476) (3.2638956,0.7289972) (3.2777845,0.70846533) (3.2916734,0.68931507) (3.3055622,0.67141983) (3.3194511,0.65466796) (3.33334,0.63896068) (3.3472289,0.62421022) (3.3611178,0.61033845) (3.3750066,0.59727552) (3.3888955,0.58495886) (3.4027844,0.57333225) (3.4166733,0.56234505) (3.4305621,0.55195152) (3.444451,0.54211026) (3.4583399,0.53278371) (3.4722288,0.52393773) (3.4861176,0.51554122) (3.5000065,0.50756577) (3.5138954,0.49998544) (3.5277843,0.49277642) (3.5416731,0.48591691) (3.555562,0.47938683) (3.5694509,0.47316772) (3.5833397,0.46724256) (3.5972286,0.46159564) (3.6111175,0.45621243) (3.6250064,0.45107949) (3.6388953,0.44618436) (3.6527841,0.44151549) (3.666673,0.43706214) (3.6805619,0.43281434) (3.6944508,0.42876281) (3.7083396,0.4248989) (3.7222285,0.42121455) (3.7361174,0.41770226) (3.7500063,0.41435499) (3.7638951,0.41116621) (3.777784,0.40812978) (3.7916729,0.40523997) (3.8055618,0.40249144) (3.8194506,0.39987916) (3.8333395,0.39739845) (3.8472284,0.3950449) (3.8611172,0.3928144) (3.8750061,0.39070309) (3.888895,0.38870737) (3.9027839,0.38682384) (3.9166728,0.38504935) (3.9305616,0.38338094) (3.9444505,0.38181581) (3.9583394,0.3803514) (3.9722283,0.37898527) (3.9861171,0.37771516) (4.000006,0.37653897) (4.0138949,0.37545474) (4.0277837,0.37446063) (4.0416726,0.37355496) (4.0555615,0.37273616) (4.0694504,0.37200279) (4.0833392,0.37135351) (4.0972281,0.37078711) (4.111117,0.37030246) (4.1250059,0.36989856) (4.1388947,0.36957449) (4.1527836,0.36932944) (4.1666725,0.36916268) (4.1805614,0.36907357) (4.1944503,0.36906158) (4.2083391,0.36912625) (4.222228,0.36926718) (4.2361169,0.36948411) (4.2500057,0.3697768) (4.2638946,0.37014513) (4.2777835,0.37058905) (4.2916724,0.37110858) (4.3055612,0.37170383) (4.3194501,0.37237497) (4.333339,0.37312227) (4.3472279,0.37394607) (4.3611167,0.37484677) (4.3750056,0.37582487) (4.3888945,0.37688095) (4.4027834,0.37801565) (4.4166722,0.3792297) (4.4305611,0.38052393) (4.44445,0.38189923) (4.4583389,0.38335659) (4.4722278,0.38489707) (4.4861166,0.38652183) (4.5000055,0.38823213) (4.5138944,0.39002931) (4.5277832,0.39191481) (4.5416721,0.39389018) (4.555561,0.39595706) (4.5694499,0.39811721) (4.5833387,0.40037247) (4.5972276,0.40272484) (4.6111165,0.4051764) (4.6250054,0.40772939) (4.6388942,0.41038613) (4.6527831,0.41314914) (4.666672,0.41602101) (4.6805609,0.41900454) (4.6944497,0.42210265) (4.7083386,0.42531842) (4.7222275,0.42865512) (4.7361164,0.43211618) (4.7500053,0.43570524) (4.7638941,0.43942611) (4.777783,0.44328282) (4.7916719,0.44727962) (4.8055607,0.45142101) (4.8194496,0.4557117) (4.8333385,0.46015668) (4.8472274,0.46476121) (4.8611162,0.46953084) (4.8750051,0.47447142) (4.888894,0.47958915) (4.9027829,0.48489054) (4.9166717,0.4903825) (4.9305606,0.49607231) (4.9444495,0.50196768) (4.9583384,0.50807675) (4.9722272,0.51440814) (4.9861161,0.52097097) (5.000005,0.5277749) (5.0138939,0.53483016) (5.0277828,0.54214759) (5.0416716,0.5497387) (5.0555605,0.55761568) (5.0694494,0.56579148) (5.0833382,0.57427985) (5.0972271,0.58309542) (5.111116,0.59225371) (5.1250049,0.60177125) (5.1388937,0.61166563) (5.1527826,0.62195559) (5.1666715,0.63266109) (5.1805604,0.64380341) (5.1944492,0.65540523) (5.2083381,0.66749078) (5.222227,0.6800859) (5.2361159,0.6932182) (5.2500047,0.70691719) (5.2638936,0.72121438) (5.2777825,0.73614348) (5.2916714,0.75174051) (5.3055603,0.768044) (5.3194491,0.78509512) (5.333338,0.80293791) (5.3472269,0.82161939) (5.3611157,0.84118979) (5.3750046,0.86170271) (5.3888935,0.88321523) (5.4027824,0.9057881) (5.4166712,0.92948581) (5.4305601,0.95437663) (5.444449,0.98053259) (5.4583379,1.0080293) (5.4722267,1.0369458) (5.4861156,1.0673639) (5.5000045,1.0993676) (5.5138934,1.133042) (5.5277822,1.1684716) (5.5416711,1.2057387) (5.55556,1.2449201) (5.5694489,1.2860837) (5.5833377,1.3292836) (5.5972266,1.3745538) (5.6111155,1.4219006) (5.6250044,1.4712922) (5.6388932,1.5226477) (5.6527821,1.5758218) (5.666671,1.6305893) (5.6805599,1.6866266) (5.6944487,1.7434926) (5.7083376,1.8006117) (5.7222265,1.8572592) (5.7361154,1.9125546) (5.7500042,1.9654674) (5.7638931,2.0148385) (5.777782,2.0594232) (5.7916709,2.0979554) (5.8055597,2.1292322) (5.8194486,2.1522094) (5.8333375,2.1660969) (5.8472264,2.1704368) (5.8611152,2.1651506) (5.8750041,2.1505463) (5.888893,2.1272822) (5.9027819,2.096298) (5.9166707,2.0587238) (5.9305596,2.0157843) (5.9444485,1.9687118) (5.9583374,1.9186758) (5.9722262,1.8667346) (5.9861151,1.8138076) (6.000004,1.7606646) (6.0138929,1.7079289) (6.0277817,1.6560882) (6.0416706,1.6055107) (6.0555595,1.556463) (6.0694484,1.509127) (6.0833372,1.4636173) (6.0972261,1.4199943) (6.111115,1.3782774) (6.1250039,1.3384546) (6.1388927,1.3004909) (6.1527816,1.2643352) (6.1666705,1.2299246) (6.1805594,1.1971894) (6.1944482,1.1660557) (6.2083371,1.1364474) (6.222226,1.1082887) (6.2361149,1.0815045) (6.2500037,1.0560218) (6.2638926,1.0317703) (6.2777815,1.0086824) (6.2916704,0.98669389) (6.3055592,0.96574377) (6.3194481,0.94577437) (6.333337,0.92673131) (6.3472259,0.90856341) (6.3611147,0.89122258) (6.3750036,0.87466369) (6.3888925,0.85884441) (6.4027814,0.84372507) (6.4166702,0.82926853) (6.4305591,0.81543997) (6.444448,0.80220682) (6.4583369,0.78953857) (6.4722257,0.77740662) (6.4861146,0.76578422) (6.5000035,0.75464626) (6.5138924,0.74396926) (6.5277812,0.73373115) (6.5416701,0.72391127) (6.555559,0.71449021) (6.5694479,0.70544976) (6.5833367,0.69677281) (6.5972256,0.6884433) (6.6111145,0.68044611) (6.6250034,0.67276704) (6.6388922,0.66539272) (6.6527811,0.65831057) (6.66667,0.65150873) (6.6805589,0.64497606) (6.6944477,0.63870202) (6.7083366,0.6326767) (6.7222255,0.62689074) (6.7361144,0.62133532) (6.7500032,0.6160021) (6.7638921,0.61088322) (6.777781,0.60597125) (6.7916699,0.60125918) (6.8055587,0.59674037) (6.8194476,0.59240857) (6.8333365,0.58825784) (6.8472254,0.5842826) (6.8611142,0.58047755) (6.8750031,0.5768377) (6.888892,0.57335831) (6.9027809,0.57003493) (6.9166697,0.56686332) (6.9305586,0.5638395) (6.9444475,0.5609597) (6.9583364,0.55822036) (6.9722252,0.55561813) (6.9861141,0.55314984) (7.000003,0.5508125) (7.0138919,0.5486033) (7.0277807,0.5465196) (7.0416696,0.54455891) (7.0555585,0.54271889) (7.0694474,0.54099734) (7.0833362,0.53939223) (7.0972251,0.53790163) (7.111114,0.53652376) (7.1250029,0.53525695) (7.1388917,0.53409966) (7.1527806,0.53305046) (7.1666695,0.53210805) (7.1805584,0.53127122) (7.1944472,0.53053888) (7.2083361,0.52991003) (7.222225,0.52938377) (7.2361139,0.52895932) (7.2500027,0.52863598) (7.2638916,0.52841315) (7.2777805,0.52829031) (7.2916694,0.52826706) (7.3055582,0.52834306) (7.3194471,0.52851807) (7.333336,0.52879194) (7.3472249,0.52916462) (7.3611137,0.52963612) (7.3750026,0.53020656) (7.3888915,0.53087613) (7.4027804,0.53164512) (7.4166692,0.5325139) (7.4305581,0.53348291) (7.444447,0.53455271) (7.4583359,0.53572393) (7.4722247,0.53699728) (7.4861136,0.53837358) (7.5000025,0.53985372) (7.5138914,0.5414387) (7.5277802,0.5431296) (7.5416691,0.54492759) (7.555558,0.54683397) (7.5694469,0.54885009) (7.5833357,0.55097744) (7.5972246,0.5532176) (7.6111135,0.55557225) (7.6250024,0.55804319) (7.6388912,0.56063232) (7.6527801,0.56334165) (7.666669,0.56617334) (7.6805579,0.56912963) (7.6944468,0.57221291) (7.7083356,0.5754257) (7.7222245,0.57877065) (7.7361134,0.58225053) (7.7500022,0.58586829) (7.7638911,0.589627) (7.77778,0.59352989) (7.7916689,0.59758035) (7.8055577,0.60178194) (7.8194466,0.60613839) (7.8333355,0.6106536) (7.8472244,0.61533167) (7.8611132,0.62017688) (7.8750021,0.6251937) (7.888891,0.63038683) (7.9027799,0.63576118) (7.9166687,0.64132187) (7.9305576,0.64707427) (7.9444465,0.65302396) (7.9583354,0.65917682) (7.9722243,0.66553895) (7.9861131,0.67211673) (8.000002,0.67891683) (8.0138909,0.68594619) (8.0277798,0.69321207) (8.0416686,0.70072203) (8.0555575,0.70848395) (8.0694464,0.71650603) (8.0833352,0.72479683) (8.0972241,0.73336524) (8.111113,0.74222051) (8.1250019,0.75137226) (8.1388907,0.76083046) (8.1527796,0.77060548) (8.1666685,0.78070804) (8.1805574,0.79114925) (8.1944463,0.80194059) (8.2083351,0.81309391) (8.222224,0.8246214) (8.2361129,0.83653561) (8.2500017,0.84884939) (8.2638906,0.8615759) (8.2777795,0.87472853) (8.2916684,0.88832088) (8.3055572,0.9023667) (8.3194461,0.91687979) (8.333335,0.93187395) (8.3472239,0.94736284) (8.3611128,0.96335984) (8.3750016,0.97987795) (8.3888905,0.99692953) (8.4027794,1.0145261) (8.4166682,1.0326782) (8.4305571,1.0513948) (8.444446,1.0706834) (8.4583349,1.0905491) (8.4722237,1.1109945) (8.4861126,1.1320192) (8.5000015,1.1536189) (8.5138904,1.1757851) (8.5277793,1.1985039) (8.5416681,1.2217557) (8.555557,1.2455138) (8.5694459,1.2697439) (8.5833347,1.2944025) (8.5972236,1.3194366) (8.6111125,1.3447817) (8.6250014,1.3703615) (8.6388902,1.3960865) (8.6527791,1.4218532) (8.666668,1.4475435) (8.6805569,1.4730243) (8.6944457,1.4981474) (8.7083346,1.5227501) (8.7222235,1.5466559) (8.7361124,1.5696769) (8.7500013,1.5916155) (8.7638901,1.6122679) (8.777779,1.6314283) (8.7916679,1.6488929) (8.8055567,1.664466) (8.8194456,1.6779648) (8.8333345,1.6892255) (8.8472234,1.6981087) (8.8611122,1.704504) (8.8750011,1.7083344) (8.88889,1.7095591) (8.9027789,1.7081749) (8.9166678,1.7042162) (8.9305566,1.6977535) (8.9444455,1.688891) (8.9583344,1.6777619) (8.9722232,1.6645244) (8.9861121,1.6493555) (9.000001,1.6324462) (9.0138899,1.6139952) (9.0277787,1.5942044) (9.0416676,1.5732739) (9.0555565,1.5513984) (9.0694454,1.5287635) (9.0833343,1.5055439) (9.0972231,1.4819012) (9.111112,1.457983) (9.1250009,1.4339223) (9.1388897,1.4098378) (9.1527786,1.3858338) (9.1666675,1.3620011) (9.1805564,1.3384175) (9.1944452,1.3151492) (9.2083341,1.2922516) (9.222223,1.26977) (9.2361119,1.2477412) (9.2500007,1.2261942) (9.2638896,1.205151) (9.2777785,1.1846279) (9.2916674,1.1646359) (9.3055563,1.1451814) (9.3194451,1.1262674) (9.333334,1.1078935) (9.3472229,1.0900566) (9.3611117,1.0727516) (9.3750006,1.0559712) (9.3888895,1.039707) (9.4027784,1.0239494) (9.4166672,1.0086876) (9.4305561,0.99391038) (9.444445,0.97960594) (9.4583339,0.96576212) (9.4722228,0.95236654) (9.4861116,0.93940675) (9.5000005,0.92687027) (9.5138894,0.91474471) (9.5277782,0.90301781) (9.5416671,0.8916775) (9.555556,0.88071192) (9.5694449,0.87010949) (9.5833337,0.85985893) (9.5972226,0.84994922) (9.6111115,0.8403697) (9.6250004,0.83111003) (9.6388893,0.82216018) (9.6527781,0.81351049) (9.666667,0.80515162) (9.6805559,0.79707456) (9.6944447,0.78927063) (9.7083336,0.78173149) (9.7222225,0.77444911) (9.7361114,0.76741577) (9.7500002,0.76062405) (9.7638891,0.75406684) (9.777778,0.74773731) (9.7916669,0.74162891) (9.8055557,0.73573536) (9.8194446,0.73005064) (9.8333335,0.72456899) (9.8472224,0.71928489) (9.8611113,0.71419305) (9.8750001,0.70928843) (9.888889,0.70456618) (9.9027779,0.70002167) (9.9166667,0.6956505) (9.9305556,0.69144842) (9.9444445,0.68741141) (9.9583334,0.68353562) (9.9722222,0.67981737) (9.9861111,0.67625315) (10,0.67283962)};
\end{axis}\end{tikzpicture}\end{center}

内区驻留增强与大正延迟标识共振, 但势全为非负, 没有负能束缚态. $\delta(0)=\delta(\infty)=0$, 满足$N=0$; 局部相移急升不等于新增束缚态.
\end{solution}
\originalsection{第5题：$z$角动量对易子（10分）}
\begin{exercise}
(a,b)求$L_z$与$x,y,z,p_x,p_y,p_z$的对易子. (c,d)证明$L_z$本征态中$x,y,p_x,p_y$均值为0.
\end{exercise}
\begin{solution}
$L_z=xp_y-yp_x$给$[L_z,x]=\ii\hbar y$,$[L_z,y]=-\ii\hbar x$,$[L_z,z]=0$; 类似$[L_z,p_x]=\ii\hbar p_y$,$[L_z,p_y]=-\ii\hbar p_x$,$[L_z,p_z]=0$. 归一化本征态中$\langle[L_z,A]\rangle=0$, 因为左右$L_z$读数相同. 分别代入上述四个非零式就得到四个均值0, 前提是矩存在和共同定义域允许操作. 不是说这些测量没有分散或粒子停在轴线上.
\end{solution}
\originalsection{第6题：球坐标算符（10分）}
\begin{exercise}
(a)求$r,\theta,\phi$对$x,y,z$的九个导数. (b)写$L_x,L_y,L_z$. (c)分别展开平方并相加得到$L^2$.
\end{exercise}
\begin{solution}
在坐标非奇点, 各行分别为$\nabla r,\nabla\theta,\nabla\phi$:
\[
\begin{pmatrix}
\sin\theta\cos\phi&\sin\theta\sin\phi&\cos\theta\\
\cos\theta\cos\phi/r&\cos\theta\sin\phi/r&-\sin\theta/r\\
-\sin\phi/(r\sin\theta)&\cos\phi/(r\sin\theta)&0
\end{pmatrix}.
\]
链式法则给$L_x=\ii\hbar(\sin\phi\,\partial_\theta+\cot\theta\cos\phi\,\partial_\phi)$,$L_y=\ii\hbar(-\cos\phi\,\partial_\theta+\cot\theta\sin\phi\,\partial_\phi)$,$L_z=-\ii\hbar\partial_\phi$. 为展示平方时对系数也要求导, 令$s=\sin\phi,c=\cos\phi,t=\cot\theta$, 写$L_x=\ii\hbar A,L_y=\ii\hbar B$:
\[
\begin{aligned}
A^2&=s^2\partial_{\theta\theta}+2sct\partial_{\theta\phi}+c^2t^2\partial_{\phi\phi}
+tc^2\partial_\theta-sc(\csc^2\theta+t^2)\partial_\phi,\\
B^2&=c^2\partial_{\theta\theta}-2sct\partial_{\theta\phi}+s^2t^2\partial_{\phi\phi}
+ts^2\partial_\theta+sc(\csc^2\theta+t^2)\partial_\phi.
\end{aligned}
\]
$L_x^2=-\hbar^2A^2,L_y^2=-\hbar^2B^2,L_z^2=-\hbar^2\partial_{\phi\phi}$, 相加混合项抵消, 用$1+\cot^2\theta=\csc^2\theta$得到第10章的$L^2$表达. 轴上的坐标奇点需要正则球面函数解释, 不能直接代$r\sin\theta=0$到导数分式.
\end{solution}

\originalchapter{8.04 第二次测试：定态与模型极限}
原件为8.04 Fall 2015 Test 2, 2015年11月17日21:30--23:00, 90分钟, 五题100分, 闭卷且禁计算器; 由Spring 2016 OCW提供. 原卷第2--4页公式表已归档, 本册理论覆盖相关公式, 不把公式表当成额外已解题. 以下为\textbf{AI独立解答}.
\originalsection{第1题：画定态波函数（20分）}
\begin{exercise}
原图势为偶函数: $|x|>a$无限高墙, 墙内两侧平底, 中央有平滑凸起; 图中$E_1<E_2$均低于中央峰顶, $E_5$高于峰顶. 画第1、2、5态, 标明对称性、节点、曲率、拐点、振幅与波长的变化.
\end{exercise}
\begin{solution}
第$n$态有$n-1$个内部节点, 奇偶交替. 基态偶且处处同号, 中央禁阻区是正的局部谷（$\psi''>0$）, 左右允许区形成两个较大的峰. 第二态奇, 中央节点$x=0$, 两侧异号; 两态尾在中央强烈重叠不足时能量很接近. 第五态偶, 有四节点; 因其能量高过峰, 墙内没有禁阻区. 由$\psi''=2m(V-E)\psi/\hbar^2$判断曲率: 允许区对正$\psi$向下弯, 禁阻区对正$\psi$向上弯, 负段反向. 在$V=E$且$\psi\neq0$处有拐点. 允许区局部波数$\sqrt{2m(E-V)}/\hbar$较小时波长较长, 慢变近似中振幅较大; 并非所有任意势的精确振幅都能逐点用WKB判定.

下图仅为满足原图定性条件的独立数值模型, 取$x/a\in[-1,1]$, 无量纲动能$-0.15\partial_{x/a}^2$, 中央$|x/a|<0.6$势为$6[1-(x/(0.6a))^2]$, 外部墙内为0. 第1、2、5能量为2.818、2.908、11.864, 前两态低于峰顶6. 各波函数分别按最大绝对值缩放, 不用于比较不同态概率大小; 原图没有提供可求精确能量的函数, 本解没有冒称这些是原题数值.
\begin{center}\begin{tikzpicture}
\begin{axis}[width=0.78\textwidth,height=3cm,xlabel={$x/a$},ylabel={$\psi_{1}$},grid=major,legend style={font=\small},xmin=-1,xmax=1,ymin=-1.2,ymax=1.2]
\addplot+[mark=none] coordinates{(-1,0) (-0.99,0.043331268) (-0.98,0.086581139) (-0.97,0.12966837) (-0.96,0.17251202) (-0.95,0.21503161) (-0.94,0.25714727) (-0.93,0.29877988) (-0.92,0.33985124) (-0.91,0.38028419) (-0.9,0.42000279) (-0.89,0.45893242) (-0.88,0.49699996) (-0.87,0.53413389) (-0.86,0.57026446) (-0.85,0.6053238) (-0.84,0.63924605) (-0.83,0.67196749) (-0.82,0.70342666) (-0.81,0.73356444) (-0.8,0.76232425) (-0.79,0.78965204) (-0.78,0.81549648) (-0.77,0.83980903) (-0.76,0.86254402) (-0.75,0.88365873) (-0.74,0.90311351) (-0.73,0.92087181) (-0.72,0.93690026) (-0.71,0.95116877) (-0.7,0.96365053) (-0.69,0.97432208) (-0.68,0.9831634) (-0.67,0.99015785) (-0.66,0.99529232) (-0.65,0.99855715) (-0.64,0.99994621) (-0.63,0.99945689) (-0.62,0.99709011) (-0.61,0.99285031) (-0.6,0.98674546) (-0.59,0.97880845) (-0.58,0.96916177) (-0.57,0.95794821) (-0.56,0.94530824) (-0.55,0.9313793) (-0.54,0.91629523) (-0.53,0.90018565) (-0.52,0.88317558) (-0.51,0.86538504) (-0.5,0.84692877) (-0.49,0.82791598) (-0.48,0.80845022) (-0.47,0.78862926) (-0.46,0.76854504) (-0.45,0.74828367) (-0.44,0.72792548) (-0.43,0.70754507) (-0.42,0.68721147) (-0.41,0.66698823) (-0.4,0.64693363) (-0.39,0.62710081) (-0.38,0.60753803) (-0.37,0.58828885) (-0.36,0.56939237) (-0.35,0.55088345) (-0.34,0.53279295) (-0.33,0.51514801) (-0.32,0.49797221) (-0.31,0.48128589) (-0.3,0.46510636) (-0.29,0.44944811) (-0.28,0.43432306) (-0.27,0.41974079) (-0.26,0.40570874) (-0.25,0.39223242) (-0.24,0.37931561) (-0.23,0.36696053) (-0.22,0.35516808) (-0.21,0.34393793) (-0.2,0.33326875) (-0.19,0.32315834) (-0.18,0.31360374) (-0.17,0.30460144) (-0.16,0.29614742) (-0.15,0.28823735) (-0.14,0.28086662) (-0.13,0.2740305) (-0.12,0.26772422) (-0.11,0.26194301) (-0.1,0.25668224) (-0.09,0.25193745) (-0.08,0.24770442) (-0.07,0.24397926) (-0.06,0.24075837) (-0.05,0.2380386) (-0.04,0.23581718) (-0.03,0.23409183) (-0.02,0.23286075) (-0.01,0.23212263) (0,0.23187667) (0.01,0.23212263) (0.02,0.23286075) (0.03,0.23409183) (0.04,0.23581718) (0.05,0.2380386) (0.06,0.24075837) (0.07,0.24397926) (0.08,0.24770442) (0.09,0.25193745) (0.1,0.25668224) (0.11,0.26194301) (0.12,0.26772422) (0.13,0.2740305) (0.14,0.28086662) (0.15,0.28823735) (0.16,0.29614742) (0.17,0.30460144) (0.18,0.31360374) (0.19,0.32315834) (0.2,0.33326875) (0.21,0.34393793) (0.22,0.35516808) (0.23,0.36696053) (0.24,0.37931561) (0.25,0.39223242) (0.26,0.40570874) (0.27,0.41974079) (0.28,0.43432306) (0.29,0.44944811) (0.3,0.46510636) (0.31,0.48128589) (0.32,0.49797221) (0.33,0.51514801) (0.34,0.53279295) (0.35,0.55088345) (0.36,0.56939237) (0.37,0.58828885) (0.38,0.60753803) (0.39,0.62710081) (0.4,0.64693363) (0.41,0.66698823) (0.42,0.68721147) (0.43,0.70754507) (0.44,0.72792548) (0.45,0.74828367) (0.46,0.76854504) (0.47,0.78862926) (0.48,0.80845022) (0.49,0.82791598) (0.5,0.84692877) (0.51,0.86538504) (0.52,0.88317558) (0.53,0.90018565) (0.54,0.91629523) (0.55,0.9313793) (0.56,0.94530823) (0.57,0.95794821) (0.58,0.96916177) (0.59,0.97880845) (0.6,0.98674546) (0.61,0.99285031) (0.62,0.99709011) (0.63,0.99945689) (0.64,0.99994621) (0.65,0.99855715) (0.66,0.99529232) (0.67,0.99015785) (0.68,0.9831634) (0.69,0.97432208) (0.7,0.96365053) (0.71,0.95116877) (0.72,0.93690026) (0.73,0.92087181) (0.74,0.90311351) (0.75,0.88365873) (0.76,0.86254402) (0.77,0.83980903) (0.78,0.81549648) (0.79,0.78965204) (0.8,0.76232425) (0.81,0.73356444) (0.82,0.70342666) (0.83,0.67196749) (0.84,0.63924605) (0.85,0.6053238) (0.86,0.57026446) (0.87,0.53413389) (0.88,0.49699996) (0.89,0.45893242) (0.9,0.42000279) (0.91,0.38028419) (0.92,0.33985124) (0.93,0.29877988) (0.94,0.25714727) (0.95,0.21503161) (0.96,0.17251202) (0.97,0.12966837) (0.98,0.086581139) (0.99,0.043331268) (1,0)};
\end{axis}\end{tikzpicture}\end{center}

\begin{center}\begin{tikzpicture}
\begin{axis}[width=0.78\textwidth,height=3cm,xlabel={$x/a$},ylabel={$\psi_{2}$},grid=major,legend style={font=\small},xmin=-1,xmax=1,ymin=-1.2,ymax=1.2]
\addplot+[mark=none] coordinates{(-1,0) (-0.99,-0.04401253) (-0.98,-0.087939762) (-0.97,-0.13169656) (-0.96,-0.17519813) (-0.95,-0.21836016) (-0.94,-0.26109901) (-0.93,-0.30333183) (-0.92,-0.34497679) (-0.91,-0.38595317) (-0.9,-0.42618156) (-0.89,-0.465584) (-0.88,-0.50408412) (-0.87,-0.54160732) (-0.86,-0.57808086) (-0.85,-0.61343406) (-0.84,-0.64759841) (-0.83,-0.68050769) (-0.82,-0.71209812) (-0.81,-0.74230849) (-0.8,-0.77108024) (-0.79,-0.79835762) (-0.78,-0.82408775) (-0.77,-0.84822077) (-0.76,-0.87070992) (-0.75,-0.8915116) (-0.74,-0.91058551) (-0.73,-0.92789467) (-0.72,-0.94340555) (-0.71,-0.95708807) (-0.7,-0.96891573) (-0.69,-0.9788656) (-0.68,-0.9869184) (-0.67,-0.99305852) (-0.66,-0.99727406) (-0.65,-0.99955685) (-0.64,-0.99990247) (-0.63,-0.99831025) (-0.62,-0.99478327) (-0.61,-0.98932837) (-0.6,-0.98195612) (-0.59,-0.97270211) (-0.58,-0.96169114) (-0.57,-0.94906806) (-0.56,-0.93497524) (-0.55,-0.91955184) (-0.54,-0.90293316) (-0.53,-0.88525012) (-0.52,-0.86662877) (-0.51,-0.84718992) (-0.5,-0.82704884) (-0.49,-0.806315) (-0.48,-0.78509195) (-0.47,-0.76347715) (-0.46,-0.74156196) (-0.45,-0.71943163) (-0.44,-0.69716531) (-0.43,-0.67483615) (-0.42,-0.65251138) (-0.41,-0.63025249) (-0.4,-0.60811534) (-0.39,-0.5861504) (-0.38,-0.5644029) (-0.37,-0.54291307) (-0.36,-0.52171632) (-0.35,-0.50084356) (-0.34,-0.48032132) (-0.33,-0.46017209) (-0.32,-0.44041446) (-0.31,-0.42106345) (-0.3,-0.40213067) (-0.29,-0.38362458) (-0.28,-0.36555068) (-0.27,-0.34791179) (-0.26,-0.33070818) (-0.25,-0.3139378) (-0.24,-0.29759648) (-0.23,-0.2816781) (-0.22,-0.26617474) (-0.21,-0.25107689) (-0.2,-0.23637353) (-0.19,-0.22205234) (-0.18,-0.20809979) (-0.17,-0.19450128) (-0.16,-0.18124123) (-0.15,-0.16830323) (-0.14,-0.15567008) (-0.13,-0.14332393) (-0.12,-0.1312463) (-0.11,-0.11941821) (-0.1,-0.10782025) (-0.09,-0.096432557) (-0.08,-0.085234978) (-0.07,-0.074207042) (-0.06,-0.06332804) (-0.05,-0.052577051) (-0.04,-0.041932986) (-0.03,-0.031374619) (-0.02,-0.020880615) (-0.01,-0.010429563) (0,1.1987668e-11) (0.01,0.010429563) (0.02,0.020880615) (0.03,0.031374619) (0.04,0.041932986) (0.05,0.052577051) (0.06,0.06332804) (0.07,0.074207042) (0.08,0.085234978) (0.09,0.096432557) (0.1,0.10782025) (0.11,0.11941821) (0.12,0.1312463) (0.13,0.14332393) (0.14,0.15567008) (0.15,0.16830323) (0.16,0.18124123) (0.17,0.19450128) (0.18,0.20809979) (0.19,0.22205234) (0.2,0.23637353) (0.21,0.25107689) (0.22,0.26617474) (0.23,0.2816781) (0.24,0.29759648) (0.25,0.3139378) (0.26,0.33070818) (0.27,0.34791179) (0.28,0.36555068) (0.29,0.38362458) (0.3,0.40213067) (0.31,0.42106345) (0.32,0.44041446) (0.33,0.46017209) (0.34,0.48032132) (0.35,0.50084356) (0.36,0.52171632) (0.37,0.54291307) (0.38,0.5644029) (0.39,0.5861504) (0.4,0.60811534) (0.41,0.63025249) (0.42,0.65251138) (0.43,0.67483615) (0.44,0.69716531) (0.45,0.71943163) (0.46,0.74156196) (0.47,0.76347715) (0.48,0.78509195) (0.49,0.806315) (0.5,0.82704884) (0.51,0.84718992) (0.52,0.86662877) (0.53,0.88525012) (0.54,0.90293316) (0.55,0.91955184) (0.56,0.93497524) (0.57,0.94906806) (0.58,0.96169114) (0.59,0.97270211) (0.6,0.98195612) (0.61,0.98932837) (0.62,0.99478327) (0.63,0.99831025) (0.64,0.99990247) (0.65,0.99955685) (0.66,0.99727406) (0.67,0.99305852) (0.68,0.9869184) (0.69,0.9788656) (0.7,0.96891573) (0.71,0.95708807) (0.72,0.94340555) (0.73,0.92789467) (0.74,0.91058551) (0.75,0.8915116) (0.76,0.87070992) (0.77,0.84822077) (0.78,0.82408775) (0.79,0.79835762) (0.8,0.77108024) (0.81,0.74230849) (0.82,0.71209812) (0.83,0.68050769) (0.84,0.64759841) (0.85,0.61343406) (0.86,0.57808086) (0.87,0.54160732) (0.88,0.50408412) (0.89,0.465584) (0.9,0.42618156) (0.91,0.38595317) (0.92,0.34497679) (0.93,0.30333183) (0.94,0.26109901) (0.95,0.21836016) (0.96,0.17519813) (0.97,0.13169656) (0.98,0.087939762) (0.99,0.04401253) (1,0)};
\end{axis}\end{tikzpicture}\end{center}

\begin{center}\begin{tikzpicture}
\begin{axis}[width=0.78\textwidth,height=3cm,xlabel={$x/a$},ylabel={$\psi_{5}$},grid=major,legend style={font=\small},xmin=-1,xmax=1,ymin=-1.2,ymax=1.2]
\addplot+[mark=none] coordinates{(-1,0) (-0.99,0.074991834) (-0.98,0.1493909) (-0.97,0.2226091) (-0.96,0.29406769) (-0.95,0.36320182) (-0.94,0.42946503) (-0.93,0.49233354) (-0.92,0.5513104) (-0.91,0.60592943) (-0.9,0.65575889) (-0.89,0.70040492) (-0.88,0.73951459) (-0.87,0.77277878) (-0.86,0.79993455) (-0.85,0.82076724) (-0.84,0.83511218) (-0.83,0.84285598) (-0.82,0.84393744) (-0.81,0.83834799) (-0.8,0.82613183) (-0.79,0.80738552) (-0.78,0.78225724) (-0.77,0.75094561) (-0.76,0.71369813) (-0.75,0.67080923) (-0.74,0.62261793) (-0.73,0.56950515) (-0.72,0.51189071) (-0.71,0.45023005) (-0.7,0.38501054) (-0.69,0.31674772) (-0.68,0.24598116) (-0.67,0.17327026) (-0.66,0.099189735) (-0.65,0.024325169) (-0.64,-0.050731675) (-0.63,-0.12538751) (-0.62,-0.19905222) (-0.61,-0.27114353) (-0.6,-0.34109158) (-0.59,-0.40835164) (-0.58,-0.47243915) (-0.57,-0.53291723) (-0.56,-0.58939156) (-0.55,-0.64151154) (-0.54,-0.68897092) (-0.53,-0.73150806) (-0.52,-0.76890577) (-0.51,-0.80099075) (-0.5,-0.82763273) (-0.49,-0.84874331) (-0.48,-0.86427456) (-0.47,-0.87421731) (-0.46,-0.87859934) (-0.45,-0.87748339) (-0.44,-0.87096499) (-0.43,-0.85917021) (-0.42,-0.84225341) (-0.41,-0.82039482) (-0.4,-0.79379822) (-0.39,-0.76268852) (-0.38,-0.72730939) (-0.37,-0.68792099) (-0.36,-0.64479766) (-0.35,-0.59822571) (-0.34,-0.54850132) (-0.33,-0.49592844) (-0.32,-0.44081693) (-0.31,-0.38348063) (-0.3,-0.32423563) (-0.29,-0.26339869) (-0.28,-0.20128564) (-0.27,-0.13820999) (-0.26,-0.074481647) (-0.25,-0.010405678) (-0.24,0.053718768) (-0.23,0.11759945) (-0.22,0.18095196) (-0.21,0.24350046) (-0.2,0.30497847) (-0.19,0.36512944) (-0.18,0.42370734) (-0.17,0.48047711) (-0.16,0.53521506) (-0.15,0.58770928) (-0.14,0.63775985) (-0.13,0.68517916) (-0.12,0.72979205) (-0.11,0.771436) (-0.1,0.80996128) (-0.09,0.84523097) (-0.08,0.87712113) (-0.07,0.90552076) (-0.06,0.93033192) (-0.05,0.9514697) (-0.04,0.96886222) (-0.03,0.98245071) (-0.02,0.99218943) (-0.01,0.99804572) (0,1) (0.01,0.99804572) (0.02,0.99218943) (0.03,0.98245071) (0.04,0.96886222) (0.05,0.9514697) (0.06,0.93033192) (0.07,0.90552076) (0.08,0.87712113) (0.09,0.84523097) (0.1,0.80996128) (0.11,0.771436) (0.12,0.72979205) (0.13,0.68517916) (0.14,0.63775985) (0.15,0.58770928) (0.16,0.53521506) (0.17,0.48047711) (0.18,0.42370734) (0.19,0.36512944) (0.2,0.30497847) (0.21,0.24350046) (0.22,0.18095196) (0.23,0.11759945) (0.24,0.053718768) (0.25,-0.010405678) (0.26,-0.074481647) (0.27,-0.13820999) (0.28,-0.20128564) (0.29,-0.26339869) (0.3,-0.32423563) (0.31,-0.38348063) (0.32,-0.44081693) (0.33,-0.49592844) (0.34,-0.54850132) (0.35,-0.59822571) (0.36,-0.64479766) (0.37,-0.68792099) (0.38,-0.72730939) (0.39,-0.76268852) (0.4,-0.79379822) (0.41,-0.82039482) (0.42,-0.84225341) (0.43,-0.85917021) (0.44,-0.87096499) (0.45,-0.87748339) (0.46,-0.87859934) (0.47,-0.87421731) (0.48,-0.86427456) (0.49,-0.84874331) (0.5,-0.82763273) (0.51,-0.80099075) (0.52,-0.76890577) (0.53,-0.73150806) (0.54,-0.68897092) (0.55,-0.64151154) (0.56,-0.58939156) (0.57,-0.53291723) (0.58,-0.47243915) (0.59,-0.40835164) (0.6,-0.34109158) (0.61,-0.27114353) (0.62,-0.19905222) (0.63,-0.12538751) (0.64,-0.050731675) (0.65,0.024325169) (0.66,0.099189735) (0.67,0.17327026) (0.68,0.24598116) (0.69,0.31674772) (0.7,0.38501054) (0.71,0.45023005) (0.72,0.51189071) (0.73,0.56950515) (0.74,0.62261793) (0.75,0.67080923) (0.76,0.71369813) (0.77,0.75094561) (0.78,0.78225724) (0.79,0.80738552) (0.8,0.82613183) (0.81,0.83834799) (0.82,0.84393744) (0.83,0.84285598) (0.84,0.83511218) (0.85,0.82076724) (0.86,0.79993455) (0.87,0.77277878) (0.88,0.73951459) (0.89,0.70040492) (0.9,0.65575889) (0.91,0.60592943) (0.92,0.5513104) (0.93,0.49233354) (0.94,0.42946503) (0.95,0.36320182) (0.96,0.29406769) (0.97,0.2226091) (0.98,0.1493909) (0.99,0.074991834) (1,0)};
\end{axis}\end{tikzpicture}\end{center}

\end{solution}
\originalsection{第2题：径向方程无量纲化（15分）}
\begin{exercise}
对$[-\hbar^2\mathrm d^2/(2m\mathrm dr^2)+\hbar^2\ell(\ell+1)/(2mr^2)-e^2/r]\psi=E\psi$, 用$a_0=\hbar^2/(me^2)$定义无量纲$u,\mathcal E$, 把二阶导数系数化为$-1$, 给坐标、能量关系及完整方程.
\end{exercise}
\begin{solution}
取$u=r/a_0$, $E=(\hbar^2/(2ma_0^2))\mathcal E$, 得$[-\partial_u^2+\ell(\ell+1)/u^2-2/u]\psi=\mathcal E\psi$. 原题把约化径向函数记作$\psi$; 它对应本册$u(r)=rR(r)$, 不是三维点密度的完整波函数. 坐标尺度也可选其他倍数, 只要三项与能量一起一致缩放.
\end{solution}
\originalsection{第3题：谐振子矩阵元（15分）}
\begin{exercise}
(a)求$\langle n|x^3|n\rangle$. (b)求$\langle0|x^{10}|10\rangle$.
\end{exercise}
\begin{solution}
$x^3$为奇算符, 定态有确定宇称, 所以(a)为0. 令$c=\sqrt{\hbar/(2m\omega)}$, $x^{10}=c^{10}(a+a^\dagger)^{10}$. 从10步到0必须每步下降, 只保留$a^{10}$路径, 即使$a,a^\dagger$不对易也没有其他净下降10的项. 因而(b)$=c^{10}\sqrt{10!}$. 不可以把非对易二项式当普通二项式展开计算一般矩阵元; 本题由唯一长度路径直接判定.
\end{solution}
\originalsection{第4题：一个深谷与两个深谷（20分）}
\begin{exercise}
偶势在远处为$V_0$, 零点附近$V\simeq\alpha x^2/2$,$\alpha>0$. (a)估计基态能量与准确条件. (b)把两个谷放在$\pm x_0$, 估计前两态能量、画波函数并给分离条件. (c)合并为$\cdots\ll\alpha\ll\cdots$.
\end{exercise}
\begin{solution}
局部频率$\omega=\sqrt{\alpha/m}$, $E_g\simeq\hbar\omega/2$. 需要$\hbar\omega\ll V_0$, 且二次近似在波函数宽度$d=(\hbar^2/(m\alpha))^{1/4}$内准确; 单给高度不能保证后一个形状条件. 两谷远隔$x_0\gg d$时, 局域基态$g_L,g_R$近正交, 前两态为$(g_L\pm g_R)/\sqrt{2(1\pm S)}$, $S=\langle g_L|g_R\rangle$小. 基态偶、第一激发奇, 能量均近$E_g$, 微小隧穿劈裂使偶态更低: $E_\pm\simeq E_g\mp|t|$. 原题未给整段势形, 不能算确切$t$. 波函数即左右两个同号Gaussian峰或异号峰, 奇态中央节点, 偶态无节点. 两个量级条件合并为
\[
\frac{\hbar^2}{mx_0^4}\ll\alpha\ll\frac{mV_0^2}{\hbar^2}.
\]
这里用$\ll$表达尺度分离, 不强调数量级条件中的2或4. 用更大$\alpha$缩窄波包有利于两谷分离, 却提高局部能量, 这正是双边限制的来源.
\begin{center}\begin{tikzpicture}[x=1cm,y=1cm]
\draw[->](-3,0)--(3,0)node[right]{$x$};
\draw[main,domain=-3:3,samples=100]plot(\x,{exp(-5*(\x+1.5)^2)+exp(-5*(\x-1.5)^2)});
\draw[dashed,domain=-3:3,samples=100]plot(\x,{exp(-5*(\x+1.5)^2)-exp(-5*(\x-1.5)^2)});
\node at(-1.5,-1.3){$-x_0$};\node at(1.5,-1.3){$x_0$};
\node at(0,1.4){实线：偶基态；虚线：奇第一激发态（示意）};
\end{tikzpicture}\end{center}
\end{solution}
\originalsection{第5题：方阱趋于\texorpdfstring{$\delta$}{δ}势（30分）}
\begin{exercise}
用宽$2a$、深$V_0$的方阱实现$-\alpha\delta(x)$,$\alpha>0$. (a)固定面积求$V_0(a)$. (b)求$z_0^2$及窄阱极限. (c)对小$a$求偶基态匹配参数$\eta,\xi$主阶. (d)推出束缚能, 检查是否依赖调节器$a$.
\end{exercise}
\begin{solution}
面积$2aV_0=\alpha$给$V_0=\alpha/(2a)$, $z_0^2=2mV_0a^2/\hbar^2=m\alpha a/\hbar^2\to0$. 深度虽发散, 相比动能尺度$\hbar^2/(ma^2)$却趋弱, 所以无量纲深度趋0并不矛盾. 偶态满足$\xi=\eta\tan\eta$,$\eta^2+\xi^2=z_0^2$. 小$z_0$时$\tan\eta\simeq\eta$, 因而$\eta\simeq z_0$, $\xi\simeq z_0^2$. 于是$E=-\hbar^2\xi^2/(2ma^2)\to-m\alpha^2/(2\hbar^2)$, 正确且$a$消去. 其衰减常数$\xi/a\to m\alpha/\hbar^2$, 波函数归一化后趋$\sqrt{m\alpha/\hbar^2}\e^{-m\alpha|x|/\hbar^2}$. 极限保持的是面积$\alpha$而非总面积1, 只有单位强度的$\delta$核自身面积为1.
\end{solution}

\originalchapter{8.04 2015期末：量子模型的连接}
原件为Fall 2015 Final Exam, 2015年12月18日13:30--16:30, 六题105分, 闭卷禁计算器, 发布于Spring 2016 OCW. PDF第2--6页公式表已保存. 以下逐题覆盖全部小问, 均为\textbf{AI独立解答}.
\originalsection{第1题：一维维里定理（15分）}
\begin{exercise}
(a)解释能量本征态的任意$\langle[H,\Omega]\rangle=0$. (b)取$\Omega=xp$, 推出$\langle T\rangle$与$V'$关系. (c)应用于谐振子.
\end{exercise}
\begin{solution}
在相关态属于运算定义域、矩收敛的条件下,$\langle H\Omega\rangle=E\langle\Omega\rangle=\langle\Omega H\rangle$. $[xp,H]=\ii\hbar(2T-xV')$, 所以$2\langle T\rangle=\langle xV'\rangle$. 对$V=m\omega^2x^2/2$, $xV'=2V$, 得$\langle T\rangle=\langle V\rangle=E/2$. “任意”算符仍需定义域条件, 对硬墙或奇异势的边界项须另查.
\end{solution}
\originalsection{第2题：氢原子的类圆轨道（15分）}
\begin{exercise}
固定$n$, $\ell=m=n-1$. (a)写波函数到无量纲归一常数. (b)求径向概率密度峰位置. (c)用$|Y_{\ell\ell}|^2\propto\sin^{2\ell}\theta$描述大$n$时电子最可能出现的区域.
\end{exercise}
\begin{solution}
径向多项式次数$n-\ell-1=0$, 所以$\psi=N a_0^{-3/2}(r/a_0)^{n-1}\e^{-r/(na_0)}Y_{n-1,n-1}(\theta,\phi)$. 径向密度$P(r)\propto r^{2n}\e^{-2r/(na_0)}$含球壳体积$r^2$, 不是点密度; 对数导数$2n/r-2/(na_0)$给峰$r=n^2a_0$. 对大$\ell$, 令$\theta=\pi/2+\epsilon$, $\sin^{2\ell}\theta\simeq\e^{-\ell\epsilon^2}$, 集中在赤道而对$\phi$均匀. 径向相对宽度约$1/\sqrt{2n}$, 因而是薄圆环附近的概率分布, 不等于一个沿确定轨迹运动的点. $n=1$时$\ell=0$方向均匀, 不适用窄环描述.
\begin{center}\begin{tikzpicture}\begin{axis}[width=.72\textwidth,height=4cm,xlabel={$\theta$},ylabel={相对角密度},xmin=0,xmax=3.14159,xtick={0,1.5708,3.14159},xticklabels={$0$,$\pi/2$,$\pi$},ymin=0,ymax=1.1]
\addplot[main,domain=0:3.14159,samples=120]{sin(deg(x))^40};
\end{axis}\end{tikzpicture}\end{center}
\end{solution}
\originalsection{第3题：出射波包（15分）}
\begin{exercise}
$x>R$定态$\psi=\e^{\ii\delta(k)}\sin(kx+\delta(k))$. (a)分为入射及出射波. (b)入射波包为
\[
\int_0^\infty f(k)\e^{-\ii kx-\ii E(k)t/\hbar}\dd k,
\]
写出射波包. (c)用驻定相位求其峰轨迹.
\end{exercise}
\begin{solution}
$\psi=[\e^{2\ii\delta}\e^{\ii kx}-\e^{-\ii kx}]/(2\ii)$, 所以相对于所给入射振幅, 反射振幅为$-\e^{2\ii\delta}$. 出射为$-\int f(k)\e^{\ii kx+2\ii\delta(k)-\ii E(k)t/\hbar}\dd k$. 若$f=|f|\e^{\ii\chi}$在$k_0$附近窄且平滑, 定相位得$x=v_g t-2\delta'(k_0)-\chi'(k_0)$,$v_g=E'(k_0)/\hbar$. 原题常以$f$相位不变为基准, 则$x=v_gt-2\delta'$. 相比自由反射, 延迟$2\delta'/v_g=2\hbar\mathrm d\delta/\mathrm dE$. 必须保留$f$中可能的初始位置相位, 不能对任意复$f$都把初始峰假定在0.
\end{solution}
\originalsection{第4题：趋近完美的无损测弹（20分）}
\begin{exercise}
$U=\left(\begin{smallmatrix}\cos\alpha&\ii\sin\alpha\\\ii\sin\alpha&\cos\alpha\end{smallmatrix}\right)$,$\alpha=\pi/(2N)$, 两侧相等距离镜面使光子重复击中分束器, 初始在左侧. (a)求$U^k$. (b)求第$k$次后左右概率及$k=N$. (c)右侧加入每次吸收光子的探测器, 求$N$次后左侧概率及被探测概率. (d)大$N$估计. (e)右侧换为有效则吸收引爆、坏则透明的EV弹, 解释找到左侧光子的结论及$N=250$爆炸概率.
\end{exercise}
\begin{solution}
记$\sigma_x=\left(\begin{smallmatrix}0&1\\1&0\end{smallmatrix}\right)$, $\sigma_x^2=I$. 题给乘法公式或$U=\e^{\ii\alpha\sigma_x}$给$U^k=\cos(k\alpha)I+\ii\sin(k\alpha)\sigma_x$. 镜面对两支路相同的反射相位只形成整体相位. 无探测器时$p_L=\cos^2(k\alpha),p_R=\sin^2(k\alpha)$, $N$次后光子完全在右. 每次右侧吸收时, 留存左振幅乘$\cos\alpha$, 所以$P_L(N)=\cos^{2N}\alpha$, $P_D=1-P_L$. 大$N$有$P_L\simeq\exp[-\pi^2/(4N)]\simeq1-\pi^2/(4N)$, $P_D\simeq\pi^2/(4N)$.

坏弹等同无探测器, 理想$N$次后不可能在左; 有效弹等同吸收器, 吸收就引爆. 因而未爆炸且找到左光子即可认证有效, 概率趋1. 爆炸概率$P_E=1-\cos^{2N}[\pi/(2N)]$. $N=250$精确约0.0098212, 首阶约0.0098696, 即约0.98\%. 极限通过多次弱分束及每次吸收条件抑制转移（量子Zeno效果）实现, 与仅把单次不对称MZ反射率推高的$1/2$上界是不同装置.
\end{solution}
\originalsection{第5题：小圆柱上的额外维度（20分）}
\begin{exercise}
粒子位于长$a$、周长$L$的圆柱表面, 用矩形坐标$0<x<a$, $0\leq y<L$且两条水平边识别. (a)分离方程及边界. (b)归一态、能量与量子数范围. (c)基态能量. (d)无偶然简并时哪些能量与一维段相同. (e)最低的新增能级. (f)$L=a/1000$, $\hbar^2/(2ma^2)=1$ eV时探测额外维度所需能量尺度.
\end{exercise}
\begin{solution}
$\psi=X(x)Y(y)$给$-\hbar^2X''/(2m)=E_xX$, $-\hbar^2Y''/(2m)=E_yY$, $E=E_x+E_y$. $X(0)=X(a)=0$, $Y(y+L)=Y(y)$且导数也周期. 解为
\[
\psi_{n\ell}=\sqrt{\frac2{aL}}\sin\frac{n\pi x}a\,\e^{2\pi\ii\ell y/L},\qquad
E_{n\ell}=\frac{\hbar^2}{2m}\left[\frac{n^2\pi^2}{a^2}+\frac{4\pi^2\ell^2}{L^2}\right],
\]
$n=1,2,\ldots$,$\ell\in\mathbb Z$. 基态$(1,0)$能量$\pi^2\hbar^2/(2ma^2)$. 与一维段相同的是全部$\ell=0$态; 最低新增能量$E_{1,\pm1}$二重, 为基态加$4\pi^2\hbar^2/(2mL^2)$. 额外方向激发差约$4\pi^2\times10^6$ eV$=39.48$ MeV, 加上约9.87 eV的基态能量可忽略. 这是给定非相对论模型的能量尺度, 若该值已不远小于粒子静能, 真正实验还须用相对论动力学修正.
\end{solution}
\originalsection{第6题：双\texorpdfstring{$\delta$}{δ}势的完全透射（20分）}
\begin{exercise}
$V=g\delta(x+a)+g\delta(x-a)$,$\lambda=mag/\hbar^2\geq0$. 左、中、右分别取$A\e^{\ii kx}+B\e^{-\ii kx}$,$C\e^{\ii kx}+D\e^{-\ii kx}$,$F\e^{\ii kx}$. (a)完全透射要求哪个系数0. (b)写四个匹配式. (c)给定共振条件$\xi\cot\xi=-2\lambda$,$\xi=2ka$, 画至$3\pi$, 求弱、强势根近似. (d)用$C/D=-1/\cos2ka$,$C=(1+\lambda/(\ii ka))A$, 求强势第一个共振的内波与概率密度图.
\end{exercise}
\begin{solution}
无损散射同侧波数相同,$T=1$等价$B=0$. 写$u=ka$, 连续与跳导数给
\[
\begin{aligned}
C+D\e^{2\ii u}&=A,& C+D\e^{-2\ii u}&=F,\\
C-D\e^{2\ii u}&=(1+2\lambda/(\ii u))A,&
C-D\e^{-2\ii u}&=(1-2\lambda/(\ii u))F.
\end{aligned}
\]
第一个导数跳跃为内减外, 第二个为外减内, 因此符号不同. $\lambda>0$时每个$\xi$根在$((n+1/2)\pi,(n+1)\pi)$,$n=0,1,\ldots$. 弱势$u\simeq(2n+1)\pi/4$; 强势$u\simeq(n+1)\pi/2$（从下侧接近）. $\lambda=0$完全无势, 所有能量透射1, 不能仅取极限离散根说只有这些能量透射.
\begin{center}\begin{tikzpicture}\begin{axis}[width=.8\textwidth,height=5cm,xlabel={$\xi$},ylabel={$\xi\cot\xi$},xmin=0,xmax=9.425,ymin=-14,ymax=8,restrict y to domain=-16:10,unbounded coords=jump]
\addplot[main,domain=.001:3.13,samples=100]{x*cot(deg(x))};
\addplot[main,domain=3.15:6.27,samples=100]{x*cot(deg(x))};
\addplot[main,domain=6.3:9.41,samples=100]{x*cot(deg(x))};
\addplot[dashed,domain=0:9.425]{-.2};\addplot[dotted,domain=0:9.425]{-10};
\node at(axis cs:1,-.9){$-2\lambda=-0.2$};\node at(axis cs:2,-11){$-2\lambda=-10$};
\end{axis}\end{tikzpicture}\end{center}
强势首根$\xi\simeq\pi-\pi/(2\lambda)$, $u\simeq\pi/2-\pi/(4\lambda)$. $\cos2u\simeq-1$, $D\simeq C\simeq-\ii\lambda/u$, 所以内波$\psi\simeq-2\ii\lambda\cos(kx)/u$. 内区概率最大$4\lambda^2/u^2\simeq16\lambda^2/\pi^2$, 很大; 外面$|\psi|^2=|A|^2=|F|^2=1$. 内部大驻波和完全透射相容, 表示粒子在双障碍间长时间储存. 下图为强势首根的主阶示意, 边界附近需要保留次阶项才能精确连接外部1.
\begin{center}\begin{tikzpicture}\begin{axis}[width=.78\textwidth,height=4cm,xlabel={$x/a$},ylabel={相对密度},xmin=-2,xmax=2,ymin=0,ymax=1.2]
\addplot[main,domain=-1:1,samples=80]{cos(deg(pi*x/2))^2};
\addplot[main,domain=-2:-1]{.02};\addplot[main,domain=1:2]{.02};
\node at(axis cs:0,1.12){内部按峰值缩放；外部很小的常值};
\end{axis}\end{tikzpicture}\end{center}
\end{solution}

\originalchapter{8.05 作业7：动力学表象与相干态}
来源为8.05 Fall 2013 Assignment 7, 2013年10月24日发布, 11月1日15时截止, 七题共65分. 题面全部小问取自OCW原件, 解答为\textbf{AI独立解答}. 第1、2题$\lambda$是磁矩与自旋的比例系数, 第5题局部同名$\lambda$是展宽系数, 与其他章的无量纲扰动参数不同.
\originalsection{第1题：振荡磁场中的自旋（10分）}
\begin{exercise}
$\boldsymbol B=B_0\cos\omega t\,\hat z$, $H=-\lambda B_0\cos\omega t\,S_z$, 初态为$S_x=\hbar/2$. (a)求演化算符, 注意各时刻$H$对易. (b)求态及方向角. (c)求$S_x=-\hbar/2$概率. (d)求允许完全翻转的最大频率.
\end{exercise}
\begin{solution}
令$\chi(t)=\lambda B_0\sin\omega t/\omega$（$\omega=0$按连续极限）. 因为所有$H$都是$S_z$的实倍数,$U=\exp[\ii\chi S_z/\hbar]$. 从$\ket{x;+}=(\ket++\ket-)/\sqrt2$演化为$(\e^{\ii\chi/2}\ket++\e^{-\ii\chi/2}\ket-)/\sqrt2$. 自旋方向$\theta=\pi/2$, $\phi=-\chi$模$2\pi$, 并非磁场振幅本身作为角速度积分后仍为正. 翻转概率为$\sin^2(\chi/2)$. 完全翻转要求$|\chi|$能达到$\pi$, 所以正频率最大为$\omega_{\max}=|\lambda B_0|/\pi$. 若磁矩或场为0, 不存在完全翻转; 可使多个周期内相位来回变化, 但不能累计突破这个振幅界.
\end{solution}
\originalsection{第2题：Heisenberg自旋算符（5分）}
\begin{exercise}
常场$H=-\lambda BS_z$, 求$S_{xH}(t),S_{yH}(t),S_{zH}(t)$.
\end{exercise}
\begin{solution}
$O_H=U^\dagger OU$, 或直接$\dot S_{xH}=\lambda BS_{yH}$,$\dot S_{yH}=-\lambda BS_{xH}$. 故$S_{xH}=S_x\cos\lambda Bt+S_y\sin\lambda Bt$,$S_{yH}=S_y\cos\lambda Bt-S_x\sin\lambda Bt$,$S_{zH}=S_z$. 算符时间变, Heisenberg态固定; 与Schrödinger表象得到相同实验均值.
\end{solution}
\originalsection{第3题：均值与Newton方程（10分）}
\begin{exercise}
(a)对$H=p^2/(2m)+V(x)$推Heisenberg方程、Ehrenfest定理及经典化条件. (b)自由初态均值$x_0,p_0$, 求$\langle x\rangle(t)$. (c)改$V=qE_0x\sin\omega t$, 求不同时刻$H$对易子并解释Ehrenfest仍成立. (d)求新位置均值.
\end{exercise}
\begin{solution}
由$\dot O_H=U^\dagger(\ii[H,O]/\hbar)U$得$\dot x_H=p_H/m$, $\dot p_H=-V'(x_H)$; 对固定初态取平均得$\ddot{\langle x\rangle}=-\langle V'(x)\rangle/m$. 若$V'$是仿射函数, 精确等于$-V'(\langle x\rangle)/m$; 窄波包且力变化缓慢时可近似. 自由时$\langle x\rangle=x_0+p_0t/m$.

含时势给$[H(t_1),H(t_2)]=\ii\hbar qE_0[\sin\omega t_1-\sin\omega t_2]p/m$, 一般非零; 题面的“不同时刻非零”有$\sin\omega t_1=\sin\omega t_2$等特例, 应按一般不对易理解. 演化方程仍是$\ii\hbar\dot U=H(t)U$, 不需假定普通无时间排序的指数, 因而均值定理仍成立. 两次积分得到
\[
\langle p\rangle=p_0+\frac{qE_0}{\omega}(\cos\omega t-1),\qquad
\langle x\rangle=x_0+\frac{p_0t}m+\frac{qE_0}{m\omega^2}(\sin\omega t-\omega t).
\]
这里题给势的力为$-qE_0\sin\omega t$, 不能改成相反号. $\omega\to0$时力也趋0, 表达式连续还原自由均值.
\end{solution}
\originalsection{第4题：三维维里定理（10分）}
\begin{exercise}
中心势$H=\boldsymbol p^2/(2m)+V(r)$, $\Omega=\boldsymbol r\cdot\boldsymbol p$. (a)求$\dot\Omega_H$. (b)解释定态下时间无显式依赖算符的均值恒定. (c)对$V=c/r^k$证$\langle T\rangle=-(k/2)\langle V\rangle$.
\end{exercise}
\begin{solution}
逐方向应用对易规则得$\dot\Omega_H=\boldsymbol p_H^2/m-\boldsymbol r_H\cdot\boldsymbol\nabla V(r_H)=2T_H-r_HV'(r_H)$. 定态的Schrödinger波函数只乘相位, 所以任意固定观测量均值不变; 表象等价给$\mathrm d\langle\Omega_H\rangle/\mathrm dt=0$. 从而$2\langle T\rangle=\langle rV'\rangle=-k\langle V\rangle$. 原题$\Omega$不是Hermitian但定态对易平均仍可用; 改对称的$(\boldsymbol r\cdot\boldsymbol p+\boldsymbol p\cdot\boldsymbol r)/2$更适合作为观测量. 对奇异势须检查正规束缚态与原点边界条件, 不是任意$k,c$都自动存在这种态.
\end{solution}
\originalsection{第5题：恒力下的展宽（10分）}
\begin{exercise}
$H=p^2/(2m)+gx$. (a)积分Heisenberg方程给$x_H,p_H$. (b)初态$\psi=N\e^{-x^2/(2\Delta^2)}$,$\Delta>0$, 求位置均值. (c)求位置方差二次增长系数, 说明其对$g$的依赖.
\end{exercise}
\begin{solution}
$p_H=p-gt$, $x_H=x+pt/m-gt^2/(2m)$, 所以$\ddot x_H=-g/m$. 初态两均值均0, $\langle x_H\rangle=-gt^2/(2m)$, 与经典同初值一致. 方差中确定的$gt^2/(2m)$消去, 得
\[
(\Delta x(t))^2=(\Delta x(0))^2+\frac{t}{m}\langle xp+px\rangle_0+\frac{t^2}{m^2}(\Delta p(0))^2
=\frac{\Delta^2}{2}+\frac{\hbar^2t^2}{2m^2\Delta^2}.
\]
实偶Gaussian对称相关矩为0, 所以$\lambda=\hbar^2/(2m^2\Delta^2)$. 均匀力只加整体动量偏移、改变中心轨迹, 不改该展宽系数; 对空间不均匀力则不成立.
\end{solution}
\originalsection{第6题：位移谐振子（10分）}
\begin{exercise}
$H=p^2/(2m)+m\omega^2x^2/2-Fx$, $F>0$. (a)求基态能量和位置均值. (b)把新基态写作$N\e^{\alpha a^\dagger}\ket0$, 求实$\alpha$.
\end{exercise}
\begin{solution}
完成平方, 令$x_c=F/(m\omega^2)$,$y=x-x_c$, 得$H=p^2/(2m)+m\omega^2y^2/2-F^2/(2m\omega^2)$. 故$E_g=\hbar\omega/2-F^2/(2m\omega^2)$, $\langle x\rangle=x_c$. 新湮灭算符$a_y=a-\sqrt{m\omega/(2\hbar)}x_c=a-\alpha$. 要求$a_y\ket{0'}=0$给$\alpha=F/\sqrt{2m\hbar\omega^3}$. 用$a\e^{\alpha a^\dagger}\ket0=\alpha\e^{\alpha a^\dagger}\ket0$检验. 若归一且相位取实正,$N=\e^{-\alpha^2/2}$.
\end{solution}
\originalsection{第7题：相干态位置波函数（10分）}
\begin{exercise}
$\ket\alpha=\exp(\alpha a^\dagger-\alpha^*a)\ket0$,$\alpha=x_0/(\sqrt2d)+\ii p_0d/(\sqrt2\hbar)$,$d=\sqrt{\hbar/(m\omega)}$. 求显式归一化波函数.
\end{exercise}
\begin{solution}
指数代入$x,p$给$\alpha a^\dagger-\alpha^*a=\ii(p_0x-x_0p)/\hbar$. 因其两项对易子为常数, BCH给$D=\e^{\ii p_0x/\hbar}\e^{-\ii x_0p/\hbar}\e^{-\ii p_0x_0/(2\hbar)}$. 平移算符将$\psi_0(x)$变为$\psi_0(x-x_0)$, 故
\[
\psi_\alpha(x)=\frac1{(\pi d^2)^{1/4}}\exp\!\left[-\frac{(x-x_0)^2}{2d^2}+\frac{\ii p_0}{\hbar}\left(x-\frac{x_0}{2}\right)\right].
\]
模方是宽$d/\sqrt2$的归一Gaussian, 平均位置$x_0$、动量$p_0$, 相位常数$-p_0x_0/2$来自题给位移算符约定, 可用BCH核对而非随意省略.
\end{solution}

\originalchapter{8.05 作业9：自旋与分子模型}
来源为8.05 Fall 2013 Assignment 9, 2013年11月8日发布, 11月15日15时截止, 六题共65分. 以下全部小问为\textbf{AI独立解答}.
\originalsection{第1、2题：方向均值与旋转（各10分）}
\begin{exercise}
第1题: 对$\ket{\boldsymbol n;+}$,$\boldsymbol n=(\sin\theta_0\cos\phi_0,\sin\theta_0\sin\phi_0,\cos\theta_0)$, 求$\langle\boldsymbol S\rangle$及另一方向$\boldsymbol n'$上的均值. 第2题: $R_\alpha(\boldsymbol n)=\exp[-\ii\alpha\boldsymbol n\cdot\boldsymbol S/\hbar]$. (a)展开$R_\epsilon^\dagger\boldsymbol SR_\epsilon$到一阶. (b)用均值说明它主动旋转态方向.
\end{exercise}
\begin{solution}
Pauli矩阵乘方向态直接给$\langle\boldsymbol S\rangle=(\hbar/2)\boldsymbol n$, 所以$\langle\boldsymbol S\cdot\boldsymbol n'\rangle=(\hbar/2)\boldsymbol n\cdot\boldsymbol n'$. 各分量算符不对易, 这个向量是三个分别进行的测量均值, 不表示一次测得三个确定分量.

展开得到$R_\epsilon^\dagger\boldsymbol SR_\epsilon=\boldsymbol S+(\ii\epsilon/\hbar)[\boldsymbol n\cdot\boldsymbol S,\boldsymbol S]+O(\epsilon^2)=\boldsymbol S+\epsilon\boldsymbol n\times\boldsymbol S+O(\epsilon^2)$. 在初始$\ket{\boldsymbol n';+}$取均值得$(\hbar/2)(\boldsymbol n'+\epsilon\boldsymbol n\times\boldsymbol n')$, 正是右手规则绕$\boldsymbol n$主动转$\epsilon$. 大角由连续小旋转积累, 不把算符共轭方向与态主动旋转的符号混淆.
\end{solution}
\originalsection{第3题：氨分子的隧穿与电场（15分）}
\begin{exercise}
在氮原子位于氢平面上下的$\ket\uparrow,\ket\downarrow$基中,$H=\left(\begin{smallmatrix}\mu\varepsilon&-\Delta\\-\Delta&-\mu\varepsilon\end{smallmatrix}\right)$,$\Delta>0$. (a)写$H=-K\boldsymbol n\cdot\boldsymbol\sigma$. (b)求能量与态. (c)讨论$\mu\varepsilon\simeq0$、正强场$\mu\varepsilon\gg\Delta$. (d)画能量. (e)从零场下$\ket\downarrow$开始演化, 然后突然开强场, 描述进一步演化及上态概率.
\end{exercise}
\begin{solution}
$K=\sqrt{\Delta^2+\mu^2\varepsilon^2}$,$\boldsymbol n=(\Delta/K,0,-\mu\varepsilon/K)$, 可取$\phi=0$,$\sin\theta=\Delta/K$,$\cos\theta=-\mu\varepsilon/K$. 基态能量$-K$, 态$\cos(\theta/2)\ket\uparrow+\sin(\theta/2)\ket\downarrow$; 激发态$+K$, 态$-\sin(\theta/2)\ket\uparrow+\cos(\theta/2)\ket\downarrow$. 零场$\theta=\pi/2$, 两态为对称/反对称叠加, 能量$\mp\Delta$. 正强场$\theta\to\pi$, 基态趋$\ket\downarrow$能量$-\mu\varepsilon$, 激发趋$-\ket\uparrow$能量$+\mu\varepsilon$; 负强场相反. 能量曲线在0处最小间隙$2\Delta$, 渐近直线$\pm|\mu\varepsilon|$.
\begin{center}\begin{tikzpicture}\begin{axis}[width=.76\textwidth,height=4cm,xlabel={$\mu\varepsilon/\Delta$},ylabel={$E/\Delta$},xmin=-3,xmax=3,ymin=-3.6,ymax=3.6]
\addplot[main,domain=-3:3,samples=80]{sqrt(1+x*x)};\addplot[main,domain=-3:3,samples=80]{-sqrt(1+x*x)};
\addplot[dashed,domain=-3:3]{abs(x)};\addplot[dashed,domain=-3:3]{-abs(x)};
\end{axis}\end{tikzpicture}\end{center}
零场态$\ket{\psi(t)}=\cos(\Delta t/\hbar)\ket\downarrow+\ii\sin(\Delta t/\hbar)\ket\uparrow$, 等效绕$x$轴按角$-2\Delta t/\hbar$旋转初始$-z$. 设开场时刻$t_s$, 忽略$\Delta$后$H\simeq\mu\varepsilon\sigma_z$, 两分量只加相反相位, $P_\uparrow\simeq\sin^2(\Delta t_s/\hbar)$固定, 不会立刻跳为0或1. 有限强场仍有小幅变化, 上界由$\Delta/|\mu\varepsilon|$控制; 这是突然开场, 不是绝热追随能量基态.
\end{solution}
\originalsection{第4题：三态二氧化碳离子（15分）}
\begin{exercise}
电子可在左氧$\ket L$、碳$\ket C$、右氧$\ket R$间跳跃, 两氧能量$E_O$、碳$E_C$, 无直接氧氧跳跃. (a)引入跳跃能量参数求谱. (b)对$E_C=E_O$求态. (c)继续取$E_C=E_O$, 从$\ket L$出发求三处概率并绘图.
\end{exercise}
\begin{solution}
选跳跃矩阵元$-\Delta$, $\Delta>0$, 则$H=\left(\begin{smallmatrix}E_O&-\Delta&0\\-\Delta&E_C&-\Delta\\0&-\Delta&E_O\end{smallmatrix}\right)$. 反对称氧态$\ket a=(\ket L-\ket R)/\sqrt2$能量$E_O$; 对称态$\ket s=(\ket L+\ket R)/\sqrt2$与碳耦合$-\sqrt2\Delta$, 另两能量为$(E_O+E_C)/2\pm\sqrt{(E_O-E_C)^2/4+2\Delta^2}$. 等能$E_O=E_C=E_*$时态为$\ket a$和$(\ket s\pm\ket C)/\sqrt2$, 后者能量分别$E_*\mp\sqrt2\Delta$.

令$\Omega=\sqrt2\Delta/\hbar$, 省略整体相位, 初态演化$\ket\psi=[\ket a+\cos\Omega t\ket s+\ii\sin\Omega t\ket C]/\sqrt2$. 因而$P_L=(1+\cos\Omega t)^2/4$,$P_R=(1-\cos\Omega t)^2/4$,$P_C=\sin^2\Omega t/2$, 和为1. $t=\pi/\Omega$电子完全到右氧, 但碳的最大概率只1/2; 反对称态始终不参与碳跳跃.
\begin{center}\begin{tikzpicture}\begin{axis}[width=.8\textwidth,height=4.5cm,xlabel={$\Omega t$},ylabel={概率},xmin=0,xmax=6.2832,ymin=0,ymax=1.1,legend pos=north east]
\addplot[main,domain=0:6.2832,samples=100]{(1+cos(deg(x)))^2/4};\addlegendentry{$P_L$}
\addplot[dashed,domain=0:6.2832,samples=100]{(1-cos(deg(x)))^2/4};\addlegendentry{$P_R$}
\addplot[dotted,domain=0:6.2832,samples=100]{sin(deg(x))^2/2};\addlegendentry{$P_C$}
\end{axis}\end{tikzpicture}\end{center}
\end{solution}
\originalsection{第5题：共振核磁自旋（5分）}
\begin{exercise}
$\boldsymbol B(t)=B_0\hat z+B_1(\hat x\cos\omega_0t-\hat y\sin\omega_0t)$,$\omega_0=\gamma B_0$, 初态$\ket+$. (a)求演化及方向角. (b)消去$t$写$\phi(\theta)$, 求首次到赤道的方向.
\end{exercise}
\begin{solution}
令$\omega_1=\gamma B_1$, 转动框架消去$z$场, 得$U=\e^{\ii\omega_0t\sigma_z/2}\e^{\ii\omega_1t\sigma_x/2}$. 所以$\ket\psi=\e^{\ii\omega_0t/2}\cos(\omega_1t/2)\ket++\ii\e^{-\ii\omega_0t/2}\sin(\omega_1t/2)\ket-$. 若$\omega_1>0$且在首次北到南区间,$\theta=\omega_1t$, $\phi=\pi/2-\omega_0t$, 故$\phi=\pi/2-(B_0/B_1)\theta$. 首次赤道$t=\pi/(2\omega_1)$, 方向$(\sin(\pi B_0/(2B_1)),\cos(\pi B_0/(2B_1)),0)$. 任意符号或更晚时间应使用无歧义向量$\boldsymbol n=(\sin\omega_1t\sin\omega_0t,\sin\omega_1t\cos\omega_0t,\cos\omega_1t)$, 再求角; 不将$\theta$无限增长当成标准极角.
\end{solution}
\originalsection{第6题：两自旋单态不随共同旋转改变（10分）}
\begin{exercise}
$\ket\Psi=(\ket+\otimes\ket--\ket-\otimes\ket+)/\sqrt2$. 已知幺正旋转$R\ket+=\ket{\boldsymbol n;+}$, (a)说明$R\ket-$为反方向态到相位. (b)证明$(R\otimes R)\ket\Psi$仅差整体相位.
\end{exercise}
\begin{solution}
幺正保持正交, 而二维中$\ket{\boldsymbol n;+}$的正交补只有一维, 所以$R\ket-=\e^{\ii\eta}\ket{\boldsymbol n;-}$. 更直接, 写$R$的列为$(a,c)^{\mathsf T},(b,d)^{\mathsf T}$, 张量展开交叉项消去, 得$(R\otimes R)\ket\Psi=(ad-bc)\ket\Psi=\det R\ket\Psi$. 一般幺正$|\det R|=1$; 物理自旋旋转$R\in SU(2)$的行列式1, 因而完全不变. 仅旋转其中一个自旋则通常改变关联, 本题是共同旋转.
\end{solution}

\originalchapter{8.05 2011测试：变分与算符}
原件为Fall 2011 Test, 2011年10月26日12:30--14:00, 五题100分, 闭卷禁计算器; 由Fall 2013 OCW提供. 以下为\textbf{AI独立解答}, 判断题加简短理由帮助复习, 原卷仅要求T/F.
\originalsection{第1题：十个判断（20分）}
\begin{exercise}
前三项势为$x<0$的$2V_0$、$0\leq x\leq a$的0、$x>a$的$V_0$, $V_0>0$. 第(4)、(5)项自旋为$1/2$. 判断: (1)$V_0<E<2V_0$每能量有两个本征态; (2)$0<E<V_0$每能量有一个正规态; (3)$E<0$没有能量本征态; (4)$z$向上自旋与$x$向上正交; (5)$S_x^2,S_y^2$对易; (6)两幺正积幺正; (7)投影可有本征值$-1$; (8)能同时对角化的两算符对易; (9)Hermitian算符指数幺正; (10)两Hermitian算符对易子Hermitian.
\end{exercise}
\begin{solution}
依次为\textbf{F,F,T,F,T,T,F,T,F,F}. (1)只有右侧一个传播通道, 左侧衰减条件与匹配留一个散射解. (2)束缚能离散, 不能每个区间能量都有态. (3)$H$非负. (4)重叠幅$1/\sqrt2$. (5)自旋$1/2$的每个平方都是$\hbar^2I/4$. (6)$(UV)^\dagger UV=I$. (7)$P^2=P$使本征值仅0或1. (8)共同基中对角矩阵积可换序. (9)实本征值指数为正数, 一般不是模1; $\exp(\ii A)$才幺正. (10)$[A,B]^\dagger=-[A,B]$, 是反Hermitian（零是特例）.
\end{solution}
\originalsection{第2题：迹与共同基（20分）}
\begin{exercise}
(a)证明$\Tr M=\sum_i\bra{a_i}M\ket{a_i}$与正交归一基无关. (b)给$A_1=\left(\begin{smallmatrix}1&0&1\\0&0&0\\1&0&1\end{smallmatrix}\right)$,$A_2=\left(\begin{smallmatrix}2&1&1\\1&0&-1\\1&-1&2\end{smallmatrix}\right)$. 在$u_1=(1,0,1)/\sqrt2,u_2=(-1,0,1)/\sqrt2,u_3=(0,1,0)$基中补全$A_2$并求其谱.
\end{exercise}
\begin{solution}
有限维时所有和有限; 无穷维须限制迹类算符及合法求和, 不能直接换序任意发散级数. 插入另一完备基得$\sum_i\sum_{jk}\braket{b_i}{a_j}\bra{a_j}M\ket{a_k}\braket{a_k}{b_i}=\sum_jM_{jj}$, 因中间求和为$\delta_{kj}$. 有限矩阵也即$\Tr(U^\dagger MU)=\Tr M$. 逐列算$A_2u_1=3u_1$,$A_2u_2=u_2-\sqrt2u_3$,$A_2u_3=-\sqrt2u_2$, 所以矩阵为$\left(\begin{smallmatrix}3&0&0\\0&1&-\sqrt2\\0&-\sqrt2&0\end{smallmatrix}\right)$, 剩下二阶特征式$\lambda^2-\lambda-2=0$, 谱$3,2,-1$. $A_1$零特征值空间内还要对角化$A_2$, 不能认为给出的每个$A_1$本征矢已经是共同矢.
\end{solution}
\originalsection{第3题：线性绝对值势的上界（20分）}
\begin{exercise}
$V=V_0|x|$, $V_0$量纲为能量/长度, 用$\psi=A\e^{-bx^2/2}$,$b>0$优化基态上界, 化为纯数乘$\hbar,m,V_0$构成的尺度.
\end{exercise}
\begin{solution}
归一后$A=(b/\pi)^{1/4}$,$\langle T\rangle=\hbar^2b/(4m)$,$\langle|x|\rangle=1/\sqrt{\pi b}$. 因而$F(b)=\hbar^2b/(4m)+V_0/\sqrt{\pi b}$. $F'=0$给$b_*=[2mV_0/(\sqrt\pi\hbar^2)]^{2/3}$, 两项在该点势能为动能2倍. 上界
\[
E_g\leq F(b_*)=\frac3{2^{4/3}\pi^{1/3}}\left(\frac{\hbar^2V_0^2}{m}\right)^{1/3}.
\]
该尺度由动能约$\hbar^2/(mL^2)$与势能$V_0L$平衡得到; 变分固定纯数使“量级”变成真正上界.
\end{solution}
\originalsection{第4题：八次势的下界（20分）}
\begin{exercise}
$V=\gamma x^8$, $\gamma>0$. (a)将基态$\langle x^8\rangle$与$\Delta x$联系. (b)用不确定性给基态能量下界并化简尺度.
\end{exercise}
\begin{solution}
偶势一维基态唯一且偶, $\langle x\rangle=0$. 连续两次Cauchy--Schwarz给$\langle x^8\rangle\geq\langle x^4\rangle^2$, $\langle x^4\rangle\geq\langle x^2\rangle^2$, 合起来给$\langle x^8\rangle\geq\langle x^2\rangle^4=(\Delta x)^8$. $\langle p\rangle=0$, 不确定性得$\langle T\rangle\geq\hbar^2/[8m(\Delta x)^2]$. 令$s=\Delta x,A=\hbar^2/(8m)$, $E_g\geq A/s^2+\gamma s^8$. 在正$s$求最小:$s^{10}=A/(4\gamma)$, 动能界项为势能界项4倍, 得
\[
E_g\geq5\,4^{-4/5}\left(\frac{\hbar^2}{8m}\right)^{4/5}\gamma^{1/5}.
\]
这里是在所有允许宽度中最小化一个\textbf{下界}, 与试探态能量的上界不同. 两个不等式不能同时饱和于真实八次势基态, 所以无需期望等号.
\end{solution}
\originalsection{第5题：绕$y$轴进动（20分）}
\begin{exercise}
自旋$1/2$的$H=-\gamma BS_y$, 初方向$\theta_0,\phi_0=0$. (a)写初态. (b)求演化. (c)给方向角函数.
\end{exercise}
\begin{solution}
初态$\cos(\theta_0/2)\ket++\sin(\theta_0/2)\ket-$. 令$b=\gamma Bt$, $U=\cos(b/2)I+\ii\sin(b/2)\sigma_y$, 其矩阵为$\left(\begin{smallmatrix}\cos(b/2)&\sin(b/2)\\-\sin(b/2)&\cos(b/2)\end{smallmatrix}\right)$. 演化为$\cos[(\theta_0-b)/2]\ket++\sin[(\theta_0-b)/2]\ket-$. 方向$\boldsymbol n=(\sin(\theta_0-b),0,\cos(\theta_0-b))$. 标准角$\theta=\arccos\cos(\theta_0-b)$, 当$x$分量非负取$\phi=0$, 负取$\phi=\pi$; 极点$\phi$无意义. 直接让$\theta_0-b$跨过负值而仍叫标准极角会混淆, 向量表达没有这个问题.
\end{solution}

\originalchapter{8.05 2012测试：对称性与三维振子}
原件为Fall 2012 Test, 2012年10月24日12:00--13:30, 五题100分, 闭卷禁计算器; Fall 2013 OCW归档. 以下为\textbf{AI独立解答}.
\originalsection{第1题：十个判断（20分）}
\begin{exercise}
判断: (1)宇称$P\ket x=\ket{-x}$既Hermitian又幺正; (2)$\ket{-x}=-\ket x$; (3)任意$N\times N$矩阵有$N$个独立本征矢; (4)有限维最高能量$E_N$为任意归一态能量均值上界. 随后三项为光滑有限实偶势恰有五束缚态: (5)最高束缚态四节点; (6)非无穷点可同时$\psi=\psi'=0$; (7)激发偶态可在0为0. 最后三项明确为自旋$1/2$, $R=\exp(\ii\pi S_x/\hbar)$: (8)幺正; (9)可对角化; (10)$R=2\ii S_x/\hbar$.
\end{exercise}
\begin{solution}
依次\textbf{T,F,F,T,T,F,F,T,T,T}, 最后三项按卷中自旋$1/2$约定. (1)$P^2=I=P^\dagger P$; (2)两个位置标签不同, 不是同一矢的负号; (3)Jordan块反例; (4)能量均值是本征值的概率加权; (5)一维节点定理; (6)二阶常微分方程零初值唯一为零解; (7)偶态必有$\psi'(0)=0$, 若$\psi(0)=0$也只能零解; (8,9)Hermitian$S_x$的虚指数为正规幺正矩阵; (10)自旋$1/2$时$R=\e^{\ii\pi\sigma_x/2}=\ii\sigma_x$. 对一般自旋$j$最后一个等式不成立, 不能脱离该课程卷的自旋$1/2$范围引用.
\end{solution}
\originalsection{第2题：算符作用的正交分量（15分）}
\begin{exercise}
对归一态与Hermitian$\Omega$, $\Delta\Omega\neq0$, 证明$\Omega\ket\Psi=\langle\Omega\rangle\ket\Psi+\Delta\Omega\ket{\Psi_\perp}$, 其中新矢归一且正交.
\end{exercise}
\begin{solution}
定义$\ket{\Psi_\perp}=(\Omega-\langle\Omega\rangle)\ket\Psi/\Delta\Omega$. 与初态内积为$(\langle\Omega\rangle-\langle\Omega\rangle)/\Delta\Omega=0$; 范数平方为$\langle(\Omega-\langle\Omega\rangle)^2\rangle/(\Delta\Omega)^2=1$. 所以离开原态方向的大小恰是测量宽度. 若宽度0, 态已是本征态, 不需要定义此垂直单位矢.
\end{solution}
\originalsection{第3题：极化恒等式（15分）}
\begin{exercise}
(a)化简$\tfrac14[\bra{u+v}T\ket{u+v}-\bra{u-v}T\ket{u-v}-\ii\bra{u+\ii v}T\ket{u+\ii v}+\ii\bra{u-\ii v}T\ket{u-\ii v}]$. (b)证明复空间中若每个$\bra wT\ket w=0$, 则$T=0$, 并解释实二维90度旋转反例.
\end{exercise}
\begin{solution}
令$A=\bra uT\ket v,B=\bra vT\ket u$, 前两个差为$2(A+B)$, 后两个差带权后为$2(A-B)$, 和除4就是$A$. 若所有对角二次形式都0, 则任意$u,v$有$\bra uT\ket v=0$, 取$u=Tv$给$\norm{Tv}^2=0$, 所以零算符. 实空间没有$\ii v$这两项, 只能得到$A+B$, 反对称90度旋转确可使$v\cdot Tv=0$却非零. 这不要求$T$Hermitian.
\end{solution}
\originalsection{第4题：高次非谐势（25分）}
\begin{exercise}
$V=\hbar^2x^{2n}/(2mL^{2n+2})$,$n\geq2$整数. (a)量纲估基态. (b)用$\psi=\e^{-b^2x^2}$,$b>0$, $c_k=\int\e^{-2y^2}y^{2k}\dd y$, 求变分函数, 不要求优化. (c)画$n\to\infty$势与预期基态能量. (d)给定最优$bL\sim\sqrt n$, 描述试探态是否越好.
\end{exercise}
\begin{solution}
唯一能量尺度$E_L=\hbar^2/(2mL^2)$, 因而$E_g=C_nE_L$. 用$y=bx$归一积分为$c_0/b$, $\langle x^{2n}\rangle=(c_n/c_0)b^{-2n}$. 动能由导数$\psi'=-2b^2x\psi$给$\langle T\rangle=4b^4\hbar^2\langle x^2\rangle/(2m)=\hbar^2b^2/(2m)$, 因$c_1/c_0=1/4$. 所以
\[
F(b)=E_L\left[(bL)^2+\frac{c_n/c_0}{(bL)^{2n}}\right].
\]
极限$|x|<L$势趋0,$|x|>L$趋无限, 端点单个值不影响Dirichlet阱极限; 基态趋宽$2L$盒的$E_L\pi^2/4$, 波函数$\cos(\pi x/(2L))$（内）. Gaussian最优宽度$1/b\sim L/\sqrt n$却缩向0, 比真实分布越来越窄, 因为Gaussian无限尾受高次势强烈惩罚. 试探族越来越差, 不是优化参数保证任何$n$都接近精确态.
\begin{center}\begin{tikzpicture}[x=1.7cm,y=.65cm,>=Stealth]
\draw[->](-1.6,0)--(1.6,0)node[right]{$x/L$};\draw[main](-1,3.4)--(-1,0)--(1,0)--(1,3.4);
\draw[main,domain=-1:1,samples=80]plot(\x,{cos(90*\x)});\node[align=center,font=\footnotesize,text width=3cm]at(0,2.6){实线：盒基态\\点线：Gaussian};
\draw[dotted,domain=-1:1,samples=100]plot(\x,{exp(-16*\x*\x)});\node[below]at(-1,0){$-1$};\node[below]at(1,0){$1$};\end{tikzpicture}\end{center}
\end{solution}
\originalsection{第5题：三维谐振子的角动量（25分）}
\begin{exercise}
三维等频振子有$a_i^\dagger,a_i,N_i$,$i=1,2,3$. (a)用产生算符列基态与前两激发组, 能量与简并. (b)令$a_L=(a_1+\ii a_2)/\sqrt2,a_R=(a_1-\ii a_2)/\sqrt2$, 写$H$, 说明$\{N_L,N_R,N_3\}$完备并求能量. (c)用新产生算符重新列三组. (d)$L_z=\hbar(N_R-N_L)$, 求$[H,L_z]$, 判断$\{H,L_z\}$完备. (e)$L_+=\sqrt2\hbar(a_3^\dagger a_L-a_R^\dagger a_3)$是否Hermitian, 求第二激发组同时被$L_z,L_+$湮灭的唯一态.
\end{exercise}
\begin{solution}
总数$N=0$有$\ket0$, 能量$3\hbar\omega/2$, 一重. $N=1$三态$a_i^\dagger\ket0$, 能量$5\hbar\omega/2$. $N=2$六态$(a_i^\dagger)^2\ket0$及$a_i^\dagger a_j^\dagger\ket0$（$i<j$）, 能量$7\hbar\omega/2$, 未归一允许. 改用$L,R,3$时三组逐项为$\ket0$; $a_L^\dagger\ket0,a_R^\dagger\ket0,a_3^\dagger\ket0$; 第二激发组六态为
\[
\begin{gathered}
(a_L^\dagger)^2\ket0,\quad(a_R^\dagger)^2\ket0,\quad(a_3^\dagger)^2\ket0,\\
a_L^\dagger a_R^\dagger\ket0,\quad a_L^\dagger a_3^\dagger\ket0,\quad a_R^\dagger a_3^\dagger\ket0.
\end{gathered}
\]. 因变换幺正维数不变. $H=\hbar\omega(N_L+N_R+N_3+3/2)$,$E=\hbar\omega(n_L+n_R+n_3+3/2)$, 三独立占有数唯一标记乘积态, 为完备组.

$H,L_z$均为数算符线性组合, 故对易, 却不完备: $a_L^\dagger a_R^\dagger\ket0$和$(a_3^\dagger)^2\ket0$都有$N=2,m=0$. $L_+$伴随为$L_-$, 非Hermitian. 在这二态的组合$[A(a_3^\dagger)^2+B a_L^\dagger a_R^\dagger]\ket0$上作用$L_+$, 得$\sqrt2\hbar(-2A+B)a_R^\dagger a_3^\dagger\ket0$. 所以$B=2A$, 归一态
\[
\ket\psi=\frac{(a_3^\dagger)^2+2a_L^\dagger a_R^\dagger}{\sqrt6}\ket0
=\frac{(a_1^\dagger)^2+(a_2^\dagger)^2+(a_3^\dagger)^2}{\sqrt6}\ket0.
\]
$m=0$且为最高权态只能是$\ell=0$, 故全角动量为0, 不是只有$z$分量为0.
\end{solution}

\originalchapter{8.05 2013测试：尺度变换与测量宽度}
原件为Fall 2013 Test, 2013年10月23日12:30--14:00, 五题100分, 闭卷禁计算器. 以下为\textbf{AI独立解答}.
\originalsection{第1题：十个判断（20分）}
\begin{exercise}
判断: (1)$\Delta Q=0$必为$Q$本征态; (2)时间无显式依赖算符的宽度恒定; (3)$N\geq2$时可有有限$X,P$满足$[X,P]=\ii\hbar I$; (4)两个对易Hermitian算符中$Q_1$的每个本征矢都是$Q_2$本征矢. 后六项为有限维复空间: (5)恒等算符矩阵依赖基; (6)谱与基无关; (7)任何基中伴随矩阵都为共轭转置; (8)向真低维子空间的正交投影在全空间中既非单射又非满射; (9)正交归一组线性独立; (10)幺正且Hermitian者平方为恒等.
\end{exercise}
\begin{solution}
依次\textbf{T,F,F,F,F,T,F,T,T,T}. (1)$(Q-\langle Q\rangle)\ket\psi$范数为0; (2)态仍演化, 除非守恒或特定态; (3)对易子迹0而右边非零; (4)简并子空间内任意矢未必对角化另一算符, 例如$Q_1=I,Q_2=\sigma_z$; (5)每个基中$I$矩阵仍$I$; (6)相似变换不改谱; (7)仅正交归一基中直接成立, 一般基有Gram矩阵; (8)核非零、像非全空间; (9)内积逐项取系数; (10)$U^\dagger=U=U^{-1}$.
\end{solution}
\originalsection{第2题：幺正尺度变换（15分）}
\begin{exercise}
$U_\lambda=\exp[-\ii\lambda(xp+px)/(2\hbar)]$. (a)幺正要求$\lambda$何条件. (b)求$U_\lambda^\dagger xU_\lambda$. (c)求动量变换, 判断是否与自由$H$对易.
\end{exercise}
\begin{solution}
正规自伴伸缩生成元$D=(xp+px)/2$要求实$\lambda$使指数反Hermitian. $[D,x]=-\ii\hbar x$,$[D,p]=\ii\hbar p$, 令$X(\lambda)=U_\lambda^\dagger xU_\lambda$, 对参数求导得$X\prime=(\ii/\hbar)U^\dagger[D,x]U=X$, 初值$x$给$X=\e^\lambda x$. 同理$P\prime=-P$, 得$U^\dagger pU=\e^{-\lambda}p$. 位置波包拉宽$\e^\lambda$时动量宽度缩小$\e^{-\lambda}$, 总概率保持. 对易$[\e^\lambda x,\e^{-\lambda}p]=\ii\hbar$保持, 但自由$H$变为$\e^{-2\lambda}H$, 一般不相等, 所以除$\lambda=0$外不是自由Hamiltonian的对称性. 幺正等谱的陈述与缩放动量能量并不矛盾, 因自由连续谱$[0,\infty)$在正缩放下集合相同.
\end{solution}
\originalsection{第3题：由长度恢复内积（15分）}
\begin{exercise}
(a)化简$\tfrac14[\norm{u+v}^2-\norm{u-v}^2+\ii\norm{u-\ii v}^2-\ii\norm{u+\ii v}^2]$, 说明保范数的线性$S$保内积. (b)从该式及$\norm{\alpha u}=|\alpha|\norm u$证明共轭对称.
\end{exercise}
\begin{solution}
展开得$\inner uv$（本册内积对第二项线性）. 线性保证$S(u\pm v)=Su\pm Sv$以及带$\ii$组合相同, 四个范数保持就使整个极化式保持. 交换$u,v$, 实部分$\norm{v-u}=\norm{u-v}$不变. 用$\norm{v-\ii u}=\norm{-\ii(u+\ii v)}=\norm{u+\ii v}$, $\norm{v+\ii u}=\norm{\ii(u-\ii v)}$使虚部交换符号, 因而$\inner vu=\inner uv^*$. 这里“内积由范数产生”包含极化式确实具有线性等内积性质的前提; 不是任意Banach范数都给内积.
\end{solution}
\originalsection{第4题：双阱势的变分估计（25分）}
\begin{exercise}
$H=p^2/(2m)+\hbar\omega-m\omega^2x^2/2+m^2\omega^3x^4/(16\hbar)$,$d_0^2=\hbar/(m\omega)$. (a)画势与极值. (b)对试探态$\psi=x\e^{-x^2/(2\lambda^2)}$求能量中四次项系数$A$. 已给$E=\hbar\omega-3m\omega^2\lambda^2/4+A\lambda^4+3\hbar^2/(4m\lambda^2)$. (c)令$\lambda=d_0\eta$, 忽略最后动能项优化再代回给上界. (d)这对哪一另一能级给更强界.
\end{exercise}
\begin{solution}
完成平方$V=m^2\omega^3(x^2-4d_0^2)^2/(16\hbar)$. 极小$x=\pm2d_0,V=0$, 局部极大$x=0,V=\hbar\omega$. 本题原卷积分常数约定$\int u^{2n}\e^{-u^2}\dd u=c_n\sqrt\pi$, $c_1=1/2,c_3=15/8$, 与上一年测试的$\e^{-2y^2}$常数不同. 试探态密度含$x^2$, 所以$\langle x^4\rangle=\lambda^4(c_3/c_1)=15\lambda^4/4$, $A=15m^2\omega^3/(64\hbar)$. 因而
\[
E(\eta)=\hbar\omega\left[1-\frac34\eta^2+\frac{15}{64}\eta^4+\frac3{4\eta^2}\right].
\]
忽略最后项只用于\textbf{选取试探参数}, 极值$\eta^2=8/5$. 将其代回完整能量而非删项能量得$E_*=139\hbar\omega/160=0.86875\hbar\omega$. 任一归一试探态的完整能量是基态上界; 删掉正动能项后的最小值并不是上界. 本试探态奇, 正交于偶基态, 展开只含奇能级, 故更直接$E_1\leq E_*$, 为第一激发态上界, 也当然高于真正基态.
\begin{center}\begin{tikzpicture}\begin{axis}[width=.72\textwidth,height=4cm,xlabel={$x/d_0$},ylabel={$V/(\hbar\omega)$},xmin=-3,xmax=3,ymin=0,ymax=2]
\addplot[main,domain=-3:3,samples=100]{(x*x-4)^2/16};\end{axis}\end{tikzpicture}\end{center}
\end{solution}
\originalsection{第5题：自旋宽度的上下限（25分）}
\begin{exercise}
自旋$1/2$: (a)推$\Delta S_x\Delta S_y$不等式. (b)最小乘积与取到的态. (c)任一分量最大宽度. (d)对一般$\ket{\boldsymbol n;+}$求三宽度和使$x,y$乘积最大的态.
\end{exercise}
\begin{solution}
$[S_x,S_y]=\ii\hbar S_z$给$\Delta S_x\Delta S_y\geq(\hbar/2)|\langle S_z\rangle|$. $S_i^2=\hbar^2I/4$, $\langle S_i\rangle=(\hbar/2)n_i$, 所以$\Delta S_i=(\hbar/2)\sqrt{1-n_i^2}$. 最小乘积0在$x$本征态$\ket{x;\pm}$或$y$本征态$\ket{y;\pm}$取到. 每个分量宽度上限$\hbar/2$, 当$n_i=0$取到. 一般角表示$n_x=\sin\theta\cos\phi,n_y=\sin\theta\sin\phi,n_z=\cos\theta$, 代入即三个公式. 最大乘积$\hbar^2/4$要求$n_x=n_y=0$, 即$\ket{z;\pm}$; 在这些态Robertson界也饱和. 测量宽度由制备决定, 不表示实验仪器的读数误差.
\end{solution}

\originalchapter{8.05 2011期末：氢、合成自旋与压缩态}
原件为Fall 2011 Final Exam, 2011年12月22日9:00--12:00, 七题100分, 闭卷禁计算器; Fall 2013 OCW归档. 以下为\textbf{AI独立解答}, 保留原卷参数歧义并说明其影响.
\originalsection{第1题：十个判断（20分）}
\begin{exercise}
判断: (1)Hermitian算符反对易子Hermitian; (2)各时刻Schrödinger $H$对易则$H_H=H_S$; (3)$a_0\alpha$为经典电子半径; (4)$a_0\alpha^2$为电子Compton波长; (5)因子分解$H^{(1)},H^{(2)}$至少一个必有正规零能态; (6)两个独立角动量之差也满足角动量代数; (7)$\boldsymbol J_1\cdot\boldsymbol J_2$与$J_{1x}+J_{2x}$对易; (8)任意自旋$j$表示各$J_i$迹0; (9)三维球对称振子第一激发组是$\ell=1$多重态; (10)第二激发组全是$\ell=2$多重态.
\end{exercise}
\begin{solution}
依次\textbf{T,T,F,F,F,F,T,T,T,F}. (1)$(AB+BA)^\dagger=AB+BA$. (2)$U$与当时$H$对易, 共轭不变. (3)$a_0\alpha$是约化Compton长度, (4)$a_0\alpha^2$才是经典半径, 非约化Compton另有$2\pi$. (5)$A\psi=0$的形式解可能不可归一, 例如全线常数超势$W=1$给指数向一侧发散, 两Hamiltonian都没有正规零能态. (6)差的自对易右侧给$J_1+J_2$, 不是差. (7)它是共同旋转标量. (8)$J_z$谱$m=-j,\ldots,j$相加0, 其他方向等谱. (9)三维第一组维数3, 三个单量子态是向量表示. (10)第二组六维分成$\ell=2$五维加$\ell=0$一维.
\end{solution}
\originalsection{第2题：Heisenberg升降算符（10分）}
\begin{exercise}
求谐振子的$a_H(t),a_H^\dagger(t),N_H(t)$.
\end{exercise}
\begin{solution}
$[H,a]=-\hbar\omega a$, 因而$\dot a_H=-\ii\omega a_H$, 解$a_H=\e^{-\ii\omega t}a$, $a_H^\dagger=\e^{\ii\omega t}a^\dagger$. 两相位抵消,$N_H=N$. 能量占有数不变, 相位旋转使位置动量均值和某些宽度仍可随时间变化.
\end{solution}
\originalsection{第3题：倒数均值不等式（15分）}
\begin{exercise}
(a)从$\langle u^{-1}(u-\langle u\rangle)^2\rangle\geq0$推出$\langle1/u\rangle$与$1/\langle u\rangle$关系,$u$谱非负. (b)以氢基态$r$验证, 可用维里定理求$\langle1/r\rangle$.
\end{exercise}
\begin{solution}
需$u^{-1}$在所取态有意义, $\langle u\rangle>0$, 有关矩有限. 展开给$\langle u\rangle^2\langle1/u\rangle-\langle u\rangle\geq0$, 即$\langle1/u\rangle\geq1/\langle u\rangle$. 零谱成分可使逆不存在或均值无穷, 不能直接除0. 氢基态径向概率$4r^2\e^{-2r/a_0}/a_0^3$给$\langle r\rangle=3a_0/2$. 维里$2T=-V$与$E=-e^2/(2a_0)$给$\langle V\rangle=2E=-e^2/a_0$, 所以$\langle1/r\rangle=1/a_0>2/(3a_0)$. 位置有分散, 故严格不等号.
\end{solution}
\originalsection{第4题：氢Hamiltonian的正性分解（15分）}
\begin{exercise}
将$H=\boldsymbol p^2/(2m)-e^2/r$写为$\sum_i(p_i+\ii\beta x_i/r)(p_i-\ii\beta x_i/r)/(2m)+\gamma$,$\beta,\gamma$实. (a)求常数. (b)证所有态$\langle H\rangle\geq\gamma$. (c)求取等号的径向波函数.
\end{exercise}
\begin{solution}
令$n_i=x_i/r$, $\sum n_i^2=1$, $\sum\partial_i n_i=2/r$（在原点外）. 展开乘积得到$p^2+\beta^2-2\beta\hbar/r$, 因而$\beta=me^2/\hbar$, $\gamma=-me^4/(2\hbar^2)$. 每项是$Q_i^\dagger Q_i$, $Q_i=p_i-\ii\beta n_i$, 所以均值为$\sum\norm{Q_i\psi}^2/(2m)+\gamma\geq\gamma$, 正规Coulomb域允许这一二次型关系. 等号要求全部$Q_i\psi=0$, 对径向$\psi$即$\hbar\psi'(r)+\beta\psi=0$, 得$\psi=(\pi a_0^3)^{-1/2}\e^{-r/a_0}$. 这既给真实基态又证明能量下界恰可取到, 不是只猜出指数试探态.
\end{solution}
\originalsection{第5题：三个自旋1的反铁磁圆环（15分）}
\begin{exercise}
三可区分自旋1有$H=\Delta(\boldsymbol S_1\cdot\boldsymbol S_2+\boldsymbol S_2\cdot\boldsymbol S_3+\boldsymbol S_3\cdot\boldsymbol S_1)/\hbar^2$,$\Delta>0$. (a)维数并改写$H$. (b)求各能量简并度. (c)在$\ket{m_1,m_2,m_3}$基中求归一基态, 从$m_1+m_2+m_3=0$的七候选态开始.
\end{exercise}
\begin{solution}
维数$3^3=27$, 总$\boldsymbol S=\sum\boldsymbol S_i$给$H=\Delta[S^2-6\hbar^2]/(2\hbar^2)$. 合成$1\otimes1\otimes1=0\oplus3(1)\oplus2(2)\oplus3$, 这里数字前系数为副本数, 最后3是$j=3$不是三个副本. 能量和总简并为
\begin{center}\begin{tabular}{cccc}\toprule $j$&副本数&$E/\Delta$&简并度\\\midrule0&1&$-3$&1\\1&3&$-2$&9\\2&2&0&10\\3&1&3&7\\\bottomrule\end{tabular}\end{center}
总和27. 唯一$j=0$态需$S_z\psi=S_+\psi=0$. 令000的系数为$g$, 六排列$+0-,+-0,0+-,0-+,-+0,-0+$的系数依次$a,b,c,d,e,f$. 除去共同升阶因子$\sqrt2\hbar$, $S_+\psi=0$在$++-,+-+,-++$上的约束为
\[
a+c=0,\qquad b+d=0,\qquad e+f=0.
\]
在$+00,0+0,00+$上的约束为$a+b+g=0,c+e+g=0,d+f+g=0$. 后三式相加并用前三式得$3g=0$; 再逐项解得$(a,b,c,d,e,f)=a(1,-1,-1,1,1,-1)$. 归一结果为
\[
\ket{\psi_g}=\frac1{\sqrt6}\bigl(\ket{+,0,-}+\ket{0,-,+}+\ket{-,+,0}
-\ket{+,-,0}-\ket{-,0,+}-\ket{0,+,-}\bigr).
\]
可逐项作用$S_+$, 每个目标排列出现一正一负, 完全抵消. 唯一副本数保证除整体相位无其他基态.
\end{solution}
\originalsection{第6题：因子分解的阶梯结构（10分）}
\begin{exercise}
假设$H_\ell=H_\ell^{(1)}+\alpha_\ell$, $H_\ell^{(1)}=A_\ell^\dagger A_\ell$, $H_\ell^{(2)}=A_\ell A_\ell^\dagger=H_{\ell+1}^{(1)}+\beta_\ell$. 每个$H_\ell^{(1)}$有已知正规零态$\phi_{0,\ell}^{(1)}$. (a)写$H_\ell$前三态及能量. (b)推广任意$n$.
\end{exercise}
\begin{solution}
恒等$H_\ell^{(1)}A_\ell^\dagger=A_\ell^\dagger H_\ell^{(2)}$把伙伴能量搬回. 因此$\phi_{0,\ell}\propto\phi_{0,\ell}^{(1)}$, $E_{0,\ell}=\alpha_\ell$; $\phi_{1,\ell}\propto A_\ell^\dagger\phi_{0,\ell+1}^{(1)}$, $E_{1,\ell}=\alpha_\ell+\beta_\ell$; $\phi_{2,\ell}\propto A_\ell^\dagger A_{\ell+1}^\dagger\phi_{0,\ell+2}^{(1)}$, $E_{2,\ell}=\alpha_\ell+\beta_\ell+\beta_{\ell+1}$. 一般
\[
\phi_{n,\ell}\propto A_\ell^\dagger A_{\ell+1}^\dagger\cdots A_{\ell+n-1}^\dagger\phi_{0,\ell+n}^{(1)},\quad
E_{n,\ell}=\alpha_\ell+\sum_{q=0}^{n-1}\beta_{\ell+q}.
\]
归一因子由每步范数平方等于所提升的伙伴能量求得. 构造需输出非零、正规且满足同一边界域; 题给链结构可产生这些态, 若要声称某任意超势的\textbf{全部}谱还须排除链外额外态.
\end{solution}
\originalsection{第7题：压缩真空的时间宽度（15分）}
\begin{exercise}
题面$S(\gamma)=\exp[-\gamma(a^{\dagger2}-a^2)/2]$, $\ket{0_\gamma}=S(\gamma)\ket0$. (a)求$S^\dagger xS,S^\dagger pS$. (b)求演化后$(\Delta x(t))^2=\hbar G/(2m\omega)$. 原题此处印$\gamma\in\mathbb C$, 但同卷公式表第5页限定$\gamma\in\mathbb R$.
\end{exercise}
\begin{solution}
必须先处理条件: 题给生成元$a^{\dagger2}-a^2$为反Hermitian, 只有\textbf{实}$\gamma$才使$S$幺正, 保证题中“压缩真空”的单位归一. 按公式表一致的实参数解, 用$S^\dagger aS=a\cosh\gamma-a^\dagger\sinh\gamma$, 得$S^\dagger xS=\e^{-\gamma}x$, $S^\dagger pS=\e^\gamma p$. Heisenberg$x_H=x\cos\omega t+p\sin\omega t/(m\omega)$, 真空交叉对称矩0, 所以
\[
G(\gamma,\omega t)=\e^{-2\gamma}\cos^2\omega t+\e^{2\gamma}\sin^2\omega t.
\]
低噪声方向随振子旋转, 能量均值仍不变.

若坚持题面任意复$\gamma$, $S^\dagger S\neq I$, 前述共轭缩放公式和归一化不能成立, 有时作用真空甚至不正规. 可在$|w|<1$, $w=-\tanh\gamma$时先将形式态正规化, 其位置宽度为$G=[1+|w|^2+2\operatorname{Re}(w\e^{-2\ii\omega t})]/(1-|w|^2)$. 这不是用原题非幺正$S^\dagger xS$直接得出的单位态公式. 通常复压缩参数应改生成元为$(\zeta^*a^2-\zeta a^{\dagger2})/2$; 本册把它列为条件修正说明, 不伪称原卷已印此式.
\end{solution}

\originalchapter{8.05 2012期末：对称性与粒子衰变}
原件为Fall 2012 Final Exam, 2012年12月18日13:30--16:30, 六题100分, 闭卷禁计算器. 原卷单独引用公式表, 题卷本身没有附入; 已取得课程公式表原件但不冒称逐字同一2012版. 以下为\textbf{AI独立解答}.
\originalsection{第1题：十个判断（20分）}
\begin{exercise}
判断: (1)三维中心势完整波函数必在原点为0; (2)球形无限阱有不能用角动量解释的能级简并; (3)$[H,L_z]=[H,L_x]=0$保证全旋转不变; (4)自旋半整数时轨道角动量也必须半整数才能使总角动量整数; (5)$\boldsymbol r\times\boldsymbol p=-\boldsymbol p\times\boldsymbol r$; (6)$\boldsymbol r\times\boldsymbol L=-\boldsymbol L\times\boldsymbol r$; (7)两个对易Schrödinger算符因显含时间而可使对应Heisenberg算符不对易; (8)$H_H$由$H_S$中每算符换成其Heisenberg版本得到; (9)谐振子$N_H=N_S$; (10)已知$U(t,t_0)$仍不足以求$H$.
\end{exercise}
\begin{solution}
依次\textbf{F,F,T,F,T,F,F,T,T,F}, 采用正规中心势/球阱与正常旋转生成元定义域. (1)$s$态例如氢基态原点非零, 消失的是约化径向$u(0)$. (2)球阱能量由球Bessel函数零点给, 简并$2\ell+1$由旋转解释, 不像Coulomb/振子另有系统隐藏简并. (3)Jacobi由$[L_z,L_x]=\ii\hbar L_y$推出$[H,L_y]=0$, 故与所有旋转生成元对易. (4)普通单值轨道只能整数$\ell$, 半整数自旋加整数仍半整数, 不能达到所述整数目标. (5)交叉分量交换$p_j,r_k$时$\delta_{jk}$被反对称张量消掉. (6)$\boldsymbol L\times\boldsymbol r=-\boldsymbol r\times\boldsymbol L+2\ii\hbar\boldsymbol r$. (7)相同时刻$[A_H,B_H]=U^\dagger[A_S(t),B_S(t)]U$总保持. (8)幺正共轭保持算符积及标量时间系数. (9)$[H,N]=0$. (10)$H(t)=\ii\hbar(\partial_tU)U^\dagger$, 包括整体时间相位携带的能量常数.
\end{solution}
\originalsection{第2题：指定进动与方向无关单态（10分）}
\begin{exercise}
(a)自旋$1/2$初始$+z$, 磁矩$\gamma\boldsymbol S$, 给常场使方向依次$z\to x\to y\to z$, 周期$T$. (b)将$[\ket{\boldsymbol n;+}\ket{\boldsymbol n;-}-\ket{\boldsymbol n;-}\ket{\boldsymbol n;+}]/\sqrt2$展开到$z$基, 证明任意$\boldsymbol n$都是总自旋单态.
\end{exercise}
\begin{solution}
取旋转轴$\boldsymbol n_0=(1,1,1)/\sqrt3$, 有向进动角速度$-\gamma\boldsymbol B=(2\pi/T)\boldsymbol n_0$, 即$\boldsymbol B=-2\pi\boldsymbol n_0/(\gamma T)$,$\gamma\neq0$. Rodrigues公式旋转$2\pi/3$把$z$变$x$, 再同角变$y$, 一周回$z$, 三方向都与轴等夹角. 自旋矢旋转一周虽自旋子得负号, 方向周期仍$T$.

令$c=\cos\theta/2,s=\sin\theta/2$, 方向态$\ket{\boldsymbol n;+}=c\ket++s\e^{\ii\phi}\ket-$,$\ket{\boldsymbol n;-}=-s\e^{-\ii\phi}\ket++c\ket-$. 交叉展开$++$和$--$项抵消,$+-$系数$c^2+s^2=1$,$-+$系数$-1$, 所以正好为$(\ket{+-}-\ket{-+})/\sqrt2$, 被总$S_z,S_\pm$湮灭.
\end{solution}
\originalsection{第3题：Runge--Lenz矢的垂直性（10分）}
\begin{exercise}
(a)检验$\boldsymbol r\cdot\boldsymbol L=0$, $\boldsymbol p\cdot\boldsymbol L=0$. (b)$\boldsymbol R=(\boldsymbol p\times\boldsymbol L-\ii\hbar\boldsymbol p)/(me^2)-\boldsymbol r/r$, 证明$\boldsymbol R\cdot\boldsymbol L=0$.
\end{exercise}
\begin{solution}
$r_iL_i=\epsilon_{ijk}r_ir_jp_k=0$, 因$r_ir_j$对称. $p_iL_i=\epsilon_{ijk}p_ir_jp_k$把$p_i$移过$r_j$得对称$p_ip_k$项与$\delta_{ij}$项, 都被$\epsilon$消掉. 对交叉积项,
$(\boldsymbol p\times\boldsymbol L)\cdot\boldsymbol L=\epsilon_{ijk}p_jL_kL_i=(1/2)\epsilon_{ijk}p_j[L_k,L_i]=\ii\hbar\boldsymbol p\cdot\boldsymbol L$. 因而与$-\ii\hbar\boldsymbol p$修正相消, 剩$-\boldsymbol r\cdot\boldsymbol L/r=0$. 量子修正不能先将非对易的$L_kL_i$当对称积略掉; 该修正保证$R$的正确Hermitian量子定义.
\end{solution}
\originalsection{第4题：压缩改变Hamiltonian参数（15分）}
\begin{exercise}
实$\gamma$的$S=\exp[-\gamma(a^{\dagger2}-a^2)/2]$幺正,$H(m,k)=p^2/(2m)+kx^2/2$. (a)求$S^\dagger HS=H(m',k')$. (b)比较谱和映射本征态. (c)压缩是不是原振子的对称性.
\end{exercise}
\begin{solution}
$S^\dagger xS=\e^{-\gamma}x$,$S^\dagger pS=\e^\gamma p$, 所以$m'=m\e^{-2\gamma}$,$k'=k\e^{-2\gamma}$. $\omega'=\sqrt{k'/m'}=\omega$, 两Hamiltonian幺正等谱; 原$H\ket E=E\ket E$对应新本征态$S^\dagger\ket E$. 对$\gamma\neq0$,$S^\dagger HS\neq H$, 因而不是固定$m,k$系统的守恒对称变换. 只是将它映到不同参数、相同谱的模型. 本征态的长度宽度同时改变, 能级相同并不代表算符相同.
\end{solution}
\originalsection{第5题：三维振子的隐藏算符（25分）}
\begin{exercise}
$T_{ij}=a_i^\dagger a_j$, $H=\hbar\omega(N+3/2)$. (a)写$H,N_i$. (b)求$[T_{ij},T_{pq}]$. (c)检验是否与$H$对易. (d)用一个$T$把$(a_3^\dagger)^3\ket0$变成$a_1^\dagger(a_3^\dagger)^2\ket0$. (e)在总数$N=n$组, 求最高及其他$\ell$, 用$a_L,a_R,a_3$构造最高多重态前三个$m$及次高多重态最高$m$. (f)找一个产生乘湮灭的$K$, 与$H$对易并把后者变为最高$\ell$最高$m$.
\end{exercise}
\begin{solution}
$N_i=T_{ii}$,$H=\hbar\omega(\sum_iT_{ii}+3/2)$. 用$[a_j,a_p^\dagger]=\delta_{jp}$展开并消四算符项得$[T_{ij},T_{pq}]=\delta_{jp}T_{iq}-\delta_{qi}T_{pj}$. 特别$[N,T_{pq}]=0$, 所以这些算符保持能量, 可在简并空间移动态. $T_{13}(a_3^\dagger)^3\ket0=3a_1^\dagger(a_3^\dagger)^2\ket0$, 所需为$T_{13}/3$.

本题圆基定义$a_L=(a_1+\ii a_2)/\sqrt2,a_R=(a_1-\ii a_2)/\sqrt2$, $L_+=\sqrt2\hbar(a_3^\dagger a_L-a_R^\dagger a_3)$, $L_-=L_+^\dagger$. 因而$L_z=\hbar(N_R-N_L)$, 最大$m=n$, 态$(a_R^\dagger)^n\ket0$且被$L_+$湮灭, 因而$\ell_*=n$. $N=n$的其余表示$\ell=n-2,n-4,\ldots$直到0或1. 对$\ell=n$, 未归一前三态可取
\[
\begin{aligned}
\ket{n,n}&\propto(a_R^\dagger)^n\ket0,\\
\ket{n,n-1}&\propto a_3^\dagger(a_R^\dagger)^{n-1}\ket0,\\
\ket{n,n-2}&\propto[(n-1)(a_3^\dagger)^2-a_L^\dagger a_R^\dagger](a_R^\dagger)^{n-2}\ket0.
\end{aligned}
\]
第三通式限$n\geq2$. 对$n=1$前三态另列$a_R^\dagger\ket0,a_3^\dagger\ket0,a_L^\dagger\ket0$, 对应$m=1,0,-1$, 不代入负次幂. 后两式由$L_-\!=\sqrt2\hbar(a_L^\dagger a_3-a_3^\dagger a_R)$连续作用, 整体符号可变. $n\geq2$次高最高态为$Q^\dagger(a_R^\dagger)^{n-2}\ket0$,$Q^\dagger=(a_3^\dagger)^2+2a_L^\dagger a_R^\dagger$, 它被$L_+$湮灭且$m=n-2$. 若$n=1$, 没有次高组, (f)亦无该映射对象; 前三态仍如上存在. 取$K=a_R^\dagger a_L$, 保持总数, 并有$KQ^\dagger(a_R^\dagger)^{n-2}\ket0=2(a_R^\dagger)^n\ket0$, 到归一因子即完成(f). $K$不是空间旋转角动量, 解释了不同$\ell$之间的额外简并.
\end{solution}
\originalsection{第6题：二体衰变的角动量限制（20分）}
\begin{exercise}
静止$X\to A+B$满足$\boldsymbol J=\boldsymbol S_A+\boldsymbol S_B+\boldsymbol L=\boldsymbol S_X$. (a)$s_X=0,s_A=s_B=1/2$允许哪些$\ell$. (b)$s_X=3/2,s_A=1/2,s_B=1$同问. (c)为何$n\to p+e^-$两体衰变不行（均自旋$1/2$）, 三体加反中微子却可. (d)未知$X$极化为$m_X=s_X$, 衰变$s_A=1/2,s_B=0$, 单一$\ell$, 测$P(m_A=1/2)=1/5$, $P(m_A=-1/2)=4/5$, 求$s_X,\ell$.
\end{exercise}
\begin{solution}
(a)两末态自旋先合为$S=0$或1, 要总$J=0$必须$\ell=S$, 所以$\ell=0,1$. 这只列角动量许可, 其他宇称或全同粒子限制仍可能排除某通道. (b)$S=1/2$给$\ell=1,2$,$S=3/2$给$\ell=0,1,2,3$, 合并$0,1,2,3$. (c)两个半整数自旋合成整数, 再加整数轨道仍整数, 不可能匹配初始半整数; 加另一半整数反中微子则总可为半整数, 与实际三体衰变相容.

(d)$J=\ell\pm1/2$. 上支最高态$m=\ell+1/2$只能$m_L=\ell,m_A=+1/2$, 与给定非平凡概率不符. 下支最高$m=\ell-1/2$写$a\ket{\ell,\ell-1}\ket++b\ket{\ell,\ell}\ket-$. $J_+$作用给$(a\sqrt{2\ell}+b)\hbar\ket{\ell,\ell}\ket+$, 必须0; 归一给$|a|^2=1/(2\ell+1)$,$|b|^2=2\ell/(2\ell+1)$. 因此$\ell=2,s_X=3/2$. 测概率足以定这两个数依赖题给单一$\ell$及极化最高态条件, 多通道混合时不能这样唯一反推.
\end{solution}

\originalchapter{8.05 2013期末：谱生成与纠缠}
原件为Fall 2013 Final Exam, 2013年12月17日9:00--12:00, 六题100分, 闭卷禁计算器. 以下全小问为\textbf{AI独立解答}. 能级图按题意大幅绘制, 图上微小劈裂作示意放大, 数值偏移由表格精确给出.
\originalsection{第1题：十个判断（20分）}
\begin{exercise}
判断: (1)$\sigma_1\otimes\sigma_1$与$\sigma_3\otimes\sigma_1$对易; (2)$\sigma_1\otimes\sigma_1$与$\sigma_3\otimes\sigma_3$对易; (3)$(T\otimes S)^\dagger=S^\dagger\otimes T^\dagger$; (4)有限维单射自映射可逆; (5)角动量守恒禁止自旋$1/2$衰变为两个自旋$1/2$; (6)$[L^2,x_i]=0$; (7)$\boldsymbol r\cdot\boldsymbol p=\boldsymbol p\cdot\boldsymbol r-3\ii\hbar$; (8)旋转矢量$\boldsymbol A$满足$\boldsymbol A\cdot\boldsymbol L=\boldsymbol L\cdot\boldsymbol A$; (9)氢Runge--Lenz矢$\boldsymbol R$本身满足角动量代数; (10)经典和量子都有$\boldsymbol R\cdot\boldsymbol L=0$.
\end{exercise}
\begin{solution}
依次\textbf{F,T,F,T,T,F,F,T,F,T}. (1)只有第一因子反对易, 整体不对易; (2)两个因子各反对易, 双负抵消; (3)伴随保持因子顺序,$T^\dagger\otimes S^\dagger$; (4)秩为维数; (5)末态合成整数不能匹配半整数; (6)位置是矢量算符, 可改变$\ell$, 例如$x_iY_{00}$产生$\ell=1$角依赖; (7)正确为$\boldsymbol r\cdot\boldsymbol p=\boldsymbol p\cdot\boldsymbol r+3\ii\hbar$; (8)$\sum_i[L_i,A_i]=\ii\hbar\sum_{ik}\epsilon_{iik}A_k=0$; (9)$[R_i,R_j]$与$HL_k$成比例, 不是$\ii\hbar\epsilon_{ijk}R_k$; (10)见上一期末第3题量子算符证明.
\end{solution}
\originalsection{第2题：位移压缩态的占有数（10分）}
\begin{exercise}
$\ket{\alpha,\gamma}=D(\alpha)S(\gamma)\ket0$,$\alpha$复、$\gamma$实, 求$\langle N\rangle$并判断自由谐振子演化中是否改变.
\end{exercise}
\begin{solution}
$D^\dagger aD=a+\alpha$, 压缩真空$\langle a\rangle=0$, $\langle a^\dagger a\rangle=\sinh^2\gamma$. 逐项展开$D^\dagger ND=(a^\dagger+\alpha^*)(a+\alpha)$, 得$\langle N\rangle=|\alpha|^2+\sinh^2\gamma$. $[H,N]=0$所以恒定. 分别代表位移的经典振动能量与压缩增加的真空涨落能量, 不随椭圆噪声方向旋转而变.
\end{solution}
\originalsection{第3题：球方阱束缚门槛（10分）}
\begin{exercise}
$V=-V_0$（$r<a$）、0（$r>a$）,$V_0>0$. 推导至少有一个三维束缚态所需深度不等式.
\end{exercise}
\begin{solution}
最容易束缚的是$\ell=0$, 因所有$\ell>0$多了非负离心势. 正规径向$u(0)=0$, 内$A\sin kr$, 外$B\e^{-\kappa r}$, 匹配$k\cot ka=-\kappa$, $k^2+\kappa^2=2mV_0/\hbar^2$. 第一束缚根要$z_0=a\sqrt{2mV_0}/\hbar>\pi/2$, 即$V_0>\pi^2\hbar^2/(8ma^2)$. 等号零能外部$u$为常数, $\psi=u/r$虽点值衰减却$\int|u|^2\dd r$发散, 不算正规束缚态. 与一维全线任意浅吸引方阱有束缚态不同, 三维有深度门槛.
\end{solution}
\originalsection{第4题：一个谱生成角动量（15分）}
\begin{exercise}
本题$\hbar=1$, $j$正整数,$J_z=-j+N$,$J_+=a^\dagger\sqrt{2j-N}$,$J_-=\sqrt{2j-N}a$. (a)在平方根合法态上验证代数. (b)求$J^2$及对易. (c)构造$j$多重态并检验作用. (d)与振子$H=N+1/2$是否对易.
\end{exercise}
\begin{solution}
在$0\leq n\leq2j$有限不变空间上,$J_+\ket n=\sqrt{(n+1)(2j-n)}\ket{n+1}$,$J_-\ket n=\sqrt{n(2j-n+1)}\ket{n-1}$. 端点向外系数0, 所以不需对负平方根运算. $J_z$作用$n-j$, 立即$[J_z,J_\pm]=\pm J_\pm$. $J_+J_-$系数$n(2j-n+1)$,$J_-J_+$系数$(n+1)(2j-n)$, 差$2(n-j)$, 因而$[J_+,J_-]=2J_z$. 用$J^2=J_z^2+(J_+J_-+J_-J_+)/2$得$j(j+1)I$, 与三分量对易. 态$\ket{j,m}=\ket{n=j+m}$,$m=-j,\ldots,j$, 给预期阶梯系数$\sqrt{j(j+1)-m(m\pm1)}$. $[H,J_z]=0$, $[H,J_\pm]=\pm J_\pm$, 后二者改变振子能量, 因而是\textbf{谱生成}而非同能量旋转对称. 对全无限振子空间没有实平方根的该定义, 必须保留有限子空间限制.
\end{solution}
\originalsection{第5题：四维振子的简并劈裂（20分）}
\begin{exercise}
四维等频$H=\hbar\omega(N+2)$. (a)列$N=0,1,2$态并计数. (b)把前两个坐标、后两个坐标各组成二维振子, 数$n_a,n_b$对应Schwinger角动量$j_1=n_a/2,j_2=n_b/2$, 列$N=0,1,2,3$的$j_1\otimes j_2$及维数. (c)加$2E_0\boldsymbol J^{(1)}\cdot\boldsymbol J^{(2)}/\hbar^2$,$0<E_0\ll\hbar\omega$, 求四组劈裂、表示与态数, 画大幅能级图.
\end{exercise}
\begin{solution}
$N=0$真空1态; $N=1$为四个$a_i^\dagger\ket0$; $N=2$为四个$(a_i^\dagger)^2\ket0$及六个$a_i^\dagger a_j^\dagger\ket0$（$i<j$）, 共10. 一般简并$\binom{N+3}3$. 各二维振子的两占有数在固定总数$n$形成$j=n/2$, 维数$n+1$. 表如下:
\begin{center}\begin{tabular}{ccc}\toprule $N$&$j_1\otimes j_2$&各表示维数\\\midrule
0&$0\otimes0$&1\\
1&$\tfrac12\otimes0,\ 0\otimes\tfrac12$&2,2\\
2&$1\otimes0,\ \tfrac12\otimes\tfrac12,\ 0\otimes1$&3,4,3\\
3&$\tfrac32\otimes0,\ 1\otimes\tfrac12,\ \tfrac12\otimes1,\ 0\otimes\tfrac32$&4,6,6,4\\\bottomrule\end{tabular}\end{center}
耦合$\boldsymbol J=\boldsymbol J^{(1)}+\boldsymbol J^{(2)}$, 偏移精确为$\delta E=E_0[j(j+1)-j_1(j_1+1)-j_2(j_2+1)]$. 因为扰动与原$H,J_1^2,J_2^2,J^2$对易, 这个线性偏移在所给模型为精确结果, 不仅是第一阶.

$N=0$一态无偏移. $N=1$两套$j=1/2$, 共4态无偏移. $N=2$: 两套$j=1$来自$1\otimes0,0\otimes1$, 共6态无偏移; $\tfrac12\otimes\tfrac12$分成$j=1$三态偏$+E_0/2$, $j=0$一态偏$-3E_0/2$. $N=3$: 两套$j=3/2$来自端点表示, 共8态无偏移; 两个中间表示各含$j=3/2$四态偏$+E_0$与$j=1/2$两态偏$-2E_0$, 总8、4态. 
\clearpage
以下图中括号为态数, 劈裂相对主间隔放大, 所有表示来源已在文字和表中逐一列出.
\begin{center}\begin{tikzpicture}[x=1cm,y=1.3cm,>=Stealth]
\draw[->](0,-.5)--(0,10.6)node[above]{$E$};
\foreach \y/\lab/\num in {0/{2\hbar\omega}/1,3/{3\hbar\omega}/4,6/{4\hbar\omega}/10,9/{5\hbar\omega}/20}{
\draw(.5,\y)--(1.6,\y);\node[left]at(-.15,\y){$\lab$};\node[above]at(1,\y){$N=\pgfmathparse{int(\y/3)}\pgfmathresult\ (\num)$};}
\draw[main](3,0)--(5,0)node[right]{$j=0$ (1)};\draw[->](1.7,0)--(2.9,0);
\draw[main](3,3)--(5,3)node[right]{$j=1/2$ (4)};\draw[->](1.7,3)--(2.9,3);
\draw[main](3,6)--(5,6)node[right]{$\delta E=0,\ j=1$ (6)};
\draw[main](3,6.35)--(5,6.35)node[right]{$+E_0/2,\ j=1$ (3)};
\draw[main](3,4.95)--(5,4.95)node[right]{$-3E_0/2,\ j=0$ (1)};
\foreach \y in {4.95,6,6.35}{\draw[->](1.7,6)--(2.9,\y);}
\draw[main](3,9)--(5,9)node[right]{$\delta E=0,\ j=3/2$ (8)};
\draw[main](3,9.7)--(5,9.7)node[right]{$+E_0,\ j=3/2$ (8)};
\draw[main](3,7.6)--(5,7.6)node[right]{$-2E_0,\ j=1/2$ (4)};
\foreach \y in {7.6,9,9.7}{\draw[->](1.7,9)--(2.9,\y);}
\node at(4,-.6){原组$\longrightarrow$耦合后（小劈裂放大示意）};
\end{tikzpicture}\end{center}
\end{solution}
\clearpage
\originalsection{第6题：纠缠的产生与局部基变换（25分）}
\begin{exercise}
$H=\hbar\omega\sigma_3\otimes\sigma_3$, $Z$基$\ket{++},\ket{+-},\ket{-+},\ket{--}$,$\ket{x;\pm}=(\ket+\pm\ket-)/\sqrt2$. (a)写$H,U$矩阵. (b)演化初态$\ket{x;+}\ket{x;+}$, 找无纠缠时刻. (c)把$\exp[\ii(K\otimes I+I\otimes K)]$分为张量积. (d)将任意$\alpha\ket{x;+}\ket{x;+}+\beta\ket{x;-}\ket{x;-}$变为$\alpha\ket{++}+\beta\ket{--}$, 用$H'=H_0\otimes I+I\otimes H_0$和$K=-H_0t'/\hbar$构造.
\end{exercise}
\begin{solution}
$H=\hbar\omega\operatorname{diag}(1,-1,-1,1)$,$U=\operatorname{diag}(\e^{-\ii\omega t},\e^{\ii\omega t},\e^{\ii\omega t},\e^{-\ii\omega t})$. 初态演化为这些四个相位各乘$1/2$. 二比特纯态系数矩阵行列式为$(\e^{-2\ii\omega t}-\e^{2\ii\omega t})/4=-(\ii/2)\sin2\omega t$, 因此无纠缠当$\omega t=k\pi/2$,$k\in\mathbb Z$; $\omega=0$则一直无纠缠. 行列式非零意味着不能写两列向量的外积, 即非乘积态.

$K\otimes I$与$I\otimes K$对易, 故指数为$\e^{\ii K}\otimes\e^{\ii K}$. 令Hadamard矩阵$F=(\sigma_x+\sigma_z)/\sqrt2$, 满足$F\ket{x;+}=\ket+$,$F\ket{x;-}=\ket-$且$F^2=I$. 取Hermitian$K=\pi(I-F)/2$, 在$F$的$\pm1$谱上$K$分别0、$\pi$, 所以$\e^{\ii K}=F$精确成立. 任取$t'>0$,$H_0=-\hbar K/t'$即完成指定转变, 包括相对相位而非仅到模方相等. 局部幺正只改变表示, 不改变两态系数$\alpha,\beta$的Schmidt值或纠缠量; 本题不是用局部操作消除纠缠.
\end{solution}

\originalchapter{8.06 作业5：隧穿电离与光的吸收}
来源为8.06 Spring 2018 Assignment 5, 2018年3月9日发布, 五题共80分. 原件没有印独立截止日期, 不补造. 以下转述全部小问并给\textbf{AI独立解答}; 电场方向按题给$-e\mathcal E x$约定, $e>0$为强度参数.
\originalsection{第1题：弱电场使束缚态成为亚稳态（15分）}
\begin{exercise}
电子（质量$m=m_e$）处于宽$d$深$V_0$的方阱加$-e\mathcal E x$, 假定$e\mathcal E d\ll\hbar^2/(2md^2)\ll V_0$. (a)用无限深阱估计基态距阱底能量, 解释实际值较低. (b)画倾斜势并解释失稳. (c)用WKB求穿透因子. (d)以经典撞壁频率求寿命. (e)$V_0=20$ eV,$d=2\times10^{-8}$ cm,$\mathcal E=7\times10^4$ V/cm, 估寿命并与宇宙年龄尺度比较. (f)说明$\exp(1/\mathcal E)$为何不能由有限阶电场微扰得到.
\end{exercise}
\begin{solution}
无限阱给$\epsilon=\pi^2\hbar^2/(2md^2)$, 总能量$E_g\simeq-V_0+\epsilon=-I$. 有限墙允许波函数外伸, 约束较少, 故实际基态更低于无限阱的这个绝对能量（距底能量更小）. 加场后右外部势$-e\mathcal E x$最终低于$-I$, 存在通向无穷的允许区, 原束缚态可以隧穿, 不再是真正离散稳定基态.

障碍从$d/2$到$x_t=I/(e\mathcal E)$, 其中$V-E_g=I-e\mathcal E x$. 题给小场条件允许忽略$e\mathcal E d/2$相对$I$, 得
\[
\mathcal T\simeq\exp\!\left[-\frac2\hbar\int_{d/2}^{I/(e\mathcal E)}\sqrt{2m(I-e\mathcal E x)}\dd x\right]
\simeq\exp\!\left[-\frac{4\sqrt{2m}I^{3/2}}{3\hbar e\mathcal E}\right].
\]
阱内速度估为$v=\sqrt{2\epsilon/m}$, 每往返$2d$撞一次右壁, 尝试频率$\nu=v/(2d)$, 寿命$\tau\sim1/(\nu\mathcal T)$. 题给数值$d=0.2$ nm,$\mathcal E=7\times10^6$ V/m, 得$\epsilon=9.40075$ eV,$I=10.59925$ eV,$\nu=4.54618\times10^{15}$ s$^{-1}$, 指数$B=33673.77$, 故$\log_{10}(\tau/\mathrm s)\simeq14608.67$. 以通常宇宙年龄约$10^{18}$ s的量级作比较, 此模型寿命大得不可观测. 该巨大数字不是精密预测: 无限阱近似能量对给定数值仅有有限准确度, 指数对$I$极敏感, 但“远长于宇宙年龄”量级结论稳健.

衰变率含$\exp(-C/\mathcal E)$,$C>0$. 在$\mathcal E\to0^+$所有Taylor系数都为0, 任何有限阶多项式微扰都不会给出非零衰变; 寿命相反按$\exp(C/\mathcal E)$增长. 微扰可描述实能量Stark位移, 不足以捕捉这个非解析的隧穿过程.
\begin{center}\begin{tikzpicture}[x=1cm,y=.4cm,>=Stealth]
\draw[->](-2.2,0)--(5,0)node[right]{$x$};\draw[->](0,-4)--(0,1)node[above]{$V$};
\draw[main](-2, .4)--(-.7,.14)--(-.7,-3)--(.7,-3.28)--(.7,-.14)--(4.7,-.94);
\draw[dashed](-2,-.8)--(4.7,-.8)node[right]{$E_g$};
\node at(3,-1.8){右侧三角障碍（示意）};\node[below]at(-.7,-3){$-d/2$};\node[below]at(.7,-3.3){$d/2$};
\end{tikzpicture}\end{center}
\end{solution}
\originalsection{第2题：一个含时两态系统（15分）}
\begin{exercise}
$H=\left(\begin{smallmatrix}E&v(t)\\v(t)&-E\end{smallmatrix}\right)$, $v$实且$\int|v|\dd t<\infty$, 初始$\ket2$在$-\infty$. (a)最低阶求$+\infty$时$\ket1$概率. (b)$E=0$时精确求概率, 比较微扰并给有效条件.
\end{exercise}
\begin{solution}
相互作用表象第一阶$c_1=-(\ii/\hbar)\int_{-\infty}^\infty v(t)\e^{2\ii Et/\hbar}\dd t$, 所以$P_{2\to1}=\hbar^{-2}|\int v\e^{2\ii Et/\hbar}\dd t|^2+O(v^4)$. 零能时$H=v(t)\sigma_x$各时刻对易, $U=\exp(-\ii A\sigma_x)$,$A=\int v\dd t/\hbar$, 因而精确$P=\sin^2A$. 微扰给$A^2$, 充分有效条件$|A|\ll1$; 更稳健的全程小扰动条件为$\int|v|\dd t/\hbar\ll1$. 若不同时间正负脉冲抵消, 积分$A=0$仍给这个特殊可对易模型的精确零转移, 不能据此概括一般不对易脉冲.
\end{solution}
\originalsection{第3题：原子与单光场模（15分）}
\begin{exercise}
$H_0=\hbar\omega_a\sigma_z/2+\hbar\omega_p(a^\dagger a+1/2)$, 张量恒等算符省略, $\delta H=\alpha(\sigma_-+\sigma_+)(a+a^\dagger)$,$\sigma_+=\ket e\bra g$. (a)求相互作用表象$\widetilde{\delta H}(t)$. (b)在$\omega_a=\omega_p=\omega$积分$0$到$t$, 解释$t\gg1/\omega$可舍项及吸收/发射意义.
\end{exercise}
\begin{solution}
由$[H_0,\sigma_+]=\hbar\omega_a\sigma_+$,$[H_0,a]=-\hbar\omega_pa$, 得
\[
\widetilde{\delta H}=\alpha[\sigma_+a\e^{\ii(\omega_a-\omega_p)t}+\sigma_-a^\dagger\e^{-\ii(\omega_a-\omega_p)t}
+\sigma_+a^\dagger\e^{\ii(\omega_a+\omega_p)t}+\sigma_-a\e^{-\ii(\omega_a+\omega_p)t}].
\]
共振时积分为$\alpha t(\sigma_+a+\sigma_-a^\dagger)+\alpha\sigma_+a^\dagger(\e^{2\ii\omega t}-1)/(2\ii\omega)+\alpha\sigma_-a(1-\e^{-2\ii\omega t})/(2\ii\omega)$. 后两项有界$O(\alpha/\omega)$, 相比随$t$增长的共振项较小. 旋转波近似还要求耦合频率$|\alpha|\sqrt{n+1}/\hbar\ll\omega$, 不能仅以$t$长证明任意强耦合下都可舍去. $\sigma_+a$消一个光子激发原子（吸收）; $\sigma_-a^\dagger$使原子降能并添光子, 振幅$\sqrt{n+1}$既含已有光子刺激也含真空$n=0$的自发发射. 单模可逆能量交换本身不是不可逆的连续谱自发衰变.
\end{solution}
\originalsection{第4题：Gaussian脉冲（10分）}
\begin{exercise}
$H_0\ket n=E_n\ket n$, $\delta H(t)=\e^{-t^2/(2\tau^2)}V/(\sqrt{2\pi}\tau)$,$V$时间无关Hermitian, $\tau>0$. 初态$\ket0$在$-\infty$, 求第一阶转移概率并讨论$\tau\to0,\infty$.
\end{exercise}
\begin{solution}
对$n\neq0$,$c_n=-(\ii/\hbar)V_{n0}\e^{-\omega_{n0}^2\tau^2/2}$,$\omega_{n0}=(E_n-E_0)/\hbar$, 因而$P_n=|V_{n0}|^2\e^{-\omega_{n0}^2\tau^2}/\hbar^2$到最低非零阶. $V$在此量纲为作用量, 因脉冲系数为每时间; 不能按普通能量矩阵直接理解. 极短脉冲频谱宽, 结果趋$|V_{n0}|^2/\hbar^2$; 极慢脉冲对非零能隙指数抑制, 简并$E_n=E_0$不受该抑制, 应在简并空间整体处理. $n=0$第一阶振幅为$1-\ii V_{00}/\hbar$, 其平方不能作为到二阶的存活概率; 一致归一化到二阶应为$P_0=1-\sum_{n\neq0}P_n$（满足微扰小概率条件）. 对角项主要给相位.
\end{solution}
\originalsection{第5题：二氧化碳振动模与选择规则（25分）}
\begin{exercise}
三个等质量$m$的一维点质量按O--C--O以两根常数$k$弹簧连接,$H_0=\sum p_j^2/(2m)+k[(x_1-x_2)^2+(x_2-x_3)^2]/2$. (a)对角化$K=(k/2)\left(\begin{smallmatrix}1&-1&0\\-1&2&-1\\0&-1&1\end{smallmatrix}\right)=R\Lambda R^{\mathsf T}$, 按特征值降序, 解释零值. (b)写$y=R^{\mathsf T}x$的势. (c)变换动量$\pi=R^{\mathsf T}p$, 验证对易与动能. (d)求以$k,m$表示的$\omega_1,\omega_2$, 并用$y,\pi$构造$a_1,a_2$, 使$H_0=\sum_{i=1}^2\hbar\omega_i(a_i^\dagger a_i+1/2)+\pi_3^2/(2m)$, 即两谐振子和一个自由模. 验证$[a_i,a_j^\dagger]=\delta_{ij}$（$i,j=1,2$）以及$[a_1,a_2]=[a_1^\dagger,a_2^\dagger]=0$. 另作补充验算$[a_1^\dagger,a_2]=0$. (e)氧各$-q$、碳$2q$, 求电偶极与能吸收振动光的模.
\end{exercise}
\begin{solution}
选正交列
\[
R=\begin{pmatrix}1/\sqrt6&1/\sqrt2&1/\sqrt3\\-2/\sqrt6&0&1/\sqrt3\\1/\sqrt6&-1/\sqrt2&1/\sqrt3\end{pmatrix},\qquad
\Lambda=\operatorname{diag}(3k/2,k/2,0).
\]
零值列三粒子同位移, 是自由质心平移, 不拉伸弹簧. 势变为$3ky_1^2/2+ky_2^2/2$, 动能因正交保持为$\sum\pi_j^2/(2m)$. 对易$[y_i,\pi_j]=\sum_{ab}R_{ai}R_{bj}[x_a,p_b]=\ii\hbar(R^{\mathsf T}R)_{ij}=\ii\hbar\delta_{ij}$, 其余同类对易0. 频率$\omega_1=\sqrt{3k/m}$,$\omega_2=\sqrt{k/m}$, $a_i=\sqrt{m\omega_i/(2\hbar)}y_i+\ii\pi_i/\sqrt{2m\hbar\omega_i}$, 给
\[
H_0=\hbar\omega_1(a_1^\dagger a_1+1/2)+\hbar\omega_2(a_2^\dagger a_2+1/2)+\pi_3^2/(2m).
\]
具体令$A_i=\sqrt{m\omega_i/(2\hbar)}$, $B_i=1/\sqrt{2m\hbar\omega_i}$, 则$a_i=A_iy_i+\ii B_i\pi_i$, $A_iB_i=1/(2\hbar)$. 从(c)的$[y_i,\pi_j]=\ii\hbar\delta_{ij}$及同类对易为0, 直接展开得
\[
[a_i,a_j^\dagger]=\hbar(A_iB_j+B_iA_j)\delta_{ij}=\delta_{ij},\qquad
[a_i,a_j]=\hbar(B_iA_j-A_iB_j)\delta_{ij}=0.
\]
因此$[a_1,a_2]=0$, 对它取伴随并交换顺序还得$[a_1^\dagger,a_2^\dagger]=0$, 完成原题的两条要求. 补充的混合对易也有$[a_1^\dagger,a_2]=-[a_2,a_1^\dagger]=-\delta_{21}=0$, 两模升降独立.  偶极$\boldsymbol d=-q(x_1-2x_2+x_3)\hat x=-\sqrt6qy_1\hat x=-q\sqrt{3\hbar/(m\omega_1)}(a_1+a_1^\dagger)\hat x$. 只有模1与均匀电场耦合, 其$\Delta n_1=\pm1$, 模2与质心模不在此电偶极近似中吸收. 该结论按题给等质量、固定部分电荷和一维近似; 真实CO$_2$有弯曲模和不同核质量, 不能把简化模型的全部频率当成实验值.
\end{solution}

\originalchapter{来源对应与覆盖边界}
本册正文为按先修顺序的中文重编, 不是逐页全文翻译. 下表给主要来源, 独立补充的证明和算例也在相应章内展开. 第15章含时微扰的基本公式为独立推导, 没有复制被暂缓打包的8.06第4章中的教材改编例题. 三门课程教师均为Barton Zwiebach. 8.04原讲义另署首版LaTeX转录者Sarah Geller、Andrew Turner; 第1、2、5、6、11、12、16、19讲署Geller, 第3、7、8、9、10、13、18讲署Turner, 第4、14--15、20--21、22讲署二人. 第17讲未另署转录者, 不代填姓名. 本册署名保留这些贡献者, 新增中文说明与题解为AI独立整理.
\begin{longtable}{p{.05\textwidth}p{.38\textwidth}p{.47\textwidth}}
\toprule 章 & 内容 & 原件资源页链接\\\midrule\endhead
1 & 复数振幅与物质波 & \href{https://ocw.mit.edu/courses/8-04-quantum-physics-i-spring-2016/resources/mit8_04s16_lecnotes1/}{8.04 LecNotes1}; \href{https://ocw.mit.edu/courses/8-04-quantum-physics-i-spring-2016/resources/mit8_04s16_lecnotes2/}{8.04 LecNotes2}; \href{https://ocw.mit.edu/courses/8-04-quantum-physics-i-spring-2016/resources/mit8_04s16_lecnotes3/}{8.04 LecNotes3}; \href{https://ocw.mit.edu/courses/8-04-quantum-physics-i-spring-2016/resources/mit8_04s16_lecnotes4/}{8.04 LecNotes4}\\
2 & 波函数与概率 & \href{https://ocw.mit.edu/courses/8-04-quantum-physics-i-spring-2016/resources/mit8_04s16_lecnotes5/}{8.04 LecNotes5}; \href{https://ocw.mit.edu/courses/8-04-quantum-physics-i-spring-2016/resources/mit8_04s16_lecnotes6/}{8.04 LecNotes6}\\
3 & 薛定谔方程与概率流 & \href{https://ocw.mit.edu/courses/8-04-quantum-physics-i-spring-2016/resources/mit8_04s16_lecnotes5/}{8.04 LecNotes5}; \href{https://ocw.mit.edu/courses/8-04-quantum-physics-i-spring-2016/resources/mit8_04s16_lecnotes6/}{8.04 LecNotes6}\\
4 & 动量表象与不确定性 & \href{https://ocw.mit.edu/courses/8-04-quantum-physics-i-spring-2016/resources/mit8_04s16_lecnotes7/}{8.04 LecNotes7}; \href{https://ocw.mit.edu/courses/8-04-quantum-physics-i-spring-2016/resources/mit8_04s16_lecnotes8/}{8.04 LecNotes8}; \href{https://ocw.mit.edu/courses/8-04-quantum-physics-i-spring-2016/resources/mit8_04s16_lecnotes9/}{8.04 LecNotes9}\\
5 & 定态与一维势阱 & \href{https://ocw.mit.edu/courses/8-04-quantum-physics-i-spring-2016/resources/mit8_04s16_lecnotes10/}{8.04 LecNotes10}; \href{https://ocw.mit.edu/courses/8-04-quantum-physics-i-spring-2016/resources/mit8_04s16_lecnotes11/}{8.04 LecNotes11}\\
6 & 反射、透射与隧穿 & \href{https://ocw.mit.edu/courses/8-04-quantum-physics-i-spring-2016/resources/mit8_04s16_lecnotes16/}{8.04 LecNotes16}; \href{https://ocw.mit.edu/courses/8-04-quantum-physics-i-spring-2016/resources/mit8_04s16_lecnotes17/}{8.04 LecNotes17}\\
7 & 谐振子与Gaussian波包 & \href{https://ocw.mit.edu/courses/8-04-quantum-physics-i-spring-2016/resources/mit8_04s16_lecnotes7/}{8.04 LecNotes7}; \href{https://ocw.mit.edu/courses/8-04-quantum-physics-i-spring-2016/resources/mit8_04s16_lecnotes14_15/}{8.04 LecNotes14\_15}\\
8 & 态空间、测量与动力学 & \href{https://ocw.mit.edu/courses/8-05-quantum-physics-ii-fall-2013/resources/lecture-notes-4/}{8.05 Chap\_04}; \href{https://ocw.mit.edu/courses/8-05-quantum-physics-ii-fall-2013/resources/lecture-notes-5/}{8.05 Chap\_05}; \href{https://ocw.mit.edu/courses/8-05-quantum-physics-ii-fall-2013/resources/lecture-notes-6/}{8.05 Chap\_06}\\
9 & 自旋与二态系统 & \href{https://ocw.mit.edu/courses/8-05-quantum-physics-ii-fall-2013/resources/lecture-notes-2/}{8.05 Chap\_02}; \href{https://ocw.mit.edu/courses/8-05-quantum-physics-ii-fall-2013/resources/lecture-notes-7/}{8.05 Chap\_07}\\
10 & 角动量与中心势 & \href{https://ocw.mit.edu/courses/8-05-quantum-physics-ii-fall-2013/resources/lecture-notes-9/}{8.05 Chap\_09}; \href{https://ocw.mit.edu/courses/8-05-quantum-physics-ii-fall-2013/resources/lecture-notes-10/}{8.05 Chap\_10}\\
11 & 氢原子的量子化 & \href{https://ocw.mit.edu/courses/8-04-quantum-physics-i-spring-2016/resources/mit8_04s16_lecnotes22/}{8.04 LecNotes22}\\
12 & 复合系统、纠缠与全同粒子 & \href{https://ocw.mit.edu/courses/8-05-quantum-physics-ii-fall-2013/resources/lecture-notes-8/}{8.05 Chap\_08}\\
13 & 变分原理与非简并微扰 & \href{https://ocw.mit.edu/courses/8-06-quantum-physics-iii-spring-2018/resources/chapter-1-non-degenerate-and-degenerate-perturbation-theory/}{8.06 ch1}\\
14 & 简并微扰与能级分裂 & \href{https://ocw.mit.edu/courses/8-06-quantum-physics-iii-spring-2018/resources/chapter-1-non-degenerate-and-degenerate-perturbation-theory/}{8.06 ch1}; \href{https://ocw.mit.edu/courses/8-06-quantum-physics-iii-spring-2018/resources/chapter-2-hydrogen-fine-structure/}{8.06 ch2}\\
15 & 含时微扰与绝热演化 & \href{https://ocw.mit.edu/courses/8-06-quantum-physics-iii-spring-2018/resources/chapter-6-adiabatic-approximation/}{8.06 ch6}\\
16 & 半经典近似与三维散射 & \href{https://ocw.mit.edu/courses/8-06-quantum-physics-iii-spring-2018/resources/chapter-3-semiclassical-approximation/}{8.06 ch3}; \href{https://ocw.mit.edu/courses/8-06-quantum-physics-iii-spring-2018/resources/chapter-7-scattering/}{8.06 ch7}\\
17 & 带电粒子与氢原子的量子对称性 & 独立推导; 未冒称缺失原件翻译\\
\bottomrule\end{longtable}
\originalsection{考核覆盖}
19份保留的完整题面考核共112道大题, 所有小问已在考核篇转述并独立解答. 本次选定课程页面没有取得官方解答, 全部解答为AI独立计算. 两张公式表作为原件提供, 不计入已解答试卷; 暂缓打包的混合文件只记录URL和校验事实, 不补造缺失的教材题面. 下表逐份列出作业与测试状态, 历史试卷年份按题面和网页分别记录. 公式表只作为原件归档, 不计为已解答试卷.
\begin{longtable}{p{.10\textwidth}p{.47\textwidth}p{.33\textwidth}}
\toprule 课程 & 文件资源页 & 中文处理状态\\\midrule\endhead
8.04 & \href{https://ocw.mit.edu/courses/8-04-quantum-physics-i-spring-2016/resources/mit8_04s16_ps1_2016/}{Problem Set 1} & 所有小问中文题面及独立解答\\
8.04 & \href{https://ocw.mit.edu/courses/8-04-quantum-physics-i-spring-2016/resources/mit8_04s16_ps2_2016/}{Problem Set 2} & 所有小问中文题面及独立解答\\
8.04 & \href{https://ocw.mit.edu/courses/8-04-quantum-physics-i-spring-2016/resources/mit8_04s16_ps3_2016/}{Problem Set 3} & 教材依赖或改编待核, 保守暂缓打包\\
8.04 & \href{https://ocw.mit.edu/courses/8-04-quantum-physics-i-spring-2016/resources/mit8_04s16_ps4_2016/}{Problem Set 4} & 所有小问中文题面及独立解答\\
8.04 & \href{https://ocw.mit.edu/courses/8-04-quantum-physics-i-spring-2016/resources/mit8_04s16_ps5_2016/}{Problem Set 5} & 所有小问中文题面及独立解答\\
8.04 & \href{https://ocw.mit.edu/courses/8-04-quantum-physics-i-spring-2016/resources/mit8_04s16_ps6_2016/}{Problem Set 6} & 教材依赖或改编待核, 保守暂缓打包\\
8.04 & \href{https://ocw.mit.edu/courses/8-04-quantum-physics-i-spring-2016/resources/mit8_04s16_ps7_2016/}{Problem Set 7} & 教材依赖或改编待核, 保守暂缓打包\\
8.04 & \href{https://ocw.mit.edu/courses/8-04-quantum-physics-i-spring-2016/resources/mit8_04s16_ps8_2016/}{Problem Set 8} & 教材依赖或改编待核, 保守暂缓打包\\
8.04 & \href{https://ocw.mit.edu/courses/8-04-quantum-physics-i-spring-2016/resources/mit8_04s16_ps9_2016/}{Problem Set 9} & 所有小问中文题面及独立解答\\
8.04 & \href{https://ocw.mit.edu/courses/8-04-quantum-physics-i-spring-2016/resources/mit8_04s16_ps10_2016/}{Problem Set 10} & 教材依赖或改编待核, 保守暂缓打包\\
8.04 & \href{https://ocw.mit.edu/courses/8-04-quantum-physics-i-spring-2016/resources/mit8_04s16_test1_fall2015/}{Mid-term Test 1 from 2015} & 所有小问中文题面及独立解答\\
8.04 & \href{https://ocw.mit.edu/courses/8-04-quantum-physics-i-spring-2016/resources/mit8_04s16_test2_fall2015/}{Mid-term Test 2 from 2015} & 所有小问中文题面及独立解答\\
8.04 & \href{https://ocw.mit.edu/courses/8-04-quantum-physics-i-spring-2016/resources/mit8_04s16_finaltest_2015/}{Final Exam from 2015} & 所有小问中文题面及独立解答\\
8.05 & \href{https://ocw.mit.edu/courses/8-05-quantum-physics-ii-fall-2013/resources/assignment-1/}{Problem Set 1} & 教材依赖或改编待核, 保守暂缓打包\\
8.05 & \href{https://ocw.mit.edu/courses/8-05-quantum-physics-ii-fall-2013/resources/assignment-2/}{Problem Set 2} & 教材依赖或改编待核, 保守暂缓打包\\
8.05 & \href{https://ocw.mit.edu/courses/8-05-quantum-physics-ii-fall-2013/resources/assignment-3/}{Problem Set 3} & 教材依赖或改编待核, 保守暂缓打包\\
8.05 & \href{https://ocw.mit.edu/courses/8-05-quantum-physics-ii-fall-2013/resources/assignment-4/}{Problem Set 4} & 教材依赖或改编待核, 保守暂缓打包\\
8.05 & \href{https://ocw.mit.edu/courses/8-05-quantum-physics-ii-fall-2013/resources/assignment-5/}{Problem Set 5} & 教材依赖或改编待核, 保守暂缓打包\\
8.05 & \href{https://ocw.mit.edu/courses/8-05-quantum-physics-ii-fall-2013/resources/assignment-6/}{Problem Set 6} & 所有小问中文题面及独立解答\\
8.05 & \href{https://ocw.mit.edu/courses/8-05-quantum-physics-ii-fall-2013/resources/assignment-7/}{Problem Set 7} & 所有小问中文题面及独立解答\\
8.05 & \href{https://ocw.mit.edu/courses/8-05-quantum-physics-ii-fall-2013/resources/assignment-8/}{Problem Set 8} & 教材依赖或改编待核, 保守暂缓打包\\
8.05 & \href{https://ocw.mit.edu/courses/8-05-quantum-physics-ii-fall-2013/resources/assignment-9/}{Problem Set 9} & 所有小问中文题面及独立解答\\
8.05 & \href{https://ocw.mit.edu/courses/8-05-quantum-physics-ii-fall-2013/resources/assignment-10/}{Problem Set 10} & 教材依赖或改编待核, 保守暂缓打包\\
8.05 & \href{https://ocw.mit.edu/courses/8-05-quantum-physics-ii-fall-2013/resources/assignment-11/}{Problem Set 11} & 教材依赖或改编待核, 保守暂缓打包\\
8.05 & \href{https://ocw.mit.edu/courses/8-05-quantum-physics-ii-fall-2013/resources/test-formula-sheet/}{Test Formula Sheet} & 公式表原件已归档\\
8.05 & \href{https://ocw.mit.edu/courses/8-05-quantum-physics-ii-fall-2013/resources/test-2011/}{2011 Test} & 所有小问中文题面及独立解答\\
8.05 & \href{https://ocw.mit.edu/courses/8-05-quantum-physics-ii-fall-2013/resources/test-2012/}{2012 Test} & 所有小问中文题面及独立解答\\
8.05 & \href{https://ocw.mit.edu/courses/8-05-quantum-physics-ii-fall-2013/resources/test-2013/}{2013 Test} & 所有小问中文题面及独立解答\\
8.05 & \href{https://ocw.mit.edu/courses/8-05-quantum-physics-ii-fall-2013/resources/final-exam-formula-sheet/}{Final Exam Formula Sheet} & 公式表原件已归档\\
8.05 & \href{https://ocw.mit.edu/courses/8-05-quantum-physics-ii-fall-2013/resources/exam-2011/}{2011 Final Exam} & 所有小问中文题面及独立解答\\
8.05 & \href{https://ocw.mit.edu/courses/8-05-quantum-physics-ii-fall-2013/resources/exam-2012/}{2012 Final Exam} & 所有小问中文题面及独立解答\\
8.05 & \href{https://ocw.mit.edu/courses/8-05-quantum-physics-ii-fall-2013/resources/exam-2013/}{2013 Final Exam} & 所有小问中文题面及独立解答\\
8.06 & \href{https://ocw.mit.edu/courses/8-06-quantum-physics-iii-spring-2018/resources/problem-set-1/}{Problem Set 1} & 所有小问中文题面及独立解答\\
8.06 & \href{https://ocw.mit.edu/courses/8-06-quantum-physics-iii-spring-2018/resources/problem-set-2/}{Problem Set 2} & 教材依赖或改编待核, 保守暂缓打包\\
8.06 & \href{https://ocw.mit.edu/courses/8-06-quantum-physics-iii-spring-2018/resources/problem-set-3/}{Problem Set 3} & 教材依赖或改编待核, 保守暂缓打包\\
8.06 & \href{https://ocw.mit.edu/courses/8-06-quantum-physics-iii-spring-2018/resources/problem-set-4/}{Problem Set 4} & 教材依赖或改编待核, 保守暂缓打包\\
8.06 & \href{https://ocw.mit.edu/courses/8-06-quantum-physics-iii-spring-2018/resources/problem-set-5/}{Problem Set 5} & 所有小问中文题面及独立解答\\
8.06 & \href{https://ocw.mit.edu/courses/8-06-quantum-physics-iii-spring-2018/resources/problem-set-6/}{Problem Set 6} & 教材依赖或改编待核, 保守暂缓打包\\
8.06 & \href{https://ocw.mit.edu/courses/8-06-quantum-physics-iii-spring-2018/resources/problem-set-7/}{Problem Set 7} & 教材依赖或改编待核, 保守暂缓打包\\
8.06 & \href{https://ocw.mit.edu/courses/8-06-quantum-physics-iii-spring-2018/resources/problem-set-8/}{Problem Set 8} & 教材依赖或改编待核, 保守暂缓打包\\
8.06 & \href{https://ocw.mit.edu/courses/8-06-quantum-physics-iii-spring-2018/resources/problem-set-9/}{Problem Set 9} & 教材依赖或改编待核, 保守暂缓打包\\
\bottomrule\end{longtable}
\originalsection{覆盖边界与许可}
正文已经补入打靶与节点排序、$\delta$势、相干与压缩、三种动力学表象、角动量合成、精细结构三项、含时跃迁与Berry相位、Airy连接推导、带电粒子、散射和全同粒子应用. 本册仍是按概念组织的选定课程学习册, 不是每页原讲义的全文翻译. 复路径变形、非终止超几何级数的增长结论、低能Dirac有效项及Runge--Lenz长算符恒等式明确作为进阶输入, 不假称已在本册从零证明. 没有包含场论、辐射修正或完整多体理论.

8.06 Spring 2018第5章原件缺失, 两次考试明确不向OCW公开, 课纲述末周不计分作业也未在作业页提供. 没有补造题面. 默认许可证不自动覆盖商业教材; 对题面依赖和明确改编信用未充分确认的混合文件, 整份保守暂缓打包. 引用、书目或改编信用本身不等于已证实的版权排除, 逐项记录区分观察事实、权利不确定性与当前打包决定. URL、页数和SHA256仍记录在清单. 其中的数学物理事实可独立推导, 本册不复制该类题面或外站图片.

课程入口: \href{https://ocw.mit.edu/courses/8-04-quantum-physics-i-spring-2016/}{8.04 Spring 2016}, \href{https://ocw.mit.edu/courses/8-05-quantum-physics-ii-fall-2013/}{8.05 Fall 2013}, \href{https://ocw.mit.edu/courses/8-06-quantum-physics-iii-spring-2018/}{8.06 Spring 2018}. 每课页面的Syllabus、Calendar、Lecture Notes、Assignments和公开Exams栏目有本地检索快照. 逐文件清单保留直接PDF网址、标题、教师、实际题面年份、页数、字节数、SHA256、明确权利声明与保守暂缓理由.

原件及本中文改编按\href{https://creativecommons.org/licenses/by-nc-sa/4.0/}{CC BY-NC-SA 4.0}署名、非商业、相同方式共享; 另参\href{https://ocw.mit.edu/pages/privacy-and-terms-of-use/}{MIT OCW使用条款}. 本册没有MIT官方中文或官方认可的含义.

\end{document}
